\documentclass[11pt]{amsart}

\usepackage[T1]{fontenc}
\usepackage[utf8]{inputenc}
\usepackage{lmodern}
\usepackage[a4paper,margin=1in]{geometry}
\usepackage{amsmath,amssymb,amsthm,mathtools,mathrsfs}
\usepackage{array,booktabs,longtable}
\usepackage{enumitem}
\usepackage{microtype}
\usepackage{xcolor}
\usepackage[colorlinks=true,linkcolor=blue!45!black,citecolor=blue!45!black,urlcolor=blue!45!black]{hyperref}
\usepackage[nameinlink,noabbrev]{cleveref}

\newtheorem{theorem}{Theorem}[section]
\newtheorem{proposition}[theorem]{Proposition}
\newtheorem{lemma}[theorem]{Lemma}
\newtheorem{corollary}[theorem]{Corollary}
\theoremstyle{definition}
\newtheorem{definition}[theorem]{Definition}
\theoremstyle{remark}
\newtheorem{remark}[theorem]{Remark}

\newcommand{\C}{\mathbf C}
\newcommand{\R}{\mathbf R}
\newcommand{\Q}{\mathbf Q}
\newcommand{\F}{\mathbf F}
\newcommand{\cO}{\mathcal O}
\newcommand{\p}{\mathfrak p}
\newcommand{\Tr}{\operatorname{Tr}}
\newcommand{\Nm}{\operatorname{N}}
\newcommand{\Gal}{\operatorname{Gal}}
\newcommand{\ph}{\operatorname{ph}}
\newcommand{\ord}{\operatorname{ord}}
\newcommand{\cond}{\operatorname{cond}}
\newcommand{\Res}{\operatorname{Res}}
\newcommand{\abs}[1]{\lvert #1\rvert}
\newcommand{\set}[1]{\left\{#1\right\}}
\newcommand{\Err}{\mathcal E}
\newcommand{\Aphase}{\mathcal A}
\newcommand{\HerbPhi}{\varphi^{\mathrm H}}
\newcommand{\HerbPsi}{\psi^{\mathrm H}}
\newcommand{\D}{\mathfrak D}
\newcommand{\eps}{\varepsilon}
\newcommand{\ceil}[1]{\left\lceil #1\right\rceil}
\newcommand{\floor}[1]{\left\lfloor #1\right\rfloor}


% Alternative macro names retained from the SML calculations.
\newcommand{\OO}{\mathcal O}\newcommand{\pp}{\mathfrak p}\newcommand{\PP}{\mathfrak P}
\newcommand{\ee}{\mathrm e}\newcommand{\echar}{\mathrm e}\newcommand{\FF}{\mathbf F}
\newcommand{\CC}{\mathbf C}\newcommand{\TT}{\mathcal T}\newcommand{\ind}{\mathbf1}
\newcommand{\pr}{\operatorname{pr}_0}\newcommand{\tr}{\operatorname{tr}}
\newcommand{\Cart}{\mathcal C}\newcommand{\muq}{\mu}
\setlength{\emergencystretch}{3em}

\title[The Second Main Lemma for Local Epsilon Factors]
{A Local Proof of Langlands's Second Main Lemma\\
for Local Epsilon Factors over Nonarchimedean Local Fields}
\author{Fukuhiro Ueda}
\address{Research Institute for Mathematical Sciences, Kyoto University, Kyoto 606--8502, Japan}
\email{fueda@kurims.kyoto-u.ac.jp}
\date{September 2026}
\subjclass[2020]{11S37, 11S40, 11R42}
\keywords{local epsilon factors, Artin $L$-functions, local constants,
Gauss sums, ramification, Lamprecht formula, biquadratic extensions,
equal characteristic}
\hypersetup{
  pdftitle={A Local Proof of Langlands's Second Main Lemma for Local Epsilon Factors over Nonarchimedean Local Fields},
  pdfauthor={Fukuhiro Ueda},
  pdfsubject={A local proof of the Second Main Lemma in mixed and equal characteristic},
  pdfkeywords={local epsilon factors, Artin L-functions, Gauss sums, ramification, Second Main Lemma},
  bookmarksdepth=2
}
\setcounter{tocdepth}{2}
\begin{document}
\begin{abstract}
We give a local proof of Langlands's Second Main Lemma for normalized
one-dimensional local epsilon factors over every nonarchimedean local
field, in both mixed and equal characteristic. The theorem asserts
independence of the inducing field in a $C_\ell\times C_\ell$ extension,
with the complete lower-field correction factors included. The
one-dimensional inputs are the author's First Main Lemma and the local
ramification and stationary-phase results established in its companion
paper. All specialized two-field calculations are given here.

The odd-prime argument constructs compatible character models and uses
weighted norm--trace estimates to reduce the comparison to a root-of-unity
ambiguity removed by the First Main Lemma. The dyadic argument instead
compares the complete residual functions, including their affine terms.
Common norm representatives coordinate the stationary values, and the
third quadratic field supplies the cancellation needed at a critical
norm image of index two. Separate local constructions treat equal
characteristic and the maximal and nonmaximal mixed-characteristic
quadratic breaks. A characteristic-uniform tame calculation and an
exhaustive ramification and conductor analysis complete the proof.

The methods retain the explicit local character of the work of Dwork,
Langlands, and Lakkis. Their two-field comparison results are not used as
premises. The remaining quoted inputs are stated classical results in
local field theory. The leading finite-field Gauss congruence and the
formal local trace--residue identity are proved in appendices; no global
functional equation is used.
\end{abstract}
\maketitle
\clearpage
\begingroup
\raggedbottom
\tableofcontents
\clearpage
\endgroup

\section{Introduction}\label{sec:introduction}
The First Main Lemma relates local constants along a cyclic extension
of prime degree. The Second Main Lemma compares two different such
routes through the same field diagram. A character on the top field
can be obtained by norm pullback from different intermediate fields;
the corresponding local constants, with their lower-field correction
factors, must give the same answer. The purpose of this paper is to
prove that compatibility locally, with every residue characteristic and
both mixed and equal characteristic included.

Let $F$ be a nonarchimedean local field with finite residue field of
characteristic $p$. Let $K/F$ be Galois with group $C_\ell\times C_\ell$,
where $\ell$ is prime. Let $\Theta:K^\times\to\CC^\times$ be a continuous
quasi-character invariant under $\Gal(K/F)$ and not of the form
$\theta_F\circ N_{K/F}$. For each degree-$\ell$ intermediate field $L$
choose $\theta_L$ with $\theta_L\circ N_{K/L}=\Theta$. Such extensions
exist by cyclic Hilbert 90 and extension of characters from the norm image.
Fix a nontrivial continuous additive character $\psi_F$, and put
$\psi_L=\psi_F\circ\Tr_{L/F}$. Write
\[
 S(L/F)=\{\nu\in\operatorname{Hom}_{\mathrm{cont}}(F^\times,\CC^\times):
             \nu\circ N_{L/F}=1\}.
\]
Here $\operatorname{Hom}_{\mathrm{cont}}$ denotes continuous group homomorphisms.
\begin{theorem}[Second Main Lemma]\label{U:main}
The complete normalized induction expression
\[
 \mathcal A_L(\theta_L)=\Delta_L(\theta_L,\psi_L)
                \prod_{\nu\in S(L/F)}\Delta_F(\nu,\psi_F)
\]
is independent of $L$. The assertion holds in mixed and equal
characteristic, for every residue characteristic and every prime $\ell$.
\end{theorem}
The proof is completed in Section~\ref{sec:completion}. An actual family of fields and
characters is the datum of the theorem; the existence of auxiliary
stationary models is proved within the relevant branches, not assumed in
its statement.


Here $\Delta_E(\theta,\psi)$ is the normalized local Gauss factor used in
the companion First Main Lemma paper~\cite{Ueda}; it is recalled in
Section~\ref{sec:local-notation}. The trivial character is included in
$S(L/F)$ and contributes the factor $1$. In a biquadratic extension the
remaining lower-field factor is the local constant of the quadratic
norm character. That factor is retained throughout the dyadic proof.

\subsection*{Historical background}
The statement is Langlands's Second Main Lemma~\cite{Langlands}.
Dwork's and Lakkis's explicit norm, trace, and Artin--Schreier
calculations provide important methodological background for its local
study~\cite{Dwork,Lakkis}. The present treatment retains that local
framework while giving the specialized two-field constructions and
comparisons within a single argument. These historical sources are
acknowledged for the problem and methods; the reader need not consult
them to supply an unstated step in the proof. No claim of historical
priority is made for the individual constructions or alternative
arguments. The leading Gauss congruence and the local residue formalism
have classical antecedents in \cite{Stickelberger,TateResidue}; the
precise results used here are proved in Appendices~\ref{app:leading-gauss}
and \ref{app:local-residues}, so these two references are historical,
not proof inputs.

\subsection*{Relation to the First Main Lemma}
The one-dimensional foundation is the author's companion
paper~\cite{Ueda}. We use its First Main Lemma, complete Lamprecht
formula, stable-twist identity, stationary duality, norm filtration,
conductor transport, and trace and different formulas. Their precise
forms are recalled in Sections~\ref{sec:local-notation}--\ref{sec:stationary}.
References to that paper use its theorem and lemma numbers. The
additional classical foundations are stated in Section~\ref{sec:ramification}.
The diamond-level norm and ramification facts are proved in
Appendix~\ref{app:diamond-foundations} from the cyclic-prime local results
and elementary algebra; neither the general local reciprocity theorem
nor the general Herbrand quotient theorem is required. Hilbert~90,
character extension, different transitivity, and the order-discriminant
bound are proved in Appendix~\ref{app:descent}. Basic algebra already
available in Mathlib is used without a redundant general development.

The First Main Lemma has a Lean formalization. The present paper is the
written mathematical proof of the Second Main Lemma; no Lean
verification of the new two-field argument is asserted. Its statements
and proofs are presented independently of a formalization workflow.

\subsection*{Contributions of the present paper}
The comparison must concern the given characters, rather than auxiliary
models whose conductors happen to agree. We therefore construct the
models, extend them to the full multiplicative groups, and prove that
the prescribed characters are obtained from them by a common descending
twist. Stationary representatives are then chosen with their complete
character formulas in force. Every use of an exact norm is supported
by an explicit construction or by a proved freedom of normalization;
a subcritical norm approximation is not promoted to an exact norm.

In wild odd degree, truncated logarithms, actual-conjugate trace
calculations, and two weighted terminal estimates control the full
four-factor comparison. The cubic boundary is treated with both of its
affine stationary corrections. In residue characteristic two, equality
of squares does not settle the comparison. We instead match the
complete residual functions, including the linear terms of their
quadratic refinements. When a critical norm map has image of index two,
the third quadratic field supplies a translation that cancels the
complementary coset in the actual sum.

\subsection*{Logical organization of the proof}
We first fix the local constant and conductor conventions, state the
one-dimensional inputs, and isolate several common comparison
principles. The tame case is treated uniformly when the extension
degree differs from the residue characteristic. The wild cases are
then separated according to inertia and, in the dyadic totally ramified
case, according to the characteristic and quadratic breaks.

For orientation, the principal branches are as follows. In the two
mixed-characteristic dyadic rows, $e=v_F(2)$ and the three entries are
the lower breaks of the quadratic intermediate fields.
\[
\begin{array}{c|l}
\text{configuration} & \text{calculation used}\\ \hline
\ell\ne p & \text{tame comparison}\\
\ell=p>2,\ |I|=p & \text{odd unramified--ramified comparison}\\
\ell=p>2,\ |I|=p^2 & \text{odd totally ramified comparison}\\
p=\ell=2,\ |I|=2 & \text{dyadic unramified--ramified comparison}\\
p=\ell=2,\ |I|=4,\ \operatorname{char}F=2
 & \text{Artin--Schreier construction}\\
(2a-1,2e,2e),\quad 1\le a\le e
 & \text{maximal quadratic-break construction}\\
(2a-1,2r-1,2r-1),\quad 1\le a\le r\le e
 & \text{nonmaximal quadratic-break construction}.
\end{array}
\]
Here $I$ denotes the inertia subgroup. The final section proves that
these possibilities exhaust the theorem. Within each branch, a
common-twist realization reduces the given characters to consecutive
minimal, intermediate, and high conductor ranges, or to a direct
comparison at all higher conductors.

\subsection*{Mixed and equal characteristic}
The odd-prime calculations are expressed by integral identities whose
mixed-characteristic error terms are kept until their valuations have
been checked. The corresponding terms vanish in equal characteristic.
The dyadic totally ramified arguments require different constructions:
local residues and Artin--Schreier equations are used in characteristic
two, while quadratic norm symbols and aligned Kummer generators are
used in mixed characteristic. Neither argument is obtained by simply
substituting into the other. In each case the elementary stationary
values and residual sums are compared together.

\subsection*{Organization}
Sections~\ref{sec:local-notation}--\ref{sec:stationary} give the common
local and stationary-phase preliminaries.
Section~\ref{sec:tame} proves the tame comparison.
Sections~\ref{sec:odd-ur}--\ref{sec:odd-calculation} give the wild
odd-prime proof, beginning with an unramified intermediate field and
then treating total ramification. Section~\ref{sec:dyadic-ur} handles
dyadic extensions with an unramified intermediate field.
Sections~\ref{sec:dyadic-equal}--\ref{sec:dyadic-nonmaximal} treat
the totally ramified dyadic cases. Section~\ref{sec:completion}
assembles the complete theorem. Appendix~\ref{app:leading-gauss} proves
the leading Gauss congruence by Jacobi sums and digit induction;
Appendix~\ref{app:local-residues} proves formal residue substitution,
trace compatibility, and the logarithmic norm identity;
Appendix~\ref{app:descent} supplies the elementary descent and
ramification arguments named above.


\section{Local constants and the elementary reduction}\label{sec:local-notation}\label{U:statement}
\subsection{Valuations, conductors, and local constants}
For a local field $E$, write $\OO_E$ for its integers, $\pp_E$ for its
maximal ideal, $k_E$ for its residue field, and $v_E(\pp_E)=1$.
Put $U_E^0=\OO_E^\times$ and $U_E^j=1+\pp_E^j$ for $j\ge1$.
The multiplicative conductor $m_E(\theta)$, also denoted $a_E(\theta)$
in the later case calculations, is the least $m\ge0$ for which
$\theta|_{U_E^m}=1$. The additive conductor $n_E(\psi)$ is specified by
$\pp_E^{-n_E(\psi)}$, the largest \emph{ideal} on which $\psi$ is trivial.
It is not generally the full kernel of $\psi$.

We use exactly the normalized local constant of \cite{Ueda}.
For a nonzero complex number $z$ put $\ph(z)=z/|z|$.
If $m=m_E(\theta)>0$, choose an admissible $\Gamma\in E^\times$ with
$v_E\Gamma=m+n_E(\psi)$ and set
\begin{equation}\label{U:delta-definition}
 \Delta_E(\theta,\psi)=\theta(\Gamma)\,
 \ph\!\left(\int_{\OO_E^\times}\psi(u/\Gamma)\theta(u)^{-1}\,du\right).
\end{equation}
The nonzero Gauss integral, or its positive-scalar finite-sum realization,
and independence of $\Gamma$ are proved in \cite{Ueda}.
For $m=0$ this convention gives
$\Delta_E(\theta,\psi)=\theta(\pi_E)^{n_E(\psi)}$.
The choice of positive Haar normalization does not change the phase.
The characters in the theorem need not be unitary; their values on
nonunits remain in the displayed multiplicative factors.

All valuations are normalized on the field named. Thus for a totally
ramified degree-$\ell$ edge $E/F$,
\[
 v_E|_F=\ell v_F,\qquad v_F(N_{E/F}x)=v_E(x).
\]
The different exponent $D_{E/F}$ is measured in $E$.
Conjugation of a character means
$\theta^\sigma(x)=\theta(\sigma^{-1}x)$.
In $A\subset B_i\subset B$, symbols $N_i,n_i$ denote the upper and lower
norms $N_{B/B_i},N_{B_i/A}$, respectively; the named domains are part of
every formula. Analogous upper and lower traces are distinguished in each
case calculation.

\subsection{The First Main Lemma and elementary identities}\label{U:interface}
For a cyclic extension $E/F$ of prime degree and a quasi-character $\chi$
of $F^\times$, the First Main Lemma is
\begin{equation}\label{U:FML}
 \Delta_E(\chi\circ N_{E/F},\psi_F\circ\Tr_{E/F})
       \prod_{\nu\in S(E/F)}\Delta_F(\nu,\psi_F)
   =\prod_{\nu\in S(E/F)}\Delta_F(\chi\nu,\psi_F).
\end{equation}
Also
\[
 \Delta_E(1,\psi)=1,\qquad
 \Delta_E(\theta,\psi(c\,\cdot))=\theta(c)\Delta_E(\theta,\psi),
\]
and for a finite-order character $\mu$,
$\Delta_E(\mu,\psi)\Delta_E(\mu^{-1},\psi)=\mu(-1)$.
These are \cite[Theorem~1.1 and Lemmas~2.8, 2.10--2.11]{Ueda}.
A change of variables in \eqref{U:delta-definition} also preserves the
local constant under a field automorphism when the additive character is
trace-invariant.


\section{Ramification and classical local inputs}\label{sec:ramification}
\subsection{Trace ideals, norm filtration, and conductors}
For a finite separable extension of ramification index $e_{E/F}$,
\begin{equation}\label{U:trace-ideal}
 \Tr_{E/F}(\pp_E^j)=\pp_F^{\lfloor(j+D_{E/F})/e_{E/F}\rfloor},
 \qquad n_E(\psi_F\Tr)=e_{E/F}n_F(\psi_F)+D_{E/F}.
\end{equation}
Here $j$ may be negative. To see the asserted generality, write
$\Tr_{E/F}(\pp_E^j)=\pp_F^c$. The inverse different is the trace-dual
of $\OO_E$. For $y\in F$, the condition
$y\Tr(\pp_E^j)\subset\OO_F$ is therefore equivalent to
$y\pp_E^j\subset\pp_E^{-D_{E/F}}$, or
$e_{E/F}v_F(y)+j\ge-D_{E/F}$. Comparing the two principal dual
ideals gives the displayed floor formula. The conductor formula follows
by applying it to the largest trivial ideal of $\psi_F$.
The cyclic-prime instances used most often below are
\cite[Theorem~3.7 and Corollary~3.8]{Ueda}.
For a totally ramified cyclic prime-degree edge of lower break $t$,
$D_{E/F}=(\ell-1)(t+1)$, and the inverse Herbrand function on
nonnegative real arguments is $s$ up to $t$, and $t+\ell(s-t)$ above $t$.
The norm filtration \cite[Theorem~3.7]{Ueda} gives the whole unit
images above the break, isomorphisms of the subcritical graded norm
quotients, and an image of index $\ell$ at the critical quotient.
Every nontrivial norm character has conductor $t+1$.


For a base character of conductor $m$, the pullback conductor $m'$ obeys
$m'-1=\psi_{E/F}(m-1)$ when $m>t+1$, and $m'=m$ when $m\le t$.
At $m=t+1$ it is at most $m$, with a drop precisely when twisting by a
norm character lowers the base conductor. These are
\cite[Corollary~3.9 and Proposition~3.10]{Ueda}.
Unramified norm pullback preserves conductor, and its unit norms are
surjective at every positive depth.

Subcritical approximation is used in the following precise form:
for $A\in F^\times$ and $0\le j\le t$ there is $x\in E^\times$ with
$v_E x=v_F A$ and $N_{E/F}x/A\in U_F^j$. This is
\cite[Lemma~6.6]{Ueda}; it makes no assertion
at the critical relative precision $t+1$.
For a wild cyclic degree-$p$ edge and $1\le i<p$, the strong
symmetric-coefficient bound is
\begin{equation}\label{U:symmetric}
 v_F(E_i(x))\ge\left\lfloor
       \frac{i v_E(x)+(p-1)(t+1)}p\right\rfloor.
\end{equation}
It includes fractional ideals (\cite[Lemma~3.4]{Ueda}).
Stationary duality always identifies coefficients from an identity on a
\emph{whole ideal}, never from the value of a character at a single point.
The high stationary tables used in Sections~\ref{sec:odd-ur}--\ref{sec:odd-calculation} are the specified rows of
\cite[Proposition~6.8]{Ueda}; their actual depth hypotheses are
checked at each application. Ordinary Hasse--Davenport is stated in Section~\ref{sec:tame}.

\subsection{Additional classical local results}\label{U:classical}
The following classical local results are the additional prerequisites.
They contain no assertion about the induction of local constants.

\paragraph{Hilbert 90 and extension of characters.}
For a finite cyclic extension with generator $\sigma$, the multiplicative
norm kernel consists of $y/\sigma(y)$ (equivalently $\sigma(y)/y$), and
the additive trace kernel consists of $\sigma(y)-y$.
Both assertions are proved in Lemma~\ref{C:hilbert90}.
The continuous character-extension statement used in realizing the
models is proved in Lemma~\ref{C:character-extension}. In particular,
agreement on overlaps must be proved before the extension lemma is
applied; existence of a desired character is never an extra assumption.

\paragraph{Norm characters in the diamond.}
The order and conductor of a cyclic-prime norm quotient are supplied by
\cite[Corollary~3.9]{Ueda}. Proposition~\ref{D:norm-separation} proves
that the lower norm kernels are distinct for every distinguished
comparison used in this paper. Lemma~\ref{D:crossed-norm} then proves
the crossed norm-character pullback by norm transitivity and the
prime order of the two quotients. In the biquadratic case,
Lemma~\ref{D:quadratic-norms} gives every pair and the product relation
between the three lower characters. These are the precise consequences
usually obtained from \cite[Chapters~XIII--XIV]{Serre}; here their proofs
are included, without the general local reciprocity theorem.

\paragraph{Ramification and differents.}
The single-edge trace, different, norm, and conductor formulas are
quoted from FML. Lemma~\ref{D:inertia} proves the cyclic unramified
quotient, the tame inertia and Frobenius action, and the subgroup rule
for lower numbering. Lemma~\ref{D:break-ledger} derives all required
prime-square break relations directly from the different formula.
Thus no general quotient ramification theorem is an additional input.
These local facts have their classical setting in
\cite[Chapters~III--V]{Serre}, retained for attribution and context.
Transitivity of the different, the order-discriminant bound, and the
leading wild ramification term (including negative valuations) are
proved in Lemmas~\ref{C:leading-ramification}--\ref{C:order-discriminant}.
The usual monogenic different formula, elementary finite Galois and
finite-field theory, and Hensel lifting are the standard algebraic
foundations; their existing Mathlib or FML formulations and the short
specialization steps are listed in the separate dependency record.
Elementary cyclic coordinates are recalled at the start of
Appendix~\ref{app:diamond-foundations}; the equal-characteristic
coefficient field is constructed in Theorem~\ref{B:trace-residue}.

\paragraph{Local residues in equal characteristic.}
For a finite separable extension of Laurent-series fields, we use
\begin{equation}\label{U:residue-input}
 \Tr_{k_E/k_F}\Res_E(\eta)=\Res_F(\Tr_{E/F}\eta).
\end{equation}
Here the trace of $x\,d\varpi_F$ is $\Tr_{E/F}(x)\,d\varpi_F$.
This is proved in Theorem~\ref{B:trace-residue}, including wild
ramification. Lemma~\ref{B:dlog-norm} proves
$d\log N(u)=\Tr(d\log u)$ separately. Neither identity is quoted from
Tate, and neither proof uses a global curve or a global residue theorem.
The Artin--Schreier symbol needed in Section~\ref{sec:dyadic-equal} is
proved there from these internal results and the explicit Cartier
operator.

\paragraph{Finite sums.}
Ordinary finite-field Hasse--Davenport is taken from \cite{Ueda}.
The leading Gauss congruence used in Section~\ref{sec:tame}, including
its leading unit and its sign at every prime, is proved in
Theorem~\ref{A:leading-gauss}. Its proof establishes the needed
Jacobi-sum relation and residue calculation directly; no further
Gauss-sum theorem is quoted. Elementary finite character orthogonality
is used throughout. All special finite identities needed for the
two-field comparison are proved below.


\section{Stationary parameters and common comparison principles}\label{sec:stationary}
\subsection{Stationary numerator, covector, and rescaled coefficient}
A stationary numerator $b\in\OO_E^\times$ relative to $\Gamma$ satisfies
$\theta(1+z)=\psi_E(bz/\Gamma)$ on the specified ideal. Its additive
covector is $b/\Gamma$. If a calculation instead uses
$\Psi_E=\psi_E(\alpha\,\cdot)$ and writes
$\theta(1-z)=\Psi_E(Cz)$, the covector in the plus convention is
$-\alpha C$. In particular its elementary Lamprecht factor is
\[
 \theta((-\alpha C)^{-1})\Psi_E(-C).
\]
One must not interchange these three descriptions without their scalar
and sign. Letters for local breaks, depths, and rescaled coefficients are
reintroduced within each case calculation. This local use of notation does not change
the definitions of conductors or normalized valuations.

\subsection{The full stationary formula and exact stable twist}
Let $m=2d+\varepsilon>1$, $\varepsilon\in\{0,1\}$. The stationary
numerator $b$ is determined modulo $\pp_E^d$ by
\[
 \theta(1+z)=\psi_E(bz/\Gamma)
                         \quad(z\in\pp_E^{d+\varepsilon}).
\]
It is a unit. If $\varepsilon=0$, Lamprecht's formula is
\[
 \Delta_E(\theta,\psi_E)=\theta(\Gamma)\theta(b)^{-1}\psi_E(b/\Gamma).
\]
If $\varepsilon=1$, choose $\delta\OO_E=\pp_E^d$ and put
\[
 H_\theta(\bar z)=\psi_E(b\delta z/\Gamma)\theta(1+\delta z)^{-1}.
\]
This is a well-defined function on $k_E$ whose sum has magnitude
$|k_E|^{1/2}$. The even formula is then multiplied by
$|k_E|^{-1/2}\sum_{k_E}H_\theta$. Its polar pairing is
\[
 H_\theta(x+y)/(H_\theta(x)H_\theta(y))
       =\psi_E(b\delta^2\widetilde{xy}/\Gamma).
\]
Changing $b$ changes both its elementary value and this full residual
function; their compensating change preserves the local constant.
These are \cite[Theorem~4.5 and Lemma~4.6]{Ueda}.
If $m_E(\nu)\le d$, the exact stable twist is
\begin{equation}\label{U:stable}
 \Delta_E(\theta\nu,\psi_E)=\nu(\Gamma/b)\Delta_E(\theta,\psi_E)
\end{equation}
(\cite[Lemma~4.7]{Ueda}). Neither a homogeneous replacement of
$H_\theta$ nor an equality of squared sums is substituted for these formulas.

\begin{lemma}[Normalized stationary factors and terminal invariance]
\label{U:stationary-factor}
Let $E$ be a local field, $\alpha\in E^\times$, and
$\Psi_E(x)=\psi_E(\alpha x)$ with largest trivial ideal $\pp_E^J$.
Write the conductor of $\theta$ as $m=2d+\varepsilon>1$, with
$\varepsilon\in\{0,1\}$. Suppose
$\theta(1-z)=\Psi_E(Zz)$ for $z\in\pp_E^{d+\varepsilon}$,
with $v_E Z=J-m$. Then
\begin{equation}\label{U:normalized-factor}
 \Delta_E(\theta,\psi_E)=
 \theta(( -\alpha Z)^{-1})\Psi_E(-Z)g_E.
\end{equation}
For even $m$, $g_E=1$. For odd $m$, choose $v_E\delta=d$ and set
\begin{equation}\label{U:normalized-residual}
 H_E(\bar x)=\Psi_E(-Z\delta x)\theta(1+\delta x)^{-1},\qquad
 g_E=|k_E|^{-1/2}\sum_{x\in k_E}H_E(x).
\end{equation}
More generally, $\Psi_E(-Zv)\theta(1+v)^{-1}$ is constant as $v$ varies
modulo $\pp_E^{d+1}$ in $\pp_E^d$ when $m$ is odd; it is identically
one on $\pp_E^d$ when $m$ is even.
\end{lemma}
\begin{proof}
Apply the full Lamprecht formula to the covector $-\alpha Z$.
For terminal invariance, if $v'-v\in\pp_E^{d+1}$, write
$1+v'=(1+v)(1+w)$ with $w\in\pp_E^{d+1}$. The quotient of the two
expressions is $\Psi_E(-Zvw)=1$, since
$v_E(Zvw)\ge J$; the linear part cancels by stationarity.
In even conductor stationarity holds already on $\pp_E^d$.
\end{proof}

\subsection{Common comparison principles}\label{U:common-principles}
The following elementary lemmas isolate the shared character-theoretic
and stationary comparison arguments. They do not assume the existence of a simultaneous stationary
origin. That existence, and the full precision of each stationary formula,
are obligations discharged within the individual branch arguments.

\begin{lemma}[Conjugate twists and the corrected common restriction]
\label{U:determinant}\label{U:conjugacy}
Let $\ell$ be prime, let $K/F$ be finite Galois with group $C_\ell^2$,
and let $L_1,L_2$ be distinct degree-$\ell$ intermediate fields
with distinct lower norm subgroups. Put
$N_i=N_{K/L_i}$, $n_i=N_{L_i/F}$ and
$\epsilon_i=\prod_{\omega\in S(L_i/F)}\omega$.
Let $\theta_iN_i=\Theta$ be primitive compatible characters.
For a generator $\sigma_i$ of $\Gal(L_i/F)$, the quotient
$\mu_i=\theta_i^{\sigma_i}/\theta_i$ generates $S(K/L_i)$.
Every character in that group is fixed by $\Gal(L_i/F)$, and the conjugates of $\theta_i$
are precisely its twists by $S(K/L_i)$. Moreover
\begin{equation}\label{U:det-character}
 D=D_\theta=\theta_i|_{F^\times}\epsilon_i
\end{equation}
is independent of $i$. For odd $\ell$, $\epsilon_i=1$; for $\ell=2$,
$\epsilon_i=\omega_i$ is the nontrivial lower norm character.
The distinct-kernel hypothesis holds for the distinguished pairs in
Proposition~\ref{D:norm-separation}, and for every biquadratic pair by
Lemma~\ref{D:quadratic-norms}.
\end{lemma}
\begin{proof}
The common pullback is invariant under both upper subgroups and hence
under $\Gal(K/F)$. Thus $\mu_i$ kills $N_iK^\times$. If it were trivial,
Hilbert~90 and character extension would make $\theta_i$ a lower norm
pullback, contradicting primitivity. The upper norm quotient has prime
order, so $\mu_i$ generates its character group. Lemma~\ref{D:crossed-norm}
identifies that group with $S(L_j/F)\circ n_i$; it is therefore invariant
under $\sigma_i$. Iterated conjugation gives all $\ell$ twists.

For $x\in L_j^\times$, compatibility gives
$\theta_i(n_jx)=\theta_j(x)^\ell$, whereas multiplication over the
conjugates gives
\[
 \theta_j(n_jx)=\theta_j(x)^\ell
       \mu_j(x)^{-\ell(\ell-1)/2}.
\]
For odd $\ell$ the last factor is one, and the characters in each
$S(L_i/F)$ pair with their inverses, so $\epsilon_i=1$.
For $\ell=2$, that factor is
$\mu_j(x)=\omega_i(n_jx)$ and $\omega_j(n_jx)=1$.
In either case the two corrected restrictions agree on $n_jL_j^\times$,
and by symmetry on $n_iL_i^\times$. These distinct index-$\ell$
subgroups generate $F^\times$, proving the result.
\end{proof}

\begin{corollary}[The odd-degree common power]\label{U:odd-power}
In Lemma~\ref{U:conjugacy}, suppose $\ell$ is odd and use trace-compatible
additive characters. Then
\[
 \Delta_{L_i}(\theta_i,\psi_{L_i})^\ell=\Delta_K(\Theta,\psi_K),
 \qquad \prod_{\omega\in S(L_i/F)}\Delta_F(\omega,\psi_F)=1.
\]
In particular, the ratio of the two complete induction expressions
has $\ell$th power one.
\end{corollary}
\begin{proof}
On any odd-prime cyclic edge, inverse pairing from FML gives
$\Delta(\omega)\Delta(\omega^{-1})=\omega(-1)=1$; the trivial factor
is one. Apply the First Main Lemma on $K/L_i$. By
Lemma~\ref{U:conjugacy}, its twists are conjugate, and their local
constants agree by an automorphism change of variables in the defining
integral. This proves the power relation as well as both product claims.
\end{proof}

\begin{proposition}[The biquadratic common square]\label{U:quadratic-square}
For a primitive compatible biquadratic family with trace-compatible
additive characters, let $\mathcal A_i=\Delta_{L_i}(\theta_i)\Delta_F(\omega_i)$. Then
\[
 \mathcal A_i^2=\Delta_K(\Theta)\prod_{j=1}^3\Delta_F(\omega_j).
\]
Thus the ratio of two such expressions belongs to $\{1,-1\}$.
\end{proposition}
\begin{proof}
By Lemma~\ref{U:conjugacy} and FML on $K/L_i$,
$\Delta_{L_i}(\theta_i)^2=\Delta_K(\Theta)\Delta_{L_i}(\omega_jn_i)$
for $j\ne i$. FML on $L_i/F$, applied to $\omega_j$, gives
$\Delta_{L_i}(\omega_jn_i)\Delta_F(\omega_i)
 =\Delta_F(\omega_j)\Delta_F(\omega_k)$, where $\{i,j,k\}=\{1,2,3\}$.
Use $\omega_i\omega_j=\omega_k$ from Lemma~\ref{D:quadratic-norms}
and multiply the two equalities.
\end{proof}

\begin{lemma}[A characteristic-free quadratic identity]\label{U:e2}
For a biquadratic extension $K/F$, $C\in K$, and any quadratic
intermediate field $L_i$,
\[
 E_2(C)=\Tr_{L_i/F}(N_iC)+n_i(S_iC),\qquad S_i=\Tr_{K/L_i},
\]
where $E_2(C)$ is the second elementary symmetric function of the four
conjugates of $C$.
\end{lemma}
\begin{proof}
Write the four conjugates, paired by the involution fixing $L_i$, as
$c_1,c_2$ and $c_3,c_4$. The two terms on the right are
$c_1c_2+c_3c_4$ and $(c_1+c_2)(c_3+c_4)$.
Their sum is exactly the sum of all six pair-products. This identity is
integral, so it holds in mixed and equal characteristic without a sign
convention or division by two.
\end{proof}

\begin{lemma}[The complete factor at a genuine common origin]\label{U:common-origin}
Use the biquadratic setting of Lemma~\ref{U:determinant}. Fix $\alpha\in F^\times$
and put $\Psi_E=\psi_E(\alpha\,\cdot)$.
Suppose all the inducing and lower norm characters under consideration
have conductor greater than one. Let $C\in K^\times$ satisfy $S_iC\ne0$,
and suppose that
\[
 Z_i=N_iC,\qquad W_i=n_i(S_iC)
\]
are actual stationary coefficients for $\theta_i$ and $\omega_i$,
respectively, in the convention
$\theta_i(1-z)=\Psi_{L_i}(Z_i z)$ and
$\omega_i(1-z)=\Psi_F(W_i z)$ on their full ordinary stationary ideals.
Let $g_i$ and $h_i$ denote the respective complete normalized residual
sums in Lamprecht's formula, with value one for even conductor. Then
\[
 \mathcal A_i(\theta_i)
   =D(-\alpha^{-1})\Theta(C)^{-1}\Psi_F(-E_2(C))\,g_i h_i.
\]
Thus equality of the complete residual products proves the desired
comparison for these characters.
\end{lemma}
\begin{proof}
The actual covectors are $-\alpha Z_i$ and $-\alpha W_i$. The two complete
Lamprecht factors are therefore
\[
 \theta_i(-\alpha^{-1})\theta_i(Z_i)^{-1}\Psi_{L_i}(-Z_i)g_i,
 \quad
 \omega_i(-\alpha^{-1})\omega_i(W_i)^{-1}\Psi_F(-W_i)h_i.
\]
Compatibility gives $\theta_i(Z_i)=\Theta(C)$. The value
$\omega_i(W_i)$ is one because $W_i$ is a literal lower norm, not merely
an element with the valuation of a norm. Multiply, apply
Lemma~\ref{U:determinant}, and use Lemma~\ref{U:e2}.
No residual sign has been discarded. In characteristic two the field
minus signs disappear, but the complex residual values need not be signs.
\end{proof}

\begin{lemma}[An identity before taking residual sums]\label{U:common-function}
Use a common origin $C$ as in Lemma~\ref{U:common-origin}. For another
$C'\in K^\times$ with all $S_iC'\ne0$, put
\[
 1+z_i=N_iC'/N_iC,\qquad 1+w_i=n_i(S_iC')/n_i(S_iC).
\]
Then the products
\[
 \Psi_{L_i}(-Z_i z_i)\theta_i(1+z_i)^{-1}
 \Psi_F(-W_i w_i)\omega_i(1+w_i)^{-1}
\]
are independent of $i$. Their common value is
\[
 \Theta(C'/C)^{-1}\Psi_F\bigl(-E_2(C')+E_2(C)\bigr).
\]
When $z_i,w_i$ belong to the respective terminal stationary ideals,
these are identities of the actual residual functions, including their
affine terms.
\end{lemma}
\begin{proof}
The upper character value is $\Theta(C'/C)$ by compatibility.
The argument of the lower character is a quotient of two nonzero norms,
so that character value is one. The additive exponent is
\[
 -\Tr_{L_i/F}(N_iC'-N_iC)-n_i(S_iC')+n_i(S_iC),
\]
which is $-E_2(C')+E_2(C)$ by Lemma~\ref{U:e2}.
The terminal-ideal interpretation follows from the full Lamprecht formula;
it is used only after the relevant ideal membership has been proved.
\end{proof}

\begin{lemma}[Cancellation of the missing norm coset]\label{U:missing-coset}
For each index $i$, let $V_i$ be a finite abelian group of the same order,
and let $H_i:V_i\to\mathbf C^\times$ have $H_i(0)=1$ and bicharacter
$B_i(x,y)=H_i(x+y)/(H_i(x)H_i(y))$.
Let $Z$ be a finite abelian group and let $p_i:Z\to V_i$ be group homomorphisms with kernel of order two and
image $I_i$ of index two. Suppose $H_i\circ p_i=Q$ for the same function
on $Z$. Suppose further that there exists $w_i\in I_i$ such that
$H_i(w_i)=1$ and $B_i(w_i,\cdot)$ is the nontrivial character of $V_i/I_i$.
Then
\[
 \sum_{x\in V_i}H_i(x)=\frac12\sum_{z\in Z}Q(z).
\]
In particular the full normalized sums are equal.
\end{lemma}
\begin{proof}
Translation by $w_i$ preserves the complementary coset $V_i\setminus I_i$.
On that coset it multiplies the summand by minus one. Thus its sum equals
its own negative and is zero in $\mathbf C$. On $I_i$, every value has
exactly two preimages under $p_i$, so summing $H_i p_i=Q$ gives the asserted
identity. The factor $1/2$ is a counting factor in $\mathbf C$, not division
in the local field. The hypothesis $H_i(w_i)=1$ concerns the actual
quadratic refinement and cannot be deduced from its polar form alone.
\end{proof}

\begin{lemma}[Norm choice within a proven normalization freedom]\label{U:normalization}
Let $L/F$ be cyclic of prime degree and let $\omega$ be a nontrivial
character with kernel $N_{L/F}L^\times$. Let $H\subset F^\times$ be a subgroup. Assume an auxiliary
normalization $\alpha$ may be changed to $\alpha u$, for every $u\in H$,
without changing any of the already fixed character formulas in an
argument. If $\omega(H)=\omega(F^\times)$, then for each $A_*\in F^\times$
one may make $A=A_*/u$ a norm while preserving the raw coefficient
$\alpha A_*$.
\end{lemma}
\begin{proof}
Choose $u\in H$ with $\omega(u)=\omega(A_*)$. Then $\omega(A_*/u)=1$,
so $A_*/u$ is a norm by the specified kernel. Moreover
$(\alpha u)(A_*/u)=\alpha A_*$ exactly. Preservation of all other formulas
is the explicit hypothesis; the branch proofs verify it on their whole
ideals before this lemma is applied. This lemma never asserts that a
previously fixed coefficient was already a norm.
\end{proof}


\section{Finite-field identities and the tame comparison}\label{O:sec:tame}\label{sec:tame}
\subsection{Finite sums and their local leading term}
For a finite field $k$, a nontrivial additive character $\psi_k$, and a
multiplicative character $\chi$, put
\[
 \tau_k(\chi)=-\sum_{x\in k^\times}\chi(x)^{-1}\psi_k(x).
\]
Thus $\tau_k(1)=1$. For a finite extension $\kappa/k$ all additive
characters are composed with trace. The two ordinary
Hasse--Davenport identities in Ueda's finite-field preliminaries are
\begin{align}\label{O:F:HD}
 \tau_\kappa(\chi\circ\Nm)&=\tau_k(\chi)^{[\kappa:k]},\notag\\
 \chi(s^s)\tau_k(\chi^s)\prod_{j=1}^{s-1}\tau_k(\eta^j)
 &=\prod_{j=0}^{s-1}\tau_k(\chi\eta^j),
\end{align}
where $\eta$ has order $s$. The second formula is used only when
$s$ is prime to the characteristic, as required for such a character.

We use the following classical leading Gauss congruence, proved in
Appendix~\ref{app:leading-gauss} by an elementary Jacobi-sum argument.
We state its local version here to fix the precise normalization. Fix an
embedding of the cyclotomic values into
$\mathbf Q_p(\mu_{p^f-1},\zeta_p)$, put $\varpi=\zeta_p-1$, and
let $\omega$ be the Teichmuller character of $\FF_{p^f}^\times$.
For the additive character $x\mapsto\zeta_p^{\Tr(x)}$ and
$0<a<p^f-1$, write $a=\sum a_i p^i$. Then
\begin{equation}\label{O:F:Stick}
 \tau_{\FF_{p^f}}(\omega^a)
 \equiv_\times \frac{\varpi^{\sum a_i}}{\prod a_i!}.
\end{equation}
Here $X\equiv_\times Y$ means $X/Y\in1+\mathfrak m$ in this local
cyclotomic field. This is a leading congruence, not merely a valuation
formula. The sign in \eqref{O:F:Stick} is positive because our $\tau$ is
minus the Gauss sum. Equivalently, in the usual plus-sum convention
$G(\omega^{-a})=\sum_{x\ne0}\omega^{-a}(x)\zeta_p^{\Tr(x)}$, the leading
term is $-\varpi^{\sum a_i}/\prod a_i!$; our definition is
$\tau(\omega^a)=-G(\omega^{-a})$. This specifies the sign without a
normalization imported from any two-field proof.
Theorem~\ref{A:leading-gauss} proves the congruence with this exact
sign, not only its valuation. It applies also at $p=2$, with
$\varpi=\zeta_2-1=-2$. No complex absolute-value argument at other places
is involved.

For later use, Wilson's theorem gives, for every nonnegative integer
$n$ with base-$p$ digits $n_i$,
\[
 n!\equiv_\times(-p)^{(n-\sum n_i)/(p-1)}\prod n_i!.
\]
Indeed separate the multiples of $p$, reduce each full nonzero residue
block to $(p-1)!\equiv-1$, and repeat for $\lfloor n/p\rfloor!$.
Since $\varpi^{p-1}\equiv_\times-p$, we may group \eqref{O:F:Stick}
into base-$q$ digits for any power $q$ of $p$:
\begin{equation}\label{O:F:groupedStick}
 a_*=\sum_{i=0}^{r-1}\gamma_iq^i
 \quad\Longrightarrow\quad
 \tau_{\FF_{q^r}}(\omega^{a_*})
 \equiv_\times\frac{\varpi^{\sum_i\gamma_i}}{\prod_i\gamma_i!}.
\end{equation}
Thus no assertion involving complex absolute values or a global
cyclotomic ideal-class theorem is needed below.

\subsection{The primitive finite-field identity}
Let $\ell$ be prime, $\ell\mid q-1$, and let $\kappa/k$ have degree
$\ell$, where $|k|=q$. Suppose
\begin{equation}\label{O:F:finitecompatible}
 \chi_\kappa^\ell=\chi_0\circ\Nm_{\kappa/k},
 \qquad \chi_0|_{\mu_\ell}\ne1.
\end{equation}
Let $\mu$ be a character of $k^\times$ of order $\ell$.

\begin{theorem}\label{O:F:finite}
With the trace-compatible additive characters,
\begin{equation}\label{O:F:finiteidentity}
 \tau_\kappa(\chi_\kappa)
 =\chi_0(\ell)\tau_k(\chi_0)
       \prod_{j=1}^{\ell-1}\tau_k(\mu^j).
\end{equation}
\end{theorem}
\begin{proof}
First take the canonical additive characters. Put
$S=(q^\ell-1)/(q-1)$ and $h=(q-1)/\ell$.
Write $\chi_0=\omega_k^a$, $1\le a\le q-2$, $\ell\nmid a$.
The exponent of $\chi_\kappa$ is $aS/\ell$ modulo
$(q^\ell-1)/\ell$. Multiplication by $q$ changes it by
$a(q^\ell-1)/\ell$; hence all $\ell$ possibilities are Frobenius
conjugates. Their Gauss sums are equal, so we may take the exponent
exactly $a_*=aS/\ell$.

Set
\[
 L=\frac{\tau_\kappa(\chi_\kappa)}
 {\chi_0(\ell)\tau_k(\chi_0)\prod_{j=1}^{\ell-1}\tau_k(\mu^j)}.
\]
The order-$\ell$ characters over $\kappa$ are the norm lifts of
$\mu^j$. Their twists of $\chi_\kappa$ are the Frobenius conjugates
just considered. The product and lifting formulas \eqref{O:F:HD}, and
$\chi_\kappa(\ell^\ell)=\chi_0(\ell)^\ell$, give
\begin{equation}\label{O:F:Lpower}
 L^\ell=1.
\end{equation}

We now determine this root of unity locally. Write $a=\ell m+k_0$,
$1\le k_0<\ell$. The unordered base-$q$ digits of $aS/\ell$ are
\begin{equation}\label{O:F:digits}
 \gamma_j=m+jh+\mathbf1_{j\ge\ell-k_0},\qquad0\le j\le\ell-1.
\end{equation}
Here ``unordered'' is sufficient for their sum and factorial product.
For completeness, multiplication of $k_0S/\ell$ by $\ell$, using
$q=\ell h+1$, gives successive residues
$i_j=((\ell-1-j)k_0\bmod\ell)$ and carries
$c_j=\lfloor(\ell-1-j)k_0/\ell\rfloor$.
The zeroth digit has the additional $k_0-c_0=1$; at a later digit the
additional $c_{j-1}-c_j$ is one exactly when $i_j\ge\ell-k_0$.
The $i_j$ permute $0,\ldots,\ell-1$, proving \eqref{O:F:digits}.
All digits lie between zero and $q-1$, since $m\le h-1$.
In particular
\begin{equation}\label{O:F:digitsum}
 \sum\gamma_j=a+\sum_{j=1}^{\ell-1}jh.
\end{equation}

We claim the multiplicative factorial congruence
\begin{equation}\label{O:F:factorials}
 \ell^a\prod_{j=0}^{\ell-1}\gamma_j!
 \equiv_\times a!\prod_{j=1}^{\ell-1}(jh)!.
\end{equation}
In fact
\[
 \frac{\prod\gamma_j!}{\prod(jh)!}
 =\prod_{j=0}^{\ell-1}\prod_{n=1}^{m}(n+jh)
      \prod_{j=\ell-k_0}^{\ell-1}(m+1+jh).
\]
For the first product,
$\ell(n+jh)=\ell n-j+jq\equiv_\times\ell n-j$:
the positive integer $\ell n-j$ is less than $q=p^f$, so its
$p$-valuation is less than $f$. The numbers $\ell n-j$ exhaust
$1,\ldots,\ell m$. For the second product the same argument applies
to $\ell(m+1)-j$, which exhausts $\ell m+1,\ldots,a$, again all
less than $q$. Multiplication proves \eqref{O:F:factorials}, even when
some of its factorials are divisible by $p$.

Apply \eqref{O:F:groupedStick} to the numerator and each factor of the
denominator of $L$. The Teichmuller value $\chi_0(\ell)$ reduces to
$\ell^a$. Equations \eqref{O:F:digitsum} and \eqref{O:F:factorials} give
$L\equiv_\times1$. Reduction is injective on the group of
$\ell$th roots of unity because $\ell\ne p$.
Together with \eqref{O:F:Lpower}, this proves $L=1$.

For an arbitrary additive character of $k$, change the canonical one
by multiplication by $c\in k^\times$. The two sides acquire factors
$\chi_\kappa(c)$ and $\chi_0(c)\prod_{j=1}^{\ell-1}\mu^j(c)$.
These agree. Indeed the restriction exponent of $\chi_\kappa$ is
$aS/\ell$ modulo $q-1$: it is $a$ for odd $\ell$, and $a+(q-1)/2$
for $\ell=2$ because $a$ is odd. This is exactly the exponent of
$\chi_0\prod\mu^j$. Thus the identity has the asserted generality.
\end{proof}

\subsection{The tame local diamond and its restriction identity}
Let $F$ have residue characteristic $p$, let $K/F$ have Galois
group $C_\ell^2$ with $\ell\ne p$, and let $U/F$ be its unique
unramified degree-$\ell$ subextension. Let $E/F$ be any other
lower field. Its tame cyclic inertia gives $\ell\mid |k_F|-1$;
$K/E$ is unramified and $K/U$ is totally ramified.
Suppose $\theta_U\Nm_{K/U}=\theta_E\Nm_{K/E}$ is invariant and
does not descend from $F$. With $S(L/F)$ as defined in
Section~\ref{sec:introduction}, put
\[
 \epsilon_L=\prod_{\nu\in S(L/F)}\nu,\quad
 \Lambda_L=\prod_{\nu\in S(L/F)}\Delta_F(\nu,\psi_F).
\]
For odd $\ell$, $\epsilon_L=1$; for $\ell=2$ it is the nontrivial
quadratic norm character. All additive characters are trace compatible.


\subsection{Conductors greater than one}
Put $m=m_U(\theta_U)$. Its conjugate quotient has conductor one, so
$m\ge1$. If $m>1$, unramified Hilbert 90 on $U_U^1$, followed by
norm surjectivity, gives a character $\lambda$ of $F^\times$ such
that
\[
 \theta_U=\chi_U(\lambda n_U),\qquad
 \theta_E=\chi_E(\lambda n_E),\qquad
 m(\chi_U)=m(\chi_E)=1,\quad m(\lambda)=m.
\]
Here $n_L=\Nm_{L/F}$. To justify the conductor-one assertion, remove
the descended restriction on $U_U^1$. The nontrivial conjugate
quotient forces the remaining $U$ character to have conductor exactly
one; all its upper-norm twists are conjugates. Ueda's minimal-orbit
conductor formula gives conductor one for its pullback to $K$;
unramified pullback from $E$ preserves conductor.

Let $g\in F^\times$ be an ordinary stationary covector of
$\lambda$, so $g=\beta/\Gamma$ and
$v_F(\Gamma)=m+n_F(\psi_F)$. The following local calculation shows that
this same $g$ is a stationary covector of both pullbacks. Their conductors
are $m$ over $U$ and $M=1+\ell(m-1)$ over $E$, and
$n_E(\psi_E)=\ell n_F(\psi_F)+\ell-1$. Hence the same
$\Gamma\in F$ is admissible in both.

\begin{lemma}[The full tame stationary covector]\label{U:tame-covector}
For $m>1$, a covector $g$ satisfying
$\lambda(1+z)=\psi_F(gz)$ on $\pp_F^{\lceil m/2\rceil}$
satisfies the corresponding ordinary stationary formula for
$\lambda\circ n_U$ and $\lambda\circ n_E$ on their full stationary
ideals. This holds in every residue characteristic different from $\ell$.
\end{lemma}
\begin{proof}
On the ramified side put $s=\lceil M/2\rceil$ and take
$x\in\pp_E^s$. For each elementary symmetric function of its $\ell$
conjugates, restriction of valuations gives
$v_F(E_j(x))\ge\lceil js/\ell\rceil$. For $j\ge2$ this is at least $m$,
since $2s\ge M>\ell(m-1)$. Thus
\[
 n_E(1+x)\equiv1+\Tr_{E/F}x\pmod{\pp_F^m}.
\]
The trace has valuation at least $\lceil s/\ell\rceil$, which is at
least $\lceil m/2\rceil$: for $m=2b$ use
$s>\ell(b-1)$, and for $m=2b+1$ use $s>\ell b$.
Both units in this congruence are principal, so their quotient belongs to
$U_F^m$. Evaluating $\lambda$ proves
$(\lambda n_E)(1+x)=\psi_E(gx)$ on the entire required ideal.
On the unramified side put $s=\lceil m/2\rceil$. All symmetric terms
of degree at least two have valuation at least $2s\ge m$, and the trace
has depth at least $s$. The same norm expansion proves the assertion.
No division by the residue characteristic, nor a half-depth formula on a
shallower group, is used.
\end{proof}

The conductor-one characters are in the stable-twist range on both
sides, including $\ell=2$ and $m=2$. Thus
\begin{equation}\label{O:F:stablesides}
 \Delta_L(\theta_L)=\chi_L(g^{-1})\Delta_L(\lambda n_L).
\end{equation}
The First Main Lemma and its exact stable-twist formula on $F$ give
\begin{equation}\label{O:F:stableFML}
 \Delta_L(\lambda n_L)\Lambda_L
 =\prod_{\nu\in S(L/F)}\Delta_F(\nu\lambda)
 =\epsilon_L(g^{-1})\Delta_F(\lambda)^\ell.
\end{equation}
All lower norm characters have conductor zero or one, hence at most
$\lfloor m/2\rfloor$. Combining \eqref{O:F:stablesides}--\eqref{O:F:stableFML}
with Lemma~\ref{U:determinant} for the compatible pair $(\chi_U,\chi_E)$
proves
\begin{equation}\label{O:F:localidentity}
 \Delta_U(\theta_U)\Lambda_U=\Delta_E(\theta_E)\Lambda_E
\end{equation}
for every $m>1$. No comparison of incomplete quadratic factors was
made; the stable-twist formula retains the complete Lamprecht factor.

\subsection{Conductor one and the sign}
Now $m_U=m_E=1$. Let $k=k_F=k_E$, $\kappa=k_U$, and denote the two
residue characters by $\chi_\kappa,\chi_0$.
Norm compatibility gives $\chi_\kappa^\ell=\chi_0\Nm_{\kappa/k}$.
The Frobenius conjugate quotient is a nontrivial character of order
$\ell$. Writing the exponents as in the proof of Theorem~\ref{O:F:finite}
shows that this is exactly $\ell\nmid a$, or
$\chi_0|_{\mu_\ell}\ne1$. In particular both residue Gauss sums are
nontrivial-character sums.

Write $n=n_F(\psi_F)$ and choose
$\Gamma=\pi_F^{n+1}\in F$. It is admissible in both fields:
$n_E=\ell n+\ell-1$, while $n_U=n$.
The conductor-one local formulas in Ueda's normalization give
\begin{align*}
 \Delta_U(\theta_U)&=-\theta_U(\Gamma)\ph\tau_\kappa(\chi_\kappa),\\
 \Delta_E(\theta_E)&=-\theta_E(\Gamma)\chi_0(\ell)
                                   \ph\tau_k(\chi_0).
\end{align*}
The factor $\chi_0(\ell)$ is present because the residual additive
character over $E$ is $x\mapsto\psi_k(\ell x)$, not $\psi_k(x)$.
For the lower norm-character products,
\[
 \Lambda_U=(-1)^{(\ell-1)n},\qquad
 \Lambda_E=(-1)^{\ell-1}\epsilon_E(\Gamma)
                 \prod_{j=1}^{\ell-1}\ph\tau_k(\mu^j).
\]
Indeed $\epsilon_U(\pi_F)=(-1)^{\ell-1}$ by the product of the
$\ell$ unramified characters, and every nontrivial member of
$S(E/F)$ has conductor one. Lemma~\ref{U:determinant} gives
$\theta_U(\Gamma)/\theta_E(\Gamma)
 =\epsilon_E(\Gamma)/\epsilon_U(\Gamma)$.
Since $\epsilon_U(\Gamma)=(-1)^{(\ell-1)(n+1)}$, all the displayed
signs cancel in the quotient of \eqref{O:F:localidentity}. What remains is
precisely the phase of the quotient in Theorem~\ref{O:F:finite}, namely
one. This proves \eqref{O:F:localidentity} at conductor one as well.

\begin{theorem}[the complete tame branch]\label{O:F:tame}
For a tame $C_\ell^2$ diamond over a local field of arbitrary residue
characteristic $p\ne\ell$, every compatible non-descending inducing pair has the
same normalized induction expression. This holds in both
characteristics, in every conductor, and includes $\ell=2$ with the
exact sign.
\end{theorem}
\begin{proof}
The unramified lower field is unique. The arguments above compare it
with every ramified lower field, so transitivity compares any two.
The higher-conductor case is \eqref{O:F:stablesides}--\eqref{O:F:stableFML};
the conductor-one case is the final calculation.
\end{proof}


\section{Wild odd degree: the unramified--ramified case}\label{O:sec:inertia}\label{sec:odd-ur}
\subsection{Statement and notation}
Let $F$ be a nonarchimedean local field of residue characteristic $p>2$.
Let $K/F$ be Galois with group $C_p\times C_p$, and suppose its inertia
subgroup has order $p$. Write $U/F$ for its unramified degree-$p$ field and
$E/F$ for any ramified degree-$p$ field. Then $K=EU$, $K/E$ is unramified,
and $E/F$ and $K/U$ are totally ramified cyclic of degree $p$, with the same
break $t\ge1$. Set
\[
 T=t+1,\qquad D_0=(p-1)T,\qquad
 q_0=\ceil{T/p},\qquad e_0=\floor{D_0/p}=T-q_0.
\]
No hypothesis $p\nmid t$ is imposed. All valuations are normalized on the
field displayed. In particular $v_K|_E=v_E$, $v_K|_U=pv_U$, and
$v_E|_F=pv_F$.
Use distinct norm symbols
\[
 n_E=\Nm_{E/F},\quad n_U=\Nm_{U/F},\quad
 N_E=\Nm_{K/E},\quad N_U=\Nm_{K/U}.
\]
Fix a nontrivial additive character $\echar_F$ and set
$\echar_L=\echar_F\circ\Tr_{L/F}$ for all four fields. Let
$\Delta_L(\theta,\echar_L)$ mean Ueda's normalized one-dimensional local
constant, with the same admissible-denominator convention.

\begin{theorem}\label{O:I:main}
Suppose $\theta_U$ and $\theta_E$ have the same norm pullback
\[
 \theta_K=\theta_U\circ N_U=\theta_E\circ N_E,
\]
and this character does not descend through $\Nm_{K/F}$. Then
\[
 \Delta_U(\theta_U,\echar_U)=\Delta_E(\theta_E,\echar_E).
\]
The cyclic lower-field products in the Second Main Lemma are both $1$, so
this is also equality of its complete normalized induction expressions.
The assertion holds in mixed and equal characteristic.
\end{theorem}

We use the one-dimensional results recalled in
Sections~\ref{sec:local-notation}--\ref{sec:stationary}, together with
the subcritical norm representatives of \cite[Lemma~6.6]{Ueda} and
the stable part of its high-conductor stationary table
\cite[Proposition~6.8]{Ueda}.
The following consequences specify the precise normalizations used here.
For a totally ramified degree-$p$ edge $L/M$ with break $t$,
\begin{align}
 \Tr_{L/M}(\pp_L^a)&=\pp_M^{\floor{(a+D_0)/p}},\label{O:I:traceideal}\\
 v_M(e_i(y))&\ge\floor{(i v_L(y)+D_0)/p}\quad(1\le i<p),\label{O:I:symbound}\\
 v_M(\Nm y)&=v_L(y).\label{O:I:normval}
\end{align}
Here $e_i(y)$ is the $i$th elementary symmetric function of the $p$
conjugates, even when $y$ lies in a proper subfield. The inequalities apply
to negative as well as positive valuations. For unramified edges, the
norm is onto every positive unit group and norm pullback preserves
conductors. If $c\in M^\times$ has valuation $a$ and $0\le r\le t$, there
is $x\in L^\times$ with
\begin{equation}\label{O:I:normrep}
 v_L(x)=a,\qquad c-\Nm(x)\in\pp_M^{a+r}.
\end{equation}
We use Hilbert 90 and extension of characters of finite abelian groups.
Every additional two-field argument is proved below.

\subsection{Truncated logarithms and the actual norm character}
Put
\[
 P(Z)=\sum_{j=1}^{p-1}\frac{Z^j}{j}.
\]
Its coefficients lie in $\mathbf Z_{(p)}$. Our sign convention is that
$P(z)$ is the truncation of $-\log(1-z)$.

\begin{lemma}[logarithm charts]\label{O:I:logchart}
The polynomial
$P(z+w-zw)-P(z)-P(w)$ has no monomial of total degree less than $p$.
Consequently, for integers $1\le r\le M$ with $pr\ge M$, the map
\[
 1-z\longmapsto P(z)
\]
is a group isomorphism $U_L^r/U_L^M\longrightarrow
\pp_L^r/\pp_L^M$. The map $P:\pp_L^r\longrightarrow\pp_L^r$ is bijective
and preserves valuations.
\end{lemma}
\begin{proof}
Compare degrees less than $p$ in the formal logarithm identity over
$\mathbf Q$. The truncated expression has coefficients in
$\mathbf Z_{(p)}$, so the assertion is valid in both characteristics.
The congruence is additive modulo $\pp_L^M$ under the stated inequality.
Finally $P(z)=z+z^2R(z)$ with $R$ integral and $P'(z)\equiv1\pmod{\pp_L}$.
Hensel's lemma solves $P(z)=w$ uniquely in $\pp_L^r$, or equivalently one
can invert successively on its successive quotients. This also proves the
valuation assertion and the group isomorphism.
\end{proof}

\begin{lemma}[full trace--norm logarithm identity]\label{O:I:normlog}
Let $L/M$ be a totally ramified degree-$p$ cyclic edge of break $t$.
For $h\ge0$ and $z\in\pp_L^{h+q_0}$,
\begin{equation}\label{O:I:normlogeq}
 P(1-\Nm(1-z))\equiv
 \Tr P(z)+\Nm(P(z))\pmod{\pp_M^{T+h}}.
\end{equation}
For an unramified degree-$p$ edge and $z\in\pp_L^{q_0}$, one instead has
\begin{equation}\label{O:I:unramlog}
 P(1-\Nm(1-z))\equiv\Tr P(z)\pmod{\pp_M^T}.
\end{equation}
\end{lemma}
\begin{proof}
Here is a universal proof, so no earlier assertion about a norm
representative is needed. Let $z_1,\ldots,z_p$ be indeterminates and $e_i$
their elementary symmetric functions. Give $e_i$ weight $i$ and put
\[
 u=1-\prod_{a=1}^p(1-z_a)=e_1-e_2+\cdots+e_p,
\quad
 \mathscr D=P(u)-\sum_aP(z_a)-\prod_aP(z_a).
\]
The polynomial $\mathscr D$ belongs to
$\mathbf Z_{(p)}[e_1,\ldots,e_p]$. Every monomial has weight at least $p$
and at least two elementary-symmetric factors. Indeed formal logarithms
kill all weights below $p$, and the weight-$p$ term before the final product
is $(\sum_a z_a^p-e_1^p)/p$. Newton's identity gives coefficient $1$ for
$e_p$ in this polynomial; the final product cancels it. There is no possible
one-factor monomial of larger weight. Moreover every coefficient of a pure
power $e_p^j$ is divisible by $p$: modulo $p$, specialize
$e_1=\cdots=e_{p-1}=0$, adjoin $w$ with $w^p=e_p$, and note that
$\sum_aP(z_a)=0$ and $\prod_aP(z_a)=P(w)^p=P(e_p)$.
The injectivity of $\F_p[e_p]\to\F_p[w]$ justifies this coefficient test.

Write $q=h+q_0$ and specialize to the conjugates of $z$. Then
$v_M(e_p)\ge q$ and, for $i<p$,
\[
 v_M(e_i)\ge\floor{(iq+D_0)/p}
 \ge \frac{iq+(p-1)t}{p},\qquad v_M(e_i)\ge e_0.
\]
A monomial containing both $e_p$ and a smaller-index factor has valuation
at least $q+e_0=T+h$. A monomial consisting of $r\ge2$ smaller-index
factors and having total weight at least $p$ has valuation at least
\[
 q+\frac{2(p-1)t}{p}\ge T+h,
\]
because $q_0+(p-2)t/p-1\ge0$.
In mixed characteristic $e=v_M(p)\ge e_0$, by
$p=\Tr_{L/M}(1)$ and \eqref{O:I:traceideal}. A pure norm-power error has
valuation at least $e+2q\ge T+h$. In equal characteristic it is zero.
This proves \eqref{O:I:normlogeq}. For an unramified edge, every monomial of
weight at least $p$ has valuation at least $pq_0\ge T$, and the norm of
$P(z)$ has that valuation also. This gives \eqref{O:I:unramlog}.
\end{proof}

Choose a nontrivial actual norm character $\tau_F$ for $E/F$. It has
conductor $T$. Lemma~\ref{O:I:logchart} and the additive trace pairing give
$\alpha\in F^\times$ such that
\begin{equation}\label{O:I:tauchart}
 \tau_F(1-z)=\echar_F(\alpha P(z))\quad(z\in\pp_F^{q_0}),
 \qquad v_F(\alpha)=-n_F(\echar_F)-T.
\end{equation}
Set
\[
 \Psi_L(z)=\echar_L(\alpha z).
\]
The largest ideal on which $\Psi_L$ is trivial is $\pp_L^T$ for
\emph{each} of $F,U,E,K$. For $E/F$, this is the different calculation
$-pT+D_0=-T$; the other edges follow by unramified base change. Thus all
ideal exponents below refer to normalized valuations, not rescaled ones.
Set $\tau_U=\tau_F\circ n_U$. It is the actual norm character for $K/U$,
and \eqref{O:I:unramlog} gives its $P$-formula on $U_U^{q_0}$ with scalar $1$
relative to $\Psi_U$.

\begin{corollary}[whole-subgroup conversion]\label{O:I:conversion}
For either ramified edge $L/M=E/F$ or $K/U$ and every
$w\in\pp_L^{q_0}$,
\begin{equation}\label{O:I:conversioneq}
 \Psi_M(\Tr_{L/M}w+\Nm_{L/M}w)=1,
 \qquad
 \Psi_L(w)=\Psi_M(-\Nm_{L/M}w).
\end{equation}
\end{corollary}
\begin{proof}
Write $w=P(z)$ with $z\in\pp_L^{q_0}$, using Lemma~\ref{O:I:logchart}.
The norm of $1-z$ belongs to $U_M^{q_0}$: the intermediate symmetric
terms have valuation at least $q_0$ by \eqref{O:I:symbound}, and so does the
endpoint. The actual norm character is trivial on this norm.
Apply its $P$-formula and Lemma~\ref{O:I:normlog} with $h=0$.
This proves both assertions. In particular this is a character identity
on a whole additive subgroup, not an inference from one test value.
\end{proof}

\subsection{A quadratic refinement of the stationary representative}
\begin{lemma}[enhanced coefficient and its exact affine correction]
\label{O:I:enhanced}
Let $\theta$ have conductor $m\ge2$ on a local field $L$ of odd residue
characteristic. Let $\Psi_L=\echar_L(\alpha\,\cdot)$ have largest trivial
ideal $\pp_L^J$, where $J\in\mathbf Z$. Put $d=\floor{m/2}$ and $s=\ceil{m/2}$.
There is $B\in L$ of valuation $J-m$, determined modulo $\pp_L^{J-d}$,
such that
\begin{equation}\label{O:I:enhancedformula}
 \theta(1-z)=\Psi_L(BP(z))\qquad(z\in\pp_L^d).
\end{equation}
If $C=B+\eta$ with $v_L(\eta)\ge J-s$, it is an ordinary stationary
coefficient in the same sign convention, and
\begin{equation}\label{O:I:completeformula}
 \Delta_L(\theta,\echar_L)=
 \theta((-\alpha C)^{-1})\Psi_L(-C)
 \Psi_L\!\left(\frac{\eta^2}{2B}\right)g_L,
 \qquad g_L\in\mu_4.
\end{equation}
For even $m$, $g_L=1$ and the correction is $1$.
If $v_L(\eta)\ge J-d$, the correction is $1$ in either parity.
\end{lemma}
\begin{proof}
Since $pd\ge m$, Lemma~\ref{O:I:logchart} identifies the character on $U_L^d$
with an additive character on $\pp_L^d/\pp_L^m$. The trace pairing gives
$B$ and its ambiguity, and the exact conductor gives $v_L(B)=J-m$.
On this depth $P(z)\equiv z+z^2/2$ at the character precision, since
$3d\ge m$. Thus
\[
 \theta(1+z)=\Psi_L(-Bz+Bz^2/2).
\]
The conventional Lamprecht linear coefficient is $-\alpha C$; the minus
sign must not be dropped. Choose any admissible denominator $\Gamma$ and
use the unit numerator $-\alpha C\Gamma$. Its elementary factor is
$\theta((-\alpha C)^{-1})\Psi_L(-C)$.
For odd $m=2d+1$, choose $\delta\OO_L=\pp_L^d$. The exact residual
function formed with $C$ is
\[
 \bar z\longmapsto
 \Psi_L(-\eta\delta z-B\delta^2z^2/2).
\]
The shift $\eta/(B\delta)$ is integral, by the assumed ordinary precision.
Completing the square produces the factor
$\Psi_L(\eta^2/(2B))$. The remaining homogeneous, nondegenerate quadratic
Gauss phase is a fourth root of unity. Indeed, writing
$G=\sum_{x\ne0}\chi_2(x)\psi_0(x)$ for a nontrivial residue additive
character, orthogonality gives $|G|^2=|k_L|$, and the substitution
$x\mapsto-x$ gives $\overline G=\chi_2(-1)G$.
Thus $G^2=\chi_2(-1)|k_L|$; scaling the quadratic coefficient only
multiplies $G$ by $\chi_2$ of that coefficient. In even conductor the ordinary
linearization already holds on $\pp_L^d$ and Lamprecht has no residual sum.
Finally $v_L(\eta^2/B)\ge J$ in the even case, and also whenever the
stronger depth $J-d$ holds. This proves every assertion.
\end{proof}

\subsection{Realizing an actual pair with an unramified Artin--Schreier element}
Let $k$ be the residue field of $F,E$ and $\kappa$ that of $U,K$.
Write $Q=|k|$ and let $\phi$ be arithmetic Frobenius in $\Gal(U/F)$.
Choose $\bar d\in\kappa$ with $\bar d^Q-\bar d=1$.
Then $\bar c=\bar d^p-\bar d\in k$ has absolute trace $1$.
Choose any lift $c\in\OO_F$ of $\bar c$. The polynomial
$X^p-X-c$ is irreducible and unramified over $F$, so we can choose
\begin{equation}\label{O:I:dchoice}
 d\in U,\qquad d^p-d=c,\qquad d\bmod\pp_U=\bar d.
\end{equation}
In mixed characteristic, Hensel's lemma at $d+1$ gives
$\phi(d)-(d+1)\in p\OO_U$. In equal characteristic this difference is zero.
Also
\begin{equation}\label{O:I:boundse0}
 e_0+q_0=T,\qquad
 \floor{T/2}\ge q_0,\qquad 2e_0\ge T,
 \qquad e_0\ge\ceil{T/2}.
\end{equation}
The elementary ceiling inequalities hold for every $T\ge2$, $p\ge3$.

\begin{proposition}[global compatible models, not merely leading classes]
\label{O:I:canonical}
With $\tau_F,\tau_U,\Psi_L$ as above, there are characters $\chi_U^0$
and $\chi_E^0$ of conductor $T$, with a common invariant primitive norm pullback,
such that
\begin{align}
 (\chi_U^0)^\phi/\chi_U^0&=\tau_U,\label{O:I:modelcomm}\\
 \chi_U^0(1-z)&=\Psi_U(d^pP(z))
 &&(z\in\pp_U^{q_0}),\label{O:I:modelU}\\
 \chi_E^0(1-z)&=\Psi_E(cP(z))
 &&(z\in\pp_E^{q_0}).\label{O:I:modelE}
\end{align}
Here $\chi^\phi(y)=\chi(\phi^{-1}y)$.
\end{proposition}
\begin{proof}
First define the right side of \eqref{O:I:modelU} on $H=U_U^{q_0}$.
It is a character by Lemma~\ref{O:I:logchart}, with exact conductor $T$.
Its conjugate quotient is $\tau_U$ on $H$ because
\[
 \phi(d^p)-d^p-1\in p\OO_U,
 \qquad v_U(p)+q_0\ge e_0+q_0=T.
\]
This uses the ideal conductor of $\Psi_U$, so it is valid for every
coefficient in $p\OO_U$, not only for rational multiples.

We now extend while preserving the \emph{exact} conjugate quotient.
Put $V=U_U^1$ and $I=\{\phi^{-1}(v)/v:v\in V\}$. Prescribe on $I$
\[
 \kappa(\phi^{-1}(v)/v)=\tau_U(v).
\]
This is well defined: the kernel of $v\mapsto\phi^{-1}(v)/v$ is
$U_F^1$, and $\tau_U(f)=\tau_F(f^p)=1$ there.
If $\phi^{-1}(v)/v\in H$, unramified Teichmuller expansions imply
$v=f h$ for $f\in U_F^1$ and $h\in H$; indeed every coefficient of a
Frobenius-fixed truncated expansion lies in $k$. Thus the prescribed
character on $I\cap H$ agrees with \eqref{O:I:modelU}.
The two characters define a character on $IH$, trivial on $U_U^T$.
Extend it across the finite abelian quotient $V/U_U^T$ to a character of
$V$. Take its value to be $1$ on the prime-to-$p$ Teichmuller subgroup
and on a fixed uniformizer of $F$. The character $\tau_U$ is trivial on
these two factors, so \eqref{O:I:modelcomm} now holds on all of $U^\times$.
This gives $\chi_U^0$ with conductor $T$ and the required full $P$-formula.

Put $\chi_K^0=\chi_U^0\circ N_U$. It is invariant under $\Gal(K/U)$
by definition and under $\Gal(K/E)$ by \eqref{O:I:modelcomm}, since $\tau_U$
is a norm character. Hilbert 90 therefore makes it trivial on
$\ker N_E$. It descends to $N_EK^\times$ and extends to a character
$\chi_E^0$ on $E^\times$; this extension only involves the valuation
factor, since $K/E$ is unramified.

For $z\in\pp_K^{q_0}$ the ramified norm of $1-z$ is in $U_U^{q_0}$.
Equations \eqref{O:I:modelU}, \eqref{O:I:normlogeq}, and \eqref{O:I:conversioneq},
using $N_U(d)=d^p$ \emph{exactly}, give
\[
 \chi_K^0(1-z)
 =\Psi_K(d^pP(z))\Psi_U(N_U(dP(z)))
 =\Psi_K((d^p-d)P(z))=\Psi_K(cP(z)).
\]
The unramified congruence \eqref{O:I:unramlog} shows that this is the pullback
of the right side of \eqref{O:I:modelE}. The norm $U_K^{q_0}\to U_E^{q_0}$
is onto, so \eqref{O:I:modelE} follows on its entire stated subgroup.
Since $c$ is a unit and the ideal conductor of $\Psi_E$ is $T$,
$\chi_E^0$ has exact conductor $T$. If the common pullback descended from
$F$, then $\chi_U^0$ would be a base norm pullback times an upper norm
character. Both are fixed by $\phi$, the latter by
Lemma~\ref{D:crossed-norm}. This contradicts its nontrivial conjugate
quotient $\tau_U$, proving primitivity.
\end{proof}


\begin{proposition}[actual common twist]\label{O:I:actualtwist}
Choose $\tau_F$ using Lemma~\ref{U:conjugacy}, and construct the models of
Proposition~\ref{O:I:canonical}. There is a character $\lambda$ of $F^\times$
such that, after changing only the unramified value of $\chi_E^0$ if needed,
\begin{equation}\label{O:I:twistpair}
 \theta_U=\chi_U^0(\lambda\circ n_U),\qquad
 \theta_E=\chi_E^0(\lambda\circ n_E).
\end{equation}
Write $m_U=T+h$. Then $h\ge0$, $m_E=T+ph$, and for $h>0$ the character
$\lambda$ has conductor $T+h$. For $h=0$ its conductor is at most $T$.
\end{proposition}
\begin{proof}
The quotient $\theta_U/\chi_U^0$ is $\phi$-invariant, so Hilbert 90
makes it a norm pullback of a character of $F^\times$ after extension
from $n_UU^\times$. Compatibility implies that
$\theta_E/(\chi_E^0(\lambda n_E))$ is an unramified norm character for
$K/E$. Absorb it into $\chi_E^0$; the full unit formula \eqref{O:I:modelE}
is unchanged.
The conductor of $\theta_U$ is at least that of its commutator, namely
$T$. If it exceeds $T$, its common twist has the same conductor because
$\chi_U^0$ has conductor $T$ and unramified norm pullback preserves
conductors. The remaining assertion about $\lambda$ follows similarly.
All twists of $\theta_U$ by $S(K/U)$ are conjugate and have the same
conductor. Therefore the minimal-orbit conductor formula in Ueda applies,
both at $m_U=T$ and above it; it gives $m_K=T+ph$.
The unramified edge $K/E$ preserves conductors, so $m_E=m_K$.
\end{proof}

\subsection{Exact norm choices at the precision actually available}
First suppose $0\le h\le t$. Set
\[
 m_U=T+h,\quad m_E=T+ph,\quad
 d_U=\floor{m_U/2},\quad d_E=\floor{m_E/2}.
\]
In particular $d_U,d_E\ge q_0$.
By the logarithm chart and the additive pairing choose $A_*\in F$ so that
\begin{equation}\label{O:I:lambdachart}
 \lambda(1-z)=\Psi_F(A_*P(z))\qquad(z\in\pp_F^{d_U}).
\end{equation}
For $h>0$, $v_F(A_*)=-h$ and its ambiguity is $\pp_F^{T-d_U}$.
For $h=0$, its class lies in $\OO_F/\pp_F^{T-d_U}$; a zero class may be
represented by $0$.

\begin{lemma}\label{O:I:chooseA}
There is $x\in E$ with $A=n_E(x)$ having the following properties.
If $0\le h<t$, one may take $A_*=A$ in \eqref{O:I:lambdachart}.
If $h=t$, one has instead
\begin{equation}\label{O:I:epsilon}
 A_*=A+\epsilon,\qquad\epsilon\in\OO_F.
\end{equation}
For $h>0$, $v_E(x)=v_F(A)=-h$. For $h=0$, $x$ is integral; the zero
coefficient is treated by $x=A=0$.
\end{lemma}
\begin{proof}
For $0<h<t$, the relative precision of the enhanced class is
$T+h-d_U=\ceil{m_U/2}\le t$. Thus \eqref{O:I:normrep} supplies an exact
norm in that class. For $h=0$, a nonzero class has valuation $a\ge0$;
its relative precision is $T-d_U-a\le\ceil{T/2}\le t$.
Again \eqref{O:I:normrep} applies. If the class is zero, use $x=0$.
At $h=t$, $m_U=2t+1$ and $d_U=t$. Use \eqref{O:I:normrep} at relative
precision $t$, not $t+1$. Its absolute error is integral, giving
\eqref{O:I:epsilon}. No exact norm representative at the critical successor
has been assumed. The valuation assertions follow from \eqref{O:I:normval}.
\end{proof}
Set $\epsilon=0$ in the nonboundary cases, and put
\begin{equation}\label{O:I:coefficients}
 B_U=A+d^p+\epsilon\in U,\qquad B_E=A-x+c\in E.
\end{equation}

\begin{proposition}[enhanced coefficients of the actual pair]\label{O:I:actualcoeff}
For the actual characters in \eqref{O:I:twistpair},
\begin{equation}\label{O:I:actualP}
 \theta_U(1-z)=\Psi_U(B_UP(z))\quad(z\in\pp_U^{d_U}),
 \qquad
 \theta_E(1-z)=\Psi_E(B_EP(z))\quad(z\in\pp_E^{d_E}).
\end{equation}
Thus these are full quadratic stationary representatives, not just
coefficients on the deeper linearization layer.
\end{proposition}
\begin{proof}
For the unramified pullback of $\lambda$, the logarithm identity through
an unramified norm is valid modulo $\pp_F^{m_U}$ at depth $d_U$, because
$pd_U\ge m_U$. It gives coefficient $A_*$ upstairs. Combine it with
\eqref{O:I:modelU} to obtain the first formula.

For $z\in\pp_E^{d_E}$ put $u=1-n_E(1-z)$. The symmetric estimates give
\[
 v_F(u)\ge\min\{d_E,\floor{(d_E+D_0)/p}\}\ge d_U.
\]
The last inequality follows from $d_E\ge d_U$ and
$d_E+D_0-pd_U=(p-1)(T+\varepsilon)/2\ge0$, where
$\varepsilon\in\{0,1\}$ is the common conductor parity.
Also $d_E-h\ge q_0$, since
$d_E-h=\floor{(T+(p-2)h)/2}\ge\floor{T/2}\ge q_0$.
For $h=0$ the same inequalities hold with $x$ integral.
Lemma~\ref{O:I:normlog} therefore gives
\[
 P(u)\equiv\Tr P(z)+n_E(P(z))\pmod{\pp_F^{T+h}}.
\]
Multiply by $A$, whose valuation is $-h$ for $h>0$ and nonnegative for
$h=0$. Conversion \eqref{O:I:conversioneq} applies to $xP(z)$, and the
\emph{exact} equality $n_E(x)=A$ gives
\[
 \Psi_F(A P(u))
 =\Psi_E(A P(z))\Psi_F(n_E(xP(z)))
 =\Psi_E((A-x)P(z)).
\]
At $h=t$, $d_E=(p+1)t/2\ge T$. It follows from the same norm estimate
that $u\in\pp_F^T$. Hence the extra coefficient $\epsilon\in\OO_F$ in
\eqref{O:I:epsilon} contributes $\Psi_F(\epsilon P(u))=1$.
Thus the displayed formula is the actual pullback of $\lambda$ at this
boundary too. Multiply by \eqref{O:I:modelE} to obtain the second formula.
\end{proof}

For $h>0$, $v_U(B_U)=-h$ and $v_E(B_E)=-ph$, with unique dominant
term $A$. For $h=0$, if $\xi=\bar x\in k$, their residues are
$\xi^p+\bar d^p$ and $\xi^p-\xi+\bar c$.
The first is nonzero because $\bar d\notin k$; the second has absolute
trace $1$, so is nonzero also. Thus the coefficients have exactly the
valuations required by their stated conductors in every case.

\subsection{Two actual norms of a single common element}
Put
\begin{equation}\label{O:I:commonC}
 C=x+d\in K^\times,\qquad Z_U=N_U(C),\qquad Z_E=N_E(C).
\end{equation}
These are literal elements in their respective fields. Because $d$ lies in
$U$, it is fixed under $\Gal(K/U)$. Because $x$ lies in $E$, it is fixed
under $\Gal(K/E)$. Hence, writing $a_i=e_i^{E/F}(x)$,
\begin{align}
 Z_U&=d^p+\sum_{i=1}^{p-1}a_i d^{p-i}+A,\label{O:I:ZU}\\
 Z_E&=x^p-x+c.\label{O:I:ZE}
\end{align}
The second identity is the exact polynomial $X^p-X-c$, evaluated at $-x$;
it uses actual conjugates in mixed characteristic as well.
Define
\[
 \eta_U=Z_U-B_U=\sum_{i=1}^{p-1}a_i d^{p-i}-\epsilon,
 \qquad \eta_E=Z_E-B_E=x^p-A.
\]

\begin{lemma}[ordinary and enhanced depths]\label{O:I:normdepths}
For $0\le h<t$, $Z_U$ and $Z_E$ are enhanced representatives in
Lemma~\ref{O:I:enhanced}; both affine corrections in \eqref{O:I:completeformula}
are $1$.
For $h=t$, they are ordinary stationary representatives. The $E$ one is
enhanced if $p>3$. In this boundary case, if $b=a_{p-1}$, then
\begin{equation}\label{O:I:etaUresidue}
 \eta_U\equiv b d-\epsilon\pmod{\pp_U},\qquad b\in\OO_F.
\end{equation}
\end{lemma}
\begin{proof}
For integral $x$, \eqref{O:I:symbound} puts all the $a_i$ in $\pp_F^{e_0}$.
The characteristic-polynomial identity gives
\begin{equation}\label{O:I:normpower}
 x^p-A=\sum_{i=1}^{p-1}(-1)^{i+1}a_i x^{p-i},\qquad
 v_E(x^p-A)\ge p v_E(x)+(p-1)t.
\end{equation}
The bound follows from $p\floor{(ia+D_0)/p}\ge ia+(p-1)t$ with
$a=v_E(x)$. Thus for $h=0$ the errors have depths at least $e_0$ and
$(p-1)t$, respectively, both at least $T-\floor{T/2}$.
If $x=0$ the assertions are immediate.

For $h>0$, \eqref{O:I:symbound} and \eqref{O:I:normpower} give
\begin{equation}\label{O:I:errorsbounds}
 v_U\!\left(\sum_{i=1}^{p-1}a_i d^{p-i}\right)
 \ge\floor{(p-1)(T-h)/p},\qquad
 v_E(\eta_E)\ge(p-1)t-ph.
\end{equation}
If $h<t$, put $r=T-h\ge2$. The first bound is
$r-\ceil{r/p}\ge\ceil{r/2}=T-d_U$.
Here $\ceil{r/p}\le\floor{r/2}$ for $r\ge2$, $p\ge3$.
For the second, $h\le t-1$ gives
\[
 2((p-1)t-ph)-(T-ph)
 =(2p-3)t-ph-1\ge(p-3)t+p-1\ge0;
\]
hence $(p-1)t-ph\ge T-d_E$. These are precisely enhanced depths.

At $h=t$, $m_U=2t+1$, $m_E=(p+1)t+1$, and the ordinary ambiguity depths
are $0$ over $U$ and $-(p-1)t/2$ over $E$.
The two bounds $v_U(\eta_U)\ge0$ and $v_E(\eta_E)\ge-t$ suffice.
If $p>3$, also $-t\ge1-(p-1)t/2=T-d_E$, so the $E$ representative is
enhanced. Finally, at this boundary $a_{p-1}$ is integral and all
$a_i$ with $i\le p-2$ have valuation at least
$\floor{(t+p-1)/p}=\ceil{t/p}\ge1$. This proves \eqref{O:I:etaUresidue}.
\end{proof}

\begin{lemma}[exact reduction to a scalar phase]\label{O:I:scalarphase}
Let
\begin{equation}\label{O:I:Sdef}
 S=\Tr_{U/F}Z_U-\Tr_{E/F}Z_E.
\end{equation}
The quotient of actual local constants satisfies
\begin{equation}\label{O:I:ratioformula}
 \frac{\Delta_U(\theta_U)}{\Delta_E(\theta_E)}
 \in\mu_4\,\Psi_F(-S)\,
 \frac{\Psi_U(\eta_U^2/(2B_U))}
      {\Psi_E(\eta_E^2/(2B_E))}.
\end{equation}
Furthermore one has the exact identity
\begin{equation}\label{O:I:Sexact}
 S=p A-\Tr_{E/F}(x^p)+p\Tr_{E/F}(x).
\end{equation}
\end{lemma}
\begin{proof}
Apply Lemma~\ref{O:I:enhanced} with the ordinary representatives $Z_U,Z_E$.
The factors at $-\alpha\in F^\times$ cancel by \eqref{U:det-character}, and
\[
 \theta_U(Z_U)=\theta_K(C)=\theta_E(Z_E).
\]
Thus all multiplicative stationary values cancel exactly, leaving
\eqref{O:I:ratioformula}, with the displayed signs forced by
\eqref{O:I:completeformula}.
Newton identities for $X^p-X-c$ give
\[
 \Tr_{U/F}(d^j)=0\ (1\le j<p-1),\quad
 \Tr_{U/F}(d^{p-1})=p-1,\quad
 \Tr_{U/F}(d^p)=pc.
\]
Take traces in \eqref{O:I:ZU} and \eqref{O:I:ZE} and subtract. This gives
\eqref{O:I:Sexact}. No approximate translated roots or discarded norm errors
enter either assertion.
\end{proof}

\subsection{Vanishing away from the last cubic row}
\begin{lemma}\label{O:I:Sdepth}
One has $S\in\pp_F^T$ if $h=0$, if $1\le h<t$, or if $p>3$ and
$h=t$.
\end{lemma}
\begin{proof}
For integral $x$, all smaller elementary coefficients and all power traces
$\Tr(x^j)$ have valuation at least $e_0$. Newton's identity at degree $p$
puts $\Tr(x^p)-pA$ in $\pp_F^{2e_0}$, and $p\Tr x$ has valuation at least
$v_F(p)+e_0\ge2e_0$. Use $2e_0\ge T$.

For $h>0$ use the following integral universal identity:
\[
 \Tr(x^p)-pA=a_1^p+p\,Q(a_1,\ldots,a_{p-1}),
\]
where every monomial of $Q$ has weight $p$ and at least two factors.
Indeed the power sum is congruent to $a_1^p$ modulo $p$ as an integral
symmetric polynomial, its $a_p$ coefficient is $p$, and there is no
one-factor monomial of weight $p$ among $a_1,\ldots,a_{p-1}$.
The estimates
\[
 v_F(a_i)\ge\frac{-ih+(p-1)t}{p},\qquad
 v_F(p)\ge\frac{(p-1)t}{p}
\]
give the respective lower bounds
\[
 v_F(pa_1)\ge\frac{2(p-1)t-h}{p},\qquad
 v_F(a_1^p)\ge(p-1)t-h,
\]
and, for every monomial of $pQ$, a lower bound
$3(p-1)t/p-h$.
All three bounds are strictly greater than $t$ when $p>3$ and $h\le t$,
or when $p=3$ and $h<t$. Since valuations are integral, they are at least
$T$. In equal characteristic the terms divisible by $p$ are zero and the
same conclusion follows. This proves the lemma.
\end{proof}

For $h<t$, both corrections in \eqref{O:I:ratioformula} are $1$ by
Lemma~\ref{O:I:normdepths}, so Lemma~\ref{O:I:Sdepth} already puts the quotient
in $\mu_4$.
At $h=t$ one has $B_U=A+d^p+\epsilon$ with $v_F(A)=-t$ and
$\eta_U\equiv b d-\epsilon\pmod{\pp_U}$. Therefore
\begin{equation}\label{O:I:Uboundarycorrection}
 \Psi_U\!\left(\frac{\eta_U^2}{2B_U}\right)
 =\Psi_U\!\left(\frac{(bd-\epsilon)^2}{2A}\right).
\end{equation}
The numerator replacement costs depth $t+1$; the inverse-denominator
replacement costs depth $2t\ge t+1$. These are exact character-kernel
estimates. If $p>3$, both $\Tr(d)$ and $\Tr(d^2)$ vanish, while the trace
of $1$ is $p$. The right side of \eqref{O:I:Uboundarycorrection} is
$\Psi_F(p\epsilon^2/(2A))=1$, since its argument has valuation at least
$v_F(p)+t\ge T$ (and is zero in equal characteristic).
The $E$ correction is $1$ by its enhanced depth. Lemma~\ref{O:I:Sdepth}
therefore also settles this last row for $p>3$.

\subsection{The cubic boundary with both affine terms retained}
It remains to treat $p=3$ and $h=t$, without any restriction on divisibility
of $t$. Write
\[
 a=\Tr_{E/F}x,\qquad b=e_2^{E/F}(x),\qquad A=n_E(x).
\]
The symmetric estimates give
\begin{equation}\label{O:I:cubicbounds}
 v_F(a)\ge\ceil{t/3}\ge1,\quad v_F(b)\ge0,
 \quad v_F(A)=-t,\quad v_E(x)=-t.
\end{equation}
The full coefficients remain $B_U=A+d^3+\epsilon$ and
$B_E=A-x+c$.
Because $\Tr_{U/F}(d^2)=2$ and $\Tr_{U/F}(d)=0$,
\eqref{O:I:Uboundarycorrection} gives
\begin{equation}\label{O:I:cubicU}
 \Psi_U\!\left(\frac{\eta_U^2}{2B_U}\right)=\Psi_F(b^2/A).
\end{equation}
The $3\epsilon^2/(2A)$ term is in $\pp_F^T$. Thus the critical ambiguity
of the approximate norm choice has disappeared by the unramified trace,
not by an assumed exact norm at depth $t+1$.

The characteristic polynomial of $x$ gives the exact identity
\begin{equation}\label{O:I:cubiceta}
 \eta_E=x^3-A=a x^2-bx.
\end{equation}
One has $v_E(\eta_E)\ge-t$, and
$v_E(1/B_E-1/A)\ge5t$. Consequently
\begin{equation}\label{O:I:cubicE}
 \Psi_E\!\left(\frac{\eta_E^2}{2B_E}\right)
 =\Psi_F\!\left(\frac{\Tr_{E/F}(\eta_E^2)}{2A}\right),
\end{equation}
since the omitted $E$-argument has valuation at least $3t\ge T$.
The trace maps $\pp_E^T$ into $\pp_F^T$ by \eqref{O:I:traceideal}.

\begin{lemma}[exact cubic completion identity]\label{O:I:cubicidentity}
Put
\[
 \Lambda=S-\frac{b^2}{A}+\frac{\Tr_{E/F}(\eta_E^2)}{2A}.
\]
Then
\begin{equation}\label{O:I:Lambdaexact}
 \Lambda=3a+a^3-\frac{b^2+b^3}{A}
          +\frac{a^2(a^2-3b)^2}{2A}.
\end{equation}
Moreover $\Psi_F(\Lambda)=1$.
\end{lemma}
\begin{proof}
Newton's identities give
\[
 \Tr x^2=a^2-2b,\quad
 \Tr x^3=a^3-3ab+3A,\quad
 \Tr x^4=a^4-4a^2b+2b^2+4aA.
\]
Using \eqref{O:I:cubiceta}, expand
\[
 \Tr\eta_E^2=a^2\Tr x^4-2ab\Tr x^3+b^2\Tr x^2
 =a^6-6a^4b+9a^2b^2-2b^3+4a^3A-6abA.
\]
Equation \eqref{O:I:Sexact} reads $S=3a-a^3+3ab$.
Substitution proves \eqref{O:I:Lambdaexact} as an exact identity in both
characteristics.

The last summand of \eqref{O:I:Lambdaexact} belongs to $\pp_F^T$:
\eqref{O:I:cubicbounds} gives $v_F(a)\ge1$, $v_F(a^2-3b)\ge1$, and
$v_F(A^{-1})=t$.
For the first remaining pair, apply \eqref{O:I:conversioneq} to the element
$a\in F\subset E$. It lies in $\pp_E^{q_0}$, since
$v_E(a)\ge3\ceil{t/3}\ge t\ge q_0$. It gives
\[
 \Psi_F(3a+a^3)=1.
\]
For the other pair apply \eqref{O:I:conversioneq} to $b/x$.
This element also has $E$-valuation at least $t\ge q_0$, and the exact
reciprocal-trace identity is
\[
 \Tr_{E/F}(b/x)=b\,\frac{e_2(x)}{n_E(x)}=b^2/A,
 \qquad n_E(b/x)=b^3/A.
\]
Hence $\Psi_F((b^2+b^3)/A)=1$. Combining these two genuine norm-character
identities with the ideal estimate proves $\Psi_F(\Lambda)=1$.
This is a phase assertion; no unjustified claim that $\Lambda$ itself
belongs to $\pp_F^T$ is made.
\end{proof}

Equations \eqref{O:I:ratioformula}, \eqref{O:I:cubicU}, and \eqref{O:I:cubicE} now give
\[
 \frac{\Delta_U(\theta_U)}{\Delta_E(\theta_E)}
 \in\mu_4\,\Psi_F(-\Lambda)=\mu_4.
\]
This finishes every conductor with $0\le h\le t$, including the affine
odd-conductor endpoint in residual characteristic $3$.

\subsection{Stable conductors and completion of the branch}
Suppose $h\ge T$, so $m=m_F(\lambda)=T+h\ge2T$.
Choose a stationary numerator $\beta\in F^\times$ for $\lambda$ with
admissible denominator $\Gamma\in F^\times$. Ueda's stable high-conductor
table permits the same $\Gamma,\beta$ for its norm pullbacks to $U$ and
$E$. The denominator assertion includes the different:
\[
 m_E(\lambda n_E)+n_E(\echar_E)
 =T+p(m-T)+p n_F(\echar_F)+(p-1)T
 =p(m+n_F(\echar_F))=v_E(\Gamma).
\]
The minimal models have conductor $T$, at most half the conductor of each
pullback. Exact stable twist gives
\[
 \Delta_L(\theta_L)=\chi_L^0(\Gamma/\beta)
       \Delta_L(\lambda\circ n_L),\qquad L=U,E.
\]
The multipliers agree by the common restriction to $F^\times$.
By Ueda's First Main Lemma on each lower edge and exact stable twist
on $F$, the product over its norm characters is
\[
 \prod_{\nu\in S(L/F)}\Delta_F(\nu\lambda)
 =\Delta_F(\lambda)^p
   \prod_{\nu\in S(L/F)}\nu(\Gamma/\beta)
 =\Delta_F(\lambda)^p.
\]
The character product is trivial in an odd cyclic group, and the product
of its individual local constants is $1$ by inverse pairing. Therefore
$\Delta_U(\lambda n_U)=\Delta_E(\lambda n_E)$ and the stable quotient
is exactly $1$.

\begin{proof}[Proof of Theorem~\ref{O:I:main}]
The stable argument covers $h\ge T$. For $0\le h\le t$, the exact common
norm representatives, the scalar phase evaluation, and the retained cubic
affine correction prove that the actual quotient belongs to $\mu_4$.
Its $p$th power is $1$ by Corollary~\ref{U:odd-power}. Since $p$ is odd,
$\mu_p\cap\mu_4=\{1\}$, and the quotient is $1$.
The lower cyclic products are $1$ by the same corollary. These cases exhaust
all possible conductors, and every step was valid in either characteristic.
\end{proof}


\section{Artin--Schreier coordinates and trace--norm estimates}\label{O:sec:as}\label{sec:odd-estimates}
\subsection{The ramification data and valuation conventions}
Let $A\subset B_1,B_2\subset B$ be a totally ramified Galois diamond with
$\operatorname{Gal}(B/A)=C_p^2$, $p>2$.  Choose the two lines so that their
breaks are $t=t_1\le t_2$, and the breaks of $B/B_2$ and $B/B_1$ are
$t$ and $t'=t+p\delta$, respectively, where $\delta=t_2-t\ge0$.
Write
\[
 N_i=\Nm_{B/B_i},\quad S_i=\Tr_{B/B_i},\qquad
 n_i=\Nm_{B_i/A},\quad T_i=\Tr_{B_i/A}.
\]
All valuations $v_B,v_i,v_A$ are normalized on their own fields.  In
particular $v_B|_{B_i}=pv_i$, $v_i|_A=pv_A$, and
$v_A(n_i y)=v_i(y)$.
The different exponents of the lower edges are
$D_i=(p-1)(t_i+1)$; that of $B/B_1$ is $D'=(p-1)(t'+1)$.

We use Ueda's proved cyclic trace ideals, norm filtration, and strong
symmetric-function estimate.  For a totally ramified degree-$p$ cyclic
extension $L/M$ with break $u$ they read
\begin{align}
 \Tr_{L/M}(\pp_L^a)&=\pp_M^{\floor{(a+(p-1)(u+1))/p}},\label{O:A:eq:trace}\\
 v_M(E_i^{L/M}(y))&\ge\floor{(i v_L(y)+(p-1)(u+1))/p},\quad 1\le i<p.
 \label{O:A:eq:strong}
\end{align}
They hold for fractional ideals as well.  Additive Hilbert 90 and the first
nonzero ramification term are standard local algebraic inputs; no local
constant comparison is used.
In mixed characteristic put $V=v_B(p)$.  Applying \eqref{O:A:eq:trace} to $1$
on $B_2/A$ gives
\begin{equation}
 V\ge p(p-1)t_2.\label{O:A:eq:pbound}
\end{equation}
In equal characteristic all expressions divisible by $p$ below are zero.

\begin{lemma}[a consequence of the characteristic polynomial]\label{O:A:lem:power}
For every $y\in L$ in the preceding cyclic extension,
\[
 v_L(y^p-\Nm_{L/M}y)\ge p v_L(y)+(p-1)u.
\]
\end{lemma}
\begin{proof}
For $1\le i<p$, \eqref{O:A:eq:strong} implies
$pv_M(E_i(y))\ge i v_L(y)+(p-1)u$.
Every intermediate term $E_i(y)y^{p-i}$ of the characteristic polynomial
therefore has the displayed valuation.  Its constant term is $-\Nm y$
because $p$ is odd.
\end{proof}

\subsection{Construction of an exact Artin--Schreier coordinate}
\begin{proposition}\label{O:A:prop:AS}
There exists $\Delta\in B$ such that
\[
 \Delta^p-\Delta=a\in B_2,\qquad v_B(\Delta)=-t,
 \quad v_2(a)=-t.
\]
In this situation $p\nmid t$ and $B=B_2(\Delta)$.
\end{proposition}
\begin{proof}
Let $\sigma$ generate $\operatorname{Gal}(B/B_2)$.  First suppose the
characteristic is zero and put $\kappa=V-(p-1)t>0$.  The trace-ideal formula
for $B/B_2$ gives
\[
 S_2(\PP_B^\kappa)=p\OO_{B_2}.
\]
Choose $\eta\in\PP_B^\kappa$ with $S_2\eta=-p$.
Since $S_2(1+\eta)=0$, additive Hilbert 90 supplies $d\in B$ with
$\sigma d-d=1+\eta$.

Subtract an element of $B_2$ from $d$ so that its valuation is maximal in
its additive $B_2$-coset.  A maximum exists because $B_2$ is closed in $B$.
This maximal valuation is not divisible by $p$: otherwise one could cancel
its leading term by an element of $B_2$ and increase it.  The first
ramification term therefore gives
\[
 v_B(\sigma d-d)=v_B(d)+t.
\]
The left side is zero, so $v_B(d)=-t$ and $p\nmid t$.

Put $a_0=d^p-d$.  The equation $\sigma d-d=1+\eta$, the binomial theorem,
and $v_B(d)=-t$ give
$v_B(\sigma a_0-a_0)\ge\kappa$.
Choose a best additive approximation $a\in B_2$ to $a_0$.
If $a_0\notin B_2$, the same first-ramification-term argument applied to
the remainder gives
\[
 v_B(a_0-a)\ge\kappa-t=V-pt>0.
\]
For $f(X)=X^p-X-a$, the shifted polynomial $f(d+Z)$ has integral
coefficients, unit linear coefficient, and constant coefficient in
$\PP_B^{V-pt}$.  Indeed
$v_B(p d^{p-1})=\kappa>0$, and the other intermediate coefficients are
integral.  Hensel's lemma gives a root
$\Delta\equiv d\pmod{\PP_B^{V-pt}}$.  It has valuation $-t$.
The valuation of $\Delta^p-\Delta$ is $-pt$, proving $v_2(a)=-t$.
Since $p\nmid t$, this root is not in $B_2$.

In characteristic $p$, additive Hilbert 90 directly solves
$\sigma d-d=1$.  The same best-approximation argument gives valuation $-t$,
and $d^p-d$ is already fixed by $\sigma$.
\end{proof}

Set
\[
 e_i=E_i^{B/B_1}(\Delta),\quad b=N_1\Delta=e_p,
 \quad s=-e_{p-1},\qquad K_0=1+pt_2.
\]
Then $N_2\Delta=a$ and $v_1(b)=-t$.
Because the polynomial over $B_2$ is exactly $X^p-X-a$,
\begin{equation}
 S_2(\Delta^j)=0\ (1\le j<p-1),\quad
 S_2(\Delta^{p-1})=p-1,\quad S_2(\Delta^p)=pa.
 \label{O:A:eq:moments}
\end{equation}

\subsection{Coefficient extraction and twisted estimates}
\begin{lemma}[coefficient extraction]\label{O:A:lem:coeff}
If $\theta\in B_1$ and $\theta=\sum_{i=0}^{p-1}Y_i\Delta^i$ with
$Y_i\in B_2$, then
\[
 v_2(Y_i)\ge v_1(\theta)+it\quad(0\le i<p).
\]
\end{lemma}
\begin{proof}
Since $p\nmid t$, the valuations $pv_2(Y_i)-it$ are distinct modulo $p$.
Each is at least $pv_1(\theta)$, proving the assertion for $i=0$ and the
preliminary bound for every $i$.
For each nonzero $\xi\in\mu_{p-1}\subset A$, a conjugate of $\Delta$ is
\[
 \Delta_\xi=\Delta+\xi+\eta_\xi,
 \qquad v_B(\eta_\xi)\ge V-(p-1)t.
\]
This follows by applying Hensel's lemma to the shifted polynomial at
$\Delta+\xi$; its constant coefficient has that valuation and its derivative
is a unit.  In equal characteristic $\eta_\xi=0$.

Induct simultaneously for all $\theta\in B_1$ on $j$.
The restriction of this conjugating automorphism to $B_1$ has break $t$,
so $v_1(\theta_\xi-\theta)\ge v_1(\theta)+t$.
The error obtained by replacing $\Delta_\xi$ by $\Delta+\xi$ in the
expansion of $\theta_\xi$ has valuation at least
\[
 pv_1(\theta)+V-(p-2)t.
\]
Its $\Delta^{j-1}$ coefficient has $v_2$ at least $v_1(\theta)+jt$,
because $V\ge p(p-1)t\ge(p-1)(j+1)t$.
The induction hypothesis applied to $\theta_\xi-\theta$ thus yields
\[
 \sum_{i=j}^{p-1}\binom{i}{j-1}\xi^{i-j+1}Y_i
 \in\PP_{B_2}^{v_1(\theta)+jt}.
\]
Multiply by $\xi^{-1}$ and sum over $\mu_{p-1}$.
The exact power sums kill every term except $(p-1)jY_j$;
this multiplier is a unit.
\end{proof}

\begin{lemma}[twisted trace and symmetric functions]\label{O:A:lem:twisted}
For $Y\in B_2$ and $1\le i<p$,
\begin{align}
 v_1(S_1(Y\Delta^i))&\ge(p-1)t_2-it+v_2(Y),\label{O:A:eq:twtrace}\\
 v_1(E_i^{B/B_1}(Y\Delta))&\ge(p-1)t_2+i(v_2(Y)-t).
 \label{O:A:eq:twsym}
\end{align}
Also $v_1(S_1Y)\ge(p-1)t_2+v_2(Y)$.
In particular
\begin{equation}
 v_1(e_i)\ge(p-1)t_2-it\ (i<p),\qquad
 v_1(s)\ge(p-1)\delta.\label{O:A:eq:ei}
\end{equation}
\end{lemma}
\begin{proof}
The assertion for $S_1Y=T_2Y\in A$ follows from \eqref{O:A:eq:trace}.
For $i\ge1$ put $c=(p-1)t_2-it+v_2Y$.
By trace duality it suffices to show $T_1(\theta S_1(Y\Delta^i))\in\OO_A$
for every $v_1(\theta)\ge-D_1-c$.
Write $\theta=\sum Y_j\Delta^j$ and use trace transitivity.
For $1\le k\le2p-2$, the only nonzero power traces over $B_2$ are
those of degrees $p-1,p,2p-2$, with values $p-1,pa,p-1$.
Thus
\[
 S_2(\theta\Delta^i)=(p-1)Y_{p-1-i}+paY_{p-i}
       +\ind_{i=p-1}(p-1)Y_{p-1}.
\]
Lemma~\ref{O:A:lem:coeff} shows that, after multiplication by $Y$, all three
terms lie in $\PP_{B_2}^{-D_2}$.  For the first this is exactly
\[
 v_2Y+(p-1-i)t-D_1-c=-D_2;
\]
for the second use $v_2(pa)\ge-t$, and for the last there is an additional
nonnegative $(p-1)t$.  This proves the trace estimate.
Newton's identity for $E_i(Y\Delta)$ involves the power traces
$S_1(Y^j\Delta^j)$, $1\le j\le i<p$.  Each product with total weight $i$
has valuation at least $(p-1)t_2+i(v_2Y-t)$, and division by $i$ is legal.
\end{proof}

\subsection{The trace--norm congruence}
\begin{theorem}\label{O:A:thm:H14}
One has
\[
 T_1(b)-T_2(a)\in\pp_A^{1+t_2}.
\]
\end{theorem}
\begin{proof}
Compare the coefficient of $Z^p$ in the exact identity
\[
 n_1(N_1(1-Z\Delta))=n_2(N_2(1-Z\Delta)).
\]
The upper polynomials are
\[
 N_1(1-Z\Delta)=1+\sum_{i=1}^{p-1}(-1)^ie_iZ^i-bZ^p,
 \qquad N_2(1-Z\Delta)=1-Z^{p-1}-aZ^p.
\]
The right coefficient is $-T_2a$.  Besides $-T_1b$, the left coefficient
contains $-n_1(e_1)$ and traces of monomials in conjugates of the $e_i$
with $i<p$, at least two factors, and total weight $p$.
Indeed the cyclic group permutes the selections from the $p$ factors;
every orbit has size $p$, except the constant selection $e_1$ from every
factor, which gives $n_1(e_1)$.
By \eqref{O:A:eq:ei}, each trace of a nonconstant selection has valuation at least
\[
 \floor{\frac{2(p-1)t_2-pt+D_1}{p}}.
\]
This is at least $1+t_2$ whenever
$(p-2)t_2-t-1\ge0$.  The same condition makes
$v_A(n_1e_1)=v_1(e_1)\ge(p-1)t_2-t\ge1+t_2$.
This covers $p\ge5$ and $p=3$, $t_2>t$.

For $p=3$, $t_2=t$, first note that $T_1e_1=T_1S_1\Delta=0$ by
\eqref{O:A:eq:moments}.  The trace map on
$\PP_{B_1}^t/\PP_{B_1}^{t+1}\longrightarrow
 \pp_A^t/\pp_A^{t+1}$ is an isomorphism: \eqref{O:A:eq:trace} gives its
surjectivity, and both spaces have dimension one over the common residue
field.  Since $v_1(e_1)\ge t$, this proves $v_1(e_1)\ge t+1$.
The norm of $e_1$ is consequently deep enough.  Every remaining weight-three
selection has one $e_1$ factor and one $e_2$ factor.  Its trace has valuation
at least $\floor{(t+1+2(t+1))/3}=t+1$.
\end{proof}

\subsection{The weighted translated-root estimate}
Put $H=pt_2+t$.  For $x\in B_2$ and $1\le k\le p$, write
\[
 X=n_2x,\qquad R=v_B(x\Delta^{k-1})=pv_2x-(k-1)t.
\]
\begin{lemma}\label{O:A:lem:translate}
If $R\ge1$, then
\[
 X T_1(b^k)\equiv X\sum_{\xi\in\TT}N_1(\Delta+\xi)^k
 \pmod{\PP_{B_1}^{pt_2+R}},\qquad\TT=\{0\}\cup\mu_{p-1}.
\]
\end{lemma}
\begin{proof}
It is exact in characteristic $p$.  In mixed characteristic fix $\xi\ne0$,
put $D=\Delta+\xi$ and write the actual conjugate as $D(1+u)$.
Polynomial division gives the exact polynomial
\[
 Q_\xi(\Delta)=\frac{(\Delta+\xi)^p-(\Delta+\xi)-(\Delta^p-\Delta)}
                         {\Delta+\xi}.
\]
It has degree $p-2$, leading coefficient $p\xi$, and every coefficient is
in $p\OO_A$.  The equation for $u$ is
\[
 u=Q_\xi+pD^{p-1}u+D^{p-1}\sum_{j=2}^p\binom pj u^j.
\]
Hence
\[
 v_Bu\ge V-(p-2)t,\qquad
 v_B(u-Q_\xi)\ge2V-(2p-3)t.
\]
For $p\ge5$, the trace of the latter error has valuation at least $H$:
subtracting $H$ from its lower bound gives
\[
 (p-3)t-\floor{(p-2)t/p}+(2p-3)\delta\ge0.
\]
The symmetric functions $E_i(u)$ with $2\le i<p$ have at least this precision,
and $v_1(N_1u)=v_Bu\ge V-(p-2)t\ge H$.  Terms of degree at most $p-3$ in $Q_\xi$ have traces of
valuation at least
$2(p-1)t_2-(p-3)t\ge H$ by Lemma~\ref{O:A:lem:twisted}.

For $p=3$ retain one further iterate:
\[
 Q_\xi=3\xi\Delta,\qquad
 u=Q_\xi+3D^2Q_\xi+E,\quad v_BE\ge3V-5t.
\]
The trace of $E$ has valuation at least $V-t+2\delta\ge H$.
Moreover
\[
 3D^2Q_\xi=9\xi\Delta^3+18\xi^2\Delta^2+9\xi^3\Delta.
\]
Using $\Delta^3=\Delta+a$, each trace on the right has valuation at least
$6t_2-2t\ge H$.  The other norm coefficients satisfy
$v_1(E_2(u))\ge2V/3+2\delta\ge H$ and
$v_1(N_1u)\ge V-t\ge H$.
Thus in all cases
\begin{equation}
 N_1(1+u)\equiv1+\xi L\pmod{\PP_{B_1}^H},\qquad
 L=pS_1(\Delta^{p-2}),\quad
 v_1L\ge pt+2(p-1)\delta.\label{O:A:eq:unit-correction}
\end{equation}
Since $2v_1L\ge H$, the same formula raised to the $k$th power has linear
term $k\xi L$.

Now $N_1(\Delta+\xi)-b$ is integral by \eqref{O:A:eq:ei}; hence
$v_1(N_1(\Delta+\xi)^k-b^k)\ge-kt+t$.
After multiplying by $X$, the error term in \eqref{O:A:eq:unit-correction}
is in $\PP_{B_1}^{pt_2+R}$ because $v_1(Xb^k)=R-t$.
In the linear term, replacing $N_1(\Delta+\xi)^k$ by $b^k$ costs at least
$R+pt+2(p-1)\delta\ge R+pt_2$.
The remaining expression vanishes after summation because
$\sum_{\xi\in\mu_{p-1}}\xi=0$ exactly.
Finally, abelianness identifies the sum over actual conjugates with
$T_1(b^k)$.
\end{proof}

\subsection{A filtered Newton-polynomial calculation}
Define
\[
 P_k=b^k+\sum_{\xi\in\mu_{p-1}}N_1(\Delta+\xi)^k,
 \qquad Q_k=S_1((\Delta^p-\Delta)^k)=T_2(a^k).
\]
\begin{lemma}\label{O:A:lem:filtered}
With $X,R$ as above and $R\ge1$,
\[
 X(P_k-Q_k)\equiv(p-1)\ind_{k=p-1}X(1+e_{p-1}^p)
       \pmod{\PP_{B_1}^{pt_2+R}}.
\]
\end{lemma}
\begin{proof}
For a monomial $e_\lambda=\prod e_i^{\lambda_i}$ put
$r=\sum\lambda_i$, $n=\sum i\lambda_i$, and $\nu=\sum_{i<p}\lambda_i$.
The exact Teichmuller power sums show that $P_k$ contains it only if
$n\equiv k\pmod{p-1}$, with coefficient
\[
 (p-1)\binom kr\frac{r!}{\prod\lambda_i!},
\]
plus the separate term $b^k$.  Write $n=(p-1)v+k$.
Newton's generating-function formula shows that its coefficient in $Q_k$
is, for $0\le v\le k$,
\[
 (-1)^{r+v}\binom kv\frac{n(r-1)!}{\prod\lambda_i!}.
\]
These are integral coefficients, even when written as factorial quotients.
The valuation estimate \eqref{O:A:eq:ei} gives
\begin{equation}
 v_1(Xe_\lambda)\ge R-t+(p-1)\bigl((r-v)t+\nu\delta\bigr).
 \label{O:A:eq:monomial}
\end{equation}
For $v\ge0$, if $r\ge v+2$, the deficiency
$pr-n=v+p(r-v)-k\ge p$ forces $\nu\ge2$.
Then \eqref{O:A:eq:monomial} is at least
$R-t+2(p-1)t_2\ge R+pt_2$.
If $r=v$, only $v=r=k$, $e_\lambda=b^k$ is possible; its coefficients
are $1+(p-1)=p$ and $p$, so they cancel.

Suppose $r=v+1\le k$.  Their coefficient difference is exactly
\[
 \frac{(r-1)!}{\prod\lambda_i!}
 \left((p-1)r\binom kr+n\binom k{r-1}\right)
 =\frac{pk\binom k{r-1}(r-1)!}{\prod\lambda_i!}.
\]
Here $r\binom kr=(k-r+1)\binom k{r-1}$ and
$n=(p-1)(r-1)+k$.  The factorial denominators are units at $p$:
if $r=p$, then $k=p$ and $n\equiv1\pmod p$, so no $\lambda_i$ can equal $p$.
Except for the pure norm monomial, one has $\nu\ge1$.
The added valuation $v_1p\ge(p-1)t_2$ then makes
\eqref{O:A:eq:monomial} at least $R+pt_2$.
The pure norm possibility is $k=p,r=1,v=0$, $e_\lambda=b$;
its coefficient difference is $p^2$, again more than sufficient.

If $v=k$ and $r=k+1$, the first coefficient is absent and the difference is
$pk\,k!/\prod\lambda_i!$.
It is divisible by $p$ except when $k=p-1$ and
$e_\lambda=e_{p-1}^p$, in which case it is exactly $p-1$.
Indeed for $k\le p-2$ all factorial denominators are units;
for $k=p-1$ the only nonunit denominator comes from that constant selection;
for $k=p$ the numerator contains at least three factors of $p$ and the
factorial denominator at most one.
The remaining monomials with $r\ge v+2$ were already dealt with.

Finally $v=-1$ contributes only the constant monomial when $k=p-1$, giving
$p-1$, or $e_1$ when $k=p$, giving $p(p-1)e_1$.
The latter is deep enough by \eqref{O:A:eq:ei} and \eqref{O:A:eq:pbound}.
These cases exhaust the possible monomials and leave exactly the two
terms displayed in the lemma.
\end{proof}

\begin{theorem}[weighted norm--trace defect]\label{O:A:thm:h}
Define
\[
 h(z)=T_1N_1(\Delta z)-T_2N_2(\Delta z).
\]
For $x\in B_2$, $0\le j\le p-1$, and $R=v_B(x\Delta^j)\ge1$,
\begin{equation}
 h(x\Delta^j)\equiv(p-1)\ind_{j=p-2}n_2(x)(1-s^p)
       \pmod{\PP_{B_1}^{pt_2+R}}.\label{O:A:eq:h}
\end{equation}
\end{theorem}
\begin{proof}
Put $k=j+1$.
Exact multiplicativity shows that replacing $h(x\Delta^{k-1})$ by
$n_2(x)(T_1b^k-T_2a^k)$ costs
$T_2((n_2x-x^p)a^k)$.
By Lemma~\ref{O:A:lem:power}, the argument has $v_2$ at least
$R-t+(p-1)t_2$.  Its trace, viewed in $B_1$, has valuation at least
$R-t+2(p-1)t_2\ge R+pt_2$.
Apply Lemmas~\ref{O:A:lem:translate} and \ref{O:A:lem:filtered}.
Since $p$ is odd and $s=-e_{p-1}$, $1+e_{p-1}^p=1-s^p$.
\end{proof}

\subsection{The linear defect without translated-root expansions}
Put
\[
 g(z)=T_1(bS_1z)-T_2(aS_2z).
\]
Expand $b=\sum_{l=0}^{p-1}Y_l\Delta^l$ over $B_2$ and
$s=\sum s_l\Delta^l$ over $B_2$.
\begin{lemma}\label{O:A:lem:bcoeff}
The coefficients satisfy
\begin{align*}
 v_2(Y_l)&\ge(p-1)t_2-(p-l)t&& (2\le l<p),\\
 v_2(Y_0-a)&\ge(p-1)t_2-t,&
 Y_1&=1-s_0+E,\quad v_2E\ge(p-1)t_2.
\end{align*}
\end{lemma}
\begin{proof}
Use $b=a+\Delta+\sum_{i=1}^{p-1}(-1)^ie_i\Delta^{p-i}$.
Expand each $e_i$ over $B_2$ by Lemma~\ref{O:A:lem:coeff}; its $\Delta^k$
coefficient has valuation at least $(p-1)t_2-it+kt$.
If $k+p-i\ge p$, reduce once using $\Delta^p=\Delta+a$.
The unreduced contribution to $Y_l$ has valuation
$(p-1)t_2-(p-l)t$; each reduced contribution has valuation at least
$(p-1)t_2+(l-1)t$.
For $l=0$ only the terms with $k=i$ contribute, with a factor $a$.
For $l=1$ the sole low term is the constant coefficient of $e_{p-1}$,
namely $-s_0$; all others have valuation at least $(p-1)t_2$.
\end{proof}

\begin{proposition}\label{O:A:prop:g}
Suppose $R=v_B(x\Delta^j)\ge1+\delta$.
If $j\ne p-2$, then $g(x\Delta^j)\in\pp_A^{1+t_2}$.
For $j=p-2$ there is a fixed $\omega\in B_2$ such that
\[
 g(x\Delta^{p-2})=T_2(x(\omega-1)),\qquad
 x\omega-xs\in\PP_B^{K_0}.
\]
\end{proposition}
\begin{proof}
For $1\le j<p$, the exact moment identities give
\[
 g(x\Delta^j)=T_2\left(x\left((p-1)Y_{p-1-j}+paY_{p-j}
              +(p-1)\ind_{j=p-1}(Y_{p-1}-a)\right)\right).
\]
For $j\le p-3$, the first term has $v_2$ at least
$v_2x+(p-1)t_2-(j+1)t\ge1+t_2$; here $v_2x\ge1$ and
$(p-2)t_2-(j+1)t\ge0$.
The $pa$ term is at least as deep by \eqref{O:A:eq:pbound} and
Lemma~\ref{O:A:lem:bcoeff}.
For $j=p-1$ combine $Y_0-a$ with the other displayed terms and use
$v_2x+(p-1)t_2-t\ge1+t_2$.
These elements are in the ideal whose trace is $\pp_A^{1+t_2}$.

For $j=0$, trace transitivity gives
\[
 g(x)=T_2(x)T_1(b)-pT_2(ax).
\]
Replace $T_1b$ by $T_2a$ using Theorem~\ref{O:A:thm:H14}.
The two resulting terms have valuations bounded below by
$(v_2x-t+2(p-1)t_2)/p>t_2$, using \eqref{O:A:eq:trace} and \eqref{O:A:eq:pbound}.

For $j=p-2$ define
$\omega=1+(p-1)Y_1+paY_2$.
The formula for $g$ is exact.  Lemma~\ref{O:A:lem:bcoeff} gives
\[
 \omega-s=p(1-s_0)+(s_0-s)+(p-1)E+paY_2.
\]
Every term except $s_0-s$ has sufficient valuation immediately after
multiplication by $x$.  For the latter, Lemma~\ref{O:A:lem:coeff} gives
\[
 v_B(s-s_0)\ge p(p-1)\delta+(p-1)t.
\]
Since $v_Bx=R+(p-2)t$, the resulting valuation exceeds or equals $K_0$;
the difference of the two lower bounds is
$(p-3)t+(p-1)^2\delta\ge0$.
\end{proof}

\subsection{A floor-corrected projection-to-norm lemma}
\begin{lemma}\label{O:A:lem:H17}
Let $x\in B_2$, $w\in B_1$, and $y\in B_2$ satisfy
\[
 xw-y\in\PP_B^{K_0},\qquad
 h_0=v_B(xw)\ge1+t_2+(p-3)t.
\]
Then
\[
 N_1(xw)-n_2(y)\in\PP_{B_1}^{K_0}.
\]
\end{lemma}
\begin{proof}
The zero case is immediate, so write $z=xw\ne0$ and expand
$z=\sum_{i=0}^{p-1}y_i\Delta^i$, with $y_i=xc_i\in B_2$.
Lemma~\ref{O:A:lem:coeff} gives
\[
 v_B(y_i\Delta^i)\ge h_0+(p-1)it.
\]
Thus $v_B(y_0)=h_0$, and every nonconstant term is strictly deeper.
The valuations of the nonzero terms in
\[
 z-y=(y_0-y)+\sum_{i=1}^{p-1}y_i\Delta^i
\]
have distinct residues modulo $p$, because $p\nmid t$.
Since their sum has valuation at least $K_0$, each term does.
In particular $v_B(y_i\Delta^i)\ge K_0$ for $i\ge1$ and
$v_2(y_0-y)\ge t_2+1$.

First keep the linear term and put $z_\ell=y_0+y_1\Delta$.
Norm multiplicativity, not scalar homogeneity, gives
\[
 N_1(z_\ell)=n_2(y_0)N_1(1+(y_1/y_0)\Delta).
\]
We have $v_1(n_2y_0)=h_0$ and $v_2(y_1/y_0)\ge t$.
The twisted symmetric estimate therefore bounds the intermediate
terms in $N_1(z_\ell)-n_2y_0$ by
\[
 h_0+(p-1)t_2+j\bigl(v_2(y_1/y_0)-t\bigr)
 \ge h_0+(p-1)t_2\ge K_0.
\]
The endpoint is $N_1(y_1\Delta)$ and has $v_1$ equal to
$v_B(y_1\Delta)\ge K_0$. Hence
$N_1(z_\ell)-n_2y_0\in\PP_{B_1}^{K_0}$.

Now put $U=\sum_{i=2}^{p-1}y_i\Delta^i$. Then
$v_B(U/z)\ge2(p-1)t$ and $v_BU\ge K_0$. The exact identity
\[
 N_1z-N_1z_\ell=N_1z\{1-N_1(1-U/z)\}
\]
is expanded on the cyclic edge $B/B_1$, whose different exponent is $D'$.
For each intermediate coefficient the strong symmetric bound gives
\[
 v_1\bigl(N_1z\,E_j(U/z)\bigr)
 \ge h_0+\floor{[2j(p-1)t+D']/p}
 \ge K_0+\floor{[(p-3)t+p-1]/p}\ge K_0.
\]
The minimum occurs at $j=1$, and the second inequality follows by inserting
$h_0\ge1+t_2+(p-3)t$ and $t'=t+p\delta$.
The endpoint is $N_1U$, which has valuation $v_BU\ge K_0$.
Thus $N_1z-N_1z_\ell\in\PP_{B_1}^{K_0}$.

Finally $y_0,y$ are integral and $v_2(y_0-y)\ge t_2+1$.
To check the last norm comparison explicitly, set $a=v_2y$.
If $a\ge t_2+1$ (or $y=0$), both norms have valuation at least $t_2+1$.
Otherwise factor out $y$, apply the strong symmetric estimate to
$(y_0-y)/y$, and note that its shallowest intermediate term has valuation
at least
\[
 a+\floor{[t_2+1-a+(p-1)(t_2+1)]/p}
 =t_2+1+\floor{(p-1)a/p}\ge t_2+1.
\]
The endpoint has valuation $v_2(y_0-y)\ge t_2+1$.
Consequently $n_2y_0-n_2y\in\pp_A^{t_2+1}$, which lies in
$\PP_{B_1}^{p(t_2+1)}\subseteq\PP_{B_1}^{K_0}$.
Combining the three comparisons proves the assertion.
\end{proof}


\section{Stationary models for totally ramified odd-degree extensions}\label{sec:odd-models}
\subsection{Construction and realization of minimal character models}\label{O:sec:models}
\subsubsection*{The local data}
Let $A\subset B_1,B_2\subset B$ be totally ramified, $\Gal(B/A)=C_p^2$,
$p$ odd. Write
\[
 t_1=t,\quad t_2=t+\delta,\quad t'=t+p\delta,\quad p\nmid t,
 \quad \kappa_2=t_2+1,\quad T'=t'+1.
\]
The upper breaks are $t'$ over $B_1$ and $t$ over $B_2$. Put
$N_i=\Nm_{B/B_i}$, $n_i=\Nm_{B_i/A}$, $S_i=\Tr_{B/B_i}$,
$T_i=\Tr_{B_i/A}$. Normalize each valuation on its own field.
Section~\ref{O:sec:as} constructs $\Delta\in B$ with
\[
 \Delta^p-\Delta=a\in B_2,\quad v_B\Delta=-t,\quad
 N_2\Delta=a,\quad b=N_1\Delta,\quad s=-E_{p-1}^{B/B_1}(\Delta).
\]
It proves, independently of the SML, the following exact inputs:
\begin{align}
 T_1b-T_2a&\in\pp_A^{\kappa_2}, &v_B(b-a)&\ge-t,                 \label{O:M:H14}\\
 v_1s&\ge(p-1)\delta, &v_B(p)&\ge p(p-1)t_2.              \label{O:M:basic}
\end{align}
For $w=\sum x_j\Delta^j$, $x_j\in B_2$, the term valuations are distinct
modulo $p$. For a coefficient $x\in B_2$ set
\[
 g(z)=T_1(bS_1z)-T_2(aS_2z),\quad
 h(z)=T_1N_1(\Delta z)-T_2N_2(\Delta z),\quad R=v_B(x\Delta^j).
\]
If $R\ge1+\delta$, $g(x\Delta^j)$ is in $\pp_A^{\kappa_2}$ except at $j=p-2$.
There is $\omega\in B_2$, independent of $x$, such that in that exception
\begin{equation}\label{O:M:ginput}
 g(x\Delta^{p-2})=T_2(x(\omega-1)),\qquad
 x\omega-xs\in\PP_B^{1+pt_2}.
\end{equation}
For $R\ge1$, the weighted norm defect is
\begin{equation}\label{O:M:hinput}
 h(x\Delta^j)\equiv(p-1)\mathbf1_{j=p-2}n_2x(1-s^p)
       \pmod{\PP_{B_1}^{pt_2+R}}.
\end{equation}
Finally, if $x\in B_2,w\in B_1,y\in B_2$,
$v_B(xw)\ge1+t_2+(p-3)t$, and $xw-y\in\PP_B^{1+pt_2}$, then
\begin{equation}\label{O:M:transferinput}
 n_2x\,w^p\equiv n_2y\pmod{\PP_{B_1}^{1+pt_2}}.
\end{equation}
These are the individual lemmas proved in Section~\ref{O:sec:as}.
All their proofs persist after an unramified extension of $A$.

Choose an actual nontrivial norm character $\tau$ for $B_2/A$.
Let $\mathrm e_L=\mathrm e_A\Tr_{L/A}$, and choose $\alpha\in A^\times$ so
\[
 \tau(1-z)=\mathrm e_A(\alpha P(z)),\qquad
 P(X)=\sum_{j=1}^{p-1}X^j/j,\quad z\in\pp_A^c,\quad c=\ceil{\kappa_2/p}.
\]
Put $\Psi_L(x)=\mathrm e_L(\alpha x)$. Their largest trivial ideals are
\[
 \pp_A^{\kappa_2},\quad\PP_{B_1}^{T'},\quad\PP_{B_2}^{\kappa_2},\quad\PP_B^{T'}.
\]
The full additive subgroup conversion proved in Corollary~\ref{O:I:conversion} gives
\begin{equation}\label{O:M:conversion}
 \Psi_A(T_2v+n_2v)=1\qquad(v\in\PP_{B_2}^c).
\end{equation}
Only Ueda's proved cyclic norm filtration, trace/different estimates,
conductor formula and local duality, plus the explicit proofs in
Sections~\ref{O:sec:inertia} and \ref{O:sec:as}, are used below. In equal characteristic expressions
containing $p$ are zero; their valuation is interpreted as $+\infty$.

\subsubsection*{The exact unit domains}
Set
\[
 m_1=1+t+t',\quad m_2=1+t+t_2,\quad
 q_i=\ceil{m_i/p},\quad q=q_1=\delta+\ceil{(2t+1)/p},
 \quad H_i=U_{B_i}^{q_i}.
\]
Define characters on these groups by
\begin{equation}\label{O:M:models}
 R_1(1-z)=\Psi_1(bPz),\qquad R_2(1-z)=\Psi_2(aPz).
\end{equation}
They are characters because $pq_i\ge m_i$ and the truncated logarithm
is additive for the multiplicative group modulo $\PP_{B_i}^{m_i}$.
They have exact conductors $m_i$ since $v_i a=v_i b=-t$ on the indicated side.

\begin{lemma}\label{O:M:domains}
One has $q_i\ge1$, $q\le t'$, $q_2\le t_2$ and
\[
 q\ge q_2,\qquad
 \floor{[q+(p-1)(t+1)]/p}\ge q_2.
\]
Consequently
\begin{equation}\label{O:M:preimage}
 N_1^{-1}(H_1)=U_B^q,\qquad N_2(U_B^q)\subset H_2.
\end{equation}
For any unramified base extension the same assertions hold.
\end{lemma}
\begin{proof}
The inequalities follow by writing $q=\delta+\lceil(2t+1)/p\rceil$ and
$q_2=\lceil(2t+\delta+1)/p\rceil$. For the floor inequality it is enough
that $q+(p-1)(t+1)\ge pq_2$; use $pq_2\le2t+\delta+p$ and
$(p-3)t+\lceil(2t+1)/p\rceil-1\ge0$.
The inequalities $q\le t'$ and $q_2\le t_2$ follow from
$2t+1\le pt$ for $p\ge3,t\ge1$ and $\delta\ge0$.
Below the upper break $t'$, the norm preserves the first nonzero unit
depth; at the critical depth its image is still in that depth.
Thus a norm in $H_1$ has a preimage of depth at least $q$, including
the case $q=t'$. The converse follows from the norm filtration.
The second assertion follows coefficientwise from the strong symmetric
bound, using the two displayed inequalities. Norm valuations first
exclude elements outside $U_B^1$ from the first preimage assertion.
\end{proof}

\subsubsection*{The monomial trace--norm phase}
\begin{lemma}\label{O:M:monophase}
For $x\in B_2$, $0\le j<p$, $v_B(x\Delta^j)\ge q$, one has
\[
 \Psi_A\bigl(g(x\Delta^j)+h(x\Delta^j)\bigr)=1.
\]
\end{lemma}
\begin{proof}
We have $q\ge1+\delta$. Outside $j=p-2$ this follows at once from
\eqref{O:M:ginput}--\eqref{O:M:hinput}: the norm defect belongs to $A$ and to
$\PP_{B_1}^{1+pt_2}$, hence to $\pp_A^{\kappa_2}$.
For $j=p-2$, $pv_2x\ge q+(p-2)t\ge \kappa_2$, so $v_2x\ge c$.
Put $K_0=1+pt_2$. The difference between $p-1$ and $-1$ in
\eqref{O:M:hinput} is in $\PP_{B_1}^{K_0}$ because
\[
 v_1p+pv_2x\ge(p-1)t_2+q+(p-2)t\ge K_0,
\]
and $s$ is integral. Also
\[
 v_B(xs)\ge q+(p-2)t+p(p-1)\delta
       \ge1+t_2+(p-3)t.
\]
Using \eqref{O:M:ginput} and \eqref{O:M:transferinput} gives
$n_2x\,s^p\equiv n_2(x\omega)\pmod{\PP_{B_1}^{K_0}}$.
The approximation $x\omega\equiv xs\pmod{\PP_B^{K_0}}$ also implies
$v_2(x\omega)\ge c$. Thus the sum in question is congruent, now between
\emph{elements of $A$}, to
\[
 -T_2x-n_2x+T_2(x\omega)+n_2(x\omega)\pmod{\pp_A^{\kappa_2}}.
\]
Both pairs have trivial phase by \eqref{O:M:conversion}.
\end{proof}

\subsubsection*{Finite-field rigidity under unramified extension}
\begin{lemma}\label{O:M:rigidity}
Let $M$ be a local field of residue characteristic $p$, and let $\Psi_M$
be a continuous additive character. Write
$f(Z)=\sum_{n=1}^{D}a_nZ^n\in M[Z]$, where $D<p^2$, and suppose
$x\mapsto\Psi_M(a_nx)$ has order dividing $p$ on $\OO_M$ for each $n$.
Let $M'/M$ be unramified of degree prime to $p$, with absolute residue
degree at least $3$, and set $\Psi_{M'}=\Psi_M\circ\Tr_{M'/M}$.
For $\lambda\in k_{M'}$, let $[\lambda]$ be its Teichmuller lift
(the coefficient-field lift in equal characteristic). Suppose
\[
 \lambda\longmapsto\Psi_{M'}(f([\lambda]))
\]
is additive on $k_{M'}$. Then
\[
 \Psi_M(f(1))=\Psi_M(a_1+a_p),
\]
where missing coefficients are zero.
\end{lemma}
\begin{proof}
A Teichmuller carry belongs to $p\OO_{M'}$, and hence is killed by
each coefficient character. Therefore these characters on Teichmuller
lifts are residue additive characters, of the form
$\psi_{\mathbf F_p}\Tr_{k_{M'}/\mathbf F_p}(\beta_n\lambda)$;
unramified trace shows that each $\beta_n$ is the scalar extension of its
base coefficient. The displayed function is consequently
\[
 \psi_{\mathbf F_p}\Tr_{k_{M'}/\mathbf F_p}
       \Bigl(\sum_n\beta_n\lambda^n\Bigr).
\]
Under the absolute trace, exponents in the same Frobenius orbit are
identified. Among positive exponents less than $p^2$, with at least three
base-$p$ digits available in the residue field, the only identifications
are $n$ with $pn$ when both occur. This follows by rotating the base-$p$
digits modulo $|k_{M'}|-1$: a two-nonzero-digit exponent cannot return below
$p^2$ except at itself. Distinct polynomial monomials of degrees less than
$|k_{M'}|$ are linearly independent as functions. An additive function can
therefore retain only the orbit of the exponents $1,p$; every other orbit
has zero trace contribution. Evaluating at $1$ proves the assertion over
$M'$. Both sides are $p$th roots of unity and the unramified degree is prime
to $p$, so the equality descends to $M$.
\end{proof}

\subsubsection*{Compatibility on the full unit subgroup}
Define on $U_B^q$
\[
 \Xi(u)=R_1(N_1u)R_2(N_2u)^{-1}.
\]
It is an actual multiplicative character. For $z\in\PP_B^q$ put
\begin{equation}\label{O:M:poly}
 F_z(Z)=T_1\bigl(bP(1-N_1(1-Zz))\bigr)
       -T_2\bigl(aP(1-N_2(1-Zz))\bigr).
\end{equation}
This is an exact polynomial of degree at most $p(p-1)<p^2$ and
$\Xi(1-Zz)=\Psi_A(F_z(Z))$ for integral base scalars $Z$.

\begin{lemma}\label{O:M:polycoeff}
Every coefficient of $F_z$ defines an additive character of order dividing
$p$ on $\OO_A$. Its degree-one coefficient is $g(z)$; its degree-$p$
coefficient is $h(z)$ modulo $\pp_A^{\kappa_2}$.
\end{lemma}
\begin{proof}
The coefficients of $1-N_i(1-Zz)$ belong to $\PP_{B_i}^{q_i}$,
by the coefficientwise proof of Lemma~\ref{O:M:domains}. Products of these
coefficients have at least that valuation. Multiplication by $b$ or $a$
loses $t$. In mixed characteristic $v_i p\ge(p-1)t_2$, and
\[
 (p-1)t_2-t+q\ge T',\qquad
 (p-1)t_2-t+q_2\ge \kappa_2.
\]
For the first inequality the difference is
$(p-3)t+\lceil(2t+1)/p\rceil-1$; the second follows from
$(p-3)t+(p-2)\delta+q_2-1\ge0$.
Thus $p$ kills the phase of every coefficient, including after any integral
base scalar. In equal characteristic this assertion is automatic.

The linear coefficient follows by differentiation at $Z=0$.
At degree $p$, the norm endpoint gives $h(z)$. All other terms contain at
least two intermediate elementary coefficients of total weight $p$.
If the upper break is $u$, their valuation after multiplication by $b$ or
$a$ is at least
\[
 q-t+2(p-1)u/p.
\]
For $u=t$ this is strictly greater than $\kappa_2-1$; for $u=t'$ it is
strictly greater than $T'-1$, using
$q=\delta+\lceil(2t+1)/p\rceil$. Integral valuations therefore put
each term in the corresponding trivial ideal. Trace to $A$ puts it in
$\pp_A^{\kappa_2}$. This argument concerns \emph{degree $p$ only}; no assertion
that the higher coefficients are individually deep is made.
\end{proof}

\begin{theorem}[full compatibility]\label{O:M:compat}
The characters \eqref{O:M:models} satisfy $R_1N_1=R_2N_2$ on $U_B^q$.
This holds for all odd $p$, in mixed and equal characteristic.
\end{theorem}
\begin{proof}
First $\Xi$ is trivial on $U_B^{m_1}$. At that depth the norm-logarithm
congruence of Lemma~\ref{O:I:normlog} applies on both upper edges, with target moduli
$\PP_{B_i}^{m_i}$. Indeed its sufficient input depths are
$t+\lceil(t'+1)/p\rceil$ for the first edge and
$t_2+\lceil(t+1)/p\rceil$ for the second, both at most $m_1$.
The two weighted norm endpoints are killed after losing $t$.
The remaining difference is
$\Psi_B((b-a)Pz)=1$ by \eqref{O:M:H14} and $m_1-t=T'$.

Proceed by descending induction, uniformly over all unramified base
extensions. If triviality on $U_B^M$, $M\le m_1$, is known, put
$s_0=\max(q,\lceil M/2\rceil)$. On $\PP_B^{s_0}$ the map
$z\mapsto\Xi(1-z)$ is additive, since every multiplicative discrepancy
has depth at least $2s_0\ge M$. In mixed characteristic a Teichmuller carry
multiplied by $z$ has depth at least $v_Bp+s_0\ge M$; the carry is zero
in equal characteristic. For each fixed $z$ in this ideal, choose an
unramified extension of degree prime to $p$ and absolute residue degree at
least $3$. Apply Lemmas~\ref{O:M:rigidity} and \ref{O:M:polycoeff} to obtain
\[
 \Xi(1-z)=\Psi_A(g(z)+h(z)).
\]
Expand $z=\sum x_j\Delta^j$. Valuation separation puts every term in
$\PP_B^{s_0}$. Additivity of the \emph{actual character} on this ideal,
and Lemma~\ref{O:M:monophase} for each term, give $\Xi(1-z)=1$.
This proves triviality on $U_B^{s_0}$. Repeated halving reaches $q$.
All conductors and calculations are stable under unramified extension,
so the simultaneous induction and prime-to-$p$ descent are legitimate.
\end{proof}

\subsubsection*{Extension to the full multiplicative groups}
\begin{theorem}\label{O:M:global}
Every extension $\theta_2$ of $R_2$ to $B_2^\times$ has a compatible
extension $\theta_1$ of $R_1$ to $B_1^\times$. Both characters have their
exact minimal conductors $m_i$, and their common pullback is primitive.
\end{theorem}
\begin{proof}
Such $\theta_2$ exists by extension of a character of an open subgroup of
an abelian group. Put $\Theta=\theta_2N_2$.
Let $\rho$ generate $\Gal(B/B_1)$, whose break is $t'$.
For \emph{every} $z\in B^\times$, not just a deep unit, the ramification
inequality gives $\rho(z)/z\in U_B^{t'}\subset U_B^q$.
Its $N_1$-norm is $1$. By Theorem~\ref{O:M:compat} and \eqref{O:M:preimage},
\[
 \Theta(\rho(z)/z)=R_2(N_2(\rho(z)/z))
                 =R_1(N_1(\rho(z)/z))=1.
\]
Hence $\Theta$ is invariant under $\Gal(B/B_1)$ on all of $B^\times$.
Hilbert 90 makes it trivial on $\ker N_1$, so it descends to the open norm
image and extends to a character $\widetilde\theta_1$ of $B_1^\times$.
On $H_1\cap N_1B^\times$ it agrees with $R_1$: any norm preimage belongs to
$U_B^q$ by \eqref{O:M:preimage}, where full compatibility holds.
Thus $\widetilde\theta_1/R_1$ on $H_1$ is a character of the image of $H_1$
in $B_1^\times/N_1B^\times$. Extend it across this finite cyclic quotient
and divide $\widetilde\theta_1$ by the resulting norm character. This gives
$\theta_1|_{H_1}=R_1$, without changing the common pullback.

The exact conductors are already visible in $R_i$ at depth $m_i-1$.
For a generator $\sigma$ of $\Gal(B_2/A)$ and $z\in\PP_{B_2}^t$,
$\sigma^{-1}(1-z)/(1-z)$ has leading depth $t+t_2=m_2-1$ and nonzero
ramification coefficient $t$. Because $p\nmid t$, its leading map on the
residue line is surjective. Therefore
$\theta_2^\sigma/\theta_2$ is nontrivial, of conductor $t+1$.
Global invariance of $\Theta$ makes this a norm character for $B/B_2$.
If $\Theta$ descended from $A$, this commutator would be trivial, a
contradiction. Here the action of the lower group of order $p$ on the
upper norm quotient of order $p$ is trivial, since
$|\operatorname{Aut}(C_p)|=p-1$.
\end{proof}

\subsubsection*{Realizing a prescribed minimal pair}
Let $(\varphi_1,\varphi_2)$ be an actual primitive compatible pair with
conductors $m_i$. Choose any initial pair from Theorem~\ref{O:M:global}.
Their second commutators are generators of $S(B/B_2)$. Choose
$j\in\{1,\ldots,p-1\}$ so that the commutator of $\varphi_2$ equals that of
$\theta_2^j$. Then $\varphi_2\theta_2^{-j}$ is invariant under
$\Gal(B_2/A)$. Hilbert 90 and character extension give a character $\lambda$
of $A^\times$ with the \emph{global} identity
\begin{equation}\label{O:M:actual2}
 \varphi_2=\theta_2^j(\lambda n_2).
\end{equation}
Put
\[
 h_0=\floor{t/p},\qquad c=\ceil{(t_2+1)/p},\qquad q_0=h_0+c.
\]
The cyclic conductor formula and $a(\lambda n_2)\le m_2$ imply
$a(\lambda)\le \kappa_2+h_0$.
The following inequalities include the case $p\mid t_2$:
\begin{equation}\label{O:M:q0}
 q_0\le q_2,\quad pq_0\ge \kappa_2+h_0,\quad
 n_2(H_2)\subset U_A^{q_0},\quad
 (p-2)t_2-ph_0+q_0-1\ge0.
\end{equation}
To verify them, write $t=ph_0+b_0$, $1\le b_0<p$.
Then $q_2=h_0+\lceil(b_0+t_2+1)/p\rceil\ge q_0$.
For the norm inclusion use the strong trace bound and
$q_2+(p-1)\kappa_2\ge pq_0$; the latter follows from
$\kappa_2\ge ph_0+b_0+1$ and the preceding ceiling formula.
The last inequality follows by replacing $t_2$ by $t$ and $q_0$ by
$h_0+\lceil(t+1)/p\rceil$; the resulting lower bound is nonnegative.

Put $\Psi'_L=\Psi_L(j\,\cdot)$. The truncated-log chart gives a coefficient
$\gamma\in A$, $v_A\gamma\ge-h_0$, such that
\[
 \lambda(1-z)=\Psi'_A(\gamma Pz)\qquad(z\in\pp_A^{q_0}).
\]
If the restriction is trivial take $\gamma=0$ and $w=0$.
Otherwise choose an \emph{approximate} norm representative by Ueda's
ordinary subcritical norm filtration:
\begin{equation}\label{O:M:approx}
 w\in B_2,\quad k=v_2w=v_A\gamma\ge-h_0,\quad
 \gamma-n_2w\in\pp_A^{k+t_2}.
\end{equation}

\begin{lemma}\label{O:M:lambdaformula}
For $x\in\PP_{B_2}^{q_2}$,
\[
 (\lambda n_2)(1-x)=\Psi'_2((w^p-w)P(x)).
\]
\end{lemma}
\begin{proof}
Set $u=1-n_2(1-x)$. It has depth at least $q_0$.
The error in replacing $\gamma$ by $n_2w$ contributes valuation at least
$k+t_2+q_0\ge \kappa_2$, so \eqref{O:M:approx} suffices.
The norm-logarithm congruence from Lemma~\ref{O:I:normlog}, now with target $\kappa_2+h_0$ and
input depth $h_0+c=q_0\le q_2$, yields
\[
 (\lambda n_2)(1-x)
 =\Psi'_2((n_2w)P(x))\Psi'_A(n_2(wP(x))).
\]
Since $k+q_2\ge c$, the norm-character conversion changes the second
factor to $\Psi'_2(-wP(x))$.
The strong symmetric estimate applied to the characteristic polynomial of
$w$ gives $v_2(w^p-n_2w)\ge pk+(p-1)t_2$. Multiplication by $P(x)$
puts this error in $\PP_{B_2}^{\kappa_2}$ by \eqref{O:M:q0}. This proves the formula.
\end{proof}

\begin{theorem}[actual realization, with the allowed conjugate explicit]
\label{O:M:realize}
There exist an actual norm-character coefficient $\alpha'=j\alpha$ and
$\Delta_\lambda\in B$ such that
\[
 \Delta_\lambda^p-\Delta_\lambda=a_\lambda\in B_2,\qquad
 v_B\Delta_\lambda=-t,
\]
for which the second full model is literally $\varphi_2|_{H_2}$ and the
first full model is the restriction of a Galois conjugate of $\varphi_1$.
Consequently both local constants are unchanged by passing to these full
models. No norm equation beyond \eqref{O:M:approx} is assumed.
\end{theorem}
\begin{proof}
Set $a_\lambda=a+w^p-w$. Its valuation is $-t$ because
$pv_2w\ge-ph_0=-t+b_0>-t$. In mixed characteristic the error of
$\Delta+w$ in $X^p-X-a_\lambda$ has valuation at least
\[
 v_B(p)-(p-1)t+pv_2w\ge v_B(p)-pt+b_0>0.
\]
Its derivative is a unit. Hensel's lemma gives the required exact root
near $\Delta+w$, still of valuation $-t$. In equal characteristic
$\Delta+w$ is already that root. It generates $B/B_2$ because $p\nmid t$.
By \eqref{O:M:actual2} and Lemma~\ref{O:M:lambdaformula}, the new second model
agrees with the actual $\varphi_2$ on \emph{all} of $H_2$.
Apply Theorem~\ref{O:M:global} to the new generator and norm character, choosing
$\varphi_2$ itself as the second extension. It supplies a first character
$\widetilde\varphi_1$ with the new full model and the same actual norm
pullback. Thus $\widetilde\varphi_1/\varphi_1\in S(B/B_1)$.
For a primitive compatible pair its nontrivial conjugate quotient generates
this order-$p$ group, on which the lower Galois group acts trivially.
Hence every such twist is a Galois conjugate. Changing variables in the
one-dimensional local integral proves equality of their local constants.
This is a character-theoretic conjugacy argument, not an SML assertion.
\end{proof}


\subsection{The second terminal congruence}\label{O:sec:r2}
\subsubsection*{Normalization}
Use $A\subset B_1,B_2\subset B$, $t_1=t$, $t_2=t+\delta$, $t'=t+p\delta$,
$p\nmid t$, and $N_i,n_i,S_i,T_i$ as in \S\ref{O:sec:as} and \S\ref{O:sec:models}. Put
\[
 T'=t'+1,\quad \kappa_2=t_2+1,\quad m_1=T'+t,\quad q_1=\ceil{m_1/p},
 \quad c_t=\ceil{(t+1)/p}=\ceil{t/p},\quad q=\delta+c_t.
\]
Let $\Delta^p-\Delta=a\in B_2$, $v_B\Delta=-t$, and
\[
 b=N_1\Delta,\quad a=N_2\Delta,\quad
 s=-E_{p-1}^{B/B_1}(\Delta),\quad H=n_1s.
\]
Choose actual norm characters $\tau_2$ for $B_2/A$ and $\tau_1$ for $B_1/A$,
with full truncated-logarithm coefficients $\alpha,\beta$ relative to the
same additive character $\mathrm e_A$. Write
\[
 \Psi_L(z)=\mathrm e_L(\alpha z),\qquad \varrho=\beta/\alpha,
 \qquad v_A \varrho=\delta,
\]
so $\Psi_A,\Psi_1,\Psi_2,\Psi_B$ have largest trivial ideals
$\pp_A^{\kappa_2},\PP_{B_1}^{T'},\PP_{B_2}^{\kappa_2},\PP_B^{T'}$.
The actual first minimal character $\chi$ is taken in the conjugacy class
supplied by Subsection~\ref{O:sec:models}, so that
\begin{equation}\label{O:R:model}
 \chi(1-y)=\Psi_1(bP(y))\qquad(y\in\PP_{B_1}^{q_1}),\qquad
 P(X)=\sum_{j=1}^{p-1}X^j/j.
\end{equation}
Its exact conductor is $m_1$. Let $\sigma$ be the restriction to $B_1$ of the
upper automorphism of $B/B_2$ carrying $\Delta$ to the root near $\Delta+1$.
Our character convention is $\chi^\sigma(u)=\chi(\sigma^{-1}u)$.

\subsubsection*{A shallow commutator evaluated by a deeper character}
\begin{lemma}\label{O:R:comm-log}
For $x\in\PP_{B_1}^{q}$,
\[
 \frac{\chi^\sigma(1-x)}{\chi(1-x)}
       =\Psi_1((\sigma b-b)P(x)).
\]
\end{lemma}
\begin{proof}
Put $d=\sigma^{-1}x-x$, so $v_1d\ge q+t$.
The quotient $(1-\sigma^{-1}x)/(1-x)$ is $1-d/(1-x)$.
We have $q+t\ge q_1$, hence \eqref{O:R:model} applies to this quotient even
though it need not apply separately to $1-x$.
As an integral formal power series in $x,d$,
\[
 P(d/(1-x))-P(x+d)+P(x)
\]
is divisible by $d$ and has total degree at least $p$. Its denominators
are products of integers $1,\ldots,p-1$ and powers of $1-x$.
Its valuation is therefore at least $(p-1)q+(q+t)=pq+t\ge T'+t=m_1$.
Multiplication by $b$ puts it in the kernel of $\Psi_1$.
Finally trace invariance gives
$\Psi_1(bP(\sigma^{-1}x))=\Psi_1((\sigma b)P(x))$.
\end{proof}

\begin{lemma}\label{O:R:shift}
Set
\[
 R=T'-q=t+1-c_t+(p-1)\delta.
\]
Then $\sigma b-b\equiv1-s\pmod{\PP_{B_1}^R}$.
\end{lemma}
\begin{proof}
Norm expansion gives
\[
 N_1(\Delta+1)-b=1-s+\sum_{i=1}^{p-2}E_i^{B/B_1}(\Delta).
\]
Each summand has valuation at least
$(p-1)t_2-(p-2)t=t+(p-1)\delta\ge R$ by Section~\ref{O:sec:as}.
For the actual translated root write
$\sigma\Delta=(\Delta+1)(1+u)$.
The elementary Hensel estimate gives
$v_Bu\ge V-(p-2)t$, where $V=v_Bp$.
For intermediate coefficients in $N_1(1+u)-1$ the strong symmetric bound
is smallest at the trace coefficient, and after multiplying by
$N_1(\Delta+1)$ it gives at least
\[
 -t+\floor{[V-(p-2)t+(p-1)(t'+1)]/p}
 =V/p-t+(p-1)\delta+c_t
 \ge(p-2)t+2(p-1)\delta+c_t\ge R.
\]
The norm endpoint has still larger valuation. In equal characteristic
$u=0$. This proves the assertion without dropping any mixed-characteristic
$p=3$ term.
\end{proof}

Thus the actual commutator $\mu=\chi^\sigma/\chi$ satisfies
\begin{equation}\label{O:R:S}
 \mu(1-x)=\Psi_1((1-s)P(x))\qquad(x\in\PP_{B_1}^{q}).
\end{equation}
It is a nontrivial norm character for $B/B_1$, of conductor $T'$.
The same is true of $\tau_2n_1$. We next identify the generator, rather than
assuming that these two nontrivial characters coincide.

\subsubsection*{The second expression uses only an approximate norm}
By the ordinary subcritical norm filtration, choose $c\in B_1^\times$ with
\begin{equation}\label{O:R:approx}
 n_1c=\varrho\varepsilon,\qquad\varepsilon\in U_A^t,\qquad v_1c=\delta.
\end{equation}
\begin{lemma}\label{O:R:T}
On $U_{B_1}^{q}$ one has
\[
 (\tau_2n_1)(1-x)=\Psi_1((1-\varrho/c)P(x)).
\]
\end{lemma}
\begin{proof}
For $u=1-n_1(1-x)$, the norm filtration gives
$v_Au\ge\lceil \kappa_2/p\rceil$, so the defining full formula for $\tau_2$
applies. The norm-logarithm congruence of Lemma~\ref{O:I:normlog} on $B_1/A$, with target
$\kappa_2$ and input depth $\delta+c_t=q$, gives
\[
 (\tau_2n_1)(1-x)=\Psi_1(P(x))\Psi_A(n_1P(x)).
\]
The error in replacing the second factor by
$\mathrm e_A(\beta n_1(P(x)/c))$ has valuation, before multiplication by
$\alpha$, at least $q+t\ge \kappa_2$. This is exactly the error
$n_1P(x)(\varepsilon^{-1}-1)$, not an exact norm equation.
Since $v_1(P(x)/c)\ge q-\delta=c_t$, the norm-character conversion for
$\tau_1$ makes this factor $\mathrm e_{B_1}(-\beta P(x)/c)$.
\end{proof}

\subsubsection*{Exact generator normalization in the one-break case}
\begin{lemma}[reciprocal trace identity]\label{O:R:reciprocal}
If $\delta=0$, then for every $Y\in\OO_A$,
\[
 \Psi_1(sY/b)=\Psi_2(Y/a)=\Psi_A(-n_1(Y/b)).
\]
\end{lemma}
\begin{proof}
Elementary symmetric functions give the \emph{exact} identities
\[
 s/b=-S_1(1/\Delta),\qquad S_2(1/\Delta)=-1/a.
\]
The second uses the coefficient $-1$ of $X$ in $X^p-X-a$.
Trace transitivity therefore gives the first equality.
As $v_2(Y/a)\ge t\ge\lceil(t+1)/p\rceil$, the actual norm character
$\tau_2$ gives $\Psi_2(Y/a)=\Psi_A(-n_2(Y/a))$.
Finally $n_1b=n_2a$ by norm transitivity and $Y\in A$, so
$n_2(Y/a)=Y^p/n_2a=n_1(Y/b)$.
\end{proof}

\begin{proposition}\label{O:R:normalize}
The actual commutator is exactly $\mu=\tau_2n_1$.
\end{proposition}
\begin{proof}
Both are nontrivial characters of the same cyclic norm quotient of order
$p$. Their restrictions to its last nonzero unit layer distinguish the
$p-1$ possible generators.
If $\delta>0$, $v_1s\ge(p-1)\delta$ and
$v_1(\varrho/c)=(p-1)\delta$ are positive. Equations \eqref{O:R:S} and
Lemma~\ref{O:R:T} agree on $U_{B_1}^{t'}$; here $P(x)\equiv x$ at the
character modulus. Hence the characters are equal.

If $\delta=0$, use $x=Y/b$, whose classes as $Y$ varies span
$\PP_{B_1}^t/\PP_{B_1}^{t+1}$. On this layer $P(x)\equiv x$.
The intermediate coefficients of $n_1(1-x)$
other than trace and norm are in $\pp_A^{t+1}$ by the strong symmetric
bound. Thus
\[
 (\tau_2n_1)(1-x)=\Psi_A(T_1x+n_1x)
                =\Psi_1(x)\Psi_A(n_1x).
\]
Lemma~\ref{O:R:reciprocal} identifies this with $\Psi_1((1-s)x)$, which is the
actual commutator by \eqref{O:R:S}. Their critical restrictions agree, and
therefore the full norm characters agree. No $R_2$ congruence was used.
\end{proof}

\subsubsection*{Whole-subgroup comparison and norm descent}
\begin{theorem}\label{O:R:R2}
For every odd $p$, in both characteristics and both break configurations,
\begin{equation}\label{O:R:strongR2}
 H\equiv \varrho^{p-1}\pmod{\pp_A^R},\qquad
 R=t+1-c_t+(p-1)\delta.
\end{equation}
Consequently if $m$ is the transition parameter with $pm>t$, then
\[
 H\equiv(\beta/\alpha)^{p-1}
       \pmod{\pp_A^{1+t_2-(p-2)m}}.
\]
\end{theorem}
\begin{proof}
Proposition~\ref{O:R:normalize}, \eqref{O:R:S} and Lemma~\ref{O:R:T} imply
\[
 \Psi_1((s-\varrho/c)P(x))=1\qquad\text{for every }x\in\PP_{B_1}^q.
\]
The polynomial $P$ is a bijection of this entire ideal, by Hensel's lemma
and $P'(0)=1$. Perfect additive duality therefore gives the absolute
coefficient congruence
\[
 s-\varrho/c\in\PP_{B_1}^{T'-q}=\PP_{B_1}^R.
\]
Put $d_s=(p-1)\delta$. Both $\varrho/c$ and $s$ have valuation $d_s$, since
$R-d_s=R_0:=t+1-c_t\ge1$. Their ratio is in $U_{B_1}^{R_0}$.
We have $1\le R_0\le t$, so the cyclic norm filtration gives a ratio of
norms in $U_A^{R_0}$. Multiplication by the norm's valuation $d_s$ yields
\[
 n_1s\equiv n_1(\varrho/c)=\varrho^{p-1}\varepsilon^{-1}
                  \pmod{\pp_A^{d_s+R_0}}=\pmod{\pp_A^R}.
\]
The remaining error has valuation at least
$d_s+t\ge d_s+t+1-c_t=R$. This proves \eqref{O:R:strongR2}, including $\delta=0$.
Finally $pm>t$ and $p\nmid t$ imply $m\ge c_t$, and
\[
 R-[1+t_2-(p-2)m]=(p-2)(\delta+m)-c_t\ge0.
\]
This proves the required terminal modulus uniformly, with no exceptional
$p=3$ strengthening and no additional critical norm representative.
\end{proof}


\subsection{Simultaneous norm parameters and stationary factors}\label{O:sec:parameters}
\subsubsection*{Truncated-logarithm coordinates}
Throughout $p$ is odd and $P(X)=\sum_{j=1}^{p-1}X^j/j$.
For a field $E$ let $\ee_E$ have largest trivial ideal $\PP_E^{-d_E}$.
A multiplicative character $\theta$ of conductor $n\ge2$ admits a
coefficient $C$ with
\begin{equation}\label{O:P:fullP}
 \theta(1-x)=\ee_E(CP(x)),\qquad
 x\in\PP_E^{\ceil{n/p}},\qquad v_E(C)=-d_E-n.
\end{equation}
This follows directly from perfect additive duality: the polynomial
$P(X+Y-XY)-P(X)-P(Y)$ has total degree at least $p$, while $P$ is a
bijection of every positive ideal, with derivative a unit. It therefore
identifies the multiplicative quotient $U_E^q/U_E^n$ with
$\PP_E^q/\PP_E^n$ when $pq\ge n$. Exact conductor gives the stated
valuation. The full coefficient is determined modulo
$\PP_E^{-d_E-q}$, not just by a half-conductor class.

\begin{lemma}[choosing a norm ratio by moving a norm-character coefficient]
\label{O:P:choose}
Let $E/F$ be ramified cyclic of odd prime degree with break $t\ge1$.
Let $\nu$ be an actual nontrivial norm character and write
\[
 \nu(1-z)=\ee_F(A P(z)),\qquad z\in\pp_F^{c},\quad
 c=\ceil{(t+1)/p},\quad v_F(A)=-d_F-t-1.
\]
For any $\gamma\in F^\times$ there is $A'$ with
\[
 A'/A\in U_F^t,\quad
 \nu(1-z)=\ee_F(A'P(z))\ (z\in\pp_F^c),\quad
 \gamma/A'\in\Nm_{E/F}(E^\times).
\]
Consequently there is $w\in E^\times$ with
$\Nm w=\gamma/A'$ and $v_E(w)=v_F(\gamma/A')$.
\end{lemma}
\begin{proof}
The character $\nu$ has exact conductor $t+1$, so its restriction to
$U_F^t$ is a surjection onto $\mu_p$.
Choose $h\in U_F^t$ with $\nu(h)=\nu(\gamma/A)$ and put $A'=Ah$.
The change in the additive coefficient on the entire prescribed
subgroup has valuation at least
\[
 v_F(A(h-1)P(z))\ge-d_F-t-1+t+c\ge-d_F.
\]
It is therefore invisible to $\ee_F$ on that whole subgroup.
Meanwhile $\nu(\gamma/A')=1$. The kernel of this nontrivial
character of the order-$p$ norm quotient is precisely the norm
subgroup. This proves exact norm membership. Total ramification gives
$v_F(\Nm w)=v_E(w)$.
\end{proof}

\begin{lemma}[Preservation of the full minimal models]\label{O:P:preserve}
Let the totally ramified notation be as in Subsection~\ref{O:sec:models}, and
suppose the actual minimal pair has full coefficients $\alpha b,\alpha a$.
If $h\in U_A^{t_2}$, then replacing $\alpha$ by $\alpha h$ leaves both
full model formulas unchanged. It also leaves the full coefficient formula
of the same norm character of $B_2/A$ unchanged.
\end{lemma}
\begin{proof}
Write $q_1=\delta+\lceil(2t+1)/p\rceil$, $q_2=\lceil(2t+\delta+1)/p\rceil$.
The old additive characters $\Psi_i=\ee_{B_i}(\alpha\,\cdot)$ have largest
trivial ideals with exponents $1+t'$ and $1+t_2$.
For $z\in\PP_{B_i}^{q_i}$ the change in the normalized exponent has
valuation at least $pt_2-t+q_i$. On $B_1$ its excess over $1+t'$ is
\[
 (p-2)t+\delta+\lceil(2t+1)/p\rceil-1\ge0,
\]
and on $B_2$ its excess over $1+t_2$ is
\[
 (p-2)t+(p-1)\delta+q_2-1\ge0.
\]
Thus neither full character formula changes. On the lower field $A$, the
change at depth $c_2=\lceil(t_2+1)/p\rceil$ has normalized valuation at
least $t_2+c_2\ge t_2+1$, proving the last assertion.
\end{proof}

\begin{remark}
This lemma permits realization of the actual minimal pair first, including
the scalar $j\in\{1,\ldots,p-1\}$ in Theorem~\ref{O:M:realize}, and the exact
norm choices afterwards. Thus an exact norm equation cannot be invalidated
by a later, unrecorded change of generator of the norm-character group.
It does not assert that a previously fixed ratio is already a norm.
\end{remark}

\subsubsection*{Full coefficients of a descending pullback}
\begin{proposition}\label{O:P:pullback}
Let $E/F$ be as in Lemma~\ref{O:P:choose}. Let $\lambda$ have conductor
$1+r$, where $r>t$, and let $\gamma$ be its full coefficient relative
to $\ee_F$. Choose $A,w$ by that lemma so that $\Nm w=\gamma/A$.
Put $h=r-t$, $n=1+t+ph$, and $q=h+\ceil{(t+1)/p}=\ceil{n/p}$.
Then the actual pullback has exact conductor $n$ and
\begin{equation}\label{O:P:pullP}
 (\lambda\Nm)(1-x)=\ee_E((\gamma-Aw)P(x)),
 \qquad x\in\PP_E^q,\quad \ee_E=\ee_F\Tr_{E/F}.
\end{equation}
Its coefficient has valuation $p\,v_F(\gamma)$.
\end{proposition}
\begin{proof}
The conductor is the high cyclic norm-transport formula, with different
exponent $(p-1)(t+1)$. One has $v_E(w)=-h$ and
\[
 v_E(Aw/\gamma)=(p-1)h>0.
\]
This proves the claimed coefficient valuation and matches conductor
$n$, since $d_E=pd_F+(p-1)(t+1)$.

For $x\in\PP_E^q$, the norm $\Nm(1-x)$ lies in the full
$\lambda$-coordinate domain $U_F^{\ceil{(r+1)/p}}$.
Indeed the norm endpoint has valuation at least $q$, and the
intermediate terms at least
$\floor{(q+(p-1)(t+1))/p}$; both are at least the required depth.
For the second inequality compare the numerator with $r+p$; its
excess is at least $\ceil{(t+1)/p}+(p-2)t-1\ge0$.
Lemma~\ref{O:I:normlog}, with $h=r-t$ and target modulus $r+1$, gives
\[
 (\lambda\Nm)(1-x)
 =\ee_E(\gamma P(x))\ee_F(\gamma\Nm(P(x))).
\]
Here its error has depth $r+1$ before multiplication by $\gamma$,
so is in $\pp_F^{-d_F}$ afterward.
Now $wP(x)$ belongs to $\PP_E^{\ceil{(t+1)/p}}$.
The actual norm-character phase, obtained by applying $\nu$ to a norm
and using Lemma~\ref{O:I:normlog} at modulus $t+1$, says
\[
 \ee_E(AwP(x))=\ee_F(-A\Nm(wP(x)))
             =\ee_F(-\gamma\Nm(P(x))).
\]
Combine the last two displays to obtain \eqref{O:P:pullP}. In the first
expression $A$ stays outside the norm. The negative sign is forced
by moving the norm contribution across a product equal to $1$.
\end{proof}

\subsubsection*{The full Lamprecht formula with its odd term retained}
\begin{lemma}[complete formula from a full $P$ coefficient]\label{O:P:lamprecht}
Under \eqref{O:P:fullP}, there is an explicitly determined
$g_E(\theta,C)\in\mu_4$ with
\begin{equation}\label{O:P:Delta}
 \Delta_E(\theta,\ee_E)
 =\theta(-C^{-1})\ee_E(-C)g_E(\theta,C).
\end{equation}
For even $n$, $g_E=1$. For odd $n=2d+1$, choose any $\delta_E$
of valuation $d$. Then
\begin{equation}\label{O:P:gauss}
 g_E=|k_E|^{-1/2}\sum_{x\in k_E}
       \ee_E(-C\delta_E^2\widetilde x^{\,2}/2).
\end{equation}
The residue phase in this display is exactly the full odd-conductor
Lamprecht term, not only its polarization.
\end{lemma}
\begin{proof}
Apply Lemma~\ref{O:I:enhanced} with $\Psi_E=\ee_E$, ideal exponent
$J=-d_E$, and enhanced coefficient $B=C$. Take the ordinary coefficient
to be the same $C$, so the displacement $\eta$ is zero. The full
$P$-formula is valid on $\pp_E^{\lfloor n/2\rfloor}$ because
$\lfloor n/2\rfloor\ge\lceil n/p\rceil$ for odd $p$ and $n\ge2$.
Its residual function is therefore exactly
$\ee_E(-C\delta_E^2\widetilde x^{\,2}/2)$, and the elementary factor is
$\theta(-C^{-1})\ee_E(-C)$. The homogeneous quadratic sum has phase in
$\mu_4$, as in the proof of that lemma. No affine displacement remains.
\end{proof}

\subsubsection*{Application to the totally ramified transition}
Let the breaks be $t=t_1\le t_2=t+\delta$, $t'=t+p\delta$, $p\nmid t$.
Let the actual pair have the form
\[
 \theta_1=\chi(\lambda n_1),\qquad
 \theta_2=\varphi(\lambda n_2),\qquad
 \chi N_1=\varphi N_2,
\]
where $a(\chi)=m_1=1+t+t'$, $a(\varphi)=m_2=1+t+t_2$.
Suppose $a(\lambda)=1+r$, $m=r-t_2$, and $pm>t$.
First realize the actual minimal pair by Theorem~\ref{O:M:realize}.
Conjugation on $B_1$ also conjugates $\theta_1$, since $\lambda n_1$ is
invariant, and hence does not change its local constant. Denote the resulting
full coefficients by $\alpha b,\alpha a$, where $b=N_1\Delta$ and
$a=N_2\Delta$. The scalar $\alpha$ is the full coefficient of an actual
nontrivial norm character for $B_2/A$.
Now choose the full coefficient $\gamma$ for $\lambda$, and choose any
actual nontrivial norm character for $B_1/A$, with coefficient $\beta$.
Apply Lemma~\ref{O:P:choose} separately to these two norm characters.
Lemma~\ref{O:P:preserve} shows that its replacement of $\alpha$ does not change
the already realized full minimal models; replacing $\beta$ affects neither
minimal model. Renaming the adjusted coefficients, the exact choices give
\begin{equation}\label{O:P:parameters}
 n_2w_2=\gamma/\alpha,\quad v_2w_2=-m,\qquad
 n_1w_1=\gamma/\beta,\quad v_1w_1=-(m+\delta),
\end{equation}
\[
 v_A\alpha=-d-1-t_2,\quad
 v_A\beta=-d-1-t,\quad v_A\gamma=-d-1-r.
\]
Both adjusted coefficients still represent the same actual lower norm
characters on their full logarithmic subgroups. The full minimal formulas
remain
\[
 \chi(1-x)=\ee_{B_1}(\alpha bP(x)),\qquad
 \varphi(1-x)=\ee_{B_2}(\alpha aP(x)),\quad
 x\in\PP_{B_i}^{\lceil m_i/p\rceil}
\]
on the appropriate side. In particular Theorem~\ref{O:R:R2} applies to these
very coefficients, with no extra norm condition on $\beta/\alpha$.
Proposition~\ref{O:P:pullback} now gives the full coefficients
\begin{equation}\label{O:P:C}
 C_1=\gamma-\beta w_1+\alpha N_1\Delta,\qquad
 C_2=\gamma-\alpha w_2+\alpha N_2\Delta.
\end{equation}
The actual conductors are
\[
 n_1^{\rm char}=1+t'+pm,
 \qquad n_2^{\rm char}=1+t_2+pm.
\]
Write $\alpha_1=\beta$, $\alpha_2=\alpha$.
The full formulas for \eqref{O:P:C} hold at depths
$\ceil{n_i^{\rm char}/p}$. Indeed $n_i^{\rm char}>m_i$,
$v_i(\alpha N_i\Delta/\gamma)=pm-t>0$, and
$v_i(\alpha_iw_i/\gamma)=(p-1)(r-t_i)>0$.
Thus no leading cancellation is possible, and the conductors and
full-coordinate domains are exactly as stated.

\begin{proposition}[Ueda-normalized orientation of the four factors]
\label{O:P:four}
Define $A_i=\gamma+\alpha N_i\Delta$ and use the corrected factors
\[
 Q_1=\frac{\chi(A_1)}{\varphi(C_2)},\quad
 Q_2=\lambda\left(\frac{n_1(A_1)}{n_2(C_2)}\right),
\]
\[
 Q_3=\ee_{B_1}(-\beta w_1)\theta_1(C_1/A_1),\qquad
 Q_4=\ee_{B_2}(\alpha w_2).
\]
Then, with the explicit $g_i$ from Lemma~\ref{O:P:lamprecht},
\begin{equation}\label{O:P:ratio}
 \frac{\Delta_{B_2}(\theta_2,\ee_{B_2})}
      {\Delta_{B_1}(\theta_1,\ee_{B_1})}
 =\frac{g_2}{g_1}\,Q_1Q_2Q_3Q_4\,
 \ee_A\left(\alpha[T_1N_1\Delta-T_2N_2\Delta]\right).
\end{equation}
In particular, after the trace--norm estimate of Theorem~\ref{O:A:thm:H14}
removes the last factor, proving $(Q_1Q_2Q_3Q_4)^{-1}=1$
implies that the actual ratio belongs to $\mu_4$.
\end{proposition}
\begin{proof}
Compatibility and odd degree give $\theta_1(-1)=\theta_2(-1)$:
apply the equal norm pullbacks to $-1\in B^\times$.
Formula \eqref{O:P:Delta} therefore gives
\[
 \frac{\Delta_{B_2}}{\Delta_{B_1}}
 =\frac{g_2}{g_1}\frac{\theta_1(C_1)}{\theta_2(C_2)}
                      \ee_{B_1}(C_1)\ee_{B_2}(-C_2).
\]
The two $\gamma$ traces both equal $p\gamma$ and cancel.
Substitute \eqref{O:P:C}; the remaining additive factors are the three
shown in \eqref{O:P:ratio}. Direct multiplication gives
\[
 Q_1Q_2=\frac{\theta_1(A_1)}{\theta_2(C_2)},\qquad
 Q_1Q_2\theta_1(C_1/A_1)=\frac{\theta_1(C_1)}{\theta_2(C_2)}.
\]
The remaining additive factors are exactly those of $Q_3Q_4$.
The same trace--norm estimate is
$T_1N_1\Delta-T_2N_2\Delta\in\pp_A^{1+t_2}$;
its weight $\alpha$ puts it in $\pp_A^{-d}$, as required.
\end{proof}


\section{The totally ramified odd-prime calculation}\label{sec:odd-calculation}
\subsection{The linear trace calculation}\label{O:sec:linear}
\subsubsection*{Notation and local estimates}
Use the totally ramified diamond $A\subset B_1,B_2\subset B$,
$\operatorname{Gal}(B/A)=C_p^2$, $p$ odd, with
\[
 t_1=t,\quad t_2=t+\delta,\quad t'=t+p\delta,\quad p\nmid t.
\]
Write $N_i=\Nm_{B/B_i}$, $n_i=\Nm_{B_i/A}$, $S_i=\Tr_{B/B_i}$,
and normalize $v=v_B$. The independently proved construction and
estimates of Section~\ref{O:sec:as} give
\[
 \Delta^p-\Delta=a\in B_2,\quad v(\Delta)=-t,\quad
 b=N_1\Delta,\quad e_i=E_i^{B/B_1}(\Delta),\quad s=-e_{p-1},
\]
\begin{equation}\label{O:W:ledger}
 v(e_i)\ge p((p-1)t_2-it),\qquad
 v(S_1(Y\Delta^i))\ge p(p-1)t_2-pit+v(Y)\quad(Y\in B_2).
\end{equation}
Put $V=v_B(p)$ in mixed characteristic; in equal characteristic every
term carrying $p$ is zero. The local trace-ideal bound is
$V\ge p(p-1)t_2$.
Let $w\in B_2$ have $n_2w=W$, put
\[
 u=W^{-1},\quad v_A(u)=m,\quad M=pm>t,\quad
 K=1+pt_2,\quad S=p(p-1)\delta.
\]
Thus $M\ge t+1$, $v(w)=-M$, $v(W)=-pM$, $v(a)=v(b)=-pt$,
and $v(s)\ge S$. All congruences below have modulus $\PP_B^K$
unless a different modulus is displayed. Set
\[
 Y=W+b,\quad D=Y^p-Y,\quad D_0=W^p-W+a^p,\quad C=1+a/W.
\]
Their denominators are nonzero: $v(Y)=-pM$ and
$v(D)=v(D_0)=-p^2M$.

\subsubsection*{The actual upper norms of the conjugates}
Let $\mathcal T=\{0\}\cup\mu_{p-1}\subset A$ and let
$\Delta_\xi$ be the root congruent to $\Delta+\xi$ at positive
relative depth. The construction of Section~\ref{O:sec:as} gives
\[
 \Delta_\xi=\Delta+\xi+\eta_\xi,\qquad
 v(\eta_\xi)\ge V-(p-1)t.
\]
These are all $B/B_2$ conjugates; in equal characteristic
$\eta_\xi=0$.

\begin{lemma}\label{O:W:actualnorm}
For each $\xi\in\mathcal T$,
\[
 v_{B_1}\bigl(N_1\Delta_\xi-N_1(\Delta+\xi)\bigr)\ge t_2+1.
\]
\end{lemma}
\begin{proof}
In mixed characteristic put $z=\eta_\xi/(\Delta+\xi)$, so
$v(z)\ge V-(p-2)t$. The different exponent for $B/B_1$ is
$D'=(p-1)(t'+1)$. The strong symmetric estimate gives
\[
 v_1(N_1(1+z)-1)\ge
 \min\left\{v(z),\left\lfloor\frac{v(z)+D'}p\right\rfloor\right\}.
\]
Here $v(z)\ge t'+1$, so the second entry is no larger than the first.
Since $v_1N_1(\Delta+\xi)=-t$, the required difference has valuation
at least
\[
 \frac Vp-t+(p-1)\delta+\left\lceil\frac tp\right\rceil
 \ge (p-2)t+2(p-1)\delta+\left\lceil\frac tp\right\rceil
 \ge t_2+1.
\]
For the last inequality, subtract $t+\delta+1$ and use $p\ge3$,
$t\ge1$, $\lceil t/p\rceil\ge1$. Equal characteristic is exact.
\end{proof}

Define $f(X)=X/(1+X/W)$. For $v(X)=-pt$ its denominator is a unit.
The exact norm polynomial gives
\[
 N_1(\Delta+\xi)=b+\xi-\xi s+R_\xi,
 \qquad v(R_\xi)\ge p((p-1)t_2-(p-2)t)\ge pt_2.
\]
The exact identity
\[
 f(X+z)-f(X)=\frac{z}{(1+X/W)(1+(X+z)/W)}
\]
and its first-order expansion show
\begin{equation}\label{O:W:normquotient}
 f(N_1\Delta_\xi)\equiv
 \frac{W(b+\xi)}{W+b+\xi}-\frac{\xi s}{(1+b/W)^2}
 \pmod{\PP_B^{pt_2}}.
\end{equation}
Indeed, the quadratic $s$-error has depth at least $pM+2S\ge pt_2$;
replacing $1+(b+\xi)/W$ by $1+b/W$ in that linear error costs at
least $pM+S\ge pt_2$. Lemma~\ref{O:W:actualnorm} is deeper still.

\subsubsection*{\texorpdfstring{The power factors, including $p=3$}{The power factors, including p=3}}
The actual trace is
\[
 W_i=S_2\left((\Delta/w)^i\frac{b}{1+b/W}\right),\qquad1\le i\le p-1.
\]
Each summand's power factor has valuation $i(M-t)\ge1$; therefore
\eqref{O:W:normquotient} may be substituted at modulus $K$.

For $p\ge5$, replacing $\Delta_\xi$ by $\Delta+\xi$ in that power
factor costs at least
\[
 V+iM-(2p+i-2)t\ge K.
\]
At the weakest values $M=t+1$, $V=p(p-1)t_2$, the excess over $K$
is $(p^2-4p+2)t+p(p-2)\delta+i-1\ge0$.

Suppose $p=3$. Hensel's equation at $\Delta+\xi$, $\xi=\pm1$,
gives the sufficient refinement
\begin{equation}\label{O:W:p3root}
 \Delta_\xi=\Delta+(1+3\Delta^2)\xi+\eta'_\xi,
 \qquad v(\eta'_\xi)\ge V-t.
\end{equation}
To see this, the initial residual is
$3\xi\Delta^2+3\xi^2\Delta$; substituting the first Hensel correction
back into the equation adds terms of valuation at least $2V-4t$,
which is at least $V-t$. Thus no coefficient beyond the displayed
$3\xi\Delta^2$ is needed. Deleting $\eta'_\xi$ in $W_i$ costs at
least $V+iM-(i+3)t\ge K$ for $i=1,2$.

Put $\lambda=1+3\Delta^2$ and
\[
 C_i(\lambda)=w^{-i}\sum_{\xi\in\mathcal T}
 (\Delta+\lambda\xi)^i\frac{W(b+\xi)}{Y+\xi}.
\]
Direct summation of the three terms gives the exact identities
\begin{align*}
 C_1(\lambda)-C_1(1)&=
 \frac{2(\lambda-1)W^2}{w(Y^2-1)},\\
 C_2(\lambda)-C_2(1)&=
 \frac{2W(\lambda-1)}{w^2(Y^2-1)}
 \bigl((\lambda+1)(bY-1)+2W\Delta\bigr).
\end{align*}
Their valuations are at least $V-2t+M$ and $V-5t+2M$, respectively,
hence at least $K$. In the signed exceptional sum the exact formulas
\[
 \sum_{\xi=\pm1}\xi(\Delta+\lambda\xi)=2\lambda,
 \qquad \sum_{\xi=\pm1}\xi(\Delta+\lambda\xi)^2=4\lambda\Delta
\]
give errors of depths $S+M+V-2t$ and $S+2M+V-3t$, also at least $K$.
Thus the unscaled Teichmuller powers may be used for every odd $p$.

\begin{proposition}[the actual split]\label{O:W:split}
For $1\le i\le p-1$,
\begin{align}\label{O:W:Wsplit}
 W_i\equiv{}&w^{-i}\sum_{\xi\in\mathcal T}
 (\Delta+\xi)^i\frac{W(b+\xi)}{Y+\xi}\notag\\
 &-\frac{s}{w^i(1+b/W)^2}
       \sum_{\xi\in\mu_{p-1}}\xi(\Delta+\xi)^i .
\end{align}
In particular, the exceptional sum has no denominator $Y+\xi$.
\end{proposition}
\begin{proof}
Combine \eqref{O:W:normquotient} with the power estimates just proved.
\end{proof}

\subsubsection*{Exact interpolation and the four replacements}
\begin{lemma}[Finite interpolation and power sums]\label{O:D:interpolation}
For $0\le i<p$ and $Z\notin-\mathcal T$,
\[
 \sum_{\xi\in\mathcal T}\frac{(\Delta+\xi)^i}{Z+\xi}
 =\frac{p\Delta^i}{Z}+(p-1)\frac{(\Delta-Z)^i}{Z^p-Z}.
\]
For $1\le i<p$, the accompanying power sums are
\begin{align*}
 \sum_{\xi\in\mathcal T}(\Delta+\xi)^i
 &=p\Delta^i+(p-1)\ind_{i=p-1},\\
 \sum_{\xi\in\mu_{p-1}}\xi(\Delta+\xi)^i
 &=(p-1)(\ind_{i=p-2}+i\Delta\ind_{i=p-1}).
\end{align*}
All three identities are exact in mixed and equal characteristic.
\end{lemma}
\begin{proof}
Regard the first identity as an equality of rational functions in $Z$.
Both sides have only simple poles at $-\mathcal T$ and vanish at infinity.
At zero the residue on each side is $\Delta^i$; at $-\xi$ with
$\xi^{p-1}=1$, it is $(\Delta+\xi)^i$, since the derivative of
$Z^p-Z$ there is $p-1$. Subtraction gives a rational function with no
poles and zero value at infinity, hence zero. This works in both
characteristics because $p-1$ is nonzero. The last two identities follow
by the binomial theorem and
$\sum_{\xi\in\mu_{p-1}}\xi^j=0$ unless $p-1$ divides $j$, when the
sum is $p-1$.
\end{proof}

Apply the lemma with $Z=Y=W+b$. This value has negative valuation,
since $W$ dominates $b$, so the denominators are nonzero. We obtain
\begin{align}\label{O:W:Wpre}
 W_i\equiv{}&\frac{pWb\Delta^i}{w^iY}
 +(p-1)\frac{W}{w^{p-1}}\ind_{i=p-1}
 -(p-1)\frac{W^2(\Delta-Y)^i}{w^iD}\notag\\
 &-(p-1)\frac{s}{(1+b/W)^2}
 \left(\frac{\ind_{i=p-2}}{w^{p-2}}
       +\frac{i\Delta\ind_{i=p-1}}{w^{p-1}}\right).
\end{align}
The first term has depth at least $V-pt+i(M-t)\ge K$.

The characteristic polynomial of $\Delta$ over $B_1$ gives two useful
forms, with all signs fixed by $s=-e_{p-1}$:
\begin{align}\label{O:W:bforms}
 b&=\Delta+a+a\omega,& v(\omega)&\ge pt_2-t,\notag\\
 b&=\Delta+a-\Delta s+\rho,&
 v(\rho)&\ge p(p-1)t_2-(p^2-2p+2)t.
\end{align}
Each follows by solving the characteristic equation for its constant
term and applying \eqref{O:W:ledger} to the remaining coefficients.

\begin{lemma}[replace the denominator]\label{O:W:den}
$W^2(\Delta-Y)^i/(w^iD)$ may be replaced by
$W^2(\Delta-Y)^i/(w^iD_0)$.
\end{lemma}
\begin{proof}
Write $q=b/W$, $a_1=a/W$, $d_1=\Delta/W$, $h_1=a\omega/W$.
Then $q=a_1+d_1+h_1$ and
\[
 D-D_0=W^p E,\qquad
 E=(1+q)^p-1-a_1^p-W^{1-p}q.
\]
Since $d_1^p=(\Delta+a)/W^p$,
$W^{1-p}q=d_1^p+W^{1-p}h_1$. Every term of $E$ is $p$-divisible,
or contains $h_1^p$, or contains $W^{1-p}h_1$. Consequently
\[
 \begin{aligned}
 v(E)&\ge\min\bigl\{V+p(M-t),\ p(pM+pt_2-(p+1)t),\\
 &\hspace{28mm}p(p-1)M+pM+pt_2-(p+1)t\bigr\}\ge K+M.
\end{aligned}
\]
The remaining factor in the difference has valuation
$(p(p-2)-i(p-1))M\ge-M$.
\end{proof}

\begin{lemma}[replace the numerator]\label{O:W:numerator}
Put $A_0=-a-W$. Modulo the required ideal,
\[
 \frac{W^2(\Delta-Y)^i}{w^iD_0}
 \equiv\frac{W^2}{w^iD_0}
 \begin{cases}
 A_0^i,&i\le p-2,\\
 A_0^{p-1}+(p-1)A_0^{p-2}\Delta s,&i=p-1.
 \end{cases}
\]
\end{lemma}
\begin{proof}
By \eqref{O:W:bforms}, $\Delta-Y=A_0+\Delta s-\rho$.
Set $Z=S-t$ and $R=p(p-1)t_2-(p^2-2p+2)t$.
For $i\le p-2$ the first perturbation after the displayed weight has
depth at least $2(p-1)M+Z\ge K$.
At $i=p-1$, deleting $\rho$ costs $(p-1)M+R\ge K$;
every term quadratic or higher in $\Delta s$ costs
$(2p-1)M+2Z\ge K$. At $p=3$, $\delta=0$, $M=t+1$ the
three excesses are $3$, $1$, $4$, respectively, so none is lost at
the boundary.
\end{proof}

\begin{lemma}[the inverse-square replacement and cancellation]\label{O:W:cancel}
The signed term $s/[w^{p-2}(1+b/W)^2]$ may use $C^{-2}$ instead,
as may its multiple $\Delta/w$. The two remaining linear
$\Delta s$ terms at $i=p-1$ cancel.
\end{lemma}
\begin{proof}
The difference of the inverse squares is
\[
 \frac{(a-b)/W\,[2+(a+b)/W]}{(1+b/W)^2(1+a/W)^2}.
\]
Since $v(a-b)\ge-t$, its weighted first error has depth
$S+2(p-1)M-t\ge K$; multiplication by $\Delta/w$ only increases it.
After extracting the common $(p-1)^2\Delta s/w^{p-1}$, the two
linear terms have bracket
\[
 -\frac1{(1+a/W)^2}
 +\frac{(1+a/W)^{p-2}}{1-W^{1-p}+(a/W)^p}.
\]
The numerator of this difference is
$W^{1-p}+\sum_{j=1}^{p-1}\binom pj(a/W)^j$.
The first weighted term has depth $S-t+(p^2-1)M\ge K$;
the others have depth at least
$S-t+(p-1)M+V+p(M-t)\ge K$.
\end{proof}

\begin{theorem}[the verified linear packet]\label{O:W:final}
For every odd $p$ and $1\le i\le p-1$,
\[
 \boxed{\begin{aligned}
 W_i\equiv{}&
 -(p-1)\frac{s}{w^{p-2}C^2}\ind_{i=p-2}
 +(p-1)\frac{W}{w^{p-1}}\ind_{i=p-1}\\
 &-(p-1)\frac{W^2(-a-W)^i}{w^i(W^p-W+a^p)}
 \pmod{\PP_B^{1+pt_2}}.
 \end{aligned}}
\]
\end{theorem}
\begin{proof}
Insert Lemmas~\ref{O:W:den}--\ref{O:W:cancel} into \eqref{O:W:Wpre}.
\end{proof}
\begin{remark}
This proves the actual packet needed in the $Q_1$ logarithm, not only
an identity for artificial translates. The delicate $p=3$ replacement
was completed before applying interpolation. The exceptional signed
sum in \eqref{O:W:Wsplit} is not changed into a rational sum at any stage.
\end{remark}


\subsection{The quadratic trace calculation}\label{O:sec:quadratic}
\subsubsection*{The character data and local estimates}
Let $A\subset B_1,B_2\subset B$ be a totally ramified $C_p^2$ diamond,
where $p$ is odd. Put
\[
 t=t_1,\quad t_2=t+\delta,\quad t'=t+p\delta,\quad
 \delta\ge0,\quad p\nmid t,
\]
\[
 N_i=\Nm_{B/B_i},\quad n_i=\Nm_{B_i/A},\quad
 S_i=\Tr_{B/B_i},\quad T_i=\Tr_{B_i/A}.
\]
Valuations are normalized on the named field; thus $v_B|_{B_i}=pv_i$
and $v_i|_A=pv_A$. Fix
\[
 \Delta^p-\Delta=a\in B_2,\quad v_B\Delta=-t,\quad
 b=N_1\Delta,\quad N_2\Delta=a,\quad
 s=-E_{p-1}^{B/B_1}(\Delta).
\]
Let $u=\alpha/\gamma\in A$, $m=v_Au$, $M=pm>t$, and choose
$w\in B_2^\times$ with $n_2w=u^{-1}$ and $v_2w=-m$.
The legitimate ordering of these choices is proved in Subsection~\ref{O:sec:parameters}.
Write
\[
 W=u^{-1},\quad C_b=1+b/W,\quad C=1+a/W,\quad
 D_0=W^p-W+a^p,
\]
\[
 q=p-1,\quad S=p(p-1)\delta,\quad K=1+pt_2,
 \quad F=1+t+t'=1+2t+p\delta.
\]
In common-field valuation $v=v_B$ one has
\[
 v(w)=-M,\quad v(W)=-pM,\quad v(a)=v(b)=-pt,
 \quad v(C)=v(C_b)=0,\quad v(D_0)=-p^2M.
\]
In mixed characteristic set $e=v_B(p)$; the cyclic ramification
bound gives $e\ge p(p-1)t_2$. In equal characteristic all
$p$-divisible terms vanish.

Normalize $\Psi_E(x)=\ee_E(\alpha x)$, with trace-compatible
$\ee_E$. Its largest trivial ideals on $A,B_1,B_2$ have exponents
$1+t_2,1+t',1+t_2$, respectively. The actual minimal first inducing
character satisfies
\begin{equation}\label{O:D:actual}
 \chi(1-x)=\Psi_{B_1}(bP(x)),\qquad
 x\in\PP_1^{\ceil{F/p}},\qquad
 P(X)=\sum_{j=1}^{p-1}X^j/j.
\end{equation}
This full formula is the conclusion of Subsection~\ref{O:sec:models}; it is not
replaced by a half-conductor assertion.

We use Lemma~\ref{O:A:lem:twisted}, with all domains visible:
\begin{align}
 v_1(E_j^{B/B_1}(Y\Delta))
   &\ge(p-1)t_2+j(v_2Y-t),\quad 1\le j<p,\label{O:D:H15s}\\
 v_1(S_1(Y\Delta^j))
   &\ge(p-1)t_2-jt+v_2Y,\quad 0\le j<p.\label{O:D:H15t}
\end{align}
Here $Y\in B_2$. These are the explicit local estimates of Section~\ref{O:sec:as},
not an appeal to Lakkis's Satz~2. We also use Ueda's cyclic estimate
\begin{equation}\label{O:D:cyclic}
 v_F(E_j^{E/F}z)\ge
 \floor{\frac{jv_Ez+(p-1)(b_0+1)}p},\quad 1\le j<p,
\end{equation}
for a cyclic edge of break $b_0$, and the corresponding trace-ideal
formula. The floor in this formula is essential.


For the quadratic packet put
\begin{equation}\label{O:D:WPdef}
 T=S_1((\Delta/w)^{p-1}),\qquad
 \mathcal Q=S_2\!\left(
 \frac{b(\Delta/w)^{p-1}T}{C_b^2}\right).
\end{equation}
Here $T\in B_1$ and $\mathcal Q\in B_2$; the calculation below does not apply a
$B_2$ character to any expression outside $B_2$.
\subsubsection*{The quadratic term: actual conjugates and one finite sum}
All congruences in this and the next subsection are in $B$ modulo
$\PP_B^K$. Set $f(X)=X/(1+X/W)^2$ and
$\mathcal T=\{0\}\cup\mu_{p-1}\subset A$.
The symmetric-coefficient estimates give
\begin{equation}\label{O:D:Tbound}
 v(T)\ge S+qM,\qquad
 v\left(S_1(\Delta^j/w^q)\right)\ge
 p(p-1)t_2-pjt+qM\quad(0\le j\le q).
\end{equation}
The characteristic polynomial of $\Delta$ over $B_1$ also gives
\begin{equation}\label{O:D:omega}
 b=\Delta+a+a\omega,\qquad v(\omega)\ge pt_2-t.
\end{equation}
Indeed the error divided by $a$ has terms of valuation at least
$(p-1)(pt_2-jt)$ for $1\le j<p$, by \eqref{O:D:H15s}; their minimum
is at least $pt_2-t$.

\begin{lemma}[uniform conjugate replacement for the quadratic term]
\label{O:D:conjugates}
In the sum defining $\mathcal Q$ one may replace every actual
$B/B_2$ conjugate of $\Delta$ by its corresponding $\Delta+\xi$,
$\xi\in\mathcal T$, modulo $\PP_B^K$. This holds for every odd $p$,
including $p=3$ in mixed characteristic.
\end{lemma}
\begin{proof}
In equal characteristic the statement is exact. In mixed
characteristic, Hensel's lemma applied at $\Delta+\xi$ gives
\[
 \Delta_\xi=\Delta+\xi+\epsilon_\xi,\qquad
 v(\epsilon_\xi)\ge e-qt.
\]
The derivative is a unit since $e-qt>0$, and the $p$ approximations
have distinct residues relative to one another. They are all the
conjugates. The outer power difference, after division by $w^q$,
has valuation at least $e-(2p-3)t+qM$.
The inner trace difference has this bound plus $q t'$, by the
trace-ideal formula for $B/B_1$. Both resulting product errors have
valuation at least
\begin{equation}\label{O:D:Henselprod}
 e+S+2qM-3qt.
\end{equation}
Here $f(N_1\Delta_\xi)$ has valuation $-pt$, and the unchanged
outer power and inner trace have bounds $q(M-t)$ and $S+qM$.

For the norm argument itself put $a_0=e-(p-2)t$. The quotient
$\epsilon_\xi/(\Delta+\xi)$ has valuation at least $a_0$.
By \eqref{O:D:cyclic},
\[
 v\{N_1\Delta_\xi-N_1(\Delta+\xi)\}
 \ge-pt+\min\{pa_0,a_0+qt'\}=e-qt+S.
\]
The equality for this lower bound uses $a_0\ge t'$, which follows
from $e\ge pq t_2$. The exact identity
\[
 f(X)-f(Y)=(X-Y)\frac{1-XY/W^2}
 {(1+X/W)^2(1+Y/W)^2}
\]
shows that $f$ does not lower this difference valuation. Multiplying
by the outer power and the inner trace gives the further bound
$e+2S+2qM-2qt$.
After $e\ge pq(t+\delta)$ and $M\ge t+1$, the excesses of the two
bounds over $K$ are at least
\[
 (p^2-3p+1)t+p(2p-3)\delta+2p-3,
\]
\[
 p(p-2)t+p(3p-4)\delta+2p-3,
\]
respectively. Both are positive already at $p=3$. Thus the entire
Hensel error is controlled, without deleting a factor $3$ formally.
\end{proof}

\begin{lemma}[reduction to a single finite rational sum]\label{O:D:onesum}
One has
\begin{equation}\label{O:D:onesumeq}
 \mathcal Q\equiv\frac{T}{w^q}
     \sum_{\xi\in\mathcal T}(\Delta+\xi)^q f(b+\xi).
\end{equation}
\end{lemma}
\begin{proof}
After Lemma~\ref{O:D:conjugates}, expand the translated inner trace as
\[
 S_1((\Delta+\xi)^q/w^q)
 =\sum_{j=0}^q\binom qj\xi^{q-j}\mathcal U_j,
 \qquad \mathcal U_j=S_1(\Delta^j/w^q).
\]
For $j\le p-3$, the product valuation is at least
\[
 pq t_2+2qM-(2p-1+pj)t\ge K.
\]
The weakest case $j=p-3$ has excess at least
$(p-1)t+p(p-2)\delta+2p-3$.
For $j=p-2$, its scalar is a multiple of $\xi$.
Subtract the constant product $f(b)\Delta^q/w^q$, whose weighted
sum is zero because $\sum_{\xi\ne0}\xi=0$.
The outer-power difference has valuation at least $qM-(q-1)t$;
also $f(N_1(\Delta+\xi))-f(b)$ has valuation at least zero.
Using \eqref{O:D:Tbound}, the two errors are bounded below by
\[
 pq t_2+2qM-(p^2-2)t,\qquad
 pq t_2+2qM-(p^2-p-1)t.
\]
Their excesses are at least $p(p-2)\delta+2p-3$ and
$(p-1)t+p(p-2)\delta+2p-3$. Thus only $\mathcal U_q=T$ remains.

The exact norm polynomial now reads
\[
 N_1(\Delta+\xi)=b+\xi-\xi s+r_\xi,
 \qquad v(r_\xi)\ge pt_2.
\]
First delete $r_\xi$ using the exact difference formula for $f$.
Next expand in $-\xi s$. Expansion of the inverse square as a
geometric power series shows that its Taylor coefficients of degree
$j\ge2$ have valuation at least $(j-1)pM$; no division by a
factorial divisible by $p$ is used. The series converges because
$v(s/W)>0$. Finally replace
$f'(b+\xi)$ by $f'(b)$; their difference has valuation at least $pM$.
After multiplication by the power and trace factors, the respective
error bounds are
\[
 S+2qM-qt+pt_2,\quad
 3S+(3p-2)M-qt,\quad
 2S+(3p-2)M-qt.
\]
At $M=t+1$, their excesses over $K$ are, respectively,
\[
 qt+S+2p-3,\quad qt+p(3p-4)\delta+3p-3,\quad
 qt+p(2p-3)\delta+3p-3,
\]
all positive. Increasing $M$ only improves them.
The only remaining signed term contains
\[
 \sum_{\xi\ne0}\xi(\Delta+\xi)^q=q^2\Delta.
\]
Its valuation is at least $2S+2qM-t$, whose excess is at least
$(p-3)t+p(2p-3)\delta+2p-3\ge0$.
This proves \eqref{O:D:onesumeq}.
\end{proof}

\subsubsection*{Exact differentiation replaces the long quadratic sum}
We differentiate the exact rational identity of
Lemma~\ref{O:D:interpolation}; no error congruence is differentiated.

\begin{theorem}[quadratic trace at the full modulus]\label{O:D:quadratic}
With the genuine $B_2$ element $\mathcal Q$ of \eqref{O:D:WPdef},
\begin{equation}\label{O:D:quadraticcong}
 \mathcal Q\equiv(p-1)\frac{aT}{w^{p-1}C^2}
                      \pmod{\PP_B^{1+pt_2}}.
\end{equation}
The right side is an expression in $B$; a $B_2$ representative is
constructed in Subsection~\ref{O:sec:transition} before its character is evaluated.
\end{theorem}
\begin{proof}
Put $L=W+b$, $A_* =\Delta-L$, and $D_L=L^p-L$.
Thus $v(L)=v(A_*)=-pM$ and $v(D_L)=-p^2M$.
Apply Lemma~\ref{O:D:interpolation} at $i=q$ and differentiate its
\emph{exact rational identity} in the variable $Z$, holding $W,b,\Delta$
constant. Writing
\[
 H(Z)=p\Delta^q/Z+q(\Delta-Z)^q/(Z^p-Z),
\]
we obtain the exact equality
\[
 \sum_\xi(\Delta+\xi)^q f(b+\xi)=W^2\{H(L)+WH'(L)\}.
\]
After collecting terms this is
\begin{align}\label{O:D:derivative}
 &\frac{qW^2A_*^{p-2}(\Delta-b)}{D_L}
 +\frac{pW^2\Delta^q b}{L^2}
 -\frac{pqW^3A_*^{p-2}}{D_L}\notag\\
 &\hspace{18mm}
 -\frac{qW^3A_*^q(pL^{p-1}-1)}{D_L^2}.
\end{align}
This calculation differentiates no congruence and divides by no $p$.
The valuations of the last three terms, splitting the last one at
$pL^{p-1}-1$, are bounded below by
\[
 e-(2p-1)t,\qquad e-pM,\qquad e-pM,\qquad p(p-2)M.
\]
The factor $T/w^q$ has valuation at least $S+2qM$.
Consequently their products have the three lower bounds
\[
 e+S+2qM-(2p-1)t,\quad e+S+(p-2)M,
 \quad S+(p^2-2)M.
\]
At $e=pq(t+\delta)$, $M=t+1$, their excesses over $K$ are,
respectively,
\begin{align*}
 &(p^2-2p-1)t+p(2p-3)\delta+2p-3,\\
 &(p^2-p-2)t+p(2p-3)\delta+p-3,\\
 &(p^2-p-2)t+p(p-2)\delta+p^2-3.
\end{align*}
All are nonnegative for $p\ge3$; increasing $e$ or $M$ preserves this.
In equal characteristic the $p$-divisible terms are simply absent.

It remains to treat the first term of \eqref{O:D:derivative}.
By \eqref{O:D:omega}, $b-\Delta=a(1+\omega)$ and
$v(a\omega)\ge p\delta-t\ge-t$.
Both $A_*/[-(W+a)]$ and $L/(W+a)$ belong to
$U_B^{pM-t}$, and
$D_L/L^p=1-L^{1-p}\in U_B^{p(p-1)M}$.
Since $p-2$ is odd,
\[
 \frac{W^2A_*^{p-2}}{D_L}
 =-C^{-2}(1+\eta_0),\qquad v(\eta_0)\ge pM-t.
\]
The first term is therefore $q a C^{-2}(1+\omega)(1+\eta_0)$.
After multiplication by $T/w^q$, the two replacement errors have
lower bounds
\[
 S+2qM+pt_2-(p+1)t,\qquad
 S+(3p-2)M-(p+1)t.
\]
Their excesses over $K$, at $M=t+1$, are
$(p-3)t+pq\delta+2p-3$ and
$(p-3)t+p(p-2)\delta+3p-3$.
The product error is deeper. Combining this with Lemma~\ref{O:D:onesum}
proves \eqref{O:D:quadraticcong}, including the full $p=3$ boundary.
\end{proof}


\subsection{\texorpdfstring{The first factor $Q_1$}{The first factor Q1}}\label{O:sec:q1}
\subsubsection*{Actual characters and the norm quotient}
Let $B/A$ be a totally ramified $C_p^2$ diamond, $p$ odd, with
$t_1=t$, $t_2=t+\delta$, $t'=t+p\delta$, and $p\nmid t$.
Write $N_i=N_{B/B_i}$, $n_i=N_{B_i/A}$, $S_i=\Tr_{B/B_i}$.
Let
\[
 \Delta^p-\Delta=a\in B_2,\quad v_B\Delta=-t,\quad b=N_1\Delta,
 \quad s=-E_{p-1}^{B/B_1}(\Delta).
\]
The actual minimal characters $\chi,\varphi$ are compatible and have full
coefficients $\alpha b,\alpha a$, with conductors
$m_1=1+t+t'$, $m_2=1+t+t_2$. Put $\Psi_L(x)=\mathrm e_L(\alpha x)$;
its largest trivial ideals over $A,B_1,B_2$ have exponents
$1+t_2,1+t',1+t_2$, respectively.
Let
\[
 u=\alpha/\gamma,
 \quad W=1/u=n_2w,\quad w=w_2\in B_2,\quad v_Au=m,\quad pm>t,
 \quad M=pm.
\]
The actual factor is
\[
 Q_1=\frac{\chi(\gamma+\alpha b)}
           {\varphi(\gamma-\alpha w+\alpha a)}.
\]

\begin{lemma}[entry from the actual factors]\label{O:Q:entry}
For $z=\Delta/w$ and $n=N_1z=ub$, one has
\[
 Q_1^{-1}=\chi\!\left(\frac{N_1(1+z)}{1+n}\right).
\]
\end{lemma}
\begin{proof}
The exact Artin--Schreier polynomial gives
\[
 N_2(1+z)=1-w^{1-p}+a/w^p.
\]
The relative difference between this and
$1-uw+ua$ is controlled by
\[
 (1-uw+ua)-N_2(1+z)
 =(uw^p-1)\frac{a-w}{w^p}.
\]
The cyclic norm--power estimate gives
$v_2(uw^p-1)\ge(p-1)t_2$. Hence the difference has valuation at least
\[
 (p-1)t_2+pm-\max(t,m)\ge1+t+t_2=m_2.
\]
For $m\le t$, subtracting $m_2$ leaves
$(p-2)t_2+pm-2t-1\ge(p-3)t$.
For $m>t$ it leaves $(p-2)t_2+(p-1)m-t-1\ge0$.
All compared elements are units, so their $\varphi$-values agree.
Compatibility gives $\varphi N_2(1+z)=\chi N_1(1+z)$.
Finally $\chi(\gamma)=\varphi(\gamma)$ by the proved common
$A^\times$ restriction, and $N_1z=ub$ exactly. This proves the formula.
\end{proof}

\subsubsection*{The full logarithm with its actual signs}
Put
\[
 e_i=E_i^{B/B_1}(z),\quad \mathcal U_i=S_1(z^i),\quad T=\mathcal U_{p-1},\quad
 E=\frac{\sum_{i=1}^{p-1}e_i}{1+n},
\]
\[
 W_i=S_2\!\left(\frac{b z^i}{1+n}\right),\qquad
 \mathcal Q=S_2\!\left(\frac{b z^{p-1}T}{(1+n)^2}\right).
\]

\begin{proposition}[the corrected logarithmic entry]\label{O:Q:logentry}
One has
\begin{equation}\label{O:Q:logentry-formula}
 Q_1^{-1}=\Psi_2\!\left(
       \sum_{i=1}^{p-1}\frac{(-1)^i}{i}W_i
       +\frac{\mathcal Q}{2(p-1)^2}\right).
\end{equation}
Every discarded term lies in the \emph{full} character kernel, before
any trace to $B_2$ is taken.
\end{proposition}
\begin{proof}
Lemma~\ref{O:A:lem:twisted} gives
\[
 v_1(e_i)\ge(p-1)t_2+i(m-t).
\]
Write $m_0=\lceil t/p\rceil$; then $m\ge m_0$ because $p\nmid t$.
All the indicated bounds are increasing in $m$, and at $m=m_0\le t$
their minimum is
$\nu_0=(p-1)(\delta+m_0)$.
Writing $t=pa+b_0$, $1\le b_0\le p-1$, checks directly that
\[
 \nu_0\ge\delta+\ceil{(2t+1)/p}=\ceil{m_1/p},\qquad
 3\nu_0\ge m_1.
\]
Thus the full $P$-formula for $\chi$ applies to $1+E$, and
\[
 \chi(1+E)=\Psi_1\bigl(b(-E+E^2/2)\bigr).
\]
This truncation is at conductor $m_1$, also when $m>t$; it does not use
the incorrect lower bound $(p-1)(\delta+m)$ for every $e_i$ in that range.

In a nonlinear Newton term, or in a product $e_i e_j$ other than
$e_{p-1}^2$, the total symmetric index is at most $2p-3$ and there
are at least two factors. After multiplication by $b$, its $B$-valuation
is at least
\[
 2p(p-1)t_2+(i+j)(M-pt)-pt.
\]
Its worst bound occurs at $M=t+1$ and $i+j=2p-3$; the excess over
$p(1+t')$ is
\[
 (p-3)t+p(p-2)\delta+p-3\ge0.
\]
Terms with additional factors have no smaller bound.
Consequently the only square term which must be retained is the top
square. Newton's identities, dividing by $1,\ldots,p-1$, give
\[
 \sum_{i=1}^{p-1}e_i
 \equiv\sum_{i=1}^{p-1}\frac{(-1)^{i-1}}{i}\mathcal U_i,
 \qquad e_{p-1}^2\equiv T^2/(p-1)^2
\]
inside this weighted character calculation. In the second replacement
the error has at least three positive-depth symmetric factors and is
killed by $3\nu_0\ge m_1$.
The remaining exponent is therefore
\[
 -\frac{b}{1+n}\sum_i\frac{(-1)^{i-1}}i \mathcal U_i
       +\frac{bT^2}{2(p-1)^2(1+n)^2}.
\]
Pull one trace $S_1$ out of each summand, then use
$\Psi_B=\Psi_2S_2$. This is precisely \eqref{O:Q:logentry-formula}.
\end{proof}

\subsubsection*{\texorpdfstring{Consequences for the final $Q_1$ expression}{Consequences for the final Q1 expression}}
Use the independently proved actual-conjugate $W_i$ formula and the
independently proved quadratic trace identity, all at $K=1+pt_2$:
\[
 W_i\equiv
 -(p-1)\frac{s\,{\bf1}_{i=p-2}}{w^{p-2}(1+a/W)^2}
 +(p-1)\frac{W\,{\bf1}_{i=p-1}}{w^{p-1}}
 -(p-1)\frac{W^2(-a-W)^i}{w^i(W^p-W+a^p)},
\]
\[
 \mathcal Q\equiv(p-1)\frac{aT}{w^{p-1}(1+a/W)^2}.
\]
Put
\[
 C=1+a/W,\qquad
 F_i=\frac{W}{W^p-W+a^p}\left(\frac{W+a}{w}\right)^i,
\]
\[
 R=W\left(-F_{p-1}+\sum_{i=1}^{p-2}\frac{F_i}{i}+w^{1-p}\right).
\]
Let $\mathscr L_1\in B_2$ denote the actual exponent in
\eqref{O:Q:logentry-formula}. Substitution gives the following congruence
\emph{in $B$}, not a character evaluated outside its domain:
\[
 \mathscr L_1\equiv R+
          \frac{s}{2w^{p-2}C^2}-\frac{aT}{2w^{p-1}C^2}
          \pmod{\PP_B^K}.
\]
With the literal projections $y_s,y_T\in B_2$ constructed in
Lemma~\ref{O:T:representatives}, this implies the well-typed identity
\[
 Q_1^{-1}=\Psi_2\bigl(R+(y_s-y_T)/2\bigr).
\]
Indeed $(-1)^i$ cancels the $(-1)^i$ in $(-a-W)^i$.
For $i<p-1$ one replaces $-(p-1)/i$ by $1/i$ only after its factor
$p$ is weighted; $i=p-1$ gives exactly $-F_{p-1}$ and is not changed.
The exceptional coefficient $(p-1)/(p-2)$ becomes $1/2$ modulo its
weighted kernel, and $1/[2(p-1)]$ becomes $-1/2$ for the quadratic term.
For $i\le p-2$, direct valuation gives
\[
 v_2(F_i)=(p-1)(p-i)m,\qquad
 v_B(pWF_i)\ge v_B(p)+(p-2)M\ge K.
\]
The coefficient of $F_{p-1}$ was never changed. The errors from the two
remaining rational-coefficient replacements are weighted by $pX_s$ and
$pX_T$; their valuations are at least $K$ by the displayed positive source
depths in Lemma~\ref{O:T:representatives}. These terms are zero in equal
characteristic. Thus the formula is valid for every $m$ with $pm>t$.
The $B_2$ representatives invoked here are constructed in
Subsection~\ref{O:sec:transition}; until then these are congruences in $B$.


\clearpage
\subsection{\texorpdfstring{The second factor $Q_2$}{The second factor Q2}}\label{O:sec:q2}
\subsubsection*{An elementary cyclic norm-quotient lemma}
Let $E/F$ be totally ramified cyclic of odd prime degree $p$, with break
$t\ge1$ and different exponent $D=(p-1)(t+1)$. Let $\lambda$ have conductor
$n\ge2$ and full coefficient $\gamma$ for a fixed additive character
$\psi_F$:
\[
 \lambda(1-z)=\psi_F(\gamma Pz),\qquad
 P(Z)=\sum_{j=1}^{p-1} Z^j/j,\qquad z\in\pp_F^{\ceil{n/p}}.
\]
Thus $v_F\gamma=-n_F(\psi_F)-n$ and
$\psi_E=\psi_F\Tr_{E/F}$.

\begin{lemma}\label{O:N:normquot}
Let $k\in\PP_E^a$ with $a\ge1$, and put
\[
 b=\floor{\frac{a+D}{p}}.
\]
If $2b\ge n$, then
\begin{equation}\label{O:N:normquot-formula}
 \lambda\!\left(\frac{N_{E/F}(1-k)}{1-N_{E/F}k}\right)
 =\psi_E\!\left(\frac{\gamma}{1-N_{E/F}k}P(k)\right).
\end{equation}
No formula for $\lambda$ on either factor separately is assumed.
\end{lemma}
\begin{proof}
Write $e_i=E_i^{E/F}(k)$ and $q_j=\Tr_{E/F}(k^j)$.
The exact norm expansion is
\[
 \frac{N(1-k)}{1-Nk}
 =1+Q,\qquad Q=\frac{\sum_{i=1}^{p-1}(-1)^i e_i}{1-Nk}.
\]
The denominator is a unit, since $v_F(Nk)=v_Ek\ge a>0$.
Ueda's strong symmetric estimate gives $v_F(e_i)\ge b$ for every
$1\le i<p$, and hence $v_FQ\ge b$.
As $2b\ge n$, one has $b\ge\ceil{n/2}\ge\ceil{n/p}$, and all
terms of degree at least two in $P(-Q)$ are killed at conductor $n$.
Consequently
\[
 \lambda(1+Q)=\psi_F(-\gamma Q).
\]
Newton's identities, dividing only by integers $j<p$, give
\[
 e_j\equiv(-1)^{j-1}q_j/j\pmod{\pp_F^{2b}},\qquad 1\le j<p.
\]
Indeed $q_i$ also has valuation at least $b$, by the trace-ideal formula,
and every non-leading term in the Newton recursion contains a product of
two such positive-depth factors. Therefore
\[
 -Q\equiv\frac{\sum_{j=1}^{p-1}q_j/j}{1-Nk}
       =\frac{\Tr_{E/F}P(k)}{1-Nk}\pmod{\pp_F^n}.
\]
Multiplication by $\gamma$ and trace compatibility prove
\eqref{O:N:normquot-formula}. The proof is valid in both characteristics,
including $p=3$.
\end{proof}

\subsubsection*{\texorpdfstring{The literal factorization of $Q_2$}{The literal factorization of Q2}}
Use the totally ramified notation $t_1=t$, $t_2=t+\delta$, $r=t_2+m$,
$pm>t$, $u=\alpha/\gamma$, $W=1/u$, $w=w_2$, $n_2w=1/u$,
$a=N_2\Delta$, $b_0=N_1\Delta$ and $H=n_1(s)$, with
$s=-E_{p-1}^{B/B_1}(\Delta)$. Put
\[
 C=1+ua,\quad P_*=n_2C,\quad q_0=u^{p-1},\quad
 Y=P_*-q_0,\quad Z=\frac{n_1(1+ub_0)}{Y},\quad
 k=\frac{uw}{C},\quad F_*=\frac{P_*}{Y}.
\]
All of $C,P_*,Y,Z$ are units, and
\[
 n_2k=q_0/P_*,\quad v_2k=(p-1)m,\quad F_*=(1-n_2k)^{-1}.
\]
The actual factor is
\[
 Q_2=\lambda\!\left(\frac{n_1(\gamma+\alpha b_0)}
                           {n_2(\gamma-\alpha w+\alpha a)}\right).
\]
Cancellation of the common $\gamma^p$, followed by multiplicativity,
gives the \emph{exact} identity
\begin{equation}\label{O:N:exactfactor}
 Q_2^{-1}=\lambda(Z^{-1})\,
       \lambda\!\left(\frac{n_2(1-k)}{1-n_2k}\right).
\end{equation}
In particular this does not require evaluation of $\lambda n_2$ at $1-k$.

\subsubsection*{\texorpdfstring{The moment comparison and the exact $Z$ modulus}{The moment comparison and the exact Z modulus}}
Set $A_0=n_2a$, $D_0=1+u^pA_0$. Write $E_i^{(1)},E_i^{(2)}$ for the
lower elementary symmetric functions of $b_0,a$, and put
$F_i=u^i(E_i^{(1)}-E_i^{(2)})$.

\begin{lemma}[weighted moments]\label{O:N:moments}
For $1\le j\le p-1$,
\[
 u^j[T_1(b_0^j)-T_2(a^j)]
 \equiv(p-1){\bf1}_{j=p-1}u^{p-1}(1-s^p)
             \pmod{\PP_1^{1+pr}}.
\]
Consequently
\[
 F_i\equiv-{\bf1}_{i=p-1}u^{p-1}(1-s^p)
             \pmod{\PP_1^{1+pr}},\qquad 1\le i<p.
\]
\end{lemma}
\begin{proof}
For $j=1$ Theorem~\ref{O:A:thm:H14}, multiplied by $u$, gives the
required modulus. For $j\ge2$, apply Theorem~\ref{O:A:thm:h} to
$w^{-j}\Delta^{j-1}$, whose $B$-valuation is $jpm-(j-1)t>0$.
Its error exponent $pt_2+jpm-(j-1)t$ is at least $1+pr$, since their
difference is $(j-1)(pm-t)-1\ge0$.
The exact discrepancy from multiplicative homogeneity is
\[
 T_2\bigl((u^j-w^{-pj})a^j\bigr).
\]
The lower cyclic norm--power estimate gives
$v_2(u^j-w^{-pj})\ge jpm+(p-1)t_2$.
After adding the different, the trace is in $\pp_A^{r+1}$, because
\[
 (j-1)pm+(p-2)t_2-jt-1
 \ge (p-2)t_2-t+j-2\ge0.
\]
This proves the moment congruence, including its exceptional term.
Apply Newton's identities to the two integral tuples $ub_0$ and $ua$.
Induction over $i<p$ introduces only integral factors and division by
$i$, a unit. At $i=p-1$ its leading sign is negative, so division of
$(p-1)u^{p-1}(1-s^p)$ by $-(p-1)$ gives the displayed $F_{p-1}$.
All other $F_i$ vanish at the same modulus.
\end{proof}

\begin{proposition}[the full $Z$ congruence]\label{O:N:Z}
One has
\[
 Z\equiv1+u^{p-1}\frac H{D_0}\pmod{\pp_A^{r+1}}.
\]
\end{proposition}
\begin{proof}
The norm endpoints of $ub_0$ and $ua$ agree by norm transitivity.
Summing Lemma~\ref{O:N:moments} over the remaining symmetric degrees gives
\[
 n_1(1+ub_0)-Y\equiv u^{p-1}s^p
             \pmod{\PP_1^{1+pr}}.
\]
There is no need to transfer $(1-s)^p$ between fields.
The ordinary lower cyclic norm--power estimate gives directly
\[
 v_1\bigl(u^{p-1}(s^p-H)\bigr)
 \ge p(p-1)m+p(p-1)\delta+(p-1)t\ge1+pr,
\]
because the excess is $p(p-2)(m+\delta)-t-1\ge0$.
Both sides after this substitution belong to $A$, so intersection with
$A$ improves the modulus to $\pp_A^{r+1}$.

It remains to replace the unit denominator $Y$ by $D_0$ with its
actual weight. One has
\[
 Y-D_0=\sum_{i=1}^{p-1}E_i^{B_2/A}(ua)-u^{p-1},
 \quad
 v_AE_i(ua)\ge L_0=
 \floor{\frac{pm-t+(p-1)(t_2+1)}p}.
\]
The weighted intermediate terms have valuation at least
$(p-1)m+(p-1)\delta+L_0\ge r+1$.
For completeness write $t=pa+b$, $1\le b\le p-1$, and use $m\ge a+1$.
At $\delta=0$ the excess is at least
\[
 (p-3)a+p-2-b+\floor{\frac{(p-2)b+p-1}{p}}\ge0.
\]
For $b\le p-2$ every indicated term is nonnegative; for $b=p-1$
the final floor is $p-2$, giving again a nonnegative result.
Increasing $\delta$ or $m$ cannot decrease the excess.
The weighted final $-u^{p-1}$ term has excess
\[
 (2p-3)m+(p-2)\delta-t-1\ge pm-t-1\ge0.
\]
All denominators are units. These estimates prove
$u^{p-1}H(Y^{-1}-D_0^{-1})\in\pp_A^{r+1}$ and hence the proposition.
\end{proof}

\begin{theorem}[the actual $Q_2$ phase]\label{O:N:Q2}
The full actual factor is
\[
 Q_2^{-1}=\Psi_A\!\left(Wu^{p-1}\frac{H}{D_0}\right)
          \Psi_2\!\left(WF_*P(k)\right),\qquad
 \Psi_L(x)=\mathrm e_L(\alpha x).
\]
\end{theorem}
\begin{proof}
The conductor of $\lambda$ is $n=r+1=t_2+m+1$.
For the second factor of \eqref{O:N:exactfactor}, Lemma~\ref{O:N:normquot}
has
\[
 a=(p-1)m,\quad D=(p-1)(t_2+1),\quad
 b=\floor{\frac{(p-1)n}{p}}.
\]
For every $n\ge2$ and odd $p$, $2\lfloor(p-1)n/p\rfloor\ge n$;
this is $n\ge2\lceil n/p\rceil$, checked at $p=3$ and then monotone
in $p$. The lemma therefore gives the second displayed phase.

For the first factor, $v_AH\ge(p-1)\delta$. Its small argument has
valuation at least $L=(p-1)(m+\delta)$, and
\[
 2L-n=(2p-3)m+(2p-3)\delta-t-1\ge pm-t-1\ge0.
\]
Hence $\lambda((1+x)^{-1})=\mathrm e_A(\gamma x)$ for this argument.
Proposition~\ref{O:N:Z}, at the full conductor, gives the first phase.
These are two different depth estimates: the floor $b$ controls the norm
quotient, while $L$ controls $Z^{-1}$. No ceiling strengthening or shallow
individual-factor linearization has been used.
\end{proof}


\subsection{The weighted first terminal estimate}\label{O:sec:weighted}
\subsubsection*{The weight that must be checked}
Use a totally ramified diamond $A\subset B_1,B_2\subset B$ with
$\operatorname{Gal}(B/A)=C_p^2$, $p$ odd. Put
\[
 t=t_1,\quad\delta=t_2-t\ge0,\quad p\nmid t,\quad
 N_i=\Nm_{B/B_i},\quad n_i=\Nm_{B_i/A},\quad S_i=\Tr_{B/B_i}.
\]
Let $\Delta^p-\Delta=a\in B_2$, $v_B\Delta=-t$, and set
\[
 s=-E_{p-1}^{B/B_1}(\Delta),\quad
 Y=w_2^{-1}\in B_2,\quad v_2(Y)=m,\quad n_2Y=u,\quad
 M=pm>t.
\]
Thus $v_B(Y)=M$, $v_A(u)=m$. Define
\[
 T=S_1((Y\Delta)^{p-1}),\quad H=n_1(s),\quad J=n_1(T),
 \quad A_0=n_2(a),\quad v_A(A_0)=-t.
\]
The first bracket in the terminal phase is multiplied by
$\alpha u^{p-2}A_0/(2D^2)$, where $D$ is a unit and
$\alpha$ gives ideal conductor $1+t_2$ on $A$.
The actual assertion needed is therefore
\begin{equation}\label{O:E:weighted}
 L_A:=u^{p-1}A_0(J+u^{p-1}H)\in\pp_A^{1+t_2}.
\end{equation}
The estimate will be proved with its weight already inserted; an
unweighted congruence is not used as a substitute for \eqref{O:E:weighted}.

\subsubsection*{A semilinear trace estimate}
Write $v=v_B$ in this subsection, and put $z=Y\Delta$.
Lemma~\ref{O:A:lem:twisted} gives
\begin{align}
 v(E_i\Delta)&\ge p(p-1)t_2-pit,\label{O:E:H15a}\\
 v(E_i z)&\ge p(p-1)t_2+i(M-pt),\qquad1\le i<p.\label{O:E:H15b}
\end{align}
The cyclic norm--power estimate, independently obtained from the
strong symmetric bound and the characteristic polynomial, is
\begin{equation}\label{O:E:normpower}
 v_E(\Nm_{E/F}x-x^p)\ge p v_E(x)+(p-1)\operatorname{break}(E/F).
\end{equation}
It holds also for negative valuations and in both characteristics.
Set
\[
 R=(p-1)M+p(p-1)\delta,\qquad
 B_0=R+(p-1)t.
\]
\begin{lemma}\label{O:E:trace}
The actual signed trace has the approximation
\begin{equation}\label{O:E:semilinear}
 v_B\bigl(T+Y^{p-1}s\bigr)\ge B_0.
\end{equation}
Moreover $v_B(T),v_B(Y^{p-1}s)\ge R$.
\end{lemma}
\begin{proof}
The last assertion is the trace estimate in Lemma~\ref{O:A:lem:twisted} and \eqref{O:E:H15a}.
Subtract $Y^p$ times the characteristic-polynomial identity for
$\Delta$ from the identity for $z$. The result is
\[
 \sum_{i=1}^{p-1}(-1)^i
  (E_i z-Y^iE_i\Delta)z^{p-i}
       -N_1\Delta\,(n_2Y-Y^p)=0.
\]
For $i\le p-2$, \eqref{O:E:H15a}--\eqref{O:E:H15b} bound the corresponding
summand below by
\[
 pM+p(p-1)t_2-[p+(p-1)i]t.
\]
Equation \eqref{O:E:normpower} for $B_2/A$ bounds the last summand below
by $pM+p(p-1)t_2-pt$. Solving for the $i=p-1$ summand and dividing
by $z$, of valuation $M-t>0$, gives
\begin{equation}\label{O:E:edifference}
 v_B(E_{p-1}z-Y^{p-1}E_{p-1}\Delta)
 \ge(p-1)M+p(p-1)t_2-(p-1)^2t=B_0.
\end{equation}
No division by $p$ occurs.

Newton's identity for degree $p-1$ says
\[
 S_1(z^{p-1})=-(p-1)E_{p-1}z+\mathcal R(E_1z,\ldots,E_{p-2}z),
\]
where $\mathcal R$ is an integral polynomial whose monomials have
weight $p-1$ and at least two elementary-symmetric factors.
By \eqref{O:E:H15b}, every such monomial has valuation at least
\[
 (p-1)M+p(p-1)(t+2\delta)\ge B_0.
\]
Using $E_{p-1}\Delta=-s$ and \eqref{O:E:edifference}, we obtain
$T\equiv(p-1)Y^{p-1}s\pmod{\PP_B^{B_0}}$.
In mixed characteristic
$v_B(p)\ge p(p-1)t_2$, so
\[
 v_B(pY^{p-1}s)\ge(p-1)M+p(p-1)(t+2\delta)\ge B_0.
\]
In equal characteristic this term is zero. Adding $Y^{p-1}s$ now
proves \eqref{O:E:semilinear}, with its plus sign forced by $s=-E_{p-1}$.
\end{proof}

\subsubsection*{Insert the weight before taking powers}
\begin{lemma}\label{O:E:power}
In $B_1$, put
\[
 L_1=u^{p-1}A_0\bigl(T^p+u^{p-1}s^p\bigr).
\]
Then
\begin{equation}\label{O:E:L1bound}
 v_B(L_1)\ge V:=2p(p-1)M-pt+p^2(p-1)\delta.
\end{equation}
\end{lemma}
\begin{proof}
Let $Z=Y^{p-1}s$ and $E=T+Z$. The identity
\[
 T^p+Z^p=E^p-\sum_{i=1}^{p-1}\binom pi T^iZ^{p-i}
\]
is exact. Lemma~\ref{O:E:trace} gives $v(E)\ge B_0$ and
$v(T),v(Z)\ge R$. The weight has valuation
\[
 v_B(u^{p-1}A_0)=p(p-1)M-p^2t.
\]
Its product with $E^p$ therefore has valuation at least $V$.
Every mixed binomial term has valuation at least
\[
 p(p-1)M-p^2t+v_B(p)+pR\ge V.
\]
In equal characteristic these mixed terms vanish.

It remains to replace $Z^p$ by $u^{p-1}s^p$.
Equation \eqref{O:E:normpower}, applied to $Y\in B_2$, gives
\[
 v_B(u/Y^p-1)\ge p(p-1)t_2,
\]
so
\[
 v_B(u^{p-1}-Y^{p(p-1)})
   \ge p(p-1)M+p(p-1)t_2.
\]
After multiplying by $u^{p-1}A_0s^p$, the bound is again at least
$V$ by \eqref{O:E:H15a}. This proves \eqref{O:E:L1bound}.
\end{proof}

\begin{theorem}[strong weighted bound]\label{O:E:main}
The actual $A$-valued expression in \eqref{O:E:weighted} satisfies
\begin{equation}\label{O:E:Abound}
 v_A(L_A)\ge2(p-1)m+(p-1)\delta-\floor{t/p}\ge1+t_2.
\end{equation}
Thus the fully weighted first terminal phase is trivial.
\end{theorem}
\begin{proof}
The symmetric-coefficient bounds give
$v_1(s)\ge(p-1)\delta$ and $v_1(T)\ge(p-1)(m+\delta)$.
Using \eqref{O:E:normpower} for $B_1/A$,
\begin{align*}
 v_1\bigl(u^{p-1}A_0(T^p-J)\bigr)
   &\ge2(p-1)M-t+p(p-1)\delta,\\
 v_1\bigl(u^{2(p-1)}A_0(s^p-H)\bigr)
   &\ge2(p-1)M-t+p(p-1)\delta.
\end{align*}
Both lower bounds are $V/p$. Lemma~\ref{O:E:power} therefore gives
$v_1(L_A)\ge V/p$.
Since $L_A\in A$, divide by $p$ again and use integrality of $v_A$:
\[
 v_A(L_A)\ge\ceil{V/p^2}
  =2(p-1)m+(p-1)\delta-\floor{t/p}.
\]
Here $\ceil{V/p^2}$ denotes the ceiling, as usual.
Write $t=pa+b$, $a\ge0$, $1\le b\le p-1$; then $m\ge a+1$.
The right side minus $1+t_2$ is at least
\[
 (p-3)a+2p-3-b+(p-2)\delta\ge p-2\ge1.
\]
This proves \eqref{O:E:Abound}. Multiplication by a unit denominator
$D^{-2}$ and the rational scalar $1/2$ preserves this ideal bound;
then multiplication by $\alpha$ reaches the actual additive conductor.
\end{proof}


\subsection{The four-factor identity and transition cancellation}\label{O:sec:transition}
\subsubsection*{The stationary data}
Use the notation of Subsections~\ref{O:sec:quadratic} and \ref{O:sec:q1}:
\[
 t=t_1,\quad t_2=t+\delta,\quad t'=t+p\delta,\quad p\nmid t,
 \quad \delta\ge0,\quad p\text{ odd},
\]
\[
 N_i=\Nm_{B/B_i},\ n_i=\Nm_{B_i/A},\ S_i=\Tr_{B/B_i},\
 T_i=\Tr_{B_i/A},\quad \Delta^p-\Delta=a\in B_2,
\]
\[
 b=N_1\Delta,\quad N_2\Delta=a,\quad v_B\Delta=-t,
 \quad s=-E_{p-1}^{B/B_1}(\Delta).
\]
The actual minimal characters $\chi,\varphi$ have equal norm pullbacks
and conductors $1+t+t'$ and $1+t+t_2$. They are the actual models of
Subsection~\ref{O:sec:models}, not auxiliary characters assumed to have the right
stationary coefficients. Put $\theta_1=\chi(\lambda n_1)$ and
$\theta_2=\varphi(\lambda n_2)$, where $a(\lambda)=1+r$,
$r=t_2+m$, $M=pm>t$. The exact choices of Subsection~\ref{O:sec:parameters} give
\[
 u=\alpha/\gamma,\quad W=u^{-1},\quad v=\beta/\alpha,
 \quad v_Au=m,\quad v_Av=\delta,
\]
\[
 n_2w=\gamma/\alpha=W,\quad v_2w=-m,\qquad
 n_1w_1=\gamma/\beta,\quad v_1w_1=-(m+\delta).
\]
Let $\Psi_E(x)=\ee_E(\alpha x)$, trace compatibly. Its largest trivial
ideal on $A$ has exponent $1+t_2$, and those on $B_1,B_2$ have
exponents $1+t',1+t_2$. The largest trivial ideal for
$x\mapsto\Psi_{B_2}(Wx)$ has exponent
\begin{equation}\label{O:T:Gammacond}
 n_\Gamma=1+t_2+pm.
\end{equation}
The largest trivial ideal of an additive character is not generally
its entire kernel; no converse valuation inference from a single
character value will be used.

Set
\[
 C=1+ua,\quad A_0=n_2a,\quad D=1+u^pA_0,
 \quad P_*=n_2C,\quad H=n_1s,
\]
\[
 T=S_1((\Delta/w)^{p-1}),\quad J=n_1T,\quad
 S=p(p-1)\delta,\quad K=1+pt_2.
\]
Thus $C,P_*,D$ are units and $v_AA_0=-t$.
The symmetric-coefficient bounds give
\begin{equation}\label{O:T:sT}
 v_1s\ge(p-1)\delta,\qquad v_1T\ge(p-1)(m+\delta).
\end{equation}
We use their coefficient version as well: for $z\in B_1$,
\begin{equation}\label{O:T:coefficient}
 z=\sum_{i=0}^{p-1}c_i\Delta^i,
 \quad c_i\in B_2\quad\Longrightarrow\quad
 v_2c_i\ge v_1z+it.
\end{equation}
The cyclic symmetric estimate is
\begin{equation}\label{O:T:sym}
 v_F(E_j^{E/F}x)\ge
 \floor{\frac{jv_Ex+(p-1)(b_0+1)}p},\quad1\le j<p,
\end{equation}
for an edge of break $b_0$. It implies
\begin{equation}\label{O:T:np}
 v_E(\Nm x-x^p)\ge p v_Ex+(p-1)b_0
\end{equation}
by Lemma~\ref{O:A:lem:power}.
All these estimates hold in both characteristics.

The trace--norm theorem in Section~\ref{O:sec:as} gives
\begin{equation}\label{O:T:H14}
 T_1b-T_2a\in\pp_A^{1+t_2}.
\end{equation}
The actual four factors of Proposition~\ref{O:P:four} are
\begin{align}\label{O:T:factors}
 Q_1&=\chi(A_1)/\varphi(C_2),&
 Q_2&=\lambda(n_1A_1/n_2C_2),\notag\\
 Q_3&=\ee_{B_1}(-\beta w_1)\theta_1(C_1/A_1),&
 Q_4&=\Psi_{B_2}(w),
\end{align}
where $A_1=\gamma+\alpha b$, $C_1=A_1-\beta w_1$ and
$C_2=\gamma-\alpha w+\alpha a$.
The full coefficient of $\theta_1$ is $C_1$, on the entire
$P$-subgroup of depth $\ceil{(1+t'+pm)/p}$.
We use the actual $Q_1$ logarithm and quadratic evaluation proved in
Subsections~\ref{O:sec:quadratic} and \ref{O:sec:q1}, together with the corrected linear $W_i$ formula of Theorem~\ref{O:W:final}.
All character arguments are in their stated fields.

\subsubsection*{Weighted denominator bounds}
\begin{lemma}\label{O:T:denominator}
One has
\begin{align}\label{O:T:weightedden}
 u^{p-2}H(P_*^{-2}-D^{-2})&\in\pp_A^{1+t_2},\notag\\
 u^{p-1}A_0J(P_*^{-2}-D^{-2})&\in\pp_A^{1+t_2}.
\end{align}
The first assertion also holds with inverse first powers instead of
inverse squares.
\end{lemma}
\begin{proof}
The norm expansion and \eqref{O:T:sym} give
\[
 v_A(P_*-D)\ge\ell:=m+
 \floor{\frac{(p-2)t+(p-1)(\delta+1)}p}.
\]
Indeed the $j$th intermediate coefficient has bound
$\floor{[j(pm-t)+(p-1)(t_2+1)]/p}$, nondecreasing in $j$.
Write $t=pa+b_0$, $1\le b_0\le p-1$, and $m=a+1+c$, $c\ge0$.
Put $f=\floor{[(p-2)b_0+(p-1)(\delta+1)]/p}$.
Using \eqref{O:T:sT}, the two weighted valuation excesses over $1+t_2$
are bounded below by
\[
 (p-3)a+(p-1)c+p-2+(p-2)\delta-b_0+f,
\]
\[
 (p-3)a+(2p-1)c+2p-2+(p-2)\delta-2b_0+f.
\]
The first is nonnegative for $b_0\le p-2$; if $b_0=p-1$, use
$f\ge p-2$. The second is nonnegative because $b_0\le p-1$.
All denominators are units; differences of their inverse powers have
valuation at least $\ell$. This proves every assertion.
\end{proof}

\subsubsection*{The rational cancellation for the actual second factor}
Use the notation of the final formula in Subsections~\ref{O:sec:quadratic} and \ref{O:sec:q1}:
\[
 F_i=\frac{W}{W^p-W+a^p}\left(\frac{W+a}{w}\right)^i.
\]
Put
\[
 k=\frac{uw}{C},\qquad q_0=u^{p-1},\qquad
 F_*=\frac{P_*}{P_*-q_0},\qquad P(X)=\sum_{j=1}^{p-1}\frac{X^j}{j}.
\]
These are exactly the variables in Theorem~\ref{O:N:Q2}; in particular
$v_2(k)=(p-1)m$. Put $Q=1+u^pa^p$ in $B_2$, and put
\[
 \eta=W/w^p,\qquad G=\frac{\eta C^p}{Q-q_0}.
\]
Both $Q-q_0$ and $P_*-q_0$ are units. Direct multiplication gives
\begin{equation}\label{O:T:Fidentity}
 F_i=G k^{p-i},\qquad G k^p=\frac{q_0}{Q-q_0}.
\end{equation}
\begin{lemma}[the two rational sums cancel to one term]\label{O:T:rational}
At the full $\gamma$-conductor \eqref{O:T:Gammacond},
\begin{equation}\label{O:T:rationalcong}
 -F_{p-1}+\sum_{i=1}^{p-2}\frac{F_i}{i}+F_*P(k)
 \equiv k(1-\eta)\pmod{\PP_2^{n_\Gamma}}.
\end{equation}
\end{lemma}
\begin{proof}
For $2\le j\le p-1$,
$1/(p-j)+1/j=p/[j(p-j)]$.
Thus \eqref{O:T:Fidentity} changes the first two terms on the left into
$-GP(k)$, with errors divisible by $pk^j$ for $j\ge2$.
In mixed characteristic their valuation is at least
$(p-1)t_2+2(p-1)m$, whose excess over $n_\Gamma$ is
$(p-2)(t_2+m)-1\ge0$. In equal characteristic the errors vanish.
The coefficient of $F_{p-1}$ has not been changed independently.

The norm--power estimate and the binomial theorem give
\begin{align*}
 v_2(P_*-C^p)&\ge(p-1)t_2,\\
 v_2(C^p-Q)&\ge(p-1)t_2+pm-t,\\
 v_2(1-\eta)&\ge(p-1)t_2.
\end{align*}
In particular $P_*-Q$ has depth at least $(p-1)t_2$.
The following identity is exact:
\begin{align}\label{O:T:FGexact}
 F_*-G-(1-\eta)={}&q_0\left(
 \frac1{P_*-q_0}-\frac\eta{Q-q_0}\right)
 -\frac{\eta(C^p-Q)}{Q-q_0}.
\end{align}
The bracket has depth at least $(p-1)t_2$, as does $F_*-G$.
For $j\ge2$, the valuation of $(F_*-G)k^j$ is at least
$(p-1)t_2+2(p-1)m\ge n_\Gamma$.
For $j=1$, the first term on the right of \eqref{O:T:FGexact}, multiplied
by $k$, has depth at least
\[
 (p^2-1)m+(p-1)t_2\ge n_\Gamma.
\]
The second has depth at least
$(p-1)t_2+pm-t+(p-1)m$; its excess is
\[
 (p-2)t_2+(p-1)m-t-1
 =(p-3)t+(p-2)\delta+(p-1)m-1\ge0.
\]
This proves \eqref{O:T:rationalcong} without deleting an unweighted
norm--power discrepancy.
\end{proof}

\begin{lemma}[the remaining elementary rational phase]\label{O:T:theta0}
There is the exact identity in $B_2$
\begin{equation}\label{O:T:Theta0}
 W[k(1-\eta)+w^{-(p-1)}]-w
 =\frac{a}{w^{p-1}C}(1-uw^p)\in\PP_2^{1+t_2}.
\end{equation}
\end{lemma}
\begin{proof}
Use $W u=1$, $\eta=1/(uw^p)$ and $C-1=ua$ to obtain the
identity by subtraction of the two fractions. The norm--power bound
gives $v_2(1-uw^p)\ge(p-1)t_2$, so the right side has valuation
at least $-t+(p-1)m+(p-1)t_2$. Its excess over $1+t_2$ is
$(p-3)t+(p-2)\delta+(p-1)m-1\ge0$.
It is its additive phase, not necessarily the element itself, which
vanishes.
\end{proof}

\subsubsection*{Representatives in $B_2$ and norm transfer}
For $z=\sum_{i=0}^{p-1}z_i\Delta^i$ with $z_i\in B_2$, let
$\pr z=z_0$. Define
\begin{equation}\label{O:T:XandY}
 X_s=\frac{s}{w^{p-2}C^2},\quad X_T=\frac{aT}{w^{p-1}C^2},
 \qquad y_s=\pr X_s,\quad y_T=\pr X_T.
\end{equation}
The first two expressions are in $B$; the last two are actual elements
of $B_2$. No $B_2$ character will be applied to $X_s$ or $X_T$.

\begin{lemma}[constant projection and its norm]\label{O:T:projection}
Let $z=xw_0$ with $x\in B_2$, $w_0\in B_1$. Suppose
\begin{equation}\label{O:T:projectdepth}
 h:=v_Bz\ge1+t',\qquad h\ge1+t_2+(p-3)t.
\end{equation}
Then $y=\pr z\in B_2$ satisfies
\begin{equation}\label{O:T:projectcong}
 z\equiv y\pmod{\PP_B^K},\qquad
 N_1z\equiv n_2y\pmod{\PP_1^K}.
\end{equation}
If $z\ne0$, then $v_By=h$.
\end{lemma}
\begin{proof}
The zero case is immediate. Write $z=\sum y_i\Delta^i$ using
\eqref{O:T:coefficient} for $w_0$. It gives
\[
 v_2y_i\ge h/p+it,\qquad
 v_B(y_i\Delta^i)\ge h+i(p-1)t.
\]
Thus all nonconstant terms are strictly deeper than $h$ and the
constant term has valuation $h$. Since
$1+t'+(p-1)t=K$, all nonconstant terms belong to $\PP_B^K$.
This proves the first congruence, with $y=y_0$.

The second bound in \eqref{O:T:projectdepth} and the congruence just
proved are precisely the hypotheses of Lemma~\ref{O:A:lem:H17}, with
$K_0=K$. Applying that lemma to $z=xw_0$ and $y=y_0$ proves the
norm congruence in \eqref{O:T:projectcong}.
\end{proof}

\begin{lemma}[all source and target depths]\label{O:T:representatives}
The elements in \eqref{O:T:XandY} satisfy
\begin{align}\label{O:T:reps}
 X_s&\equiv y_s\pmod{\PP_B^K},&
 X_T&\equiv y_T\pmod{\PP_B^K},\notag\\
 n_2y_s&\equiv u^{p-2}H/D^2\pmod{\pp_A^{1+t_2}},&
 n_2y_T&\equiv u^{p-1}A_0J/D^2\pmod{\pp_A^{1+t_2}}.
\end{align}
Moreover $y_s,y_T\in\PP_2^{c_2}$, where
$c_2=\ceil{(1+t_2)/p}$.
\end{lemma}
\begin{proof}
By \eqref{O:T:sT},
\[
 h_s:=v_BX_s\ge(p-2)M+S,\qquad
 h_T:=v_BX_T\ge2(p-1)M-pt+S.
\]
At $M\ge t+1$, their excesses over $1+t'$ are at least
\[
 (p-3)(t+1)+p(p-2)\delta,\qquad
 (p-3)t+2p-3+p(p-2)\delta,
\]
and their excesses over $1+t_2+(p-3)t$ are at least
\[
 p-3+[p(p-1)-1]\delta,\qquad
 2p-3+[p(p-1)-1]\delta.
\]
All are nonnegative, including $p=3,\delta=0$.
Lemma~\ref{O:T:projection} therefore applies to both products and proves
the first two assertions of \eqref{O:T:reps}. It also gives
$v_By_s,v_By_T\ge1+t_2$, hence the stated depth $c_2$ in $B_2$.
The norm part of that lemma gives
\[
 n_2y_s\equiv\frac{u^{p-2}s^p}{P_*^2},\qquad
 n_2y_T\equiv\frac{u^{p-1}A_0T^p}{P_*^2}
 \pmod{\PP_1^K}.
\]
Here $N_1$ fixes $B_1$ with norm the $p$th power, and its restriction
to $B_2$ is $n_2$; these two domains have not been interchanged.
By \eqref{O:T:np} on $B_1/A$,
\begin{align*}
 v_1\bigl(u^{p-2}(s^p-H)\bigr)
 &\ge(p-2)M+p(p-1)\delta+(p-1)t,\\
 v_1\bigl(u^{p-1}A_0(T^p-J)\bigr)
 &\ge2(p-1)M+p(p-1)\delta-t.
\end{align*}
The first excess over $K$ is at least
$(p-3)(t+1)+p(p-2)\delta$, and the second at least
$(p-3)t+2p-3+p(p-2)\delta$.
Thus the two powers may be replaced by $H,J$. Now both sides lie
in $A$, so $A\cap\PP_1^K=\pp_A^{1+t_2}$ applies.
Lemma~\ref{O:T:denominator} replaces $P_*^{-2}$ by $D^{-2}$ at precisely
that weighted modulus, proving the remaining assertions.
\end{proof}

\begin{lemma}[actual norm character and the rational half]\label{O:T:actualnorm}
One has
\begin{align}\label{O:T:normhalves}
 \Psi_{B_2}(y_s/2)&=\Psi_A(-u^{p-2}H/(2D^2)),\notag\\
 \Psi_{B_2}(-y_T/2)&=\Psi_A(u^{p-1}A_0J/(2D^2)).
\end{align}
\end{lemma}
\begin{proof}
The actual norm character of $B_2/A$ has full coefficient $\alpha$.
For any $y\in\PP_2^{c_2}$ choose $x$ with $P(x)=y$, using
bijectivity of $P$ on that positive ideal. The cyclic coefficient
estimate places $n_2(1-x)$ in the full $P$-domain of that character.
Its value on this norm is $1$. The norm--logarithm identity of
Lemma~\ref{O:I:normlog}, with break $t_2$ and $h=0$, gives at modulus $\pp_A^{1+t_2}$
\[
 P(1-n_2(1-x))\equiv T_2P(x)+n_2(P(x)).
\]
Consequently $\Psi_{B_2}(y)\Psi_A(n_2y)=1$; this is the actual
norm-character phase of Corollary~\ref{O:I:conversion}, not an auxiliary model equality.
For the half, the product
$\Psi_{B_2}(y/2)\Psi_A(n_2y/2)$ has square $1$ and has $p$-primary
order. Indeed the image of a continuous additive character of a
local field of residue characteristic $p$ is $p$-primary: sufficiently
large $p$-powers of each argument lie in an open trivial ideal.
As $p$ is odd, this product is $1$. This argument does not move
$1/2$ inside the norm, and does not extend phase vanishing under an
arbitrary local-unit multiplier.
Apply it to $y_s,y_T$ and use Lemma~\ref{O:T:representatives}.
The errors after multiplication by $1/2$ still lie in the whole
trivial ideal $\pp_A^{1+t_2}$.
\end{proof}

\begin{proposition}[the actual $Q_1Q_2Q_4$ product]\label{O:T:Q124}
The genuine character values satisfy
\begin{equation}\label{O:T:Q124phase}
 (Q_1Q_2Q_4)^{-1}
 =\Psi_A\left(
 Wu^{p-1}\frac HD-\frac{u^{p-2}H}{2D^2}
                   +\frac{u^{p-1}A_0J}{2D^2}\right).
\end{equation}
\end{proposition}
\begin{proof}
The end of Subsections~\ref{O:sec:quadratic} and \ref{O:sec:q1} expresses the actual $Q_1$ exponent, a member
of $B_2$, modulo $\PP_B^K$ as
\[
 -WF_{p-1}+W\sum_{i=1}^{p-2}F_i/i+W/w^{p-1}
       +\frac{p-1}{p-2}X_s+\frac{1}{2(p-1)}X_T.
\]
Replace $X_s,X_T$ by $y_s,y_T$ using Lemma~\ref{O:T:representatives}.
The resulting expression belongs to $B_2$, so its difference from
the genuine exponent lies in
$B_2\cap\PP_B^K=\PP_2^{1+t_2}$ and has trivial $\Psi_{B_2}$ phase.
This supplies the field-valued representatives required by the
calculation, rather than silently applying $\Psi_{B_2}$ to $B$.

The two coefficient replacements are
\[
 \frac{p-1}{p-2}-\frac12=\frac p{2(p-2)},\qquad
 \frac1{2(p-1)}+\frac12=\frac p{2(p-1)}.
\]
Their errors have common-field valuations at least
$v_B(p)+h_s$ and $v_B(p)+h_T$, respectively, and hence at least $K$.
They vanish identically in equal characteristic. Thus the surviving
exceptional phase is $\Psi_{B_2}((y_s-y_T)/2)$.
Theorem~\ref{O:N:Q2} and Lemma~\ref{O:T:rational} combine the two rational
sums. Their error is at the full $\gamma$ modulus, not only the
$\alpha$ modulus. Including $W/w^{p-1}$ and $Q_4^{-1}=\Psi_{B_2}(-w)$,
Lemma~\ref{O:T:theta0} kills the remaining elementary phase.
The product is therefore exactly
\[
 \Psi_A(Wu^{p-1}H/D)\,\Psi_{B_2}((y_s-y_T)/2).
\]
Finally apply \eqref{O:T:normhalves}. This proves \eqref{O:T:Q124phase}
with every scalar outside the appropriate norm.
\end{proof}

\subsubsection*{The third factor $Q_3$}
Put $h_0=m+\delta$ and $z=\beta w_1/A_1$. Then
\[
 v_1z=(p-1)h_0,\qquad
 n_1^{\mathrm{char}}=1+t'+pm=1+t+ph_0
\]
is the conductor of $\theta_1$.
\begin{lemma}[direct $Q_3$ evaluation]\label{O:T:Q3}
One has
\begin{equation}\label{O:T:Q3phase}
 Q_3^{-1}=\Psi_A\left(-\frac W2\frac{(vu)^{p-1}}D\right).
\end{equation}
\end{lemma}
\begin{proof}
The inequalities $p(p-1)h_0\ge n_1^{\mathrm{char}}$ and
$3(p-1)h_0\ge n_1^{\mathrm{char}}$ follow from $pm\ge t+1$:
for the latter the excess is $(2p-3)h_0-t-1\ge0$, and the former
is no smaller. Thus the full $P$-formula of $\theta_1$, whose
coefficient is $C_1=A_1(1-z)$, applies at $z$ with all terms of
degree at least three in the following expression beyond its conductor:
\[
 Q_3=\ee_{B_1}\bigl(A_1[(1-z)P(z)-z]\bigr)
     =\ee_{B_1}(-A_1z^2/2).
\]
The identity used here is exact before deleting the higher terms;
for $p=3$ the cubic term is $-z^3/2$ and the same displayed
inequality includes it. Consequently
\[
 Q_3^{-1}=\ee_{B_1}\left(\frac{\beta^2w_1^2}{2A_1}\right).
\]
Set
\[
 y=\frac{(\beta/\gamma)w_1^2}{1+ub}.
\]
Its valuation is $(p-2)h_0\ge\rho:=\ceil{t/p}$.
The actual norm character of $B_1/A$, with coefficient $\beta$,
has the norm phase on all of this ideal. The proof is exactly the
argument of Lemma~\ref{O:T:actualnorm} on that edge; here
$\rho=\ceil{(t+1)/p}$ because $p\nmid t$.
It gives
\[
 \ee_{B_1}(\beta y/2)=\ee_A(-\beta n_1y/2).
\]
The exact norm choice $n_1w_1=\gamma/\beta$ now implies
\begin{equation}\label{O:T:Q3beforeden}
 Q_3^{-1}=\ee_A\left(-\frac\gamma2
       \frac{(\beta/\gamma)^{p-1}}{n_1(1+ub)}\right).
\end{equation}
No norm is replaced by a $p$th power in this calculation.

The endpoint of the norm expansion in the denominator is
$n_1(ub)=u^pA_0$, so \eqref{O:T:sym} on the lower edge of break $t$ gives
\[
 v_A(n_1(1+ub)-D)\ge
 m+\floor{\frac{(p-2)t+p-1}p}.
\]
The numerator weight $(\beta/\gamma)^{p-1}$ has valuation
$(p-1)(m+\delta)$. After subtracting the $\gamma$-conductor $r+1$,
the weighted denominator error has excess at least
\[
 (p-1)m+(p-2)\delta+
 \floor{\frac{(p-2)t+p-1}p}-t-1\ge0.
\]
To verify the last inequality write $t=pa+b_0$,
$1\le b_0\le p-1$, and $m=a+1+c$. The expression becomes
\[
 (p-3)a+(p-1)c+p-2-b_0+
 \floor{\frac{(p-2)b_0+p-1}p}+(p-2)\delta.
\]
It is nonnegative for $b_0\le p-2$; for $b_0=p-1$ the floor is
$p-2$. This includes the equality boundary $p=3,b_0=2$.
Both denominators are units. Hence \eqref{O:T:Q3beforeden} permits
replacement by $D$ at the actual additive modulus, yielding
\eqref{O:T:Q3phase} because $\gamma=\alpha W$ and $\beta/\gamma=vu$.
\end{proof}

\subsubsection*{The terminal phase and its cancellation}
\begin{theorem}[actual transition phase]\label{O:T:phase}
The actual four character factors satisfy
\begin{equation}\label{O:T:phaseeq}
 (Q_1Q_2Q_3Q_4)^{-1}=\Psi_A(\phi),
\end{equation}
where the expression on the right has now been derived from those
factors and is
\begin{equation}\label{O:T:phidef}
 \phi=-\frac W2\frac{(vu)^{p-1}}D
       +Wu^{p-1}\frac HD
       -\frac{u^{p-2}H}{2D^2}
       +\frac{u^{p-1}A_0J}{2D^2}.
\end{equation}
\end{theorem}
\begin{proof}
Combine Proposition~\ref{O:T:Q124} and Lemma~\ref{O:T:Q3}. In particular,
\eqref{O:T:phaseeq} is not inferred from vanishing of a candidate terminal
expression or from the desired equality of local constants.
\end{proof}

\begin{corollary}[transition cancellation]\label{O:T:cancel}
For the actual model supplied by Subsections~\ref{O:sec:models} and \ref{O:sec:parameters} and the local inputs
specified above, $Q_1Q_2Q_3Q_4=1$ in both characteristics.
\end{corollary}
\begin{proof}
Since $Wu=1$ and $D=1+u^pA_0$, direct expansion gives the exact
algebraic identity
\begin{equation}\label{O:T:factorphase}
 \phi=\frac{L_A}{2D^2}
       +\frac{u^{p-2}}{2D}(H-v^{p-1}),\qquad
 L_A=u^{p-1}A_0(J+u^{p-1}H).
\end{equation}
The independently proved, fully weighted estimate of Theorem~\ref{O:E:main} is
\[
 v_A L_A\ge2(p-1)m+(p-1)\delta-\floor{t/p}\ge1+t_2.
\]
The actual-character $R_2$ of Theorem~\ref{O:R:R2} is
\[
 H-v^{p-1}\in\pp_A^{L+(p-1)\delta},\qquad
 L=1+t-\ceil{t/p}.
\]
The second term of \eqref{O:T:factorphase} has depth at least $1+t_2$
because, with $\rho=\ceil{t/p}$ and $m\ge\rho$,
\[
 (p-2)m+L+(p-1)\delta-(1+t_2)
 =(p-2)(m+\delta)-\rho\ge0.
\]
The unit denominator and rational half preserve these whole ideal
bounds. Both terms therefore lie in the largest trivial ideal of
$\Psi_A$, and Theorem~\ref{O:T:phase} proves the assertion.
Neither an unweighted substitute for the first estimate nor an exact
auxiliary norm for $\beta/\alpha$ is used.
\end{proof}

\begin{corollary}[actual local factors in this transition row]\label{O:T:actual-local-factors}
Under the same explicit inputs, the two actual normalized local
factors agree in the totally ramified transition row $pm>t$.
\end{corollary}
\begin{proof}
Subsection~\ref{O:sec:parameters} derives the full Ueda-normalized Lamprecht formula for
these actual characters, with its complete odd-conductor term:
\[
 \frac{\Delta_{B_2}(\theta_2,\ee_{B_2})}
      {\Delta_{B_1}(\theta_1,\ee_{B_1})}
 =\frac{g_2}{g_1}\,Q_1Q_2Q_3Q_4\,
   \Psi_A(T_1b-T_2a),\qquad g_i\in\mu_4.
\]
The $g_i$ are the full evaluated residue factors, not merely their
polar forms. Equation \eqref{O:T:H14} and Corollary~\ref{O:T:cancel} put
this actual quotient in $\mu_4$. Independently, Ueda's First Main
Lemma and the odd cyclic lambda products give its $p$th power equal
to $1$, as proved for the actual pair in Corollary~\ref{U:odd-power}.
Since $p$ is odd, $\mu_p\cap\mu_4=\{1\}$, proving equality.
No equality of local constants was used to establish any field
congruence or character phase above.
\end{proof}


\subsection{Conductor reduction and completion of the odd-prime case}\label{O:sec:total}
\subsubsection*{Local notation}
Let $B/A$ be totally ramified with Galois group $C_p^2$, $p$ odd.
Let $B_1,B_2$ be the distinguished lower fields, with breaks
\[
 t=t_1\le t_2=t+\delta,\quad t'=t+p\delta,\quad p\nmid t.
\]
The upper breaks are $u_1=t'$ and $u_2=t$.
Set $N_i=\Nm_{B/B_i}$, $n_i=\Nm_{B_i/A}$ and
$T_i=\Tr_{B_i/A}$ and $\ee_E=\ee_A\Tr_{E/A}$.
Suppose $\psi_1N_1=\psi_2N_2=\psi_B$ is non-descending from $A$.
It is invariant under all of $\Gal(B/A)$, because it is invariant
under each of the two distinct upper cyclic subgroups.


\subsubsection*{The exact conductor reduction}
Let $a(\theta)$ denote a multiplicative conductor exponent. For a generator
$\tau_i$ of $\Gal(B_i/A)$ put $\mu_i=\psi_i^{\tau_i}/\psi_i$.
By Lemma~\ref{U:conjugacy}, $\mu_i$ generates the upper norm-character
group, the restrictions $\psi_1|_A,\psi_2|_A$ agree, and
Corollary~\ref{U:odd-power} supplies the common $p$th power.
\begin{lemma}[primitive lower bound]\label{O:G:lower}
For $i=1,2$,
\[
 a(\psi_i)\ge1+t_i+u_i.
\]
\end{lemma}
\begin{proof}
The nontrivial $\mu_i$ of Lemma~\ref{U:conjugacy} has conductor $u_i+1$.
Choose $y\in U_i^{u_i}$ with $\mu_i(y)\ne1$.
Write $y=1+z$, with $v_i z\ge u_i$. The ramification expansion gives
$v_i(\tau_i^{-1}z-z)\ge t_i+v_i z$, so
$\tau_i^{-1}y/y\in U_i^{t_i+u_i}$. Its value under $\psi_i$ is
$\mu_i(y)\ne1$. Thus $a(\psi_i)>t_i+u_i$, as asserted.
No equality in the ramification estimate is needed.
\end{proof}

\begin{lemma}[deep norm-kernel descent]\label{O:G:deep}
Let $E/A$ be a ramified cyclic edge with break $s$, and let $M/E$ be
the other upper edge, of break $u$ with $p\nmid u$. If $\theta$ is
an endpoint of a primitive compatible pair, then
\[
 \theta(x)=1\quad\text{for }x\in U_E^{1+s+u},\quad\Nm_{E/A}x=1.
\]
\end{lemma}
\begin{proof}
Hilbert 90 gives $x=\sigma(y)/y$. If $x=1$ there is nothing to prove.
Multiply $y$ by an element of $A^\times$ so that its valuation is
in $\{0,\ldots,p-1\}$. If this valuation is nonzero, the first
ramification term of $\sigma(y)/y$ has depth exactly $s$:
its uniformizer contribution is nonzero modulo $p$, while the unit
contribution has depth at least $s+1$. This contradicts the depth of
$x$. Thus $y$ is a unit, and a residue-field lift from $A$ makes
$y\in U_E^1$.

Choose the representative of $yA^\times$ with maximal possible unit
depth $q$. Such a finite maximum exists: otherwise representatives
from the closed subfield $A$ approximate $y$ arbitrarily closely,
forcing $y\in A^\times$ and $x=1$.
One has $p\nmid q$. In fact
\[
 U_E^{pk}=U_A^k U_E^{pk+1}\quad(k\ge1),
\]
because the image of $\pp_A^k$ in
$\PP_E^{pk}/\PP_E^{pk+1}$ is the entire common residue field.
If $q=pk$, this would improve the representative's depth, a contradiction.
The first ramification term for a depth-$q$ unit is consequently
nonzero and has depth $q+s$. Since $x\in U_E^{1+s+u}$, it follows
that $q\ge1+u$.
The commutator $\theta^{\sigma}/\theta$ has conductor $u+1$.
It is trivial on this $y$, so $\theta(\sigma(y)/y)=1$, with either
consistent convention for the generator. This proves the lemma.
\end{proof}

\begin{proposition}\label{O:G:reduce}
Every primitive compatible pair can be written
\[
 \psi_1=\chi(\lambda n_1),\qquad
 \psi_2=\varphi(\lambda n_2),
\]
where $\chi N_1=\varphi N_2$ is primitive and
\[
 a(\chi)=m_1=1+t+t',\qquad a(\varphi)=m_2=1+t+t_2.
\]
If $a(\psi_2)>m_2$, then $a(\lambda)=1+r$ with $r>t_2$, and
$\,p(r-t_2)>t$. Thus this is precisely the transition range of
Subsection~\ref{O:sec:parameters}, without a missing intermediate conductor row.
\end{proposition}
\begin{proof}
By Lemma~\ref{O:G:deep}, $\psi_1$ on $U_1^{m_1}$ factors through
$n_1(U_1^{m_1})$. The factored character is continuous, because the
unit groups are compact and the norm map is a quotient onto its
image. Extend it to a continuous character $\lambda$ of $A^\times$;
the image is open, and divisibility of $\mathbf C^\times$ extends a
character across the resulting discrete quotient.
Remove $\lambda n_i$ from both endpoints.
The pair remains compatible and primitive, and its first conductor is
at most $m_1$. Lemma~\ref{O:G:lower} makes it exactly $m_1$.

Both endpoint conductors are greater than their upper breaks plus
one, by Lemma~\ref{O:G:lower}. Therefore the high cyclic conductor formula
for their common norm pullback gives
\[
 p\,a(\chi)-(p-1)(t'+1)
 =p\,a(\varphi)-(p-1)(t+1).
\]
It follows that $a(\varphi)=m_2$.
If $a(\psi_2)>m_2$, then $a(\lambda n_2)=a(\psi_2)>m_2$.
The lower conductor formula forces $a(\lambda)=1+r>1+t_2$ and
\[
 a(\psi_2)=1+t_2+p(r-t_2)>1+t+t_2,
\]
which is equivalent to the asserted range.
\end{proof}

\subsubsection*{The minimal-conductor row}
\begin{theorem}\label{O:G:minimal}
If the primitive compatible pair has conductors
$m_1,m_2$, its two complete Second-Main-Lemma factors are equal.
\end{theorem}
\begin{proof}
Choose the actual $\Delta$ furnished by Theorem~\ref{O:M:realize} and write
$b=N_1\Delta$, $a=N_2\Delta$. Its full coefficients are
$C_1=\alpha b$, $C_2=\alpha a$, where
$\ee_A(\alpha\cdot)$ has ideal conductor $1+t_2$.
The full Lamprecht formula of Subsection~\ref{O:sec:parameters}, including its evaluated odd
residual factor, gives
\[
 \frac{\Delta_{B_2}(\varphi)}{\Delta_{B_1}(\chi)}
 =\frac{g_2}{g_1}\frac{\chi(\alpha b)}{\varphi(\alpha a)}
  \ee_A\left(\alpha[T_1b-T_2a]\right),\qquad g_i\in\mu_4.
\]
Lemma~\ref{U:determinant} cancels the two values at $\alpha$.
Compatibility, evaluated on $\Delta\in B^\times$, cancels
$\chi(b)/\varphi(a)$. The trace--norm bound of Theorem~\ref{O:A:thm:H14}
$T_1b-T_2a\in\pp_A^{1+t_2}$ kills the last factor.
Thus the ratio lies in $\mu_4$.
Corollary~\ref{U:odd-power} gives its $p$th power equal to $1$; since $p$ is
odd, the ratio is $1$. That lemma also identifies each lower cyclic
norm-character product with $1$, so these are the complete target
factors, not just a quotient missing lambda constants.
\end{proof}
\begin{theorem}[The totally ramified branch]\label{O:G:totalbranch}
For a totally ramified $C_p^2$ diamond over a local field of odd residue
characteristic, every compatible non-descending pair gives equal complete
Second Main Lemma factors. This holds in mixed and equal characteristic.
\end{theorem}
\begin{proof}
First compare the distinguished smallest-break field with an arbitrary
other lower field, using the notation fixed above. If the second
conductor is $m_2$, the conductor equality in
Proposition~\ref{O:G:reduce} gives $m_1$ on the first side, and
Theorem~\ref{O:G:minimal} applies. Otherwise Proposition~\ref{O:G:reduce} writes the
pair as a common base twist of an actual minimal pair, with
$p(r-t_2)>t$. The simultaneous choices of Subsection~\ref{O:sec:parameters}
realize that pair without changing either local constant. The actual
four-factor calculation of Subsection~\ref{O:sec:transition}, using
Subsections~\ref{O:sec:q1} and \ref{O:sec:q2}, puts the quotient in $\mu_4$.
Corollary~\ref{U:odd-power} puts it in $\mu_p$, so it is $1$. The lower cyclic
products are also $1$ by that corollary. For a pair not containing the
distinguished field, extend its invariant common pullback to that field
by Hilbert~90 and character extension; the two distinguished comparisons
then give the desired equality by transitivity. The preceding proofs
include both characteristics and $p=3$.
\end{proof}


\section{Dyadic extensions with an unramified intermediate field}\label{sec:dyadic-ur}
\subsection{The local configuration and notation}
Let $F$ be a nonarchimedean local field of residue characteristic $2$,
in mixed or equal characteristic. Let $K/F$ be biquadratic. For a quadratic
intermediate field $L$, let $\omega_L$ be the nontrivial character of
$F^\times/\Nm_{L/F}L^\times$. A primitive compatible pair means characters
$\theta_i$ of $L_i^\times$ with the same pullback $\Theta$ to $K^\times$,
where $\Theta$ does not descend from $F$ through $\Nm_{K/F}$.
Fix $\psi_F$ and use $\psi_L=\psi_F\Tr_{L/F}$. The target is equality of
\begin{equation}\label{D:UR:Adef}
 A_L(\theta_L)=\Delta_L(\theta_L,\psi_L)\Delta_F(\omega_L,\psi_F).
\end{equation}
The omitted trivial-character factor is $1$.
Write $a_L(\theta)$ for the multiplicative conductor, $n_L$ for the additive
conductor in the convention that $\pp_L^{-n_L}$ is the largest trivial
\emph{ideal}, and $v_L(\pp_L)=1$. On a ramified quadratic edge $L/F$ of
break $t$, put $T=t+1$. Its different exponent is $T$, and
\begin{equation}\label{D:UR:localledger}
 n_L=2n_F+T,\qquad
 \Tr_{L/F}(\pp_L^j)=\pp_F^{\floor{(j+T)/2}},\qquad
 v_F(\Nm_{L/F}x)=v_L(x).
\end{equation}
For an unramified edge, conductors and normalized valuations restrict
without a ramification multiplier.

We use the one-dimensional results recalled in
Sections~\ref{sec:local-notation}--\ref{sec:stationary}, and the local
norm-character identifications of Appendix~\ref{app:diamond-foundations}.

Put $N_i=\Nm_{K/L_i}$, $n_i=\Nm_{L_i/F}$, and let $\nu_i$ denote
the upper norm character. Lemma~\ref{U:conjugacy} gives
$\theta_i^{\sigma_i}=\theta_i\nu_i$ and the common restriction
$D=\theta_i|_F\omega_i$. Proposition~\ref{U:quadratic-square} leaves a
possible sign; the residual calculations below determine that sign.

\subsection{The quadratic norm-character phase}
Let $E/F$ be ramified quadratic of break $t\ge1$, put $T=t+1$ and
$q=\ceil{T/2}$, and let $\tau$ be its actual norm character. By additive
duality on $U_F^q/U_F^T$, choose $\alpha$ such that
\begin{equation}\label{D:UR:tauchart}
 \tau(1-z)=\Psi_F(z),\quad z\in\pp_F^q,\qquad
 \Psi_L(w):=\psi_L(\alpha w),\qquad v_F\alpha=-n_F-T.
\end{equation}
Both $\Psi_F$ and $\Psi_E$ have largest trivial ideal with exponent $T$.
The linear chart is valid because $2q\ge T$; there is no coefficient $1/2$.
In mixed characteristic put $e=v_F(2)$; in equal characteristic put
$e=+\infty$. The trace of $1$ in \eqref{D:UR:localledger} gives
\begin{equation}\label{D:UR:ebound}
 e\ge\floor{T/2},\qquad e+q\ge T.
\end{equation}
\begin{lemma}[Actual dyadic norm phase]\label{D:UR:normphase}
For every $z\in\pp_E^q$,
\begin{equation}\label{D:UR:phase}
 \Psi_E(z)=\Psi_F(\Nm_{E/F}z).
\end{equation}
\end{lemma}
\begin{proof}
The exact quadratic identity is
\[
 \Nm(1-z)=1-\Tr z+\Nm z.
\]
Both $\Tr z$ and $\Nm z$ have valuation at least $q$, by
\eqref{D:UR:localledger}. Evaluate \eqref{D:UR:tauchart} at $\Tr z-\Nm z$ and use
$\tau(\Nm(1-z))=1$. This gives
$\Psi_F(\Tr z-\Nm z)=1$, which is \eqref{D:UR:phase}.
\end{proof}
\begin{remark}
The sign in \eqref{D:UR:phase} is positive. The odd-prime expression with a
negative norm term on the right must not simply be imported at $p=2$.
This proof is exact on its entire stated ideal; it does not assert
\eqref{D:UR:phase} on shallower ideals.
\end{remark}
\begin{lemma}[The odd quadratic conductor]\label{D:UR:oddT}
If $T$ is odd, then $\operatorname{char}F=0$ and $T=2e+1$. In that case,
with $k$ the residue field,
\begin{equation}\label{D:UR:critical4}
 \Psi_F(4z)=(-1)^{\Tr_{k/\F_2}(\bar z)},\qquad z\in\OO_F.
\end{equation}
\end{lemma}
\begin{proof}
Write $T=2r+1$, so $r\ge1$, and $e\ge r$ by \eqref{D:UR:ebound}.
If $e>r$, including equal characteristic, the square
$(1+\pi^rz)^2$ is $1+\pi^{2r}z^2$ modulo $\pp_F^T$.
Since $\tau$ has order $2$, this would make its nontrivial last-layer
additive character trivial on every residue square. Frobenius is onto
$k$, a contradiction. Hence $e=r$.
Now $v_F4=2e=T-1$, so $z\mapsto\Psi_F(4z)$ is a nontrivial residue
additive character. The equality
$(1+2z)^2=1+4(z+z^2)$ and the order of $\tau$ show that it kills
$\bar z+\bar z^2$. The image of this Artin--Schreier map is exactly the
absolute-trace-zero hyperplane: its kernel has two elements and its image
has zero trace. The unique nontrivial character of the quotient is the
right side of \eqref{D:UR:critical4}.
\end{proof}

\subsection{Minimal characters and common twists}
Assume from now through Section~\ref{D:UR:sec:critical} that $U/F$ is unramified
quadratic, $E/F$ ramified quadratic of break $t$, and $K=EU$. Write
\[
 N_U=\Nm_{K/U},\quad N_E=\Nm_{K/E},\quad
 n_U=\Nm_{U/F},\quad n_E=\Nm_{E/F}.
\]
Let $\eta=\omega_U$ and $\tau=\omega_E$. The norm character on $K/U$ is
$\tau_U=\tau n_U$. All four $\Psi_L=\psi_L(\alpha\,\cdot)$ have largest
trivial ideal $\pp_L^T$. Let $k\subset\kappa$ be the residue fields of
$F\subset U$, and put $Q=|k|$.
Choose $c\in\OO_F$ whose residue has absolute trace $1$, and choose
\begin{equation}\label{D:UR:dchoice}
 d\in\OO_U,\qquad d^2-d=c.
\end{equation}
This polynomial gives the unramified quadratic extension. For its
nontrivial automorphism $\sigma$,
\[
 \sigma(d)=1-d,\quad \Tr_{U/F}d=1,\quad n_U(d)=-c.
\]
In particular $c,d$ are units and $\bar d\notin k$.

\begin{theorem}[Compatible minimal characters on the full multiplicative groups]\label{D:UR:models}
There exist characters $\chi_U^0,\chi_E^0$, both of conductor $T$, with
a primitive common norm pullback, such that
\begin{align}
 (\chi_U^0)^\sigma/\chi_U^0&=\tau_U,\label{D:UR:modelcomm}\\
 \chi_U^0(1-z)&=\Psi_U(d^2z) &&(z\in\pp_U^q),\label{D:UR:modelU}\\
 \chi_E^0(1-z)&=\Psi_E(cz) &&(z\in\pp_E^q).\label{D:UR:modelE}
\end{align}
These are identities on the full specified subgroups, with actual
characters on the entire multiplicative groups.
\end{theorem}
\begin{proof}
The right side of \eqref{D:UR:modelU} is a character on $H=U_U^q$, since
$2q\ge T$. Its exact conductor is $T$. Unramified norm expansion shows
\[
 \tau_U(1-z)=\Psi_U(z)\quad(z\in\pp_U^q),
\]
because $v_F n_Uz\ge2q\ge T$. Further,
$\sigma(d^2)-d^2=1-2d$; the $2d$ error is killed on this whole ideal by
\eqref{D:UR:ebound}. Thus the prescribed character on $H$ has the required
conjugate quotient on $H$.

Here is the extension to all units, with the commutator retained.
Put $V=U_U^1$ and $I=\{\sigma(v)/v:v\in V\}$. Prescribe
$\kappa_I(\sigma(v)/v)=\tau_U(v)$. This is well defined because
$\tau_U(f)=\tau(f^2)=1$ on $U_F^1$. If $\sigma(v)/v\in H$, then $v$
modulo $\pp_U^q$ is fixed, hence is represented by a unit of $F$; write
$v=fh$ with $f\in U_F^1$, $h\in H$. (For unramified extensions this
follows coefficient by coefficient in the uniformizer expansion.)
The characters prescribed on $I$ and $H$ consequently agree on the
intersection. Extend their product character from $IH/U_U^T$ across the
finite abelian group $V/U_U^T$. Make it trivial on the odd-order
Teichmuller group and on a uniformizer of $F$. The character $\tau_U$ is
trivial on both these factors. The resulting $\chi_U^0$ has
\eqref{D:UR:modelcomm} on all of $U^\times$ and has exact conductor $T$.

Its pullback to $K$ is invariant under both quadratic automorphisms,
since $\tau_U$ kills $N_UK^\times$. Hilbert 90 makes this pullback trivial
on $\ker N_E$; it therefore descends and extends to $\chi_E^0$ on
$E^\times$. The unramified norm is onto units, so only an unramified value
is left to choose in this extension.
For $z\in\pp_K^q$, the exact ramified quadratic norm gives
\[
 \chi_U^0(N_U(1-z))
 =\Psi_K(d^2z)\Psi_U(-d^2N_Uz)
 =\Psi_K((d^2-d)z)=\Psi_K(cz).
\]
The middle equality uses Lemma~\ref{D:UR:normphase} on $dz$ and
$N_U(dz)=d^2N_Uz$, an exact equality. The unramified norm onto $U_E^q$
and its expansion modulo $\pp_E^T$ now give \eqref{D:UR:modelE}. Since $c$ is
a unit, its conductor is exactly $T$. Nontriviality of \eqref{D:UR:modelcomm}
excludes descent of the common pullback from $F$.
\end{proof}

\begin{proposition}[Realizing any given pair]\label{D:UR:actual}
For every primitive compatible pair $(\theta_U,\theta_E)$ there is a
character $\lambda$ of $F^\times$, after an unramified adjustment of
$\chi_E^0$ if necessary, such that
\begin{equation}\label{D:UR:twistpair}
 \theta_U=\chi_U^0(\lambda n_U),\qquad
 \theta_E=\chi_E^0(\lambda n_E).
\end{equation}
There is an integer $h\ge0$ with
\begin{equation}\label{D:UR:conductors}
 a_U(\theta_U)=T+h,\qquad a_E(\theta_E)=T+2h.
\end{equation}
For $h>0$, $a_F(\lambda)=T+h$; for $h=0$, $a_F(\lambda)\le T$.
\end{proposition}
\begin{proof}
Lemma~\ref{U:conjugacy} identifies the actual conjugate quotient over $U$ with
$\tau_U$, with no choice of power in a group of order $2$.
Thus $\theta_U/\chi_U^0$ is invariant. Hilbert 90 and character extension
make it $\lambda n_U$. The other quotient is a norm character for the
unramified edge $K/E$, so it is absorbed by the unramified adjustment.
The conductor of $\theta_U$ is at least $T$, the conductor of its
commutator. If it is greater, the common twist has this conductor, since
$\chi_U^0$ has conductor $T$ and unramified pullback preserves conductor.
At the endpoint it has conductor at most $T$. Both twists of $\theta_U$
by $\tau_U$ are conjugate and have the same conductor. Ueda's norm-pullback
conductor formula, including its no-drop criterion at $T$, gives upstairs
conductor $T+2h$. The unramified edge $K/E$ preserves it.
\end{proof}

\subsection{Common norm representatives}
For $0\le h\le t$ put
\[
 s_U=\ceil{(T+h)/2},\quad s_E=\ceil{T/2}+h,\quad
 d_U=\floor{(T+h)/2}.
\]
Choose a coefficient $A_*$ of the base twist on its half-conductor layer:
\begin{equation}\label{D:UR:lambdachart}
 \lambda(1-z)=\Psi_F(A_*z)\quad(z\in\pp_F^{s_U}).
\end{equation}
For $h>0$, $v_FA_*=-h$; at $h=0$ the coefficient can be chosen integral.
Its ambiguity is $\pp_F^{T-s_U}$.
\begin{lemma}\label{D:UR:choosex}
There is $x\in E$, with $A=n_Ex$, such that $A$ can replace $A_*$ in
\eqref{D:UR:lambdachart}. If $h>0$, $v_Ex=-h$; at $h=0$, $x$ is integral.
The zero coefficient class may be represented by $x=0$.
\end{lemma}
\begin{proof}
For $h>0$ the required relative precision is
$T-s_U+h=d_U\le t$, since $T+h\le2T-1$. Thus the subcritical norm
representative theorem applies. For an integral coefficient of valuation
$a<T-s_U$ use relative precision $T-s_U-a\le t$. If the integral class
is zero, choose zero. No approximation beyond the break is requested.
\end{proof}
Put
\begin{equation}\label{D:UR:BCdefs}
 B_U=A+d^2,\quad B_E=A-x+c,\quad C=x+d,
 \quad Z_U=N_UC,\quad Z_E=N_EC,\quad b=\Tr_{E/F}x.
\end{equation}
\begin{lemma}[The actual coefficients]\label{D:UR:coeffs}
For the given characters in \eqref{D:UR:twistpair},
\begin{equation}\label{D:UR:actualcoeffs}
 \theta_U(1-z)=\Psi_U(B_Uz)\ (z\in\pp_U^{s_U}),\qquad
 \theta_E(1-z)=\Psi_E(B_Ez)\ (z\in\pp_E^{s_E}).
\end{equation}
Their coefficient valuations are $-h$ and $-2h$, respectively.
\end{lemma}
\begin{proof}
For the unramified pullback of $\lambda$, the extra norm term has phase
valuation at least $-h+2s_U\ge T$; use \eqref{D:UR:modelU}. On the ramified side,
$u=\Tr z-n_Ez$ belongs to $\pp_F^{s_U}$ for $z\in\pp_E^{s_E}$, because
\[
 s_E\ge s_U,\qquad \floor{(s_E+T)/2}\ge s_U.
\]
Then \eqref{D:UR:lambdachart} and Lemma~\ref{D:UR:normphase}, applied to $xz$ of
valuation at least $\ceil{T/2}$, give
\[
 \lambda(n_E(1-z))=\Psi_F(A\Tr z-A n_Ez)
                  =\Psi_E((A-x)z).
\]
Together with \eqref{D:UR:modelE} this proves the formulas. For $h>0$, $A$
is the unique leading term in each coefficient. At $h=0$, write
$\xi=\bar x\in k$. The residues are $\xi^2+\bar d^2$ and
$\xi^2-\xi+\bar c$. The first is nonzero since $\bar d\notin k$, and
the second has absolute trace $1$.
\end{proof}
The norms of the \emph{same} element $C$ have the exact expressions
\begin{align}
 Z_U&=A+b d+d^2,\label{D:UR:ZU}\\
 Z_E&=x^2+x-c,\label{D:UR:ZE}\\
 S:=\Tr_{U/F}Z_U-\Tr_{E/F}Z_E
     &=1+4c-\mathcal D_x,\qquad
       \mathcal D_x=b^2-4A=(2x-b)^2.\label{D:UR:Sidentity}
\end{align}
Indeed $\Tr d=1$, $\Tr d^2=1+2c$, and $\Tr x^2=b^2-2A$.
These identities do not use translated roots in characteristic zero.

The normalized factors and residual functions are those of
Lemma~\ref{U:stationary-factor}, with $J=T$. We denote them by $H_L,g_L$.

\subsection{The minimal-conductor comparison}
\begin{theorem}\label{D:UR:minimal}
Every primitive compatible pair with
$a_U(\theta_U)=a_E(\theta_E)=T$ satisfies
\begin{equation}\label{D:UR:URtarget}
 \Delta_U(\theta_U)\Delta_F(\eta)
 =\Delta_E(\theta_E)\Delta_F(\tau).
\end{equation}
This includes the odd-conductor maximal-break case in mixed characteristic
and all equal-characteristic minimal pairs.
\end{theorem}

\subsubsection*{Even \texorpdfstring{$T$}{T}}
Here $T=2r$ and $x$ in Lemma~\ref{D:UR:choosex} is integral. The differences
$Z_U-B_U=bd$ and
\[
 Z_E-B_E=bx-2A+2x-2c
\]
have respective valuations at least $r$ in $U$ and $2r$ in $E$, by
\eqref{D:UR:localledger} and $v_F2\ge r$. Thus $Z_U,Z_E$ are legitimate
stationary coefficients. All three conductors $T$ are even, so there are
no residual sums.
Put $c_0=-\alpha^{-1}$. Compatibility gives
$\theta_U(Z_U)=\theta_E(Z_E)$; Lemma~\ref{U:determinant} gives
$\theta_U(c_0)/\theta_E(c_0)=\eta(c_0)\tau(c_0)$.
Using \eqref{U:normalized-factor} and
$\Delta_F(\tau)=\tau(c_0)\Psi_F(-1)$ therefore gives
\begin{equation}\label{D:UR:evenratio}
 \frac{\Delta_U(\theta_U)\Delta_F(\eta)}
      {\Delta_E(\theta_E)\Delta_F(\tau)}
 =(-1)^T\Psi_F(1-S).
\end{equation}
Here $\Delta_F(\eta)=(-1)^{n_F}$ and
$\eta(c_0)=(-1)^{n_F+T}$; this explains the sign instead of discarding it.
Now $2x-b$ is trace zero and has $E$-valuation at least $2r$.
Lemma~\ref{D:UR:normphase} gives $\Psi_F(\mathcal D_x)=1$ since
$n_E(2x-b)=-\mathcal D_x$. Also $v_F(4c)\ge2r=T$.
Equation \eqref{D:UR:Sidentity} proves that \eqref{D:UR:evenratio} is $1$.

\subsubsection*{Odd \texorpdfstring{$T$}{T}: the elementary sign}
By Lemma~\ref{D:UR:oddT}, $T=2e+1$, $\operatorname{char}F=0$ and $e\ge1$.
We still have integral $x$, $b\in\pp_F^e$ and legitimate stationary
coefficients $Z_U,Z_E$: their errors from \eqref{D:UR:actualcoeffs} have depths
at least $e$ and $2e$, which exceed their ambiguity depth $e$ as required
(with equality allowed in $U$). In addition put
\begin{equation}\label{D:UR:Wdefs}
 P_C=\Tr_{K/E}C=2x+1,\quad W_E=n_E(P_C),\quad
 Q_C=\Tr_{K/U}C=b+2d,\quad W_U=n_U(Q_C).
\end{equation}
The element $P_C$ is a unit and
$W_E=1+2b+4A\equiv1\pmod{\pp_F^{2e}}$. Thus $W_E$ is a legitimate
stationary coefficient for $\tau$ in the normalized $\Psi_F$ convention,
and $\tau(W_E)=1$ because it is an actual norm.

Lemma~\ref{U:e2}, in the present notation, gives
\begin{equation}\label{D:UR:e2}
 \Tr_{U/F}N_UY+n_U(\Tr_{K/U}Y)
 =\Tr_{E/F}N_EY+n_E(\Tr_{K/E}Y).
\end{equation}
Use the three stationary coefficients $Z_U,Z_E,W_E$ in \eqref{U:normalized-factor}.
Compatibility, \eqref{U:det-character}, and \eqref{D:UR:e2} give
\begin{equation}\label{D:UR:oddratio}
 \frac{\Delta_U(\theta_U)\Delta_F(\eta)}
      {\Delta_E(\theta_E)\Delta_F(\tau)}
 =(-1)^T\Psi_F(W_U)\frac{g_U}{g_Eg_\tau}.
\end{equation}
The prefactor on the right is $1$. Indeed $a=b/2\in\OO_F$ and
\[
 W_U=4n_U(d+a),\qquad
 \overline{n_U(d+a)}=\bar a^2+\bar a-\bar c.
\]
Its absolute trace is $1$. Lemma~\ref{D:UR:oddT} gives $\Psi_F(W_U)=-1$,
which cancels $(-1)^T=-1$.

\subsubsection*{Odd \texorpdfstring{$T$}{T}: an identity between the three actual residual functions}
Choose a uniformizer $\Pi$ of $E$ and put $\pi=n_E\Pi$; this is a
uniformizer of $F$ and $U$. In \eqref{U:normalized-residual} use critical elements
$\pi^e$ for $U,F$ and $\Pi^e$ for $E$.
For $z\in\OO_K$ set
\[
 Y_z=C(1+\Pi^e z).
\]
All traces whose norms are used below are nonzero. The trace over $E$ stays
a unit. The trace over $U$ stays at valuation $e$, as will follow below.
The three normalized critical residues are
\begin{equation}\label{D:UR:critmaps}
 \begin{aligned}
 \frac{N_UY_z/Z_U-1}{\pi^e}&\equiv\bar z^2 &&\text{in }\kappa,\\
 \frac{N_EY_z/Z_E-1}{\Pi^e}&\equiv\Tr_{\kappa/k}\bar z &&\text{in }k,\\
 \frac{n_E(\Tr_{K/E}Y_z)/W_E-1}{\pi^e}
    &\equiv\bigl(\Tr_{\kappa/k}(\bar C\bar z)\bigr)^2 &&\text{in }k.
 \end{aligned}
\end{equation}
To check the first row, the trace of $\Pi^ez$ on the ramified edge has
$U$-valuation at least $\floor{(e+T)/2}\ge e+1$; its norm has leading
term $\pi^e\bar z^2$. For the second row the unramified norm's trace term
has leading coefficient $\Tr\bar z$, and its quadratic term has depth
$2e\ge e+1$. For the third, the relative trace change has leading
$E$-coefficient $\Tr(\bar C\bar z)$ at depth $e$, since
$\overline{P_C}=1$; apply the first ramified calculation to its norm.
These verifications include the smallest case $e=1$.

Moreover
\begin{equation}\label{D:UR:extraUconstant}
 n_U(\Tr_{K/U}Y_z)-W_U\in\pp_F^T.
\end{equation}
In fact $v_UQ_C=e$, because $Q_C=2(d+b/2)$ and $\bar d\notin k$.
The trace change $\Tr_{K/U}(C\Pi^ez)$ has valuation at least
$\floor{(e+T)/2}\ge e+1$. Expanding its unramified quadratic norm proves
\eqref{D:UR:extraUconstant}.

Apply \eqref{D:UR:e2} to $Y_z$ and to $C$, and subtract. The common
multiplicative character values cancel by compatibility. The two
$\tau$-values of the norms of the $E$-traces are both $1$.
Equation \eqref{D:UR:extraUconstant} removes precisely the remaining
$U$-trace-norm phase. Lemma~\ref{U:stationary-factor} and \eqref{D:UR:critmaps} give the
\emph{pointwise identity of actual residual functions}
\begin{equation}\label{D:UR:actualHidentity}
 H_U(z^2)=H_E(\Tr_{\kappa/k}z)
           H_\tau\bigl((\Tr_{\kappa/k}(\bar C z))^2\bigr),
 \qquad z\in\kappa.
\end{equation}
No identification of polar forms alone has been used.

The map
\[
 z\longmapsto\bigl(\Tr z,(\Tr(\bar C z))^2\bigr)
        :\kappa\longrightarrow k\times k
\]
is bijective. Write $z=u+v\bar d$ and $\bar x=\xi\in k$; its coordinates
are $v$ and $(u+(\xi+1)v)^2$. Frobenius on $k$ is bijective.
Since Frobenius on $\kappa$ is bijective too, summing
\eqref{D:UR:actualHidentity} gives
\[
 \sum_{\kappa}H_U
     =\left(\sum_kH_E\right)\left(\sum_kH_\tau\right),
 \qquad g_U=g_Eg_\tau.
\]
Together with \eqref{D:UR:oddratio}, this proves Theorem~\ref{D:UR:minimal} in the
odd row as well. Notice that the local norm maps themselves produced the
affine phases and their coordinates, before any finite sum was evaluated.

\subsection{Positive critical conductors}
\label{D:UR:sec:critical}
The naive pair of norms in \eqref{D:UR:ZU}--\eqref{D:UR:ZE} need not both be
stationary near the critical boundary. We replace the second norm by a
conjugate multiple. This changes its character value by an explicitly
known sign, not by an uncontrolled character error.

Throughout this subsection let
\[
 1\le h\le t=T-1,\qquad m_U=T+h,\quad m_E=T+2h.
\]
Use the actual model and base twist of Proposition~\ref{D:UR:actual}, and choose
$A=n_E x$ in the stationary coefficient class by Lemma~\ref{D:UR:choosex}.
In particular $v_E x=-h$, so $x\ne0$. Let $x'$ be its nontrivial
$E/F$ conjugate, and put
\begin{equation}\label{D:UR:newrep}
 \rho=\frac{x'}x,\qquad
 \widehat Z_U=A+bd+d^2,\qquad
 \widehat Z_E=\rho(x^2+x-c),\qquad b=\Tr_{E/F}x.
\end{equation}

\subsubsection*{The trace parity lemma}
\begin{lemma}[An exact last-layer sign]\label{D:UR:traceparity}
Put $\epsilon=(T-h)\bmod2$ and $\mathfrak u=b^2/A$.
If $\epsilon=0$, then $v_F\mathfrak u\ge T$, with $v_F0=+\infty$.
If $\epsilon=1$, then
\[
 v_F\mathfrak u=T-1,\qquad
 \Psi_F(\mathfrak u z)=(-1)^{\Tr_{k/\F_2}(\bar z)}
 \quad(z\in\OO_F).
\]
Consequently, for the fixed $c$ in \eqref{D:UR:dchoice},
\begin{equation}\label{D:UR:paritysign}
 (-1)^{T+h}\Psi_F(-c\mathfrak u)=1.
\end{equation}
This holds in both mixed and equal characteristic.
\end{lemma}
\begin{proof}
The trace-ideal formula gives
\[
 v_Fb\ge B:=\floor{(T-h)/2}.
\]
If $T-h$ is even, $2B+h=T$, proving the first assertion.
If $T-h$ is odd, the induced map
\[
 \pp_E^{-h}/\pp_E^{1-h}
 \longrightarrow \pp_F^B/\pp_F^{B+1}
\]
is a nonzero $k$-linear map between one-dimensional $k$-spaces.
Indeed the trace images of the two consecutive source ideals are exactly
$\pp_F^B$ and $\pp_F^{B+1}$. Thus an element of exact valuation $-h$
has trace of exact valuation $B$. Hence $b\ne0$ and
$v_F\mathfrak u=2B+h=T-1$.

For $z\in\OO_F$, quadratic norm expansion gives the exact identity
\begin{equation}\label{D:UR:traceparitynorm}
 n_E\!\left(1+\frac bA xz\right)
       =1+\mathfrak u(z+z^2).
\end{equation}
Here $v_E((b/A)x)=2B+h=T-1>0$, so the argument is a unit.
Since $T-1\ge q$, the actual norm-character formula
\eqref{D:UR:tauchart} applies to the right side. It follows that the residue
additive character $z\mapsto\Psi_F(\mathfrak u z)$ annihilates
$z+z^2$. It is nontrivial, because $\Psi_F$ has largest trivial ideal
$\pp_F^T$ and $v_F\mathfrak u=T-1$. The image of $z\mapsto z+z^2$
in $k$ is the absolute-trace-zero hyperplane. The character is therefore
exactly $(-1)^{\Tr_{k/\F_2}z}$. Since $\bar c$ has absolute trace one,
$\Psi_F(-c\mathfrak u)=-1$ in this parity and is $1$ in the other.
Finally $T-h$ and $T+h$ have the same parity, proving \eqref{D:UR:paritysign}.
\end{proof}

\subsubsection*{Stationarity and the exact elementary factor}
\begin{lemma}[Corrected common representatives]\label{D:UR:correctedreps}
The two elements $\widehat Z_U,\widehat Z_E$ in \eqref{D:UR:newrep} are actual
ordinary stationary coefficients for $\theta_U,\theta_E$ in the
normalization of \eqref{U:normalized-factor}, for every $1\le h<T$. Moreover
\begin{align}
 \widehat Z_U-B_U&=bd,\label{D:UR:newerrU}\\
 \widehat Z_E-B_E&=b(1-c/x),\label{D:UR:newerrE}\\
 \widehat S:=\Tr_{U/F}\widehat Z_U-
                   \Tr_{E/F}\widehat Z_E&=1+c\mathfrak u,\label{D:UR:newS}\\
 \frac{\theta_E(\widehat Z_E)}{\theta_U(\widehat Z_U)}
                 &=(-1)^h.\label{D:UR:newmult}
\end{align}
\end{lemma}
\begin{proof}
The first error is immediate. Since $xx'=A$ and $x+x'=b$,
\[
 \rho(x^2+x-c)=A+x'-c\rho,
\]
whose difference from $A-x+c$ is
$b-c(x'/x+1)=b(1-c/x)$. As $v_E x=-h<0$, the latter parenthesis is a unit.
With $B=\floor{(T-h)/2}$, the error valuations are at least $B$ over $U$
and $2B$ over $E$. The required ambiguity depths are
\[
 T-s_U=B,\qquad T-s_E=\floor{T/2}-h.
\]
For $\epsilon=(T-h)\bmod2$,
\[
 2B-(T-s_E)=\ceil{T/2}-\epsilon\ge0.
\]
Both errors therefore lie in the correct coefficient ambiguity ideals.
The actual coefficients have valuations $-h,-2h$; the corrected ones do
also, either by these estimates or directly by their unique leading terms.

The traces of $\rho$ and of the corrected coefficients are
\[
 \Tr_{E/F}\rho=\frac{b^2-2A}{A}=\mathfrak u-2,
\]
\[
 \Tr_{E/F}\widehat Z_E=2A+b+2c-c\mathfrak u,
 \qquad \Tr_{U/F}\widehat Z_U=2A+b+1+2c.
\]
This proves \eqref{D:UR:newS}. The uncorrected $Z_U,Z_E$ are the two norms of
$C=x+d$, so compatibility gives $\theta_U(Z_U)=\theta_E(Z_E)$.
Lemma~\ref{U:conjugacy} on $E$ gives
\[
 \theta_E(\rho)=\frac{\theta_E(x')}{\theta_E(x)}
      =\eta(n_E x)=(-1)^{v_FA}=(-1)^h,
\]
proving \eqref{D:UR:newmult}. No equation $n_E x=A_*$ was assumed: the exact
norm $A=n_E x$ was chosen inside its justified coefficient class.
\end{proof}

\begin{remark}[A literal common element remains available]
The correction can also be recorded using a common element and a base
scalar. Put $Y=x'(x+d)\in K$ and $\lambda_0=A^{-1}\in F$.
Then $\lambda_0N_UY=\widehat Z_U$ and
$\lambda_0N_EY=\widehat Z_E$. The symbol $\lambda_0$ is a field element,
not the twisting character $\lambda$. We use \eqref{D:UR:newmult} directly
rather than suppressing the character value introduced by this scalar.
\end{remark}

\begin{proposition}[Reduction to actual residual factors]\label{D:UR:criticalratio}
Use the corrected coefficients in the two Lamprecht formulas, and use
coefficient $1$ for $\tau$. Let $g_U,g_E,g_\tau$ be their full residual
factors, with $g=1$ in even conductor. Then
\begin{equation}\label{D:UR:newratio}
 \frac{\Delta_U(\theta_U)\Delta_F(\eta)}
      {\Delta_E(\theta_E)\Delta_F(\tau)}
  =(-1)^{T+h}\Psi_F(-c\mathfrak u)
                      \frac{g_U}{g_Eg_\tau}
  =\frac{g_U}{g_Eg_\tau}.
\end{equation}
\end{proposition}
\begin{proof}
Let $c_0=-\alpha^{-1}$. Lemma~\ref{U:determinant} gives
$\theta_U(c_0)/\theta_E(c_0)=\eta(c_0)\tau(c_0)$.
The product $\eta(c_0)\Delta_F(\eta)$ is $(-1)^T$, since
$v_Fc_0=n_F+T$ and $\Delta_F(\eta)=(-1)^{n_F}$.
The quotient of the multiplicative stationary values is \eqref{D:UR:newmult}.
The additive phase from \eqref{U:normalized-factor} and the norm-character factor is
$\Psi_F(1-\widehat S)=\Psi_F(-c\mathfrak u)$.
These exact identities give the first equality in \eqref{D:UR:newratio}.
Lemma~\ref{D:UR:traceparity} gives the second. In particular the elementary
sign has been evaluated, not inferred from the square relation.
\end{proof}

\subsubsection*{The unramified residual sum}
\begin{lemma}\label{D:UR:resU}
If $T+h$ is odd and $h>0$, the actual residual function for
$\theta_U$ with coefficient $\widehat Z_U$ is $1$ on the residue subfield
$k\subset\kappa$. Its normalized sum is therefore $g_U=1$.
\end{lemma}
\begin{proof}
Set $a=(T+h-1)/2$ and choose the critical element
$\delta_U=\pi^a\in F$, for any uniformizer $\pi$ of $F$.
For $z\in\OO_F$ put $v=\delta_Uz$. One has
\[
 2a-s_E=\floor{T/2}-1\ge0.
\]
Also $a\ge q$: for even $T$ this follows from odd $h\ge1$; for odd $T$
the present parity requires even $h\ge2$. Thus both $\theta_E(1+v)$ and
$\tau(1+v)$ are evaluated by their legitimate stationary formulas.
The determinant identity and triviality of $\eta$ on units give
\[
 \theta_U(1+v)=\theta_E(1+v)\tau(1+v).
\]
Substitution into the actual residual function yields
\[
 H_U(\bar z)=\Psi_F((1-\widehat S)v)
            =\Psi_F(-c\mathfrak u\,\delta_Uz)=1.
\]
Indeed $v_F\mathfrak u\ge T-1$ and $a\ge1$.

For completeness, the complete Lamprecht function is a nondegenerate
quadratic phase on the additive group of $\kappa$. Its restriction to $k$
is $1$, so $k$ is isotropic for its polar pairing. Since $|k|^2=|\kappa|$,
nondegeneracy gives $k^\perp=k$. Summation on a coset of $k$ is zero unless
the coset is $k$ itself, and the sum on $k$ is $|k|$. Thus
$|\kappa|^{-1/2}\sum_\kappa H_U=1$. This computes the full refinement,
including any linear term.
\end{proof}

\subsubsection*{Reciprocal residual functions at odd \texorpdfstring{$T$}{T}}
If $T$ is even, the two remaining factors $g_E,g_\tau$ are already $1$.
For odd $T$ write $T=2e+1$, where $e=v_F(2)\ge1$ by Lemma~\ref{D:UR:oddT}.
We prove their exact cancellation, without dividing a characteristic-two
phase by two or guessing a square root.

\begin{lemma}[Critical quadratic norm reciprocity]\label{D:UR:resreciprocity}
For $w\in\pp_E^e$ one has
\begin{equation}\label{D:UR:taureciprocity}
 \tau(1+n_Ew)=\Psi_E(w).
\end{equation}
In particular, writing the actual norm-character residual function as
\[
 \mathcal H_\tau(z)=\Psi_F(-z)\tau(1+z)^{-1},\qquad z\in\pp_F^e,
\]
one has
\begin{equation}\label{D:UR:Htaucritical}
 \mathcal H_\tau(n_Ew)^{-1}=\Psi_E(w)\Psi_F(n_Ew).
\end{equation}
\end{lemma}
\begin{proof}
The trace of $w$ has valuation at least
$\floor{(e+T)/2}\ge e+1$, and $v_F n_Ew\ge e$.
The exact identity
\[
 n_E(1+w)=(1+n_Ew)
             \left(1+\frac{\Tr_{E/F}w}{1+n_Ew}\right)
\]
can therefore be evaluated by $\tau$, using its linear formula on the
second factor only. Since the left side is a norm,
\[
 \tau(1+n_Ew)
   =\Psi_F\!\left(\frac{\Tr w}{1+n_Ew}\right)
   =\Psi_E\!\left(\frac{w}{1+n_Ew}\right).
\]
The difference between $w/(1+n_Ew)$ and $w$ has $E$-valuation at least
$3e\ge2e+1=T$. This proves \eqref{D:UR:taureciprocity}, including $e=1$.
Equation \eqref{D:UR:Htaucritical} follows from the definition of the actual
residual function. It is a statement at the critical depth $e$;
it was not obtained by extending Lemma~\ref{D:UR:normphase} below its domain.
\end{proof}

\begin{lemma}[The ramified residual factor is the reciprocal]\label{D:UR:resE}
For every $1\le h<T$ when $T$ is odd,
\[
 g_E=g_\tau^{-1}.
\]
\end{lemma}
\begin{proof}
The critical $E$-ideal is $\pp_E^{e+h}$. For $v$ in this ideal, both the
core formula for $\chi_E^0$ and the base-twist formula apply. The first is
valid since $e+h\ge e+1=q$. For the second, the two terms in
$n_E(1+v)-1=\Tr v+n_Ev$ have $F$-valuations at least
\[
 \floor{(3e+h+1)/2}\quad\hbox{and}\quad e+h,
\]
respectively, both at least
$s_U=e+\ceil{(h+1)/2}$. It follows that
\begin{equation}\label{D:UR:Ecriticalchar}
 \theta_E(1+v)=\Psi_E(-(A+c)v)\Psi_F(-A n_Ev).
\end{equation}
The difference between $\widehat Z_E$ and $A-x+c$ is $b(1-c/x)$.
Its product with $v$ has valuation at least
\[
 2\floor{(T-h)/2}+e+h
       =3e+(h\bmod2)\ge T.
\]
Consequently the actual residual function at $v$ is
\begin{align*}
 \mathcal H_E(v)
   &=\Psi_E(-\widehat Z_Ev)\theta_E(1+v)^{-1}\\
   &=\Psi_E(xv)\Psi_F(A n_Ev)\\
   &=\mathcal H_\tau(n_E(xv))^{-1},
\end{align*}
where the last equality is Lemma~\ref{D:UR:resreciprocity}, since
$v_E(xv)\ge e$ and $n_E(xv)=A n_Ev$ exactly.

Choose a uniformizer $\Pi$ of $E$ and put $\pi=n_E\Pi$.
Write $x=\Pi^{-h}u$ with $u\in\OO_E^\times$, and use $\Pi^{e+h}$ and
$\pi^e$ as the two critical elements. On the two residue fields $k$ the
map $v\mapsto n_E(xv)$ induces
\[
 z\longmapsto\bar u^2 z^2.
\]
This is a bijection. Terminal representative invariance therefore gives
\[
 g_E=|k|^{-1/2}\sum_{z\in k}H_\tau(z)^{-1}
      =\overline{g_\tau}=g_\tau^{-1}.
\]
The final equality uses the unit magnitude of the full normalized
Lamprecht sum. No individual value of either Gauss factor was assumed.
\end{proof}

\begin{theorem}[Every positive critical unramified--ramified row]
\label{D:UR:criticalcomplete}
Every primitive compatible pair with $1\le h<T$ satisfies
\eqref{D:UR:URtarget}, in mixed and equal characteristic.
\end{theorem}
\begin{proof}
Proposition~\ref{D:UR:criticalratio} leaves $g_U/(g_Eg_\tau)$.
If $T+h$ is even, $g_U=1$ by the even-conductor formula; if it is odd,
Lemma~\ref{D:UR:resU} gives $g_U=1$. If $T$ is even, $g_E=g_\tau=1$.
If $T$ is odd, Lemma~\ref{D:UR:resE} gives $g_Eg_\tau=1$.
Thus the actual full ratio is $1$ in every case.
\end{proof}

\subsection{Sufficiently ramified common twists}
\begin{lemma}[Comparison of high stationary covectors]\label{D:UR:highcov}
Let $L/F$ be ramified quadratic with different exponent $\delta$, and let
$\lambda$ have conductor $n>\delta$. Let $b_F$ be a stationary covector
for $\lambda$, so $\lambda(1+z)=\psi_F(b_Fz)$ for
$z\in\pp_F^{\ceil{n/2}}$. A stationary covector $b_L$ for $\lambda\circ\Nm_{L/F}$
satisfies
\begin{equation}\label{D:UR:highclose}
 b_L/b_F\in U_L^{n-\delta}.
\end{equation}
For an unramified quadratic edge $b_F$ itself is a stationary covector
for the pullback.
\end{lemma}
\begin{proof}
On the ramified edge the pullback conductor is $M=2n-\delta$, and
$n_L(\psi_L)=2n_F+\delta$, so $v_Lb_L=v_Lb_F=-2(n+n_F)$.
For $z\in\pp_L^n$, its trace belongs to $\pp_F^{\ceil{n/2}}$ because
$\floor{(n+\delta)/2}\ge\ceil{n/2}$, and its norm belongs to $\pp_F^n$.
Exact norm expansion and the downstairs linear formula give
$(\lambda\circ\Nm_{L/F})(1+z)=\psi_L(b_Fz)$ on this \emph{whole} ideal.
The stationary formula for $b_L$ also applies there. Perfect additive
duality gives
\[
 b_L-b_F\in\pp_L^{-2n_F-\delta-n}.
\]
Divide by $b_F$ to obtain \eqref{D:UR:highclose}. On an unramified edge at depth
$s=\ceil{n/2}$, the norm endpoint has depth $2s\ge n$ and the trace has
depth at least $s$, so the same expansion gives the final assertion
at the full stationary depth.
\end{proof}

\begin{theorem}[Exact stable comparison]\label{D:UR:stable}
Let $(\chi_1,\chi_2)$ be any primitive compatible pair in a biquadratic
diagram, and put $a_i=a_{L_i}(\chi_i)$ and
$\delta_i=a_F(\omega_{L_i})$. Let $\lambda$ have conductor $n\ge2$.
Assume for every ramified $L_i/F$ that
\[
 n>\delta_i,\qquad n\ge a_i+\delta_i,\qquad n\ge2\delta_i,
\]
and for every unramified $L_i/F$ that $n\ge2a_i$.
For the actual twisted characters $\theta_i=\chi_i(\lambda n_i)$,
\begin{equation}\label{D:UR:stableidentity}
 A_{L_i}(\theta_i)=D(c_F)\Delta_F(\lambda)^2,
 \qquad c_F=b_F^{-1},\quad D=\chi_i|_F\,\omega_{L_i}.
\end{equation}
In particular these full expressions are independent of $i$.
\end{theorem}
\begin{proof}
The inequalities ensure that $\chi_i$ is within the exact stable-twist
range of $\lambda n_i$. On a ramified edge its stationary depth is
$n-\ceil{\delta_i/2}\ge a_i$; the unramified assertion follows from
$n\ge2a_i$. Ueda's exact stable twist gives
\[
 \Delta_{L_i}(\chi_i(\lambda n_i))
 =\chi_i(c_i)\Delta_{L_i}(\lambda n_i),\qquad c_i=b_{L_i}^{-1}.
\]
Lemma~\ref{D:UR:highcov} and $n-\delta_i\ge a_i$ show
$\chi_i(c_i)=\chi_i(c_F)$. On the unramified edge $c_i=c_F$ may be used.
The First Main Lemma on $L_i/F$ is
\[
 \Delta_{L_i}(\lambda n_i)\Delta_F(\omega_{L_i})
 =\Delta_F(\lambda)\Delta_F(\lambda\omega_{L_i}).
\]
A downstairs stable twist is valid since $n\ge2\delta_i$ (and
$\delta_i=0$ on the unramified edge). Its right side is
$\omega_{L_i}(c_F)\Delta_F(\lambda)^2$.
Multiply the two formulas and use Lemma~\ref{U:determinant}.
This argument retains all quadratic norm-character factors and all
residual factors; stable twisting cancels the latter exactly.
\end{proof}
\begin{corollary}\label{D:UR:URstable}
In the unramified--ramified diagram, every primitive pair with $h\ge T$
in \eqref{D:UR:conductors} satisfies the Second Main comparison.
\end{corollary}
\begin{proof}
Use the actual minimal pair from Theorem~\ref{D:UR:models} and
Proposition~\ref{D:UR:actual}. Its two conductors are $T$, the lower different
exponents are $0,T$, and $a_F(\lambda)=T+h\ge2T$.
All hypotheses of Theorem~\ref{D:UR:stable} hold.
\end{proof}

\subsection{Completion of the unramified--ramified case}
\begin{theorem}[Complete dyadic unramified--ramified comparison]\label{D:UR:URall}
Let $K/F$ be biquadratic over a local field of residue characteristic two,
and assume it has an unramified quadratic intermediate field $U$.
For every other quadratic intermediate field $E$ and every primitive
compatible pair $(\theta_U,\theta_E)$, one has
\[
 \Delta_U(\theta_U,\psi_U)\Delta_F(\omega_U,\psi_F)
 =\Delta_E(\theta_E,\psi_E)\Delta_F(\omega_E,\psi_F).
\]
The complete induction expression is therefore independent of all three
quadratic intermediate fields. The assertion holds in mixed and equal
characteristic, for every conductor.
\end{theorem}
\begin{proof}
The other quadratic fields are ramified, since an unramified local Galois
group is cyclic. For a comparison with either such $E$,
Proposition~\ref{D:UR:actual} gives conductors $T+h,T+2h$ with $h\ge0$.
The case $h=0$ is Theorem~\ref{D:UR:minimal}; the cases $1\le h<T$ are
Theorem~\ref{D:UR:criticalcomplete}; and every $h\ge T$ is
Corollary~\ref{D:UR:URstable}. This exhausts the conductors. Comparing each
ramified field with $U$ gives independence of all three by transitivity.
The possible extensions of a fixed upstairs character differ by the
upper norm character and are conjugate, by Lemma~\ref{U:conjugacy}; their local
constants agree by an automorphism change of variables. Thus the
conclusion does not depend on those choices either.
\end{proof}

The ordinary common norm $N_E(x+d)$ was not used as stationary
outside its proved range. Its replacement is $\rho N_E(x+d)$, with
$\rho=x'/x$, and the nontrivial character value of $\rho$ was retained.
The sign from the last trace row was derived from an actual norm identity,
and the two odd-$T$ residual factors were compared pointwise, not merely
through their squares or polar forms.


\section{Totally ramified extensions in characteristic two}\label{sec:dyadic-equal}
\subsection{Statement and normalization}
Let $F=k((\varpi))$, where $k$ is a finite field of characteristic $2$.
Let $K/F$ be totally ramified and Galois with group $C_2^2$.
Write its three nontrivial automorphisms as $g_1,g_2,g_3$, with
$g_1g_2=g_3$, and put
\[
 L_i=K^{\langle g_i\rangle},\qquad
 N_i=\Nm_{K/L_i},\quad S_i=\Tr_{K/L_i},\quad
 n_i=\Nm_{L_i/F},\quad T_i^{\rm tr}=\Tr_{L_i/F}.
\]
We reserve $T_i$ without a superscript for a different/conductor exponent.
Let $\omega_i$ be the nontrivial character of
$F^\times/n_iL_i^\times$.
\begin{theorem}\label{D:EQ:main}
Suppose $\Theta$ is $\Gal(K/F)$-invariant and is not a norm pullback
from $F$. Choose characters $\theta_i$ of $L_i^\times$ with
$\theta_iN_i=\Theta$. For any nontrivial additive character $\psi_F$,
put $\psi_{L_i}=\psi_FT_i^{\rm tr}$. Then
\begin{equation}\label{D:EQ:Adef}
 \mathcal A_i(\theta_i)=
 \Delta_{L_i}(\theta_i,\psi_{L_i})\Delta_F(\omega_i,\psi_F)
\end{equation}
is independent of $i$. Continuous quasi-characters are allowed.
\end{theorem}
The local constants and their admissible denominators have Ueda's
normalization. In particular, the omitted trivial-character factor is $1$.
Only restrictions to compact unit groups enter the finite calculations.

We use the First Main Lemma, Lamprecht's formula, exact stable
twisting, additive scaling and duality, norm filtration, and the
trace and conductor formulas from~\cite{Ueda}, recalled in
Sections~\ref{sec:local-notation}--\ref{sec:stationary}.
The subcritical norm representatives are those of
\cite[Lemma~6.6]{Ueda}.
On a ramified quadratic edge $E/D$ of break $t$, these give
\begin{equation}\label{D:EQ:edge}
 \Tr_{E/D}(\pp_E^a)=\pp_D^{\floor{(a+t+1)/2}},\qquad
 v_D(\Nm x)=v_E(x),\qquad a_D(\omega_{E/D})=t+1.
\end{equation}
If $a_D(\lambda)>t+1$, its norm pullback has conductor
$2a_D(\lambda)-(t+1)$. For $0\le j\le t$, an arbitrary $A\in D^\times$
has a norm approximation with relative error in $U_D^j$ and the correct
valuation. No approximation modulo $U_D^{t+1}$ is assumed.
We also use Hilbert 90, extension of abelian characters, elementary
Artin--Schreier theory, and the local trace--residue identity for
separable Laurent-series fields. The residue formula for the norm
characters needed here is proved next. There is no higher-dimensional
induction theorem among the inputs.

Throughout this section $K/F$ is totally ramified.
Fix a uniformizer $\pi$ of $K$ and take
\begin{equation}\label{D:EQ:compatiblepi}
 \Pi_i=N_i\pi,\qquad \varpi=\Nm_{K/F}\pi=n_i\Pi_i.
\end{equation}
These are compatible uniformizers; the coefficient field $k$ is contained
in all four fields. Write $\tr=\Tr_{k/\F_2}$.
We first prove the theorem for
\begin{equation}\label{D:EQ:canonicalpsi}
 \psi_F(x)=(-1)^{\tr\Res_F(x\,d\varpi)}.
\end{equation}
Its largest trivial ideal is $\OO_F$. The passage to an arbitrary
additive character is given at the end.

\subsection{Local Artin--Schreier characters and a residue identity}
For a Laurent differential in a uniformizer $v$, define
\[
 \Cart\left(\sum_n a_nv^n\,dv\right)
       =\sum_m a_{2m+1}^{1/2}v^m\,dv.
\]
Coefficient square roots are unique because the residue field is perfect.
The identities used below follow from this formula:
\begin{equation}\label{D:EQ:Cartier}
 \Cart(f^2\eta)=f\Cart(\eta),\qquad
 \Cart(dv/v)=dv/v,\qquad
 \tr\Res\Cart(\eta)=\tr\Res\eta.
\end{equation}
Also $\Cart(du/u)=du/u$ for every nonzero Laurent series $u$.
One direct verification writes a principal unit, successively in its
coefficients, as a convergent product of factors $1-cv^m$.
Factors with even $m$ have zero logarithmic derivative. For odd $m$,
\[
 d\log(1-cv^m)=\sum_{j\ge1}c^jv^{mj-1}\,dv,
\]
and the displayed formula for $\Cart$ retains $j=2j'$ and returns the
same series. The valuation and constant factors of $u$ give the other
parts of the assertion. Every residue calculation uses only finitely
many terms of these convergent formal series.

We use the local trace--residue identity proved in
Theorem~\ref{B:trace-residue}
\begin{equation}\label{D:EQ:residue-trace}
 \Tr_{k_E/k_D}\Res_E(\eta)
       =\Res_D\bigl(\Tr_{E/D}\eta\bigr)
\end{equation}
for a finite separable extension of equal-characteristic local fields.
Here the trace of $x\,d\varpi_D$ is
$\Tr_{E/D}(x)\,d\varpi_D$. Appendix~\ref{app:local-residues}
proves this formal identity by finite free Laurent-series expansions
and coefficient specialization; wild ramification is included.
The logarithmic norm identity needed with it is
Lemma~\ref{B:dlog-norm}. Thus both ingredients have proofs here,
independently of a residue theorem on a global curve.

\begin{lemma}[The actual norm character]\label{D:EQ:ASsymbol}
Let $z^2+z=f\in F$ define a nontrivial quadratic extension $E/F$.
Then its actual norm character is
\begin{equation}\label{D:EQ:symbol}
 \omega_f(u)=(-1)^{\tr\Res_F(f\,du/u)}.
\end{equation}
If $f=f_0+\sum_{j>0,\ j\text{ odd}}f_j\varpi^{-j}$ is reduced and
its largest pole is $t>0$, this character has conductor $t+1$.
\end{lemma}
\begin{proof}
The formula is multiplicative in $u$. Adding $w^2+w$ to $f$ does not
change it: by \eqref{D:EQ:Cartier} and the logarithmic identity, the residues
of $w^2du/u$ and $wdu/u$ have the same absolute trace.
Every Artin--Schreier class has a reduced representative of the stated
form. Remove a negative even power by adding $w^2+w$ and repeat; remove
the positive-depth part using the bijectivity of $w\mapsto w^2+w$ on
the maximal ideal. Constants are taken modulo the same map on $k$.

For $u=\Nm_{E/F}v$, \eqref{D:EQ:residue-trace} makes its exponent equal to
the absolute trace of
$\Res_E((z^2+z)\,dv/v)$, which is zero by \eqref{D:EQ:Cartier}.
Thus the formula is trivial on norms.
For a reduced representative with pole $t$, it is trivial on $U_F^{t+1}$,
where $f\,du/u$ has nonnegative valuation. At
$u=1+c\varpi^t$ its last-layer exponent is $\tr(f_tc)$, which is
nontrivial as $c$ varies. If instead the class is an unramified nonzero
constant, evaluation at a uniformizer is nontrivial.
Consequently the formula is nontrivial for every nonzero
Artin--Schreier class. The quadratic norm quotient has order two, so
this nontrivial character trivial on norms is precisely its norm
character. The conductor assertion follows as well.
\end{proof}

\begin{lemma}[A local Cartier calculation]\label{D:EQ:rationalCartier}
For $f,w\in F$, the denominator $1+\varpi w^2$ is nonzero and
\begin{equation}\label{D:EQ:Cartier-rational}
 \tr\Res_F\left(\frac{\varpi f^2}{1+\varpi w^2}\,d\varpi\right)
 =\tr\Res_F\left(\frac{f}{1+\varpi w^2}\,d\varpi\right).
\end{equation}
\end{lemma}
\begin{proof}
Nonvanishing follows from the parity of valuations. Put
$H=1+\varpi w^2$ and write
\[
 \frac{\varpi f^2}{H}
   =\varpi(f/H)^2+\varpi^2(fw/H)^2.
\]
Applying \eqref{D:EQ:Cartier} gives $(f/H)\,d\varpi$, and taking absolute
residue traces proves the assertion.
\end{proof}

\subsection{A simultaneous origin and its ramification data}
Choose the labeling so that $L_1/F$ has the smallest break, and take
reduced representatives
\[
 z_i^2+z_i=f_i\quad(i=1,2),\qquad z_3=z_1+z_2,\quad f_3=f_1+f_2.
\]
Because $K/F$ is totally ramified, all three nonzero classes have a pole.
Write their breaks as $t_1\le t_2=t_3$, all odd, and put
\begin{equation}\label{D:EQ:rledger}
 T_1=t_1+1=2r_1,\qquad T_2=t_2+1=2r_2,\qquad
 \delta=r_2-r_1\ge0.
\end{equation}
In the one-break case the two leading pole coefficients are distinct;
otherwise their sum would have smaller conductor and would be the
first field. In the two-break case $t_2-t_1=2\delta$.

Write $f_i=f_{i,0}+f_i^-$, where $f_i^-$ is the negative-power part.
Define
\begin{equation}\label{D:EQ:dbeta}
 \beta_i=\varpi^{t_2}f_i^-=d_i^2\in\OO_F,\qquad d_3=d_1+d_2.
\end{equation}
The square roots exist in $F$: every exponent $t_2-j$ is even, and
$k$ is perfect. Thus $v_Fd_1=\delta$ and $d_2,d_3$ are units.
Set
\begin{equation}\label{D:EQ:origin}
 Y=d_1z_2+d_2z_1,\quad
 B=\beta_1f_2+\beta_2f_1
       =\beta_1f_{2,0}+\beta_2f_{1,0}\in\OO_F,\quad
 \kappa_i=d_jd_k\quad(\{i,j,k\}=\{1,2,3\}).
\end{equation}
The automorphisms are labeled so that $g_1$ changes $z_2$ by $1$ and
fixes $z_1$, while $g_2$ changes $z_1$ by $1$ and fixes $z_2$.
Exact identities in the fields are
\begin{align}
 g_iY-Y&=d_i,&S_iY&=d_i,&
 a_i:=N_iY&=\kappa_i z_i+B.                         \label{D:EQ:exact-origin}
\end{align}
For example $Y^2=\beta_1z_2+\beta_2z_1+B$, which proves every formula.
In particular,
\begin{equation}\label{D:EQ:originval}
 v_KY=-t_1,\qquad v_{L_i}a_i=-t_1,\qquad
 \mathcal N:=\Nm_{K/F}Y=n_i a_i,\quad v_F\mathcal N=-t_1.
\end{equation}
For $i=2,3$, $v_{L_i}\kappa_i=2\delta=t_2-t_1$; for $i=1$ its
valuation is zero. These observations prove the second assertion of
\eqref{D:EQ:originval}, and the first follows by the norm valuation formula.

The upper breaks are $u_1=2t_2-t_1$ and $u_2=u_3=t_1$.
For completeness, a ramification automorphism $g$ changes an element
of odd valuation $v$ by an element of exact valuation $v+u_g$:
expand in a uniformizer and use the first nonzero ramification term.
Apply this to $Y$ in \eqref{D:EQ:exact-origin}. The valuation of $d_1$ in $K$
is $4\delta$, and those of $d_2,d_3$ are zero, proving these upper
breaks. Thus the upper different exponents are
\begin{equation}\label{D:EQ:upperdifferent}
 D'_1=2T_2-T_1,\qquad D'_2=D'_3=T_1.
\end{equation}

Put
\begin{equation}\label{D:EQ:Psis}
 \alpha=\varpi^{-T_2},\qquad
 \Psi_F(x)=\psi_F(\alpha x),\qquad
 \Psi_{L_i}(x)=\Psi_F(T_i^{\rm tr}x).
\end{equation}
The largest trivial ideals for these characters have exponents
\begin{equation}\label{D:EQ:Jledger}
 J_F=T_2,\quad J_1=2T_2-T_1,\quad J_2=J_3=T_2.
\end{equation}
Do not confuse $J_2=T_2$ with the upper different $D'_2=T_1$.
Since $T_2$ is even, $\Psi_F(x^2)=1$ for every $x\in F$.
Lemma~\ref{D:EQ:ASsymbol}, integration by parts for a formal residue, and
$df_i=\alpha\beta_i\,d\varpi$ give
\begin{equation}\label{D:EQ:lowerchart}
 \omega_i(1-z)=\Psi_F(\beta_i z)\qquad(z\in\pp_F^{r_i}),
 \qquad r_3=r_2.
\end{equation}
Indeed replacing $(1-z)^{-1}$ by $1$ in $f_i\,dz/(1-z)$ costs no
residue at depth $r_i$: the error has valuation at least
$-t_i+r_i+(r_i-1)=0$.
Notice the exact simultaneous norm property
\begin{equation}\label{D:EQ:loweractualnorm}
 \beta_i=n_i(S_iY)=n_i(d_i),\qquad \omega_i(\beta_i)=1.
\end{equation}
These are full stationary coefficients, not just their leading residues.
Finally exact quadratic norm expansion yields the phase conversion
\begin{equation}\label{D:EQ:normphase}
 \Psi_{L_i}(\beta_i z)=\Psi_F(\beta_i n_i z)
                      \qquad(z\in\pp_{L_i}^{r_i}).
\end{equation}
Both trace and norm of $z$ lie in $\pp_F^{r_i}$; evaluate
\eqref{D:EQ:lowerchart} on $n_i(1-z)=1+T_i^{\rm tr}z+n_i z$ and use
triviality on norms. This establishes the conversion on its whole
specified ideal.

\subsection{Construction of compatible characters}
Define the minimal conductors and stationary depths
\begin{equation}\label{D:EQ:coreledger}
 \begin{aligned}
 m_1&=2T_2-1=2t_2+1,&s_1&=T_2,\\
 m_2=m_3&=T_1+T_2-1=t_1+t_2+1,&s_2=s_3&=r_1+r_2.
 \end{aligned}
\end{equation}
Thus $m_i=J_i+t_1$, $m_i=2s_i-1$, and the formulas
\begin{equation}\label{D:EQ:Rmodel}
 R_i(1-z)=\Psi_{L_i}(a_i z),\qquad z\in\pp_{L_i}^{s_i}
\end{equation}
define characters of $U_{L_i}^{s_i}$ with exact conductor $m_i$.
Multiplicativity follows since $2s_i-t_1\ge J_i$.

\begin{lemma}[The required full commutator compatibility]\label{D:EQ:commmodel}
Let $\sigma$ be the nontrivial automorphism of $L_1/F$ and
$\nu_1=\omega_2n_1$, the actual norm character for $K/L_1$.
For every $v\in L_1^\times$ with $\sigma(v)/v\in U_{L_1}^{s_1}$,
\begin{equation}\label{D:EQ:commmodel-eq}
 R_1(\sigma(v)/v)=\nu_1(v).
\end{equation}
\end{lemma}
\begin{proof}
The description of $\nu_1$ follows from Lemmas~\ref{D:crossed-norm}
and \ref{D:quadratic-norms}, whose characteristic-two proof uses only
\eqref{D:EQ:symbol} and \eqref{D:EQ:residue-trace}.
Multiplication of $v$ by an element of $F^\times$
does not affect either side. Reduce its valuation to $0$ or $1$.
In valuation $1$ its conjugate quotient has exact depth $t_1<s_1$,
by the first ramification term. Thus only units occur.
Write $v=a+bz_1$, with $a,b\in F$. The two summands have valuations of
opposite parity, since $t_1$ is odd. Consequently $a$ is a unit and
$v_Fb\ge r_1$. The conjugate quotient is $1+b/v$; the assumed depth
$s_1=2r_2$ forces $v_Fb\ge r_2$. Put $w=b/a$, so $v_Fw\ge r_2$.
The case $b=0$ is immediate. Exact quadratic trace calculation gives
\begin{equation}\label{D:EQ:comm-rational}
 T_1^{\rm tr}\left(a_1\frac b v\right)
   =\frac{\kappa_1ab+B b^2}{n_1v}
   =\frac{\kappa_1w+B w^2}{1+w+w^2f_1}.
\end{equation}

Here is an independent evaluation of the other side. Put
$R=w^2f_1$, $H=1+R$ and $V=H+w$. These are units, since
$v_FR\ge2r_2-t_1\ge1$. Factor $V=H(1+w/H)$.
The residue formula evaluates $\omega_2(H)$ exactly, and the full
linear chart \eqref{D:EQ:lowerchart} applies to the second factor. They give
\[
 \omega_2(V)=
 \Psi_F\left(\frac{\beta_1f_2w^2+\beta_2w}{H}\right)
 =\Psi_F\left(\frac{\beta_2R+\beta_2w}{H}\right).
\]
In the last equality $B w^2/H$ was discarded at depth at least $T_2$,
using $\beta_1f_2+\beta_2f_1=B\in\OO_F$.
Write $R_0=\varpi W^2$, where $W=\varpi^{-r_2}d_1w\in\OO_F$.
Then $R-R_0=f_{1,0}w^2\in\pp_F^{T_2}$, so it can be discarded in
all these unit denominators. Lemma~\ref{D:EQ:rationalCartier}, with
$f=\varpi^{-r_2}d_2W$, proves
\[
 \Psi_F\left(\frac{\beta_2R_0}{1+R_0}\right)
   =\Psi_F\left(\frac{d_1d_2w}{1+R_0}\right).
\]
It follows that $\omega_2(V)=\Psi_F(\kappa_1w/H)$ because
$\beta_2+d_1d_2=d_2d_3=\kappa_1$.
Replacing $H$ by $V$ here costs only $\kappa_1w^2/(HV)$, of depth
at least $T_2$. The same is true of the $B w^2/V$ term in
\eqref{D:EQ:comm-rational}. Since $n_1v=a^2V$ and $\omega_2(a^2)=1$,
\eqref{D:EQ:comm-rational} proves \eqref{D:EQ:commmodel-eq}.
Every discarded expression was evaluated at its actual $F$-character
modulus $\pp_F^{T_2}$.
\end{proof}

\begin{lemma}[Compatibility on the full common unit subgroup]\label{D:EQ:Rcompat}
For $z\in\pp_K^{T_2}$ all $N_i(1-z)$ lie in the domains of $R_i$, and
$R_i(N_i(1-z))$ is independent of $i$.
\end{lemma}
\begin{proof}
The norm term has $L_i$-valuation at least $T_2$.
The upper trace terms have respective depths at least
$3r_2-r_1$ and $r_1+r_2$; these are at least $s_1$ and $s_2$.
The identity of Lemma~\ref{U:e2} is
\begin{equation}\label{D:EQ:E2}
 E_2(C)=T_i^{\rm tr}N_iC+n_iS_iC.
\end{equation}
Apply it to $Y(1-z)$ and $Y$. Put $h_i=S_i(Yz)/d_i$.
The upper trace formula and \eqref{D:EQ:originval} give
\[
 v_{L_1}S_1(Yz)\ge r_1+3\delta,\qquad
 v_{L_2}S_2(Yz),v_{L_3}S_3(Yz)\ge r_2.
\]
After division by $d_1$, of $L_1$-valuation $2\delta$, the first
bound is also $r_2$. Thus $v_{L_i}h_i\ge r_i$ on every side.
Since $S_i(Y(1-z))=d_i(1-h_i)$, its lower norm divided by $\beta_i$
is an actual norm. Formula \eqref{D:EQ:lowerchart} applied to that norm
quotient shows that the additive phase of
$n_iS_i(Y(1-z))-\beta_i$ is $1$.
The remaining term in \eqref{D:EQ:E2} is
$T_i^{\rm tr}(a_i(S_iz+N_iz))$, exactly the exponent of
$R_i(N_i(1-z))$. Therefore this value is
$\Psi_F(E_2(Y(1-z))-E_2(Y))$, independent of $i$.
\end{proof}

\begin{theorem}[Existence of actual minimal models]\label{D:EQ:core-exists}
There are characters $\chi_i$ on $L_i^\times$, with exact conductors
$m_i$, primitive common norm pullback, and
\begin{equation}\label{D:EQ:corechart}
 \chi_i(1-z)=\Psi_{L_i}(a_i z)
                    \qquad(z\in\pp_{L_i}^{s_i}).
\end{equation}
\end{theorem}
\begin{proof}
Let $I=\{\sigma(v)/v:v\in L_1^\times\}$ and prescribe
$\kappa_I(\sigma(v)/v)=\nu_1(v)$. This is well defined because
$\nu_1$ is trivial on $F^\times$: it is $\omega_2(n_1x)=\omega_2(x^2)$
there. Lemma~\ref{D:EQ:commmodel} proves agreement with $R_1$ on
$I\cap U_{L_1}^{s_1}$. Their product character on
$I U_{L_1}^{s_1}$ is trivial on $U_{L_1}^{m_1}$.
Extend it across the abelian group $L_1^\times/U_{L_1}^{m_1}$.
One can first extend on its finite unit subgroup and then choose the
uniformizer value. The resulting $\chi_1$ is continuous, has exact
conductor $m_1$, and satisfies $\chi_1^\sigma/\chi_1=\nu_1$ globally.

The character $\chi_1N_1$ is invariant under $G$: its other conjugate
quotient is $\nu_1N_1=\omega_2\Nm_{K/F}=1$.
Hilbert 90 therefore makes it trivial on $\ker N_i$ for $i=2,3$;
extend its descent from the norm image to $\chi_i$ on $L_i^\times$.
Above the upper break $t_1$, norm filtration gives
\[
 N_i(U_K^{T_2+1})=U_{L_i}^{s_i}\qquad(i=2,3),
\]
since $2s_i-t_1=T_2+1$. Lemma~\ref{D:EQ:Rcompat} then proves
\eqref{D:EQ:corechart} on those entire groups, not just on a last layer.
Their conductors are at least $m_i$. The common pullback has conductor
$2m_1-D'_1=t_1+2t_2+1$; since $m_i>D'_i$, the high pullback-conductor
formula forces their conductors to equal $m_i$.
Finally, descent of the common pullback from $F$ would make $\chi_1$
a base pullback times its upper norm character. Both are invariant
under $\sigma$, contradicting the nontrivial commutator $\nu_1$.
\end{proof}

\subsection{Common-twist realization and conductor ranges}


Write $\nu_i$ for the nontrivial norm character of $K/L_i$.
The common restriction $D_\theta=\theta_i|_{F^\times}\omega_i$
is the one in Lemma~\ref{U:conjugacy}.

\begin{proposition}[Actual common-twist realization]\label{D:EQ:realize}
For every primitive compatible family there is a core family satisfying
\eqref{D:EQ:corechart}, and a character $\lambda$ of $F^\times$, such that
\begin{equation}\label{D:EQ:actual-family}
 \theta_i=\chi_i(\lambda n_i).
\end{equation}
Put $n=a_F(\lambda)$ and $h=n-T_2$.
If $n\le T_2+r_1-1$, the conductors are exactly $m_i$.
If $n\ge T_2+r_1$, they are $M_i=2n-T_i$, all even.
These alternatives exhaust all primitive conductors.
\end{proposition}
\begin{proof}
Start with Theorem~\ref{D:EQ:core-exists}. By Lemma~\ref{U:conjugacy},
$\theta_1/\chi_1$ is invariant under $\Gal(L_1/F)$.
Hilbert 90 and abelian character extension make it $\lambda n_1$.
For $i=2,3$, the remaining quotient is a power of $\nu_i$.
Absorb it into $\chi_i$. This does not change its common pullback or
\eqref{D:EQ:corechart}, because $a_{L_i}(\nu_i)=T_1\le s_i$.

Primitivity forces $a_{L_i}(\theta_i)\ge t_i+u_i+1=m_i$:
on the last unit row of the nontrivial commutator, the first
ramification term in $\sigma(v)/v$ has exact depth $t_i+u_i$.
Its coefficient is nonzero since $u_i$ is odd; a character of smaller
conductor would make the commutator trivial on that row.
For $n>T_i$, the pullback $\lambda n_i$ has conductor $2n-T_i$;
for $n\le T_i$ it has conductor at most $T_i$.
The number $2n-T_i$ is even, whereas $m_i$ is odd, so cancellation
at equal conductors is impossible. The same threshold
$n=T_2+r_1$ is obtained for all three fields. This proves the assertions.
\end{proof}

\subsection{Moving the common origin by a norm from the third field}
We compare fields $1$ and $2$ and use $L_3$ for the adjustment. The same
construction compares $1$ and $3$ with their roles exchanged.
For $X\in L_3$, put
\begin{equation}\label{D:EQ:adjust}
 A=n_3X,\quad a=T_3^{\rm tr}X,\quad C=Y+X,\quad
 Q_i=aY+d_iX\in L_i\quad(i=1,2).
\end{equation}
The following are exact algebraic identities:
\begin{align}
 N_iC&=a_i+A+Q_i,&S_iC&=d_i+a,                    \label{D:EQ:commonC}\\
 n_iQ_i&=\beta_iA+E,&
 E&=a(a a_3+\kappa_3X)\in F.                    \label{D:EQ:commonE}
\end{align}
For example, write $X=p+qz_3$. Then $a=q$,
$Q_1=d_1p+qd_3z_1$ and $Q_2=d_2p+qd_3z_2$; taking the two quadratic
norms proves \eqref{D:EQ:commonE}, with $E=q^2B+q\kappa_3p$.

If $v_{L_3}X=-h$, trace ideals and \eqref{D:EQ:originval} give
\begin{align}
 v_F a&\ge r_2-\ceil{h/2},                       \label{D:EQ:a-bound}\\
 v_{L_1}Q_1&\ge2\delta-h,&
 v_{L_2}Q_2&\ge-h,                              \label{D:EQ:Q-bound}\\
 v_F E&\ge r_2+\delta-h.                        \label{D:EQ:E-bound}
\end{align}
These statements include negative $h$ and vanishing traces.
For the last one, the invariant element $a a_3+\kappa_3X$ has
$L_3$-valuation at least $2\delta-h$, hence $F$-valuation at least
$\delta-\floor{h/2}$; combine this with \eqref{D:EQ:a-bound}.
For the $Q_i$ bounds one can use directly the upper trace of
$Yg_iX$, with valuation $-t_1-2h$, and the upper differents
\eqref{D:EQ:upperdifferent}.

Suppose $\lambda(1-z)=\Psi_F(Az)$ on a downstairs ideal containing
both trace and norm of a variable $z\in L_i$.
Whenever $(Q_i/\beta_i)z\in\pp_{L_i}^{r_i}$, equations
\eqref{D:EQ:normphase} and \eqref{D:EQ:commonE} give
\begin{equation}\label{D:EQ:pullback-adjust}
 \Psi_{L_i}(Q_iz)
 =\Psi_F\bigl((A+E/\beta_i)n_i z\bigr).
\end{equation}
Thus, if $(E/\beta_i)n_i z$ is in $\pp_F^{T_2}$, the full pulled-back
character has coefficient $A+Q_i$. No exact equation
$n_i(Q_i/\beta_i)=A/\beta_i$ is asserted.

\subsection{Stationary origins at minimal conductors}
Let the family be realized as in \eqref{D:EQ:actual-family}, with
$n\le T_2+r_1-1$. Put $q_0=r_1+r_2$.
For $z\in\pp_{L_i}^{s_i}$, both $T_i^{\rm tr}z$ and $n_i z$ lie in
$\pp_F^{q_0}$. For $i=1$ the trace bound is exactly $q_0$;
for $i=2,3$ its bound is at least $q_0$.
If $n\le q_0$, take $X=0$ and $C=Y$ below.
Otherwise put $h=n-T_2$, so
\begin{equation}\label{D:EQ:min-h}
 1-\delta\le h\le r_1-1.
\end{equation}
Since $2q_0\ge n$, additive duality supplies
\[
 \lambda(1-z)=\Psi_F(A_*z)\quad(z\in\pp_F^{q_0}),\qquad
 v_FA_*=-h,
\]
with ambiguity $\pp_F^{T_2-q_0}=\pp_F^\delta$.
Choose $X\in L_3$ by norm approximation so that
\begin{equation}\label{D:EQ:min-approx}
 A=n_3X\equiv A_*\pmod{\pp_F^\delta},\qquad v_{L_3}X=-h.
\end{equation}
The relative precision requested is $\delta+h\le r_2-1\le t_2$,
so this choice is legitimate. The exact norm $A$ may replace $A_*$
on the whole indicated ideal.

\begin{lemma}[The full stationary properties]\label{D:EQ:min-origin}
The element $C$ just constructed has valuation $-t_1$ in $K$.
For $i=1,2$, its actual upper norms and lower trace norms satisfy
\begin{align}
 \theta_i(1-z)&=\Psi_{L_i}(N_iC\,z)
                            &&(z\in\pp_{L_i}^{s_i}),       \label{D:EQ:min-core}\\
 \omega_i(1-z)&=\Psi_F(n_iS_iC\,z)
                            &&(z\in\pp_F^{r_i}).          \label{D:EQ:min-lower}
\end{align}
If $\delta=0$, the same single $C$ satisfies these identities for all
three indices.
\end{lemma}
\begin{proof}
The case $X=0$ follows from the core formulas. Otherwise
$-2h>-t_1$ by \eqref{D:EQ:min-h}, so $C$ has the valuation of $Y$.
For $i=1,2$, \eqref{D:EQ:Q-bound} gives
\[
 v_{L_1}((Q_1/\beta_1)z)\ge T_2-2\delta-h=2r_1-h\ge r_1,
 \quad
 v_{L_2}((Q_2/\beta_2)z)\ge r_1+r_2-h\ge r_2.
\]
The error in \eqref{D:EQ:pullback-adjust}, at these depths, has valuations
at least $T_2+r_1-h$ and $3r_2-h$, respectively; both are at least
$T_2$. Therefore exact norm expansion evaluates the actual pullback
$\lambda n_i$ with coefficient $A+Q_i$. Combining it with the core
coefficient gives $N_iC$ by \eqref{D:EQ:commonC}.
Moreover $n_iS_iC=(d_i+a)^2=\beta_i+a^2$.
The bound
\[
 2v_Fa\ge2r_2-2\ceil{h/2}\ge2r_2-r_1=T_2-r_1
\]
proves \eqref{D:EQ:min-lower} on both full lower stationary ideals.
It also shows $v_Fa>\delta$, so the traces $d_i+a$ are nonzero with
the original valuations.

When $\delta=0$, $r_1=r_2=r$ and $h\le r-1$.
On the third side the exact norm-phase conversion gives the
pullback coefficient $A+d_3X$, because
$n_3(X/d_3)=A/\beta_3$. Its argument has depth at least $2r-h\ge r$.
But
\[
 N_3C=a_3+X^2+d_3X
       =a_3+A+d_3X+aX,
\]
and $v_{L_3}(aX)\ge2r-2h-(h\bmod2)\ge0$.
The last term is invisible on $\pp_{L_3}^{2r}$.
Finally $S_3C=d_3$, so its lower norm coefficient is unchanged.
This proves the assertion for all three fields in the one-break case.
\end{proof}

\subsection{The minimal sign and the critical norm image}
Use Lemma~\ref{U:stationary-factor} with stationary coefficient $A_L$
and critical element $\Pi_L^{\lfloor m/2\rfloor}$, writing its residual
function and sum as $H_L,g_L$. In characteristic two the field minus
signs in that lemma disappear; complex phases are unchanged.

For the $C$ of Lemma~\ref{D:EQ:min-origin}, use $N_iC$ for the upper
coefficient and $n_iS_iC$ for the lower one. The latter is an actual
norm and the lower conductor $T_i$ is even. Compatibility and
\eqref{D:EQ:E2} therefore give
\begin{equation}\label{D:EQ:minfull}
 \mathcal A_i(\theta_i)
  =D_\theta(\alpha^{-1})\Theta(C)^{-1}
                       \Psi_F(E_2(C))g_i.
\end{equation}
All character values and elementary phases in this identity have been
retained.

For $z\in\OO_K$ vary $C$ to $C_z=C(1+\pi^{t_2}z)$.
The relative upper norms are on the terminal stationary rows, of
depths
\[
 \ell_1=t_2,\qquad \ell_2=\ell_3=(t_1+t_2)/2.
\]
The changes $S_i(C\pi^{t_2}z)$ have valuations at least
$r_1+3\delta$ and $r_2$ on sides $1$ and $2$, respectively.
After division by $S_iC$, both relative changes have valuation at
least $r_2\ge r_i$. The actual lower norm quotient can thus be evaluated
by \eqref{D:EQ:min-lower}; its additive phase is $1$.
Subtracting \eqref{D:EQ:E2} for $C_z,C$, and using compatibility and
terminal-coset invariance, proves
\begin{equation}\label{D:EQ:commonH-general}
 H_i(p_i(z))=Q(z),
\end{equation}
where $p_i$ is the actual graded upper norm map at this row and $Q$
is independent of the compared index. These are identities between
full functions, not just between polar forms.

\begin{proposition}[The two-break minimal row]\label{D:EQ:min-two}
If $t_1<t_2$, every minimal family satisfies $\mathcal A_1=\mathcal A_2$.
\end{proposition}
\begin{proof}
On side $1$, $t_2$ is below the upper break $2t_2-t_1$.
The trace term is at least one level deeper than the norm term, so
$p_1(z)=z^2$, with compatible uniformizers \eqref{D:EQ:compatiblepi}.
On side $2$, the trace induces an isomorphism of the one-dimensional
residue lines at source depth $t_2$ and target depth $(t_1+t_2)/2$:
\[
 \floor{(t_2+T_1)/2}=(t_1+t_2)/2,\qquad
 \floor{(t_2+1+T_1)/2}=(t_1+t_2)/2+1.
\]
The norm endpoint is strictly deeper. Hence $p_2$ is a nonzero
$k$-linear map. Both $p_i$ are bijections of $k$.
Summing \eqref{D:EQ:commonH-general} gives $g_1=g_2$, and
\eqref{D:EQ:minfull} proves equality. The same comparison with $L_3$
proves independence of every field in this configuration.
\end{proof}

\begin{proposition}[The one-break minimal row]\label{D:EQ:min-one}
If $t_1=t_2=t$, every minimal family satisfies
$\mathcal A_1=\mathcal A_2=\mathcal A_3$.
\end{proposition}
\begin{proof}
Lemma~\ref{D:EQ:min-origin} supplies one $C$ on all three sides.
Put $u=\overline{C\pi^t}\in k^\times$ and
$c_i=\overline{g_iC-C}\in k^\times$.
The three $c_i$ are distinct, with $c_1+c_2=c_3$.
Indeed $g_iY-Y=d_i$ and the possible adjustment changes the first two
by $a\in\pp_F$, leaving their residues unchanged.
The first ramification coefficient on $\pi$ is $c_i/u$:
expand the leading term $u\pi^{-t}$ of $C$ and use oddness of $t$.
Thus the upper critical norm maps are
\begin{equation}\label{D:EQ:one-pi}
 p_i(z)=z^2+(c_i/u)z.
\end{equation}
Write $P=c_1c_2c_3$ and $k_i=c_jc_k$.
For the lower edge, the difference of the two conjugates of $N_iC$
has residue $k_i$, while its leading term is $u^2\Pi_i^{-t}$.
Consequently its lower critical norm polynomial has coefficient
$\eta_i=k_i/u^2$, namely $z^2+\eta_i z$.

The coefficient of the actual lower norm character in
\eqref{D:EQ:min-lower} has residue $c_i^2$.
Triviality on its critical norms gives
\[
 \tr\bigl(c_i^2(z^2+\eta_i z)\bigr)=0\qquad(z\in k).
\]
Nondegeneracy of the residue trace pairing forces
$c_i+c_i^2\eta_i=0$, hence
\begin{equation}\label{D:EQ:onecalibration}
 \eta_i=c_i^{-1},\qquad u^2=P.
\end{equation}
This determines the unit coefficients; it is not a valuation-only claim.
The polar pairing of the actual upper residual function is therefore
\begin{equation}\label{D:EQ:onepolar}
 H_i(x+z)/(H_i(x)H_i(z))=(-1)^{\tr(k_i xz)}.
\end{equation}
Here tracing $\Pi_i^t$ gives the coefficient $\eta_i$ and the leading
coefficient of $N_iC$ contributes $u^2$.

The image $V_i$ of \eqref{D:EQ:one-pi} is
$\{x:\tr(u^2x/c_i^2)=0\}$, an index-two hyperplane.
For $j\ne i$, put $w_i=k_i/u^2=p_i(c_j/u)$.
The common actual function in \eqref{D:EQ:commonH-general} gives
$H_i(w_i)=H_j(0)=1$ and $w_i\in V_i\setminus\{0\}$.
By \eqref{D:EQ:onecalibration}, its polar character in \eqref{D:EQ:onepolar}
is $x\mapsto(-1)^{\tr(u^2x/c_i^2)}$.
Lemma~\ref{U:missing-coset}, applied with ambient group $(k,+)$
and image $V_i$, gives $\sum_kH_i=\frac12\sum_kQ$ for every $i$.
The normalization $|k|^{-1/2}$ is common, so
\eqref{D:EQ:minfull} proves the result.
\end{proof}

\subsection{Intermediate common twists}
Suppose now $n\ge T_2+r_1$, and write $h=n-T_2\ge r_1$.
The actual conductors and stationary depths are
\begin{equation}\label{D:EQ:twistdepths}
 M_i=2n-T_i,\qquad b_i=M_i/2=n-r_i;
 \quad b_1=T_2+h-r_1,\quad b_2=r_2+h.
\end{equation}
For this subsection assume
\begin{equation}\label{D:EQ:transitionrange}
 r_1\le h\le r_1+2r_2-2.
\end{equation}
Put
\[
 q_F=r_2+\floor{(h+r_1)/2}.
\]
The trace and norm of each variable $z\in\pp_{L_i}^{b_i}$, $i=1,2$,
belong to $\pp_F^{q_F}$. Moreover
$q_F\ge\ceil{n/2}$ and $b_i\ge s_i$.
Choose the actual coefficient $A_*$ of $\lambda$ on this whole
base ideal. It has valuation $-h$ and ambiguity $\pp_F^{T_2-q_F}$.
A norm from $L_3$ can represent this class, because its requested
relative precision is
\begin{equation}\label{D:EQ:transitionprecision}
 h+T_2-q_F=r_2+\ceil{(h-r_1)/2}\le2r_2-1=t_2.
\end{equation}
Thus choose $X\in L_3$ with $v_{L_3}X=-h$ and $A=n_3X$ in that
class, and use \eqref{D:EQ:adjust}.

\begin{theorem}[Exact intermediate comparison]\label{D:EQ:intermediate}
Every family in \eqref{D:EQ:transitionrange} satisfies
$\mathcal A_1(\theta_1)=\mathcal A_2(\theta_2)$.
\end{theorem}
\begin{proof}
The elements $(Q_i/\beta_i)z$ in \eqref{D:EQ:pullback-adjust} have depths
at least $r_1$ and $r_2$, respectively, because
\[
 -h-2\delta+b_1=r_1,\qquad -h+b_2=r_2.
\]
The error terms there have depths at least
\[
 (r_2+\delta-h)-2\delta+b_1=T_2,\qquad
 (r_2+\delta-h)+b_2=T_2+\delta.
\]
Thus the actual pulled-back coefficient is $A+Q_i$ on its full
stationary ideal. Combining it with \eqref{D:EQ:corechart} proves that
$N_iC=a_i+A+Q_i$ is the stationary coefficient for $\theta_i$.
Here $v_KX=-2h<-t_1=v_KY$, so $C\ne0$ and
$v_{L_i}N_iC=-2h$, exactly the required $J_i-M_i$.

Both upper and lower conductors are even. Use $N_iC$ for the upper
coefficient, but retain the original actual coefficient $\beta_i$
for $\omega_i$. It is unnecessary, and generally unjustified, to
claim that $n_iS_iC$ is still stationary for $\omega_i$.
The full Lamprecht formula instead gives
\[
 \mathcal A_i(\theta_i)
   =D_\theta(\alpha^{-1})\Theta(C)^{-1}
                    \Psi_F(T_i^{\rm tr}N_iC+\beta_i).
\]
By \eqref{D:EQ:E2} and \eqref{D:EQ:commonC}, its last exponent is
\[
 E_2(C)+(d_i+a)^2+d_i^2=E_2(C)+a^2,
\]
the same for $i=1,2$. This proves exact equality without a residual
sign argument. In fact $\Psi_F(a^2)=1$, but this extra simplification
is not needed. Relabeling the third field proves the other comparison.
\end{proof}

\subsection{Higher twists and their common correction}
The remaining range is
\begin{equation}\label{D:EQ:highrange}
 h\ge h_0:=r_1+2r_2-1,\qquad n=T_2+h.
\end{equation}
Let $A\in F^\times$ be an actual stationary coefficient of $\lambda$
on $\pp_F^{\ceil{n/2}}$, in the normalized $\Psi_F$ convention.
Thus $v_FA=-h$, and put $c_F=(\alpha A)^{-1}$.
For each $i=1,2$, choose only an approximate norm
\begin{equation}\label{D:EQ:highapprox}
 n_i x_i=(A/\beta_i)\varepsilon_i,\qquad
 \varepsilon_i\in U_F^{t_i},\qquad
 v_{L_i}x_i=-h-(T_2-T_i).
\end{equation}
Such $x_i$ exist by subcritical norm approximation at its stated
precision $t_i$.

\begin{lemma}[The high stationary coefficient and core value]\label{D:EQ:highcoef}
A stationary coefficient of $\lambda n_i$ is $A+\beta_ix_i$.
If $c_i=(\alpha(A+\beta_ix_i))^{-1}$, then
\begin{equation}\label{D:EQ:highcorrection}
 \chi_i(c_i)=\chi_i(c_F)\Psi_F(\mathcal N/A).
\end{equation}
The correction is identical for both fields.
\end{lemma}
\begin{proof}
For $z\in\pp_{L_i}^{b_i}$, $b_i=n-r_i$, both its trace and norm lie
in the base stationary ideal. Also $x_iz$ has valuation at least
$r_i$, so \eqref{D:EQ:normphase} applies. The error from
\eqref{D:EQ:highapprox}, after multiplication by $n_i z$, has normalized
valuation at least
\[
 -h+t_i+b_i=T_2+r_i-1\ge T_2.
\]
Exact quadratic norm expansion proves the first assertion on the
whole stationary ideal. The leading term is $A$, since
\[
 v_{L_i}(\beta_ix_i/A)=h+T_2-T_i>0.
\]
Set $v_i=\beta_ix_i/A$. The lower bound \eqref{D:EQ:highrange} places
$v_i$ in the core domain $\pp_{L_i}^{s_i}$, hence
\[
 \chi_i(c_i/c_F)=\chi_i((1+v_i)^{-1})
                        =\Psi_{L_i}(a_i v_i).
\]
The last equality keeps the actual core character; its displayed
value has order at most two. To convert it, the variable
$a_ix_i/A$ has valuation
\[
 h-t_1-(T_2-T_i)\ge r_i.
\]
For $i=1$ this inequality is precisely $h\ge h_0$; for $i=2$ it is
weaker. Apply \eqref{D:EQ:normphase}, using $n_i a_i=\mathcal N$, to get
\[
 \Psi_{L_i}(\beta_i a_ix_i/A)
                      =\Psi_F(\mathcal N\varepsilon_i/A).
\]
The error on deleting $\varepsilon_i$ has valuation at least
$h-t_1+t_i$, which is at least $T_2$ on both sides by
\eqref{D:EQ:highrange}. This proves \eqref{D:EQ:highcorrection}.
\end{proof}

\begin{theorem}[Exact high comparison]\label{D:EQ:highcomparison}
For every family in \eqref{D:EQ:highrange},
\begin{equation}\label{D:EQ:highanswer}
 \mathcal A_i(\theta_i)
   =D_\chi(c_F)\Psi_F(\mathcal N/A)\Delta_F(\lambda,\psi_F)^2
                    \qquad(i=1,2).
\end{equation}
In particular the full induction expressions are equal.
\end{theorem}
\begin{proof}
The core conductor lies within the upstairs stable range:
\[
 b_1=T_2+h-r_1\ge 2T_2-1=m_1,\qquad
 b_2=r_2+h\ge T_1+T_2-1=m_2.
\]
Ueda's exact stable twist gives
$\Delta_{L_i}(\theta_i)=\chi_i(c_i)\Delta_{L_i}(\lambda n_i)$.
The First Main Lemma on $L_i/F$ gives
\[
 \Delta_{L_i}(\lambda n_i)\Delta_F(\omega_i)
                =\Delta_F(\lambda)\Delta_F(\lambda\omega_i).
\]
Since $n\ge2T_2\ge2T_i$, the downstairs stable twist is valid and
the right side is $\omega_i(c_F)\Delta_F(\lambda)^2$.
Combine this with Lemma~\ref{D:EQ:highcoef} and the common determinant
restriction of the core family. No arbitrary exact norm or comparison
of the desired local constants has been used to obtain the correction.
\end{proof}

\subsection{Completion in characteristic two}
\begin{theorem}\label{D:EQ:total-complete}
Theorem~\ref{D:EQ:main} holds when $K/F$ is totally ramified.
\end{theorem}
\begin{proof}
Proposition~\ref{D:EQ:realize} realizes the actual family by the stated core
and a common base twist. If $n\le T_2+r_1-1$, use
Propositions~\ref{D:EQ:min-two} and \ref{D:EQ:min-one}.
Otherwise $h=n-T_2\ge r_1$. The interval
$r_1\le h\le r_1+2r_2-2$ is Theorem~\ref{D:EQ:intermediate}, and the
remaining integers are Theorem~\ref{D:EQ:highcomparison}.
Their endpoints are consecutive, including $r_1=r_2=1$.
In the two-break case repeat the comparison of the distinguished first
field with the third field. This proves independence of all three.
Changing an inducing extension by an upper norm character conjugates it
and preserves its local constant by an automorphism change of variables
in the defining local integral. The passage to an arbitrary additive
character is made once for all branches at the end.
\end{proof}


\begin{proof}[Proof of Theorem~\ref{D:EQ:main}]
For the canonical additive character this is Theorem~\ref{D:EQ:total-complete}.
Every nontrivial additive character of $F=k((\varpi))$ is a scaling
$x\mapsto\psi_F(cx)$ of \eqref{D:EQ:canonicalpsi}, with $c\in F^\times$.
Ueda's additive-scaling formula multiplies the complete induction
expression by $\theta_i(c)\omega_i(c)=D_\theta(c)$, independent of $i$
by Lemma~\ref{U:conjugacy}. This proves the stated arbitrary-character result.
\end{proof}


\section{Totally ramified mixed characteristic: the maximal break}\label{sec:dyadic-maximal}
\subsection{Statement and ramification data}
Let $F/\mathbf Q_2$ be finite, with residue field $k$, normalized
valuation $v_F$, and $e=v_F(2)$. Suppose $K/F$ is totally ramified,
$\Gal(K/F)=C_2^2$. Label its three quadratic fields so that
\begin{equation}\label{D:MX:breaks}
 t_1=2a-1,\qquad t_2=t_3=2e,\qquad 1\le a\le e.
\end{equation}
Put $T_1=2a$, $T_2=T_3=T=2e+1$. Write $g_i$ for the nontrivial
upper automorphism fixing $L_i$, and
\[
 N_i=\Nm_{K/L_i},\quad S_i=\Tr_{K/L_i},\quad
 n_i=\Nm_{L_i/F},\quad \operatorname{tr}_i=\Tr_{L_i/F}.
\]
Let $\omega_i$ be the actual quadratic norm character for $L_i/F$.
Fix a nontrivial additive $\psi_F$ and compose with trace on all fields.
\begin{theorem}\label{D:MX:main}
For every primitive compatible family $\theta_iN_i=\Theta$,
where $\Theta$ is not a pullback from $F$, the full expressions
\begin{equation}\label{D:MX:target}
 \mathcal A_i(\theta_i)=
 \Delta_{L_i}(\theta_i,\psi_{L_i})\Delta_F(\omega_i,\psi_F)
\end{equation}
are equal. Continuous quasi-characters are allowed. The conclusion
holds for every conductor under \eqref{D:MX:breaks}.
\end{theorem}
We prove the comparison of $L_1,L_2$ using $L_3$ as auxiliary field.
Relabeling the two maximal fields then proves the other comparison.

The one-dimensional inputs are the results of~\cite{Ueda}: the full Lamprecht formula,
stationary additive duality, norm filtration and norm-character conductor,
norm-pullback conductor, trace/different formulas, and the trivial-character
local constant. Hilbert~90 and character extension are proved in
Appendix~\ref{app:descent}; the crossed norm characters and ramification
relations are proved in Appendix~\ref{app:diamond-foundations}.
We give the additional quadratic-symbol calculation explicitly.
No two-inducing-field comparison or global functional equation is used.

For a ramified quadratic edge $E/F$ of different exponent $D$,
\begin{equation}\label{D:MX:trace}
 \Tr_{E/F}(\pp_E^j)=\pp_F^{\floor{(j+D)/2}},\quad
 v_F(\Nm x)=v_E x,\quad n_E(\psi_E)=2n_F(\psi_F)+D.
\end{equation}
The upper different exponents in our diagram are
\begin{equation}\label{D:MX:upper}
 D'_1=2T-T_1=4e-2a+2,\qquad D'_2=D'_3=2a.
\end{equation}
Indeed the upper breaks are $2t_2-t_1$ and $t_1$, by the usual
break calculation of Lemma~\ref{D:break-ledger}.
All valuations below name their field;
restriction of $v_K$ to $F$ is $4v_F$.

\subsection{A local symbol identity and the aligned Kummer origin}
For $A,B\in F^\times$ let $(A,B)$ be the quadratic norm symbol: it is
$1$ exactly when $B$ is a norm from $F(\sqrt A)$, with the split case
interpreted as $1$. It is a character in $B$. Symmetry follows from the
symmetric conic equation $z^2=A x^2+B y^2$, and hence it is bilinear in
both square classes. For $A\ne1$, one also has $(A,1-A)=1$, since
$1-A$ is the nonzero norm of $1+\sqrt A$. If $A+B\ne0$, the same conic, under
$u=Ax+By$, $v=x-y$, gives
\begin{equation}\label{D:MX:symbol-transform}
 (A,B)=(A+B,-AB),\qquad
 (A+B)z^2=u^2+ABv^2.
\end{equation}
The linear change in $(x,y)$ is invertible. Thus this equality is a
local quadratic-norm identity, not a local-constant identity.

\begin{lemma}[A sufficient unit depth]\label{D:MX:unit-symbol}
If $v_Fu=m\ge1$, $v_Fv=n\ge1$, and $m+n>2e$, then
$(1+u,1+v)=1$.
\end{lemma}
\begin{proof}
If $1+u$ is a square, the assertion is immediate; this includes
$m>2e$ by Hensel's lemma. Otherwise put
$r=\floor{m/2}\le e$. The integral element
$(\sqrt{1+u}-1)/\varpi^r$ generates an order of discriminant valuation
$2e-2r$. The field discriminant divides the order discriminant, so the
quadratic norm character has conductor at most $2e-2r$; the unramified
case also satisfies this bound. If $m$ is even, $n\ge2e-2r+1$;
if $m$ is odd, $n\ge2e-2r$. In both cases the character kills $1+v$.
\end{proof}

Put $b=2e-2a+1$, an odd positive integer, and $q=\varpi^b$.
There are Kummer generators $s,R_0$ with
\[
 s^2=1+q u_0,\qquad R_0^2=q w_0,\qquad u_0,w_0\in\OO_F^\times.
\]
Here is a verification at the stated breaks. For a uniformizer $\Pi_1$
of $L_1$, its quadratic minimal polynomial and different imply
$v_F\Tr\Pi_1=a$: the term $2\Pi_1$ has $L_1$-valuation $2e+1>2a$.
Thus $s=1-2\Pi_1/\Tr\Pi_1$ has the required square defect of valuation
$b$. For a maximal quadratic field, a uniformizer $\Pi$ has
$v_F\Tr\Pi\ge e+1$, and $\Pi-\Tr\Pi/2$ is a trace-zero uniformizer.
Multiplying it by $\varpi^{e-a}$ gives $R_0$.

\begin{lemma}[Alignment beyond the leading residue]\label{D:MX:align}
The generators can be replaced by $S=s(1+\delta)$ and $R=R_0v$,
with $v$ a base unit, so that, writing $\epsilon=S^2-1$,
\begin{equation}\label{D:MX:aligned}
 v_F\epsilon=v_FR^2=b,\qquad
 \epsilon/R^2\in U_F^a.
\end{equation}
\end{lemma}
\begin{proof}
Take $\delta=\varpi^{e-a+1}w$, $w\in\OO_F$. Modulo $\pp_F^a$,
\[
 \frac{S^2-1-R^2}{q}
 \equiv u_0-w_0v^2+\varpi(1+q u_0)w^2.
\]
The omitted term $2\delta/q$ has valuation at least $a$.
Because $a\le e$, $\OO_F/\pp_F^a$ has characteristic two and is
$k[\varpi]/(\varpi^a)$ after choosing its coefficient field. Such a
coefficient field is obtained by the unique lifts of the roots of
$X^{|k|}-X$, whose derivative is a unit.
In this truncated ring solve
\[
 w_0v^2+\varpi(1+q u_0)w^2=u_0.
\]
Successively in the degree of $\varpi$, an even degree determines the next
coefficient of $v^2$, with unit leading coefficient $\bar w_0$; an odd
degree determines the next coefficient of $w^2$, with leading coefficient
$1$. Contributions of earlier coefficients are already known. Frobenius
on $k$ supplies the unique required square roots. The constant coefficient
of $v$ is nonzero. Lift $v,w$ to $F$. The displayed congruence proves
\eqref{D:MX:aligned}, and the positive-depth correction to $u_0$ preserves its
unit leading coefficient.
\end{proof}

Set $L_1=F(S)$, $L_2=F(R)$, $L_3=F(SR)$, and define
\begin{equation}\label{D:MX:origin}
 Y=(S+R-1)/2,\qquad a_i=N_iY,\qquad h_i=S_iY,\qquad \beta_i=n_i h_i.
\end{equation}
For
$B_0=(\epsilon-R^2)/4$ and $C_0=(\epsilon+R^2)/4$ one has
\begin{align}\label{D:MX:origin-table}
 a_1&=(1-S)/2+B_0,&h_1&=S-1,&\beta_1&=-\epsilon,\notag\\
 a_2&=-R/2-B_0,&h_2&=R-1,&\beta_2&=1-R^2,\notag\\
 a_3&=-SR/2-C_0,&h_3&=-1,&\beta_3&=1.
\end{align}
The alignment gives $v_FB_0,v_FC_0\ge1-a$; for the plus sign use also
$v_F(2R^2)\ge b+a$. Consequently
\begin{equation}\label{D:MX:valuations}
 v_KY=-t_1,\quad v_{L_i}a_i=-t_1,\quad
 v_{L_1}h_1=b,\quad v_{L_2}h_2=v_{L_3}h_3=0,
 \quad\Tr_{K/F}Y=-2.
\end{equation}
For instance $a_2$ has the unique leading term $-R/2$ of valuation
$1-2a$, while $B_0$ has $L_2$-valuation at least $2-2a$.
The value of $Y$ follows by norm valuation; the other entries follow
from \eqref{D:MX:origin-table} and $(S-1)(S+1)=\epsilon$.

\subsection{Simultaneous coefficients for the norm characters}
Choose the normalization from the \emph{third} maximal character:
\begin{equation}\label{D:MX:normalization}
 \omega_3(1-z)=\Psi_F(z)\quad(z\in\pp_F^{e+1}),\qquad
 \Psi_L(x)=\psi_L(\alpha x),\quad v_F\alpha=-n_F(\psi_F)-T.
\end{equation}
The largest trivial ideals of $\Psi_F,\Psi_{L_1},\Psi_{L_2},\Psi_{L_3}$
have respective exponents
\begin{equation}\label{D:MX:J}
 T,\quad J_1=2T-2a,\quad J_2=J_3=T.
\end{equation}
The existence and precision of \eqref{D:MX:normalization} are the full linear
chart from additive duality; $2(e+1)\ge T$.

\begin{lemma}[Full lower character formulas]\label{D:MX:lower}
With $q_1=a$, $q_2=q_3=e+1$, one has
\begin{equation}\label{D:MX:lowerformula}
 \omega_i(1-z)=\Psi_F(\beta_i z)\quad(z\in\pp_F^{q_i}).
\end{equation}
Each $\beta_i$ is a genuine norm from $L_i$, and on the whole stated source
ideal the resulting phase conversion is
\begin{equation}\label{D:MX:normphase}
 \Psi_{L_i}(\beta_i x)=\Psi_F(\beta_i n_i x)
                    \quad(x\in\pp_{L_i}^{q_i}).
\end{equation}
\end{lemma}
\begin{proof}
For $z\in\pp_F^a$ put $C=1+\epsilon z$. The Steinberg identity and
\eqref{D:MX:symbol-transform} give
\[
 (1+\epsilon,1-z)
 =(1+\epsilon,-\epsilon(1-z))
 =(C,\epsilon(1-z)(1+\epsilon)).
\]
Here $v_F(C-1)\ge b+a$, and $(b+a)+a=2e+1$.
Lemma~\ref{D:MX:unit-symbol} removes $1-z$ from the second argument and replaces
$\epsilon$ by $R^2$, since $\epsilon/R^2\in U_F^a$. By symmetry the result
is $(R^2(1+\epsilon),C)=\omega_3(C)$.
Also $b+a=2e-a+1\ge e+1$, so \eqref{D:MX:normalization} gives
$\omega_1(1-z)=\Psi_F(-\epsilon z)$ on the \emph{entire} ideal.
On $U_F^{e+1}$ the product $\omega_2=\omega_1\omega_3$ therefore has
coefficient $1-\epsilon$. Replacing it by $1-R^2$ costs valuation at least
$b+a+e+1=3e-a+2\ge T$. This proves \eqref{D:MX:lowerformula}; the third formula
is the chosen normalization.
The norm assertions are \eqref{D:MX:origin-table}. Finally for
$x\in\pp_{L_i}^{q_i}$ both its trace and norm lie in $\pp_F^{q_i}$ by
\eqref{D:MX:trace}. Evaluate \eqref{D:MX:lowerformula} at
$n_i(1-x)=1-\operatorname{tr}_i x+n_i x$ and use triviality on norms.
This proves \eqref{D:MX:normphase}, including its positive norm sign.
\end{proof}

\subsection{Construction and realization of the minimal characters}
Put
\begin{equation}\label{D:MX:core-depths}
 m_1=4e+1,\quad m_2=m_3=2e+2a,\qquad
 s_1=2e+1,\quad s_2=s_3=e+a.
\end{equation}
Thus $m_i=J_i+t_1$ and $s_i=\ceil{m_i/2}$.
The formulas
\begin{equation}\label{D:MX:R}
 R_i(1-z)=\Psi_{L_i}(a_i z),\qquad z\in\pp_{L_i}^{s_i},
\end{equation}
define characters of their indicated unit groups with exact conductor
$m_i$: $2s_i-t_1\ge J_i$, and the last-layer coefficient is nonzero.

\begin{lemma}[Compatibility before extension]\label{D:MX:Rcompat}
For $z\in\pp_K^{2e+1}$, all $N_i(1-z)$ lie in the domains of $R_i$,
and $R_i(N_i(1-z))$ is independent of $i$.
\end{lemma}
\begin{proof}
The norm endpoints have valuation at least $2e+1$ on the target field.
The upper traces have depths at least $3e-a+1$ and $e+a$, hence are in
the required ideals. Use the exact identity
\begin{equation}\label{D:MX:E2}
 E_2(W)=\operatorname{tr}_i N_iW+n_iS_iW
\end{equation}
for the six pair-products of the four conjugates of $W$.
Apply it to $Y(1-z)$ and $Y$.
The quantities $w_i=S_i(Yz)/h_i$ have valuation at least $e+1$ on all
three fields: the first upper trace bound is $3e-2a+2$, from which one
subtracts $b$, and the other two are $e+1$.
Thus $n_iS_i(Y(1-z))/\beta_i=n_i(1-w_i)$ is an actual norm in the domain
of \eqref{D:MX:lowerformula}. Its additive phase implies
\[
 \Psi_F(n_iS_i(Y(1-z))-\beta_i)=1.
\]
The remaining difference in \eqref{D:MX:E2} identifies
$R_i(N_i(1-z))=\Psi_F(-E_2(Y(1-z))+E_2(Y))$, independent of $i$.
\end{proof}

\begin{theorem}[Actual global core models]\label{D:MX:cores}
There are primitive compatible characters $\chi_i$ on $L_i^\times$ of
conductors $m_i$, whose full stationary formulas are \eqref{D:MX:R}.
\end{theorem}
\begin{proof}
First extend $R_2$ across the abelian group $L_2^\times/U_{L_2}^{m_2}$;
extend on the finite unit subgroup and choose a uniformizer value.
Call the character $\chi_2$. Its lower conjugate quotient is determined
by $R_2$ on every element of $L_2^\times$, because every ratio
$\sigma(v)/v$ has depth at least $t_2=2e\ge s_2$.
For $v=N_2 Z$, this ratio is $N_2(g_1Z/Z)$.
The element $g_1Z/Z$ has depth at least $4e-2a+1\ge2e+1$ and its
$N_1$-norm is $1$. Lemma~\ref{D:MX:Rcompat} therefore proves that the conjugate
quotient of $\chi_2$ kills $N_2K^\times$.
It is nontrivial: on $U_{L_2}^{t_1}$ the first nonzero lower ramification
term maps onto the last row $U_{L_2}^{t_1+t_2}/U_{L_2}^{t_1+t_2+1}$,
since $t_1$ is odd. Formula \eqref{D:MX:R} is nontrivial on this row.
Thus the quotient is the actual upper norm character $\nu_2$.

Consequently $\chi_2N_2$ is invariant under all of $G$; Hilbert 90 makes
it trivial on $\ker N_1$ and $\ker N_3$. Extend its descents to
$\widetilde\chi_1,\chi_3$. On side $3$, norm filtration gives
$N_3(U_K^{2e+1})=U_{L_3}^{e+a}$, since
$2(e+a)-t_1=2e+1$. Lemma~\ref{D:MX:Rcompat} gives the full $R_3$ formula.
On side $1$, every preimage of $U_{L_1}^{s_1}$ under $N_1$ lies in
$U_K^{s_1}$: the norm has nonzero graded map at every smaller positive
depth, all below the upper break $4e-2a+1$, and its residue map is
Frobenius. Norm valuation excludes nonunits. Hence
$\widetilde\chi_1$ agrees with $R_1$ on
$U_{L_1}^{s_1}\cap N_1K^\times$. Their quotient extends across the order-two
norm quotient; divide by this norm character to obtain $\chi_1$ with the
full $R_1$ formula and the same pullback. The extensions are trivial at
depth $m_i$ and nontrivial at $m_i-1$, by their prescribed charts and norm
filtration, so their conductors are exactly $m_i$.
If the common pullback descended from $F$, $\chi_2$ would be a base
pullback times an upper norm character, both invariant under its lower
automorphism. This contradicts its nontrivial commutator.
\end{proof}

\begin{lemma}[Conjugacy, determinant, and realization]\label{D:MX:realize}
For every primitive compatible family,
\[
 \theta_i^\sigma/\theta_i=\nu_i=\omega_jn_i\quad(j\ne i),\qquad
 D_\theta=\theta_i|_{F^\times}\omega_i
\]
is independent of $i$. For comparison of $1$ and $2$ we may, without
changing either local constant or the common pullback, write
\begin{equation}\label{D:MX:twisted}
 \theta_i=\chi_i(\lambda n_i)\qquad(i=1,2)
\end{equation}
with cores as in Theorem~\ref{D:MX:cores}. If $n=a_F(\lambda)$, the conductors
are $m_i$ for $n\le2e+a$; for $n\ge2e+a+1$ they are
\begin{equation}\label{D:MX:higher-conductors}
 M_1=2n-2a,\qquad M_2=2n-T.
\end{equation}
\end{lemma}
\begin{proof}
Lemma~\ref{U:conjugacy} gives the stated conjugate quotient and
corrected common restriction, for the given family and the core family.
The quotient $\theta_2/\chi_2$ is invariant and is therefore $\lambda n_2$.
The discrepancy on side $1$ is either $1$ or $\nu_1$.
Conjugate $\theta_1$ in the latter case. Its local constant is unchanged
by a field-automorphism change of variables in the defining integral;
its norm pullback is unchanged because $\nu_1N_1=1$.
This proves \eqref{D:MX:twisted}; it is not an application of SML.
Finally the conductor formula for norm pullback gives the alternatives:
below the displayed threshold both twists have smaller conductor than
their cores; above it they have larger conductor. On side $1$ the core
is odd and the higher pullback even; on side $2$ the core is even and the
higher pullback odd. No equal-conductor cancellation is omitted.
\end{proof}

\subsection{Normalization and the exact norm choice}
\begin{lemma}[Preservation and an exact third-field norm]\label{D:MX:freedom}
Replacing $\alpha$ by $\alpha u$, with $u\in U_F^{2e}$, leaves
\eqref{D:MX:lowerformula} and the full actual core formulas \eqref{D:MX:R} unchanged.
For any nonzero coefficient $A_*$ of a base character relative to $\Psi_F$,
such a replacement can be made so that $A=A_*/u$ is an exact norm from
$L_3/F$. It remains the coefficient of the same base character relative
to the new $\Psi_F$.
\end{lemma}
\begin{proof}
For the lower formula the change has valuation at least
$2e+v_F\beta_i+q_i\ge T$. For the upper core formula it has
$L_i$-valuation at least $4e-t_1+s_i\ge J_i$.
Thus these formulas for the already fixed characters remain valid on
whole ideals, not merely on their final residue rows.
The actual $\omega_3$ has conductor $2e+1$ and is nontrivial on
$U_F^{2e}$. Apply Lemma~\ref{U:normalization} to this subgroup.
It gives the claimed norm without changing the raw coefficient
$\alpha A_*$. Norm valuation fixes $v_{L_3}X=v_FA$.
\end{proof}
We henceforth use the new normalization when needed.


For comparison of $1,2$ take $X\in L_3$, and put
\[
 A=n_3X,\qquad p_X=\operatorname{tr}_3X,\quad C=Y-X,\quad
 Q_i=S_i(Yg_iX)\ (i=1,2).
\]
Then, exactly,
\begin{align}\label{D:MX:adjust-identities}
 N_iC&=a_i+A-Q_i,&S_iC&=h_i-p_X,\notag\\
 n_iQ_i&=\beta_i A+\mathcal E,&
 \mathcal E&=(X'-X)(a_3X'-a_3'X)\in F,\\
 n_iS_iC&=\beta_i+\Delta,&\Delta&=p_X^2+2p_X.\label{D:MX:Delta}
\end{align}
Primes in the middle line denote the automorphism of $L_3/F$.
These identities follow by multiplying the four conjugates, with $X$
fixed by $g_3$; the last uses $\operatorname{tr}_ih_i=\Tr_KY=-2$.

\begin{lemma}[Weighted common error]\label{D:MX:errors}
If $v_{L_3}X=-h$, for any integer $h$, then
\begin{equation}\label{D:MX:error-bounds}
 \begin{split}
 v_Fp_X&\ge e-\floor{h/2},\qquad
 v_F\mathcal E\ge T-a-h,\\
 v_{L_1}Q_1&\ge b-h,\qquad v_{L_2}Q_2\ge-h.
 \end{split}
\end{equation}
\end{lemma}
\begin{proof}
Let $r_3=SR/\varpi^{e-a}$, a trace-zero uniformizer of $L_3$.
Write $X=P+Qr_3$; parity of the two valuations gives
$v_FP\ge-\floor{h/2}$ and $v_FQ\ge-\ceil{h/2}$.
The trace is $2P$, proving its bound. Write $a_3=A_3+B_3r_3$, with
$v_FA_3\ge1-a$ and $v_FB_3=-a$ by \eqref{D:MX:origin-table}.
Now $X'-X=-2Qr_3$ and
$a_3X'-a_3'X=2(B_3P-A_3Q)r_3$.
Their product has $F$-valuation at least $2e+1-a-h=T-a-h$.
The upper trace formula applied to $Yg_iX$, of $K$-valuation
$-t_1-2h$, gives respectively $b-h$ and $-h$ for $Q_i$.
\end{proof}
The conversion \eqref{D:MX:normphase} and \eqref{D:MX:adjust-identities} imply,
whenever $(Q_i/\beta_i)z\in\pp_{L_i}^{q_i}$,
\begin{equation}\label{D:MX:Qphase}
 \Psi_{L_i}(Q_i z)
   =\Psi_F\bigl((A+\mathcal E/\beta_i)n_i z\bigr).
\end{equation}
This retains the actual norm error with its weight.

\subsection{Minimal conductors: upper and lower stationary origins}
Suppose $n\le2e+a$. Both trace and norm of a variable in
$\pp_{L_i}^{s_i}$ lie in $\pp_F^{e+a}$, for $i=1,2$.
If $n\le e+a$, the twist is trivial there, and use $X=0$, $C=Y$.
Otherwise put $h=n-T$, so
\begin{equation}\label{D:MX:minimal-range}
 a-e\le h\le a-1.
\end{equation}
Additive duality on $U_F^{e+a}$ supplies a nonzero base coefficient
$A_*$ of valuation $-h$, since $2(e+a)\ge n$.
Use Lemma~\ref{D:MX:freedom} to make its new value $A=n_3X$ an exact norm.

\begin{lemma}[Full stationary formulas for the given characters]\label{D:MX:minimal-origin}
The resulting $C$ has $v_KC=-t_1$. For $i=1,2$,
\begin{align}\label{D:MX:minimal-stationary}
 \theta_i(1-z)&=\Psi_{L_i}(N_iC\,z)
                                &&(z\in\pp_{L_i}^{s_i}),\notag\\
 \omega_i(1-z)&=\Psi_F(n_iS_iC\,z)
                                &&(z\in\pp_F^{q_i}).
\end{align}
Both $S_iC$ are nonzero and have the same valuations as $h_i$.
\end{lemma}
\begin{proof}
For $X\ne0$, $-2h>-t_1$, so $Y$ dominates in $C$.
The two arguments of \eqref{D:MX:Qphase} have depths at least
$2a-h\ge a+1$ and $e+a-h\ge e+1$.
The error terms $(\mathcal E/\beta_i)n_i z$ have depths at least
$T+a-h$ and $T+e-h$, both at least $T$.
Exact quadratic norm expansion evaluates the actual $\lambda n_i$ on its
full stationary ideal as $\Psi_{L_i}((A-Q_i)z)$, giving the first formula.
For the second, \eqref{D:MX:error-bounds} and $h\le a-1$ give
\[
 v_F\Delta\ge\min(2e-2\floor{h/2},\ 2e-\floor{h/2}).
\]
Adding $a$ to either bound gives at least $T$; adding $e+1$ does also.
Thus \eqref{D:MX:Delta} preserves both full lower formulas.
Moreover $v_{L_1}p_X\ge2e-2\floor{h/2}>b$, while $p_X$ has positive
valuation on $L_2$. The traces retain their nonzero leading terms.
\end{proof}

Use the full stationary formula and terminal invariance of
Lemma~\ref{U:stationary-factor}, with $c_0=-\alpha^{-1}$ and
residual factors denoted $g_\theta$.

\begin{theorem}[Exact minimal comparison]\label{D:MX:minimal}
Every family with $n\le2e+a$ satisfies $\mathcal A_1=\mathcal A_2$.
\end{theorem}
\begin{proof}
Use the coefficients of Lemma~\ref{D:MX:minimal-origin} in \eqref{U:normalized-factor}.
Their norm-character values are $1$, and compatibility cancels the two
upper multiplicative values. Formula \eqref{D:MX:E2} gives
\begin{equation}\label{D:MX:minimal-full}
 \mathcal A_i=D_\theta(c_0)\Theta(C)^{-1}
                    \Psi_F(-E_2(C))g_{\theta_i}g_{\omega_i}.
\end{equation}
Only $\theta_1$ and $\omega_2$ have odd conductors, with critical depths
$2e$ and $e$ respectively.
Choose a top uniformizer $\pi$ and put $C_z=C(1+\pi^{2e}z)$ for integral
$z\in K$. On side $1$ the relative upper norm has leading term its norm
endpoint at depth $2e$; its trace has depth at least $3e-a+1\ge2e+1$.
Its residue map is a nonzero scalar times Frobenius, hence bijective.
On side $2$ the relative upper norm is in the even stationary ideal:
its trace has depth $e+a$, and its norm endpoint depth $2e\ge e+a$.

The trace change $S_i(C\pi^{2e}z)$ has first-side depth at least
$3e-2a+1$; after division by $S_1C$ its depth is at least $e\ge a$.
Its lower norm quotient is therefore on the even stationary ideal of
$\omega_1$. On side $2$, the upper trace of an element of $K$-valuation
$b=2e-2a+1$ induces an isomorphism from its residue line to the
$L_2$ row of depth $e$: the consecutive trace-ideal exponents are $e$
and $e+1$. Division by the unit $S_2C$ preserves this isomorphism.
The lower maximal norm then has leading endpoint at depth $e$ and trace
at depth at least $\floor{(3e+1)/2}\ge e+1$.
Thus its residue map is again bijective.

Subtract \eqref{D:MX:E2} for $C_z,C$. The multiplicative character values
cancel by the actual common pullback, and both lower arguments are ratios
of actual norms. The complete residual expressions therefore give
\[
 H_{\theta_1}(p(z))=H_{\omega_2}(q(z))\qquad(z\in k),
\]
where both $p,q:k\to k$ are bijective maps just constructed.
All even expressions are $1$; all arguments of multiplicative characters
are nonzero, since the relative changes have positive depth.
Summing gives $g_{\theta_1}=g_{\omega_2}$. Equation \eqref{D:MX:minimal-full}
proves the exact sign.
\end{proof}

\subsection{Higher conductors: two stationary origins}
Now $n\ge2e+a+1$, and put $h=n-T\ge a$.
The upper stationary depths are
\begin{equation}\label{D:MX:high-depths}
 b_1=T+h-a,\qquad b_2=e+1+h,
\end{equation}
with $M_1=2b_1$ even and $M_2=2b_2-1$ odd.
Choose an actual coefficient $A_*$ of $\lambda$ on its ordinary base
stationary ideal. Both trace and norm of $z\in\pp_{L_i}^{b_i}$ lie in
that base ideal, by \eqref{D:MX:trace}. Its valuation is $-h$.
Lemma~\ref{D:MX:freedom} supplies, after its harmless normalization change,
$A=n_3X$ \emph{exactly}, with $v_{L_3}X=-h$. There is no upper bound on $h$.

\begin{lemma}\label{D:MX:higher-coeff}
For the actual given characters, $N_iC$ is a full stationary coefficient
on $\pp_{L_i}^{b_i}$, $i=1,2$. Its valuation is $-2h$.
\end{lemma}
\begin{proof}
The arguments $(Q_i/\beta_i)z$ of \eqref{D:MX:Qphase} have respective depths
at least $a$ and $e+1$. The weighted errors there have depths
\[
 (T-a-h)-b+b_1=T,\qquad
 (T-a-h)+b_2=3e+2-a\ge T.
\]
Consequently exact quadratic norm expansion yields the pulled-back
coefficient $A-Q_i$ on the whole indicated ideals. These ideals are
contained in the core domains, so adding $a_i$ gives $N_iC$.
The element $X$ dominates $Y$ because $-2h<-t_1$; hence
$v_KC=-2h$ and $v_{L_i}N_iC=-2h$, as required by the actual conductors.
\end{proof}

Retain $\beta_i=n_i h_i$, rather than $n_iS_iC$, for the lower characters.
No assertion that the latter remain stationary in this range is needed.
Equations \eqref{D:MX:E2}, \eqref{D:MX:Delta}, and \eqref{U:normalized-factor} give
\begin{equation}\label{D:MX:higher-full}
 \mathcal A_i=D_\theta(c_0)\Theta(C)^{-1}
              \Psi_F(-E_2(C)+\Delta)g_{\theta_i}g_{\omega_i}.
\end{equation}
The prefactor is already the same for $i=1,2$, even if $\Delta$ is shallow.
Both residual factors on side $1$ are $1$; both on side $2$ are odd.

\begin{theorem}[Exact comparison at every higher conductor]\label{D:MX:higher}
For every $h\ge a$, one has $\mathcal A_1=\mathcal A_2$.
\end{theorem}
\begin{proof}
Put
\[
 V=Y\pi^{2e}z,\qquad C_z=C+V,\qquad Y_z=Y+V\quad(z\in\OO_K).
\]
The upper character rows use $C_z/C$, but the lower norm-character rows
use $S_iY_z/h_i$. These are \emph{two} origins, a distinction necessary
when $n_iS_iC$ is not stationary.
The element $V$ has $K$-valuation at least $b$ and $V/C$ has depth
$b+2h$. On side $1$, its upper trace depth is at least
$3e-2a+1+h\ge b_1$, and its norm endpoint depth $b+2h\ge b_1$.
Thus the upper residual expression is $1$ on its even stationary ideal.
On side $2$ the consecutive upper trace-ideal exponents at source depths
$b+2h,b+2h+1$ are $e+h,e+h+1$.
The norm endpoint is strictly deeper, since
$b+2h-(e+h)=e-2a+1+h\ge1$.
It follows that the actual upper odd residue map $p:k\to k$ is bijective.

For the lower rows, $S_1V/h_1$ has depth at least $e\ge a$, so its norm
quotient is in the even stationary ideal of $\omega_1$.
The map $z\mapsto S_2V/h_2$ has a nonzero graded trace at depth $e$,
by the consecutive trace exponents at $b,b+1$. Its lower norm has
leading endpoint at depth $e$ and trace of depth at least $e+1$.
Hence the actual lower odd residue map $q:k\to k$ is bijective as well.
All traces in these lower character arguments stay nonzero.

Here the elementary phase equality is exact, not an error estimate.
Since $S_iC=h_i-p_X$, quadratic norm expansion gives
\[
 [n_iS_iC_z-n_iS_iC]-[n_iS_iY_z-\beta_i]=-p_X\Tr_{K/F}V.
\]
Using \eqref{D:MX:E2}, we obtain
\begin{equation}\label{D:MX:two-origins}
 \operatorname{tr}_i(N_iC_z-N_iC)+(n_iS_iY_z-\beta_i)
  =E_2(C_z)-E_2(C)+p_X\Tr_{K/F}V,
\end{equation}
independent of $i$.
Upper multiplicative values are the common $\Theta(C_z/C)$ and lower
ones are $1$, because their arguments are norms. Thus \eqref{D:MX:two-origins}
and the complete residual definition give
\[
 H_{\theta_2}(p(z))H_{\omega_2}(q(z))=1\qquad(z\in k).
\]
Both residue maps are additive bijections. Therefore
$\sum H_{\theta_2}=\overline{\sum H_{\omega_2}}$, after reindexing.
Each normalized sum has absolute value $1$ by Ueda's nondegenerate
finite quadratic-sum formula. Consequently
$g_{\theta_2}g_{\omega_2}=1$.
Equation \eqref{D:MX:higher-full} proves the assertion, for arbitrarily large
$h$, without a separate stability or norm-approximation limit.
\end{proof}

\subsection{Completion of the maximal-break case}
\begin{proof}[Proof of Theorem~\ref{D:MX:main}]
The explicit aligned origin and Lemma~\ref{D:MX:lower} give full actual lower
coefficients. Theorem~\ref{D:MX:cores} constructs actual cores, and
Lemma~\ref{D:MX:realize} realizes the given comparison, using at most a
conjugation that preserves the individual local constant and pullback.
Every base conductor is either $n\le2e+a$, covered by
Theorem~\ref{D:MX:minimal}, or $n\ge2e+a+1$, covered by
Theorem~\ref{D:MX:higher}. Their endpoints are consecutive.
Repeat the construction with the two maximal fields interchanged to
compare the distinguished $L_1$ with $L_3$. This proves independence
of all three expressions for the original given family.
\end{proof}


\section{Totally ramified mixed characteristic: nonmaximal breaks}\label{sec:dyadic-nonmaximal}
\subsection{Statement and ramification data}
Let $F/\mathbf Q_2$ be finite, with residue field $k$, uniformizer
$\varpi$, and $e=v_F(2)$. Normalize all valuations on their named fields.
Let $K/F$ be totally ramified with group $C_2^2$. Label its lower
quadratic fields $L_i=K^{\langle g_i\rangle}$ so that
\begin{equation}\label{D:NM:breaks}
 t_1=2a-1,\qquad t_2=t_3=2r-1,\qquad 1\le a\le r\le e.
\end{equation}
Put $\delta=r-a$, $T=2r$, and $T_1=2a$, $T_2=T_3=T$.
These $T_i$ are the lower different exponents and the conductors of the
actual quadratic norm characters $\omega_i$. Write
\[
 N_i=\Nm_{K/L_i},\quad S_i=\Tr_{K/L_i},\quad
 n_i=\Nm_{L_i/F},\quad \operatorname{tr}_i=\Tr_{L_i/F}.
\]
For an arbitrary nontrivial additive character $\psi_F$, use
$\psi_{L_i}=\psi_F\operatorname{tr}_i$. Set $n_F=n_F(\psi_F)$, so
$\pp_F^{-n_F}$ is its largest trivial \emph{ideal}.
\begin{theorem}\label{D:NM:main}
Suppose $\theta_iN_i=\Theta$ is a primitive compatible family: $\Theta$
is invariant under $\Gal(K/F)$ and does not descend by $\Nm_{K/F}$.
Then the three expressions
\begin{equation}\label{D:NM:Adef}
 \mathcal A_i(\theta_i)=
 \Delta_{L_i}(\theta_i,\psi_{L_i})\Delta_F(\omega_i,\psi_F)
\end{equation}
are equal. This holds for every continuous quasi-character conductor
under \eqref{D:NM:breaks}.
\end{theorem}
A prescribed compatible primitive pair extends to such a family by
Hilbert 90 and character extension from the third norm image; its two
prescribed members are not changed. The trivial-character factor omitted
from \eqref{D:NM:Adef} is one.

The one-dimensional inputs are the results of~\cite{Ueda}: full Lamprecht
(\cite[Theorem~4.5]{Ueda}), stationary additive duality, quadratic norm
filtration and norm-character conductor, norm-pullback conductor
(\cite[Proposition~3.10]{Ueda}), trace/different formulas, and
nondegeneracy and magnitude of the residual quadratic sum.
For a ramified quadratic edge $E/D$ of different exponent $d$ these give
\begin{equation}\label{D:NM:edge}
 \Tr_{E/D}(\pp_E^j)=\pp_D^{\floor{(j+d)/2}},\qquad
 v_D(\Nm x)=v_Ex,\qquad n_E(\psi_E)=2n_D(\psi_D)+d.
\end{equation}
A base character of conductor $n>d$ pulls back with conductor $2n-d$;
for $n\le d$ its pullback has conductor at most $d$.
The descent and order-discriminant facts are proved in
Appendix~\ref{app:descent}, and the required crossed norm-character and
ramification facts in Appendix~\ref{app:diamond-foundations}.
In particular Lemma~\ref{D:break-ledger} gives the upper
quadratic different exponents
\begin{equation}\label{D:NM:upper}
 D'_1=4r-2a,\qquad D'_2=D'_3=2a.
\end{equation}
Every additional two-field calculation is proved below.

\subsection{Alignment of the quadratic generators}
Choose actual generators $x\in L_1$, $y\in L_2$ with
\begin{equation}\label{D:NM:xy}
 x^2+x=f,\qquad y^2+y=g,\qquad
 v_Ff=1-2a,\quad v_Fg=1-2r.
\end{equation}
Here is a construction. If a uniformizer of a quadratic field has
minimal polynomial $X^2-cX+b$, its different is the valuation of
$2\Pi-c$. The two possible valuations $2e+1$ and $2v_Fc$ have different
parities. Different $2a<2e+1$ therefore forces $v_Fc=a$.
Then $x=-\Pi/c$ satisfies \eqref{D:NM:xy}; use the same argument for $y$.
All conjugations are $x\mapsto-1-x$, $y\mapsto-1-y$.

\begin{lemma}\label{D:NM:alignment}
The generator $x$ may be changed, without changing its field or the
valuations in \eqref{D:NM:xy}, and $d\in F$ may be chosen so that
\begin{equation}\label{D:NM:align}
 v_Fd=\delta,\qquad E:=f+d^2g\in\pp_F^{1-a}.
\end{equation}
Moreover $k_0:=1+d$ is a unit, including when $a=r$.
\end{lemma}
\begin{proof}
Start with $x_0,f_0$ and put $S_0=1+2x_0$, $U_0=S_0^2=1+4f_0$.
Let $b=2e-2a+1$ and $q=\varpi^b$. Replace
\[
 S_0\quad\hbox{by}\quad S=S_0(1+\varpi^{e-a+1}w),\qquad x=(S-1)/2,
 \quad d=\varpi^{r-a}v.
\]
The two base variables $v,w$ will be integral, with $v$ a unit.
Modulo $\pp_F^a$,
\[
 \frac{S^2-1+4d^2g}{q}
 \equiv u_0+\varpi U_0w^2+v_0v^2,
 \qquad u_0=4f_0/q,\quad v_0=4\varpi^{2r-2a}g/q.
\]
The omitted linear term has valuation at least $a$; $u_0,v_0,U_0$ are
units. Since $a\le e$, the ring $\OO_F/\pp_F^a$ has characteristic two
and is $k[\varpi]/(\varpi^a)$. Its coefficient field can be constructed
as the unique lifts of the roots of $Z^{|k|}-Z$, using its unit derivative.
Solve $v_0v^2+\varpi U_0w^2=u_0$ successively in powers of $\varpi$.
At an even degree the next coefficient of $v^2$ is determined, and at
an odd degree the next coefficient of $w^2$ is determined; earlier
contributions are known and the leading coefficients are units.
Frobenius on $k$ supplies the required residue square roots. The constant
coefficient of $v$ is nonzero. Lifting $v,w$ to $F$ proves
$v_F E\ge b+a-2e=1-a$. The perturbation of $4f_0$ is strictly deeper
than $b$, so $v_Ff=1-2a$ is unchanged. This did not request a square
root of a fixed arbitrary unit modulo $\pp_F^a$.

If $r>a$, then $d\in\pp_F$. Suppose $r=a$ and $\bar d=1$.
Then $f+g$ has valuation at least $2-2a$, and so does
$f+g+4fg$, since $e\ge a$. The generator
$z=x+y+2xy\in L_3$ satisfies $z^2+z=f+g+4fg$.
The integral element $\varpi^{a-1}z$ generates an order whose
polynomial discriminant has valuation $2a-2$. The field different is
at most this value, contradicting $T_3=2a$. Thus $1+d$ is a unit.
\end{proof}

Put
\begin{equation}\label{D:NM:origin}
 z=x+y+2xy\in L_3,\qquad Y=x+dy,\qquad
 a_i=N_iY,\quad h_i=S_iY,\quad \beta_i=n_i h_i.
\end{equation}
The following are exact identities, not characteristic-two approximations:
\begin{align}
 z^2+z&=f+g+4fg,                                      \label{D:NM:zpoly}\\
 h_1&=2x-d,&h_2&=2dy-1,&h_3&=-k_0,                  \label{D:NM:htable}\\
 \beta_1&=d^2+2d-4f,&\beta_2&=1+2d-4d^2g,&\beta_3&=k_0^2,\label{D:NM:betatable}\\
 a_1&=2f-E-k_0x,&a_2&=E-2f-dk_0y,&a_3&=-E-dz.     \label{D:NM:atable}
\end{align}
The right side of \eqref{D:NM:zpoly} has valuation $1-2r$. If $r>a$, $g$
is its unique leading term; if $r=a$, Lemma~\ref{D:NM:alignment} excludes
cancellation of the leading terms of $f+g$. The term $4fg$ is strictly
deeper in both cases. Thus $v_{L_3}z=1-2r$.
The constant terms in \eqref{D:NM:atable} have base valuation at least $1-a$,
whereas its nonconstant terms have $L_i$-valuation $1-2a$. Hence
\begin{equation}\label{D:NM:originvals}
 v_KY=v_{L_i}a_i=-t_1,
 \quad v_{L_1}h_1=2\delta,\quad v_{L_2}h_2=v_{L_3}h_3=0,
 \quad \Tr_{K/F}Y=-2k_0.
\end{equation}
For $h_1$, the depth of $2x$ is $2e-2a+1>2\delta$; the other traces
are units. In particular all trace norms in \eqref{D:NM:origin} are nonzero,
$v_F\beta_1=2\delta$, and $\beta_2,\beta_3$ are units.

\subsection{The lower stationary coefficients}
Use the quadratic norm symbol, the conic identity
\eqref{D:MX:symbol-transform}, and the unit-depth estimate of
Lemma~\ref{D:MX:unit-symbol}. Their hypotheses depend only on $F$ and
$e=v_F(2)$, not on maximality of the breaks.

Choose full ordinary stationary covectors $b_1,b_2$ for the actual
lower characters:
\[
 \omega_1(1-u)=\psi_F(b_1u)\quad(u\in\pp_F^a),\qquad
 \omega_2(1-u)=\psi_F(b_2u)\quad(u\in\pp_F^r).
\]
They exist by additive duality, with valuations $-n_F-2a$ and
$-n_F-2r$. The product character $\omega_3$ has coefficient $b_1+b_2$
on $\pp_F^r$. Define
\begin{equation}\label{D:NM:normalization}
 \alpha=b_2/\beta_2,\qquad \Psi_L(u)=\psi_L(\alpha u),
 \qquad v_F\alpha=-n_F-T.
\end{equation}
The largest trivial ideals for $\Psi_F,\Psi_{L_1},\Psi_{L_2},\Psi_{L_3}$
have exponents
\begin{equation}\label{D:NM:J}
 T,\qquad J_1=4r-2a,\qquad J_2=J_3=T.
\end{equation}

\begin{lemma}[All three full lower formulas]\label{D:NM:lower}
With $q_1=a$, $q_2=q_3=r$, one has
\begin{equation}\label{D:NM:lowerformula}
 \omega_i(1-u)=\Psi_F(\beta_i u)\qquad(u\in\pp_F^{q_i}).
\end{equation}
Moreover $\omega_i(\beta_i)=1$, and
\begin{equation}\label{D:NM:normphase}
 \Psi_{L_i}(\beta_i v)=\Psi_F(\beta_i n_i v)
                              \qquad(v\in\pp_{L_i}^{q_i}).
\end{equation}
\end{lemma}
\begin{proof}
Write $U=1+4f$, $V=1+4g$. For $u\in\pp_F^a$, put
$C=1+4fu$ and $w=(f/g)u$. Then
\[
 v_F(C-1)\ge2e-a+1,\qquad v_Fw\ge2r-a\ge r.
\]
Multiplying the second symbol argument by the norm $-4f$ and using
\eqref{D:MX:symbol-transform} gives
\[
 \omega_1(1-u)=(C,f(1-u)U).
\]
The ratio $f/(-d^2g)$ lies in $U_F^a$, by \eqref{D:NM:align}.
Lemma~\ref{D:MX:unit-symbol} removes this ratio and $1-u$ from the second
argument. The result is $(C,-gU)$. Applying the same transformation
with $g,w$ gives
$\omega_2(1-w)=(C,g(1-w)V)=(C,gV)$.
The factor $-1$ in the product is also killed by
Lemma~\ref{D:MX:unit-symbol}, since $v_F(-2)=e$.
Consequently the exact identity on this whole ideal is
\begin{equation}\label{D:NM:character-relation}
 \omega_1(1-u)=\omega_2(1-(f/g)u)\,\omega_3(1+4fu).
\end{equation}
All uses of the lemma have depth sum greater than $2e$;
in particular $(2e-a+1)+a=2e+1$. The last argument has depth at
least $2e-a+1\ge r$, so all three full linear formulas apply.
Equation \eqref{D:NM:character-relation} gives, by full-ideal duality,
\begin{equation}\label{D:NM:ratio-covectors}
 b_1-b_2 R\in\pp_F^{-n_F-a},\qquad
 R=\frac{f(1-4g)}{g(1+4f)}.
\end{equation}
Here multiplication of the variable $u$ by the unit $1+4f$ preserves
its whole ideal; it is not an inference from one trivial phase.

We check that the explicit norm coefficients have exactly this ratio.
Direct expansion gives
\begin{align}\label{D:NM:lower-error}
 f(1-4g)\beta_2-g(1+4f)\beta_1
 ={}&E(1+2d+16fg)\\[-6pt]
 &-8fgd(d+2)-2gd(d^2+d+1).\notag
\end{align}
Every term has valuation at least $1-a$. For the middle term,
$v_F(d+2)=\delta<e$, so its valuation is at least
$3e+2-4a\ge1-a$; for the last it is at least
$e-r-a+1\ge1-a$. The first follows from \eqref{D:NM:align} and
$v_F(16fg)\ge2$. Division by $g(1+4f)\beta_2$, of valuation $1-2r$,
gives $R-\beta_1/\beta_2\in\pp_F^{2r-a}$.
Together with \eqref{D:NM:ratio-covectors} this proves the first lower formula.
The second is \eqref{D:NM:normalization}. Finally
\[
 \beta_3-\beta_1-\beta_2=-2d+4E\in\pp_F^r,
\]
which proves the third formula on $\pp_F^r$.
The $\beta_i$ are the actual norms in \eqref{D:NM:origin}, so their character
values are one. For $v\in\pp_{L_i}^{q_i}$ both its trace and norm lie
in $\pp_F^{q_i}$. Evaluate \eqref{D:NM:lowerformula} at the exact norm
$n_i(1-v)=1-\operatorname{tr}_i v+n_i v$. Its character is one,
which proves \eqref{D:NM:normphase} with the positive sign shown.
\end{proof}

\subsection{Construction of the minimal characters}
Put
\begin{equation}\label{D:NM:cores}
 m_1=4r-1,\quad m_2=m_3=2a+2r-1,\qquad
 s_1=2r,\quad s_2=s_3=a+r.
\end{equation}
Then $m_i=J_i+t_1=2s_i-1$. The expressions
\begin{equation}\label{D:NM:Rcharts}
 R_i(1-v)=\Psi_{L_i}(a_i v)\qquad(v\in\pp_{L_i}^{s_i})
\end{equation}
define characters with exact conductor $m_i$ on the indicated groups:
$2s_i-t_1\ge J_i$, and the last-layer coefficient is nonzero.

\begin{lemma}[The shallow commutator is evaluated, not assumed]\label{D:NM:commmodel}
Let $\sigma$ be the lower involution on $L_1$ and
$\nu_1=\omega_2 n_1$, the actual norm character for $K/L_1$.
For every $v\in L_1^\times$ with $\sigma(v)/v\in U_{L_1}^{2r}$,
\begin{equation}\label{D:NM:commmodel-eq}
 R_1(\sigma(v)/v)=\nu_1(v).
\end{equation}
\end{lemma}
\begin{proof}
Multiplication by a base element changes neither side. Odd valuation
would give conjugate quotient of exact depth $t_1<2r$, by the first
ramification term, so it is excluded. Otherwise scale to a unit and
write $v=P+Qx$. Parity of the two summand valuations makes $P$ a unit
and $v_FQ\ge a$. The element $1+2x$ is a unit of $L_1$. Thus the
assumed conjugate-quotient depth forces $w=Q/P\in\pp_F^r$.
Put
\[
 V_0=1-w-fw^2=n_1(1+wx),\quad
 Q_0=1-dw-d^2gw^2=n_2(1+dwy),\quad H_0=1+d^2gw^2.
\]
All three are units. Write $c=f-d^2g=E-2d^2g$, so $v_Fc\ge1-a$.
The difference $Q_0-V_0=(1-d)w+cw^2$ lies in $\pp_F^r$.
Since $Q_0$ is a norm from $L_2$, the actual lower formula gives
\begin{equation}\label{D:NM:comm-norm}
 \nu_1(v)=\omega_2(V_0)=
 \Psi_F\bigl(\beta_2(Q_0-V_0)/Q_0\bigr).
\end{equation}
On the other hand $\sigma(v)/v=1-w(1+2x)/(1+wx)$.
Using $a_1=f-d^2g-k_0x$, an exact quadratic trace calculation yields
\begin{equation}\label{D:NM:comm-trace}
 R_1(\sigma(v)/v)=\Psi_F(B_0/V_0),\qquad
 B_0=- (1+4f)(cw^2+k_0w).
\end{equation}

Here is the full modulus check comparing these expressions.
Both numerators in \eqref{D:NM:comm-norm}--\eqref{D:NM:comm-trace} lie in
$\pp_F^r$. Also
\[
 V_0-H_0=-w-Ew^2\in\pp_F^r,\quad
 Q_0-H_0=-dw-2d^2gw^2\in\pp_F^r.
\]
Thus replacing either unit denominator by $H_0$ costs an element of
$\pp_F^{2r}$, exactly the trivial ideal of $\Psi_F$.
With $U=1+4f$, the remaining numerator difference is
\[
 B_0-\beta_2(Q_0-V_0)
 =-(U+\beta_2)cw^2-(Uk_0+\beta_2(1-d))w.
\]
The identities
\[
 U+\beta_2=2k_0+4c,\qquad
 Uk_0+\beta_2(1-d)=2(1+d-d^2)+4Ek_0-8d^2g
\]
show that the first coefficient has valuation at least $e$, and the
second at least $e\ge r$. The first summand therefore has valuation
at least $2r+e+1-a\ge2r$ and the second at least $2r$.
Equations \eqref{D:NM:comm-norm}--\eqref{D:NM:comm-trace} now prove the assertion.
In particular no character formula on an unjustified shallow ideal was
used to evaluate $\nu_1(v)$.
\end{proof}

For any $Z\in K$ we will repeatedly use the exact identity
\begin{equation}\label{D:NM:E2}
 E_2(Z)=\operatorname{tr}_i N_iZ+n_iS_iZ.
\end{equation}
Both sides are the sum of the six pair-products of its four conjugates.
\begin{lemma}\label{D:NM:compat}
For $v\in\pp_K^{2r}$ all $N_i(1-v)$ belong to the domains in
\eqref{D:NM:Rcharts}, and $R_i(N_i(1-v))$ is independent of $i$.
\end{lemma}
\begin{proof}
The upper traces have depths at least $3r-a$ and $r+a$, and the norm
endpoints have depth $2r$. These meet all three domain requirements.
Apply \eqref{D:NM:E2} to $Y(1-v)$ and $Y$. The upper trace changes
$S_i(Yv)$ have depths at least $3r-2a$ and $r$, respectively.
After division by $h_i$, each relative change has depth at least $r$,
which is at least $q_i$. Its lower norm is an actual norm, and
Lemma~\ref{D:NM:lower} makes the phase of
$n_iS_i(Y(1-v))-\beta_i$ trivial. The other part of \eqref{D:NM:E2} is
$\operatorname{tr}_i(a_i(-S_iv+N_iv))$.
It follows that
\[
 R_i(N_i(1-v))=\Psi_F(-E_2(Y(1-v))+E_2(Y)),
\]
independent of $i$.
\end{proof}

\begin{theorem}[Full global core models]\label{D:NM:core-exists}
There exist characters $\chi_i$ of the full multiplicative groups
$L_i^\times$, with primitive common norm pullback, conductors $m_i$,
and the full formulas \eqref{D:NM:Rcharts}.
\end{theorem}
\begin{proof}
Let $I=\{\sigma(v)/v:v\in L_1^\times\}$ and prescribe
$\kappa_I(\sigma(v)/v)=\nu_1(v)$. This is well defined since
$\nu_1$ is trivial on $F^\times$. Lemma~\ref{D:NM:commmodel} proves agreement
with $R_1$ on $I\cap U_{L_1}^{s_1}$. Their product character factors
through $I U_{L_1}^{s_1}/U_{L_1}^{m_1}$; extend it across
$L_1^\times/U_{L_1}^{m_1}$, first on its finite unit subgroup and then
on a uniformizer. The resulting $\chi_1$ has exact conductor $m_1$
and satisfies $\chi_1^\sigma/\chi_1=\nu_1$ globally.

Its norm pullback to $K$ is invariant under both involutions, since
$\nu_1N_1=1$. Hilbert 90 makes it trivial on each other norm kernel;
extend its descents to $\chi_2,\chi_3$. Norm filtration gives
\[
 N_i(U_K^{2r+1})=U_{L_i}^{a+r}\qquad(i=2,3),
\]
since these upper edges have break $2a-1$ and
$2(a+r)-(2a-1)=2r+1$. Lemma~\ref{D:NM:compat} gives the full $R_i$ formulas
on these groups. Their conductors are at least $m_i>D'_i$.
The common upstairs conductor, computed from $\chi_1$, is
$2m_1-D'_1=4r+2a-2$. The high norm-pullback conductor formula then
forces the other conductors to be exactly $m_i$.
If the common pullback descended from $F$, $\chi_1$ would be a base
pullback times its upper norm character, both invariant under $\sigma$.
This contradicts its nontrivial commutator. Thus the family is primitive.
\end{proof}

\subsection{Common-twist realization and norm choices}
\begin{lemma}\label{D:NM:realize}
For every primitive compatible family,
\[
 \theta_i^\sigma/\theta_i=\nu_i=\omega_jn_i\ (j\ne i),\qquad
 D_\theta=\theta_i|_{F^\times}\omega_i
\]
is independent of $i$. There are cores with the formulas of
Theorem~\ref{D:NM:core-exists} and a base character $\lambda$ such that
\begin{equation}\label{D:NM:family}
 \theta_i=\chi_i(\lambda n_i).
\end{equation}
If $n=a_F(\lambda)$, their conductors are $m_i$ for $n\le2r+a-1$;
for $n\ge2r+a$ they are $M_i=2n-T_i$, all even.
\end{lemma}
\begin{proof}
Lemma~\ref{U:conjugacy} gives the stated conjugate quotient and
corrected common restriction, for the given family and the core family.
The quotient $\theta_1/\chi_1$ is invariant, so Hilbert 90 and character
extension give $\lambda n_1$. The remaining discrepancies on sides
2 and 3 are powers of $\nu_i$. Absorb them into the cores: their
conductors $2a$ are at most $s_i=a+r$, so the full core charts do not
change. This proves \eqref{D:NM:family} for the given family.
The pullback of a base character with $n>T_i$ has conductor $2n-T_i$;
otherwise it has conductor at most $T_i<m_i$. Since $m_i$ is odd and
$2n-T_i$ is even, equal-conductor cancellation does not occur.
The common threshold is $n=2r+a$, giving precisely the alternatives.
\end{proof}

\begin{lemma}[Normalization freedom after realization]\label{D:NM:freedom}
Replacing $\alpha$ by $\alpha u$ with $u\in U_F^{T-1}$ preserves all
full core and lower formulas for the already fixed characters.
For any nonzero coefficient $A_*$ of $\lambda$ relative to $\Psi_F$,
such a replacement makes $A=A_*/u$ an exact norm from $L_3/F$.
The raw coefficient $\alpha A_*$ and the base character itself are
unchanged.
\end{lemma}
\begin{proof}
On each lower ideal the normalized error has depth at least
$T-1+v_F\beta_i+q_i\ge T$. On the upper core ideals it has depth
at least $2(T-1)-t_1+s_i\ge J_i$. These inequalities include $a=r=1$.
Thus the full formulas are preserved. The actual $\omega_3$ is
onto $\{1,-1\}$ on $U_F^{T-1}$, so Lemma~\ref{U:normalization}
applies with this subgroup and gives the exact norm while keeping
$\alpha A_*$ fixed.
\end{proof}

\subsection{Adjustment through the third quadratic field}
For comparison of fields 1 and 2, take $X\in L_3$ and put
\begin{equation}\label{D:NM:adjust}
 A=n_3X,\quad p_X=\operatorname{tr}_3X,\quad C=Y-X,\quad
 Q_i=S_i(Yg_iX)\quad(i=1,2).
\end{equation}
For $X=0$ all expressions below have their evident zero values.
With primes denoting conjugation in $L_3/F$, direct multiplication gives
\begin{align}
 N_iC&=a_i+A-Q_i,&S_iC&=h_i-p_X,                       \label{D:NM:commonC}\\
 n_iQ_i&=\beta_i A+\mathcal E,&
 \mathcal E&=(X'-X)(a_3X'-a_3'X)\in F,               \label{D:NM:common-error}\\
 n_iS_iC&=\beta_i+\Delta,&
 \Delta&=p_X^2+2k_0p_X.                              \label{D:NM:common-Delta}
\end{align}
The last identity uses $\operatorname{tr}_i h_i=-2k_0$, the same for both
fields. It is valid even when a trace of $C$ vanishes; no character is
evaluated at such a zero.

\begin{lemma}\label{D:NM:bounds}
If $v_{L_3}X=-h$ for any integer $h$, then
\begin{equation}\label{D:NM:bounds-eq}
 \begin{aligned}
 v_Fp_X&\ge r-\ceil{h/2},& v_F\mathcal E&\ge T-a-h,\\
 v_{L_1}Q_1&\ge2\delta-h,&v_{L_2}Q_2&\ge-h.
 \end{aligned}
\end{equation}
\end{lemma}
\begin{proof}
Write $X=P+Qz$, with $z$ as in \eqref{D:NM:zpoly}. Parity of valuations gives
$v_FP\ge-\floor{h/2}$ and $v_FQ\ge r-\ceil{h/2}$.
The trace $2P-Q$ has the asserted bound because $e\ge r$.
Using $a_3=-E-dz$, the exact common error is
\[
 \mathcal E=Q(dP-EQ)(1+4(f+g+4fg)).
\]
Its last factor is a unit. The first product term has depth at least
$r+\delta-h=T-a-h$; the second at least
$1-a+2r-2\ceil{h/2}\ge T-a-h$.
Finally $Yg_iX$ has $K$-valuation $-t_1-2h$.
The upper trace formula \eqref{D:NM:edge}, with \eqref{D:NM:upper}, gives the two
$Q_i$ bounds exactly as stated.
\end{proof}
Whenever $(Q_i/\beta_i)v\in\pp_{L_i}^{q_i}$, the actual norm phase gives
\begin{equation}\label{D:NM:Qphase}
 \Psi_{L_i}(Q_iv)
   =\Psi_F\bigl((A+\mathcal E/\beta_i)n_i v\bigr).
\end{equation}
If the error term is in $\pp_F^T$, exact quadratic norm expansion
therefore evaluates $\lambda n_i$ with coefficient $A-Q_i$ on its
whole stationary ideal. Adding the core coefficient gives $N_iC$.
The shared error, not an unjustified individual exact norm equation,
is what matters in this step.

\subsection{Common stationary origins at minimal conductors}
Assume $n\le2r+a-1$ in \eqref{D:NM:family}. Both trace and norm of a variable
in the core stationary ideal $\pp_{L_i}^{s_i}$ lie in
$\pp_F^{q_F}$, where $q_F=a+r$. If $n\le q_F$, take $X=0$.
Otherwise $h=n-T$ satisfies
\begin{equation}\label{D:NM:minimal-h}
 1-\delta\le h\le a-1.
\end{equation}
Since $2q_F\ge n$, the base character has a linear formula
$\lambda(1-v)=\Psi_F(A_*v)$ on $\pp_F^{q_F}$, with $v_FA_*=-h$.
Apply Lemma~\ref{D:NM:freedom} and choose $A=n_3X$ exactly, with
$v_{L_3}X=-h$. Use the adjusted $\alpha$ henceforth.

\begin{lemma}\label{D:NM:minimal-origin}
The element $C$ so constructed has $v_KC=-t_1$. For $i=1,2$,
\begin{align}
 \theta_i(1-v)&=\Psi_{L_i}(N_iC\,v)
                              &&(v\in\pp_{L_i}^{s_i}),  \label{D:NM:upper-stationary}\\
 \omega_i(1-v)&=\Psi_F(n_iS_iC\,v)
                              &&(v\in\pp_F^{q_i}).      \label{D:NM:lower-stationary}
\end{align}
The traces $S_iC$ are nonzero with the valuations of $h_i$.
If $a=r$, this same $C$ satisfies both formulas for all three indices.
\end{lemma}
\begin{proof}
For $X\ne0$, the inequality $h\le a-1$ gives $-2h>-t_1$, so $Y$
dominates in $C$. On the two upper ideals the arguments of
\eqref{D:NM:Qphase} have respective depths at least $2a-h\ge a$ and
$a+r-h\ge r$. Its weighted errors have depths at least
$T+a-h$ and $3r-h$, respectively, both at least $T$.
The trace and norm variables lie in the full base ideal already specified.
Thus exact norm expansion and \eqref{D:NM:commonC} prove
\eqref{D:NM:upper-stationary} for the actual characters.

By Lemma~\ref{D:NM:bounds},
\[
 v_F\Delta\ge\min\{2r-2\ceil{h/2},\ e+r-\ceil{h/2}\}.
\]
Adding $a$ to both bounds gives at least $T$, since $h\le a-1$;
adding $r$ does also. Equation \eqref{D:NM:common-Delta} therefore proves
\eqref{D:NM:lower-stationary}. Moreover
$2r-2\ceil{h/2}>2\delta$ and $r-\ceil{h/2}>0$, so the trace
perturbations cannot cancel the leading terms of $h_1,h_2$.

When $a=r$, one has $s_3=T$, $q_3=r$, and $h\le r-1$.
On side 3 the exact norm phase applied to $(X/k_0)v$ gives the
pullback coefficient $A-k_0X$, since $\beta_3=k_0^2$.
But
\[
 N_3C=a_3+X^2+k_0X,
 \qquad (X^2+k_0X)-(A-k_0X)=p_XX-2A+2k_0X.
\]
The three terms on the right have $L_3$-valuations at least
$2r-2h-(h\bmod2)$, $2e-2h$, and $2e-h$, all nonnegative.
They are invisible on $\pp_{L_3}^{T}$, proving the third upper formula.
Also $S_3C=-k_0-2X$ is a unit and
\[
 n_3S_3C-\beta_3=2k_0p_X+4A.
\]
Both terms have base valuation at least $r$, giving the third lower
formula. The case $X=0$ is already contained in the core and lower
formulas.
\end{proof}

\subsection{The exact minimal sign}
Write $c_0=-\alpha^{-1}$ and use the full factors and terminal
invariance of Lemma~\ref{U:stationary-factor}. Their residual functions
and normalized sums will be denoted by $H_i,g_i$.

Use $N_iC$ for the upper coefficient and $n_iS_iC$ for the lower one
in Lemma~\ref{D:NM:minimal-origin}. The latter is an actual norm and all
$T_i$ are even. By compatibility, the determinant restriction, and
\eqref{D:NM:E2}, the full expression is
\begin{equation}\label{D:NM:minimal-factor}
 \mathcal A_i=D_\theta(c_0)\Theta(C)^{-1}
                         \Psi_F(-E_2(C))g_i.
\end{equation}
Here $g_i$ is the full upper residual factor.

Choose a top uniformizer $\pi$ and compatible lower uniformizers
$\Pi_i=N_i\pi$, $\varpi=n_i\Pi_i$. These choices need not be the ones
used to construct the generators earlier. Vary
$C$ to $C_v=C(1+\pi^{t_2}v)$ for integral $v\in K$.
The upper critical depths are $t_2$ on side 1 and $a+r-1$ on sides 2,3.
The changes $S_i(C\pi^{t_2}v)$ have depths at least $3r-2a$ and $r$.
After division by $S_iC$ they have depth at least $r\ge q_i$.
Their lower norm quotients are actual norms on the lower stationary
ideals. Subtracting \eqref{D:NM:E2} at $C_v,C$, using the full lower formulas
and canceling the common multiplicative value, gives
\begin{equation}\label{D:NM:actual-functions}
 H_i(p_i(\bar v))=Q(\bar v),
\end{equation}
where $p_i$ is the actual upper graded norm map and $Q$ is independent
of the compared index. All these assertions concern full functions,
not merely their polar forms. Every trace used as a multiplicative
argument stays nonzero, since its relative change has positive depth.

\begin{proposition}[Two unequal breaks]\label{D:NM:minimal-two}
If $a<r$, the minimal family satisfies $\mathcal A_1=\mathcal A_2$.
\end{proposition}
\begin{proof}
The source depth $t_2=2r-1$ is below the first upper break
$4r-2a-1$. Thus $p_1(v)=v^2$, with its trace term deeper.
On the second edge the consecutive trace ideals at source depths
$2r-1,2r$ have depths $a+r-1,a+r$. Hence the trace induces a nonzero
$k$-linear map of residue lines, and the norm endpoint is strictly deeper
because $r>a$. Thus $p_2$ is also bijective.
Summing \eqref{D:NM:actual-functions} proves $g_1=g_2$.
Equation \eqref{D:NM:minimal-factor} gives the assertion. The relabeled
construction compares the first with the third field.
\end{proof}

\begin{proposition}[One equal break]\label{D:NM:minimal-one}
If $a=r$, every minimal family satisfies
$\mathcal A_1=\mathcal A_2=\mathcal A_3$.
\end{proposition}
\begin{proof}
Now $t=t_1=t_2=2r-1$ and Lemma~\ref{D:NM:minimal-origin} provides the same
$C$ on all three sides. Put
\[
 u=\overline{C\pi^t}\in k^\times,\qquad
 c_i=\overline{g_iC-C}\in k^\times,\qquad P=c_1c_2c_3,
 \qquad k_i=c_jc_k\quad(\{i,j,k\}=\{1,2,3\}).
\]
The residues are $\bar d,1,1+\bar d$, respectively. Indeed the origin
formulas give these residues for $Y$, and the adjustment $X$ changes
none: its conjugate differences have positive valuation by
$h\le r-1$ and the ramification bound. Thus the $c_i$ are distinct and
$c_1+c_2=c_3$.

Expansion of the odd-valuation leading term $u\pi^{-t}$ shows that the
first ramification coefficient on $\pi$ is $c_i/u$.
The exact quadratic critical norm polynomial consequently is
\begin{equation}\label{D:NM:critical-polynomials}
 p_i(v)=v^2+(c_i/u)v.
\end{equation}
One can fix the trace coefficient without a convention choice: the norm
endpoint has leading coefficient one, and $g_i\pi/\pi$ is the nontrivial
critical kernel element with residue $c_i/u$.
The lower conjugate difference of $N_iC$ has residue $k_i$.
For $Y$ this follows immediately from \eqref{D:NM:atable}; the change to
$C$ has conjugate difference of positive valuation. Its leading term is
$u^2\Pi_i^{-t}$. Therefore the lower critical norm polynomial is
$v^2+\eta_i v$ with $\eta_i=k_i/u^2$.

Write the actual last-layer additive character as
\[
 \Psi_F(\varpi^t v)=(-1)^{\Tr_{k/\F_2}(\gamma\bar v)},\qquad
 \gamma\in k^\times.
\]
The lower stationary norm coefficient $n_iS_iC$ has residue $c_i^2$:
$2C$ is of positive $K$-valuation. Triviality of $\omega_i$ on its
actual critical norms and full residue duality give
\[
 \Tr_{k/\F_2}(\gamma c_i^2(v^2+\eta_i v))=0\quad(v\in k),
 \quad \eta_i=(\sqrt\gamma c_i)^{-1},\quad u^2=\sqrt\gamma P.
\]
Every unit scalar is thus fixed. Tracing the leading term of
$N_iC\,\Pi_i^{2t}$ in the Lamprecht polar formula gives the polar
pairing of the actual $H_i$:
\begin{equation}\label{D:NM:polar}
 \frac{H_i(x+y)}{H_i(x)H_i(y)}
       =(-1)^{\Tr_{k/\F_2}(\gamma k_i xy)}.
\end{equation}
The omitted next term in that trace lies in $\pp_F^{t+1}$ by
\eqref{D:NM:edge}. No normalized coefficient $\gamma=1$ has been presumed.

The image of \eqref{D:NM:critical-polynomials} is the index-two hyperplane
\[
 V_i=\{x:\Tr_{k/\F_2}(u^2x/c_i^2)=0\}.
\]
For $j\ne i$, put $w_i=p_i(c_j/u)=k_i/u^2\ne0$.
The common \emph{actual} function \eqref{D:NM:actual-functions} gives
$H_i(w_i)=H_j(0)=1$. The polar character of $w_i$ in \eqref{D:NM:polar}
is the character defining $V_i$, since $u^4=\gamma P^2$.
All hypotheses of Lemma~\ref{U:missing-coset} now hold on $(k,+)$,
with image $V_i$. It gives $\sum_kH_i=\frac12\sum_kQ$, independently
of $i$. The common normalization and \eqref{D:NM:minimal-factor}
therefore identify the complete factors.
\end{proof}

\subsection{All higher conductors}
Now $n=T+h\ge2r+a$, so $h\ge a$. The upper conductors and stationary
depths are
\begin{equation}\label{D:NM:higher-depths}
 M_i=2n-T_i,\qquad b_1=2r+h-a,\quad b_2=r+h,
 \qquad M_i=2b_i.
\end{equation}
For a variable in either $\pp_{L_i}^{b_i}$, $i=1,2$, its trace and norm
belong to the ordinary base stationary ideal
$\pp_F^{\ceil{n/2}}$. For the first trace this follows from
$r+\floor{(h+a)/2}\ge r+\ceil{h/2}$; for the second use $r\ge1$.
Choose the actual coefficient $A_*$ of $\lambda$ there, of valuation
$-h$. Apply Lemma~\ref{D:NM:freedom} to make $A=n_3X$ exact, and use
\eqref{D:NM:adjust}. There is no upper bound on $h$.

\begin{theorem}[Exact comparison for every higher family]\label{D:NM:higher}
For every $h\ge a$, one has $\mathcal A_1=\mathcal A_2$.
\end{theorem}
\begin{proof}
On the full stationary ideals the arguments $(Q_i/\beta_i)v$ of
\eqref{D:NM:Qphase} have depths at least $a$ and $r$, respectively:
\[
 (2\delta-h)-4\delta+b_1=a,\qquad -h+b_2=r.
\]
The weighted errors have base depths at least
\[
 (T-a-h)-2\delta+b_1=T,\qquad (T-a-h)+b_2=3r-a\ge T.
\]
Thus the actual pullbacks $\lambda n_i$ have full coefficient $A-Q_i$,
and the core charts, valid because $b_i\ge s_i$, give coefficient $N_iC$
for the actual $\theta_i$. Moreover $X$ dominates $Y$, as
$-2h<-t_1$. Hence $C\ne0$ and $v_{L_i}N_iC=-2h=J_i-M_i$.

Retain the original lower coefficients $\beta_i$, which are actual norms.
It is unnecessary, and not asserted, that $n_iS_iC$ are still lower
stationary coefficients in this range. All upper and lower conductors
are even, so every residual factor is one. The complete formula gives
\[
 \mathcal A_i=D_\theta(c_0)\Theta(C)^{-1}
               \Psi_F(-\operatorname{tr}_i N_iC-\beta_i)
 =D_\theta(c_0)\Theta(C)^{-1}\Psi_F(-E_2(C)+\Delta).
\]
The last identity is \eqref{D:NM:E2} and \eqref{D:NM:common-Delta}.
Every factor is the same for $i=1,2$. In particular a possibly shallow
$\Delta$ is retained as a common phase rather than discarded.
Relabel the third field for the other comparison.
\end{proof}

\subsection{Completion of the nonmaximal case}
\begin{proof}[Proof of Theorem~\ref{D:NM:main}]
Lemmas~\ref{D:NM:alignment} and \ref{D:NM:lower} construct an actual common origin
and full norm-character coefficients. Lemmas~\ref{D:NM:commmodel} and
\ref{D:NM:compat} prove the extension and compatibility conditions on their
whole domains, giving the actual cores of Theorem~\ref{D:NM:core-exists}.
Lemma~\ref{D:NM:realize} realizes the prescribed family by a common base twist.
For $n\le2r+a-1$ use Lemma~\ref{D:NM:minimal-origin} and
Propositions~\ref{D:NM:minimal-two}--\ref{D:NM:minimal-one}.
For $n\ge2r+a$ use Theorem~\ref{D:NM:higher}.
These two conductor intervals are consecutive. The one-break and two-break
ramification possibilities in \eqref{D:NM:breaks} are both included.
No step uses an equality of the desired local constants to deduce a
field congruence or a norm equation.
\end{proof}


\section{Completion of the Second Main Lemma}\label{sec:completion}
\subsection{Ramification possibilities}\label{U:exhaustion}
\begin{lemma}[Selection of fields in wild odd degree]\label{U:odd-selection}
Suppose $p$ is odd and $K/F$ is totally ramified with group $C_p^2$.
One can choose a degree-$p$ lower field $B_1$ so that, for every other
lower field $B_2$, their lower breaks are $t\le t_2=t+\delta$, and the
breaks of $K/B_2,K/B_1$ are $t,t'=t+p\delta$, respectively.
Here $\delta\ge0$ is an integer and $p\nmid t$.
\end{lemma}
\begin{proof}
Lemma~\ref{D:break-ledger} gives the distinguished field $B_1$,
the lower breaks $t,t_2$, and the upper-edge breaks
$t'=t+p(t_2-t),t$. Its proof uses only the subgroup definition and
transitivity of the different, not a quotient ramification theorem.
In particular $\delta=t_2-t$ is a nonnegative integer.
In mixed characteristic, the trace of $1$ on $B_2/F$ and
\eqref{U:trace-ideal} imply $v_K(p)\ge p(p-1)t_2$.
The proof of Proposition~\ref{O:A:prop:AS}, before using any assumption
$p\nmid t$, constructs an exact Artin--Schreier coordinate by additive
Hilbert~90 and best approximation and proves precisely that divisibility
assertion. Its hypotheses are the edge breaks and the displayed
$p$-valuation bound, so it applies here without circularity.
The equal-characteristic argument of that proposition proves the same
assertion. There is consequently no extra maximal-break case missing from
the odd total-ramification calculation.
\end{proof}

\begin{lemma}[Quadratic conductors in mixed characteristic]\label{U:quadratic-types}
Let $F/\mathbf Q_2$ be finite and $e=v_F(2)$. A ramified quadratic norm
character has conductor $2a$, $1\le a\le e$, or $2e+1$.
Any two such characters of conductor $2e+1$ have the same restriction to
$U_F^{2e}$, so their product has conductor at most $2e$.
\end{lemma}
\begin{proof}
The trace of $1$ gives $e\ge\lfloor D/2\rfloor$ for the quadratic
different exponent $D$, so $2\le D\le2e+1$. The even possibilities
are as stated. Lemma~\ref{D:UR:oddT} proves that an odd $D$ is
$2e+1$ and fixes the last-layer value in coordinates $1+4z$ as
$(-1)^{\Tr_{k_F/\F_2}\bar z}$. This value is independent of the
maximal character; hence the product of two is trivial on $U_F^{2e}$.
\end{proof}

\begin{lemma}[Exhaustive local classification]\label{U:dispatch}
Every $C_\ell^2$ local diamond belongs to one of the proof branches in
Sections~\ref{sec:tame}--\ref{sec:dyadic-nonmaximal}.
\end{lemma}
\begin{proof}
By Lemma~\ref{D:inertia} the finite residue Galois group is cyclic;
a $C_\ell^2$ extension cannot be unramified. Its inertia has order $\ell$ or $\ell^2$.
If $\ell\ne p$, wild inertia is trivial and tame inertia is cyclic,
so inertia has order $\ell$. Abelianness and the residue-Frobenius
conjugation action imply $\ell\mid |k_F|-1$. There is a unique unramified
lower field and every other lower field is tamely ramified. These are
exactly the hypotheses of Section~\ref{sec:tame}.

Suppose $\ell=p$. Inertia of order $p$ gives the unramified--ramified
configuration, treated in Sections~\ref{sec:odd-ur}--\ref{sec:odd-calculation} for odd $p$ and Section~\ref{sec:dyadic-ur} for $p=2$.
If inertia has order $p^2$, the extension is totally ramified.
For odd $p$, Lemma~\ref{U:odd-selection} provides all the hypotheses of
the totally ramified proof in Sections~\ref{sec:odd-estimates}--\ref{sec:odd-calculation}.

It remains to classify totally ramified biquadratic extensions.
In characteristic exactly two, the reduced Artin--Schreier representatives
of Section~\ref{sec:dyadic-equal} have positive odd breaks and the two largest equal, including
the one-break case. Theorem~\ref{D:EQ:main} treats all of them.
In mixed characteristic the three nontrivial lower norm characters obey
$\omega_3=\omega_1\omega_2$ by Lemma~\ref{D:quadratic-norms},
and all are ramified. The product of two
characters of unequal conductors has the larger conductor, so the maximum
of their three conductors occurs at least twice. If it is $2e+1$,
Lemma~\ref{U:quadratic-types} makes the third smaller, and the lower breaks
are $(2a-1,2e,2e)$ with $1\le a\le e$, precisely Section~\ref{sec:dyadic-maximal}.
Otherwise their conductors are $(2a,2r,2r)$ after relabeling, with
$1\le a\le r\le e$, giving the breaks $(2a-1,2r-1,2r-1)$ of Section~\ref{sec:dyadic-nonmaximal}.
The upper different formulas in those parts follow from
Lemma~\ref{D:break-ledger}. No further parity or divisibility case is
assumed away.
\end{proof}

\subsection{Proof of the main theorem}
\begin{proof}[Proof of Theorem~\ref{U:main}]
For $\ell\ne p$, Theorem~\ref{O:F:tame} compares the unique unramified
lower field with every other lower field. This includes $\ell=2$ and
its exact tame sign when $p$ is odd.
For $\ell=p$ odd and inertia of order $p$, use
Theorem~\ref{O:I:main}; its conductor dispatch includes the final cubic
affine correction. If inertia has order $p^2$, choose the distinguished
lower field as in Lemma~\ref{U:odd-selection} and apply
Theorem~\ref{O:G:totalbranch} to its comparison with every other field.

For $p=\ell=2$ and inertia of order two, use Theorem~\ref{D:UR:URall}.
For total ramification in characteristic two, use
Theorem~\ref{D:EQ:main}. For total ramification in mixed characteristic,
Theorem~\ref{D:MX:main} or Theorem~\ref{D:NM:main} applies according to
Lemma~\ref{U:dispatch}. These branches exhaust every character conductor:
each has an actual common-twist realization and consecutive minimal,
intermediate, and high conductor ranges, or a direct all-high-conductor
comparison. All quadratic correction factors and full residual functions
are retained.

In every case the given compatible characters have precisely the common
norm pullback required by the branch theorem, and primitivity is the
non-descent hypothesis in those theorems. Where realization conjugates
an inducing character, it is an upper norm-character twist. An
automorphism change of variables in the defining local integral preserves
its local constant and its common pullback. This uses no two-field
comparison. Comparisons with a distinguished field give independence of
all the intermediate fields by transitivity of equality.
The canonical additive character in the equal-characteristic calculation
was removed by its proved scaling argument and the common determinant
restriction. All other arguments already use the given trace-compatible
additive characters. The result is therefore the full normalized expression
of Theorem~\ref{U:main}, for the given quasi-characters in either
characteristic.
\end{proof}


\appendix
\addtocontents{toc}{\protect\setcounter{tocdepth}{1}}
\section{The leading Gauss congruence, with proof}\label{app:leading-gauss}
This appendix proves the precise leading congruence used in
Section~\ref{sec:tame}. Only finite-field character orthogonality, the
binomial theorem, and the elementary local lifting of roots of unity are
needed. In particular, neither a factorization theorem for Gauss sums nor
a formula involving the $p$-adic gamma function is an input.

\subsection{The coefficient field and notation}
Let $k=\F_q$, $q=p^f$, and choose an unramified extension $K_0/\Q_p$
with residue field $k$. For clarity, this elementary lift can be
constructed as follows. Lift a monic irreducible polynomial of degree $f$
over $\F_p$ to a monic polynomial over $\mathbf Z_p$, and take its
quotient ring $R$. It is finite free and $p$-adically complete over
$\mathbf Z_p$, with $R/pR=k$. Every element with nonzero residue is a
unit: lift an inverse modulo $p$ and invert the remaining element of
$1+pR$ by its convergent geometric series. Every nonzero element is
$p^n$ times such a unit, since $R$ is finite free and separated over
$\mathbf Z_p$. Thus $R$ is a discrete valuation ring with uniformizer
$p$, and its fraction field is the required $K_0$.

Hensel lifting of $X^{q-1}-1$ gives a unique multiplicative lift
$x\mapsto[x]$ from $k^\times$ into $R^\times$: its derivative is a
unit at every nonzero residue root. The lifts multiply because their
products are again roots with the prescribed residue. Choose a primitive
$p$th root $\zeta$ and put
\[
 \varpi=\zeta-1,\qquad K_1=K_0(\zeta),\qquad \OO_1=\OO_{K_1}.
\]
For $p>2$, the polynomial
\[
 \frac{(1+X)^p-1}{X}=p+\binom p2X+\cdots+X^{p-1}
\]
is Eisenstein; hence $\varpi$ is a uniformizer of $K_1$ and its residue
field is $k$. For $p=2$, $K_1=K_0$ and $\varpi=-2$ is already a
uniformizer. All root-of-unity values below are viewed in $K_1$; the
identities are algebraic identities among these values.

Let $\operatorname{tr}_k=\Tr_{k/\F_p}$ and
$\psi(x)=\zeta^{\operatorname{tr}_k(x)}$. For $0<a<q-1$, write
\[
 a=\sum_{i=0}^{f-1}a_ip^i,\qquad
 s_p(a)=\sum_i a_i,\qquad a!_p=\prod_i a_i!.
\]
Define the plus-sign Gauss sum
\[
 G_a=\sum_{x\in k^\times}[x]^{-a}\psi(x).
\]
The convention of Section~\ref{sec:tame} is
$\tau_k(\omega^a)=-G_a$, where $\omega(x)=[x]$.

\subsection{A first-order calculation and a Jacobi sum}
\begin{lemma}\label{A:gauss-first}
If $q>2$, then $G_1\equiv-\varpi\pmod{\varpi^2\OO_1}$.
Moreover $G_{p^j}=G_1$ whenever $0\le j<f$ and $p^j<q-1$.
\end{lemma}
\begin{proof}
For each $x$, let $t_x\in\{0,\ldots,p-1\}$ represent
$\operatorname{tr}_k(x)$. Then
$\psi(x)=(1+\varpi)^{t_x}\equiv1+t_x\varpi\pmod{\varpi^2}$.
The exact identity $\sum_{x\ne0}[x]^{-1}=0$ follows by summing a
nontrivial character of the cyclic group $k^\times$. Therefore
\[
 G_1\equiv\varpi\sum_{x\ne0}[x]^{-1}t_x\pmod{\varpi^2}.
\]
The residue of the last sum is
\[
 \sum_{x\ne0}x^{-1}\operatorname{tr}_k(x)
   =\sum_{i=0}^{f-1}\sum_{x\ne0}x^{p^i-1}=-1\quad\text{in }k.
\]
Indeed the $i=0$ inner sum is $q-1=-1$, while every other exponent is
strictly between $0$ and $q-1$ and its power sum is zero. A power sum
vanishes by multiplying its variable by an element on which the
corresponding character is nontrivial. Finally, Frobenius permutes
$k^\times$ and preserves its absolute trace. Reindexing by the inverse
Frobenius shows that $G_{pa}=G_a$, with indices reduced modulo $q-1$.
Iteration gives the last assertion.
\end{proof}

\begin{lemma}[Jacobi sums when the digits add without carries]
\label{A:jacobi-unit}
Suppose $a,b>0$, $a+b<q-1$, and $a_i+b_i<p$ for every digit. Put
\[
 J(a,b)=\sum_{x\in k\setminus\{0,1\}}[x]^{-a}[1-x]^{-b}.
\]
Then
\begin{equation}\label{A:jacobi-product}
 G_aG_b=J(a,b)G_{a+b},
\end{equation}
and $J(a,b)$ is a unit with
\begin{equation}\label{A:jacobi-reduction}
 J(a,b)\equiv-\prod_{i=0}^{f-1}\binom{a_i+b_i}{a_i}
                  \pmod{\varpi\OO_1}.
\end{equation}
\end{lemma}
\begin{proof}
Expand the product of the Gauss sums and separate $t=x+y$. When $t=0$
the contribution is zero because the character $[x]^{-(a+b)}$ is
nontrivial. When $t\ne0$, write $x=tu$, $y=t(1-u)$. The two sums
separate, giving \eqref{A:jacobi-product}.

For the reduction of the Jacobi sum, both exponents below are positive,
so the terms at $0$ and $1$ may be included. In $k$ the reduction is
\[
 \sum_{x\in k}x^{q-1-a}(1-x)^{q-1-b}
 =\sum_{j=0}^{q-1-b}(-1)^j\binom{q-1-b}{j}
                          \sum_{x\in k}x^{q-1-a+j}.
\]
The exponent is positive and smaller than $2(q-1)$. Its only possible
positive multiple of $q-1$ occurs at $j=a$. This index occurs because
$a+b<q-1$. The reduction is consequently
$-(-1)^a\binom{q-1-b}{a}$.

To evaluate this binomial coefficient, use in $\F_p[X]$ the identity
\[
 (1+X)^{q-1-b}
    =\prod_{i=0}^{f-1}(1+X^{p^i})^{p-1-b_i}.
\]
Comparison of coefficients gives
$\binom{q-1-b}{a}=\prod_i\binom{p-1-b_i}{a_i}$ modulo $p$;
each exponent on the right has a unique digit expansion. Since
$a_i+b_i<p$, ordinary factorials give
\[
 \binom{p-1-b_i}{a_i}
       \equiv(-1)^{a_i}\binom{a_i+b_i}{a_i}\pmod p.
\]
For odd $p$, $a\equiv\sum_i a_i\pmod2$, so the signs cancel.
For $p=2$ the same sign equality holds in the residue field.
This proves \eqref{A:jacobi-reduction}. None of its factorials is
zero modulo $p$, proving that $J(a,b)$ is a unit.
\end{proof}

\subsection{The complete leading unit}
\begin{theorem}[Leading Gauss congruence]\label{A:leading-gauss}
For every prime $p$, every $q=p^f$, and $0<a<q-1$,
\begin{equation}\label{A:leading-plus}
 G_a\equiv-\frac{\varpi^{s_p(a)}}{a!_p}
                       \pmod{\varpi^{s_p(a)+1}\OO_1}.
\end{equation}
In particular $G_a$ has $\varpi$-valuation exactly $s_p(a)$, and
\begin{equation}\label{A:leading-minus}
 \tau_k(\omega^a)\equiv_\times
       \frac{\varpi^{s_p(a)}}{\prod_i a_i!}.
\end{equation}
Here $X\equiv_\times Y$ means $X/Y\in1+\varpi\OO_1$.
\end{theorem}
\begin{proof}
If $q=2$ there is no index in the specified interval. Otherwise use
induction on $s_p(a)$. For digit sum one, $a=p^j$, and
Lemma~\ref{A:gauss-first} proves the assertion.
For larger digit sum choose $j$ with $a_j>0$ and put $c=a-p^j$.
Then $c>0$, and $c+p^j=a<q-1$ is an addition without carries.
Lemma~\ref{A:jacobi-unit} gives $J(c,p^j)\equiv-a_j\pmod\varpi$.
It is a unit, so the exact product formula and the induction hypothesis
give
\[
 G_a=\frac{G_cG_{p^j}}{J(c,p^j)}
 \equiv-\frac{\varpi^{s_p(a)}}{
                a_j(a_j-1)!\prod_{i\ne j}a_i!}
                \pmod{\varpi^{s_p(a)+1}}.
\]
Division by the unit does not lower the error precision. This proves
\eqref{A:leading-plus}, including its nonzero leading coefficient.
Changing from $G_a$ to $-G_a$ proves \eqref{A:leading-minus} with the
positive sign required in \eqref{O:F:Stick}. The argument applies
unchanged when $p=2$ and $\varpi=-2$.
\end{proof}

Thus the leading unit used to resolve the root of unity in the tame
comparison is proved here, not imported as an additional Gauss-sum
theorem. The base-$q$ grouping and factorial manipulations following
\eqref{O:F:Stick} are consequences of this theorem and the proofs
already given in Section~\ref{sec:tame}.


\section{Formal local residues and their trace compatibility}\label{app:local-residues}
This appendix proves the local identity used in the characteristic-two
argument. A differential on $k((t))$ means a \emph{formal differential}
$h(t)\,dt$, with termwise differentiation of Laurent series, and
\[
 \Res_t\left(\sum_n h_nt^n\,dt\right)=h_{-1}.
\]
No global curve or residue theorem is used. We first prove substitution
and trace compatibility algebraically, including in wild ramification.

\subsection{Substitution and a finite free expansion}
\begin{lemma}[Formal substitution]\label{B:substitution}
Let $k$ be a field and let $u(t)\in t^e k[[t]]$ have exact order
$e\ge1$. For every $H(s)\in k((s))$,
\begin{equation}\label{B:subst}
 \Res_t\bigl(H(u(t))u'(t)\,dt\bigr)
                         =e\Res_s(H(s)\,ds).
\end{equation}
In particular the residue is unchanged by a change of uniformizer.
\end{lemma}
\begin{proof}
It suffices to treat $H(s)=s^n$: only finitely many of its Laurent
monomials can contribute to the residue after substitution. Write
$u=t^e v$ with $v\in k[[t]]^\times$. If $n=-1$, then
$u'/u=e/t+v'/v$ and $v'/v\in k[[t]]$, proving the assertion.
For $n\ne-1$ one must not divide by $n+1$ in characteristic dividing
$n+1$. Instead, perform the calculation first over the torsion-free
ring
\[
 R=\mathbf Z[c_0,c_0^{-1},c_1,c_2,\ldots],\qquad
 v=c_0+c_1t+c_2t^2+\cdots.
\]
The coefficient giving the residue uses only finitely many $c_i$ and
lies in this ring, also for negative $n$. The residue of a derivative
is zero, so differentiating $u^{n+1}$ gives
$(n+1)\Res_t(u^n u'\,dt)=0$. Torsion-freeness gives residue zero in
$R$. Specializing the coefficients to any field proves the desired
identity without division in that field. Linearity now proves
\eqref{B:subst}; for a change of uniformizer, $e=1$.
\end{proof}

\begin{lemma}[A basis compatible with coefficients]\label{B:free-basis}
Let $R$ be a commutative ring and
$u(t)=t^e(c_0+c_1t+\cdots)$, where $c_0\in R^\times$.
Under $s\mapsto u(t)$, the module $R[[t]]$ is free over $R[[s]]$
with basis $1,t,\ldots,t^{e-1}$. After inverting $s$ it is
$R((t))$, free over $R((s))$ with the same basis. These expansions
and the resulting multiplication matrices commute with specialization
of the coefficients.
\end{lemma}
\begin{proof}
For $n=ej+i$, $0\le i<e$, the series $s^jt^i$ has leading term
$c_0^j t^n$. Given a power series, cancel its lowest coefficient
using the corresponding multiple of $s^jt^i$ and continue successively.
This constructs an expansion $\sum_{i=0}^{e-1}H_i(s)t^i$, with
$H_i\in R[[s]]$. It is unique: the first nonzero terms of the
nonzero $H_i(s)t^i$ have different degrees modulo $e$, and hence
cannot cancel. The algorithm uses only ring operations and $c_0^{-1}$.
Every output coefficient is determined after finitely many operations,
which proves compatibility with specialization.
Since $s=t^e v(t)$ with $v(t)$ a unit, inverting $s$ is equivalent
to inverting $t$. This proves the assertions for Laurent series as well.
In particular, the algebraic trace of multiplication by an element of
$R((t))$ is an element of $R((s))$, obtained by the trace of its finite
matrix, and it commutes with specialization.
\end{proof}

\subsection{Trace compatibility}
\begin{lemma}[Totally ramified substitution]\label{B:trace-substitution}
For a field $k$ and $s=u(t)$ of positive order $e$,
\begin{equation}\label{B:trace-subst}
 \Res_s\left(\Tr_{k((t))/k((s))}(h(t))\,ds\right)
       =\Res_t\left(h(t)u'(t)\,dt\right)
                 \qquad(h\in k((t))).
\end{equation}
The trace here is the algebraic trace on the finite free extension.
The identity remains valid even when $u'=0$.
\end{lemma}
\begin{proof}
First suppose that $k$ has characteristic zero. After a finite extension
of the coefficient field, we can take an $e$th root of the leading
coefficient of $u$ and include all $e$th roots of unity. This scalar
extension preserves the trace matrix in Lemma~\ref{B:free-basis}, so an
identity proved there descends. The binomial series then gives a formal
parameter $z=t(c_0+c_1t+\cdots)^{1/e}$ with $s=z^e$.
Write $h=\sum_n b_nz^n$. The extension $k((z))/k((s))$ is now
Galois with automorphisms $z\mapsto\xi z$, $\xi^e=1$. Thus
\[
 \Tr(h)=e\sum_{j\in\mathbf Z}b_{ej}s^j,
 \qquad h\,ds=e\sum_n b_nz^{n+e-1}\,dz.
\]
Both residues equal $e b_{-e}$. The change from $t$ to $z$ preserves
residue by Lemma~\ref{B:substitution}. The trace is continuous for the
Laurent-series topology, either from the displayed automorphisms or
from its finite multiplication matrix, so these calculations apply to
infinite Laurent series, not just Laurent polynomials.

Now take a field of arbitrary characteristic. Lift every coefficient
of $u(t)$ and $h(t)$ to an independent indeterminate in a polynomial
ring over $\mathbf Z$, inverting the nonzero leading coefficient of
$u$. This ring is a characteristic-zero integral domain. The two
residues in \eqref{B:trace-subst} belong to this ring:
Lemma~\ref{B:free-basis} gives this for the trace, and the right side
is an ordinary coefficient of a product and derivative. Each residue
uses only finitely many of the indeterminates. Over the fraction field
of this ring the characteristic-zero proof gives equality. It is
therefore an equality in the ring itself. Specialize the coefficients
to their values in $k$, using Lemma~\ref{B:free-basis} to specialize
the trace. This proves \eqref{B:trace-subst} in every characteristic.
This coefficient argument does not invert the possibly zero integer
$e$ in $k$ and does not assume a tame extension.
\end{proof}

\begin{theorem}[Local trace--residue identity]\label{B:trace-residue}
Let $E/F$ be a finite separable extension of equal-characteristic local
fields with finite residue fields. For every formal differential
$\eta$ on $E$,
\begin{equation}\label{B:residue-trace}
 \Tr_{k_E/k_F}\Res_E(\eta)=\Res_F(\Tr_{E/F}\eta).
\end{equation}
Here if $s$ is a uniformizer of $F$ and $\eta=x\,ds$, the trace of
the differential is $\Tr_{E/F}(x)\,ds$.
\end{theorem}
\begin{proof}
We explain the Laurent-series presentations used in the proof. In an
equal-characteristic complete discrete valuation field with finite
residue field $k$ of size $q$, each residue has a unique lift solving
$X^q-X=0$, by Hensel's lemma with derivative $-1$. In characteristic
$p$ these lifts form a field: sums and products are again roots with
the expected residues. Choosing a uniformizer $t$ and successively
subtracting residue coefficients gives a unique convergent Laurent
expansion. Thus the field is $k((t))$. In a finite extension, these
coefficient fields are compatible, because the lifts of elements of
the smaller residue field are the same unique roots.

Write $F=k_F((s))$ and $E=k_E((t))$. The intermediate field
$F'=k_E((s))$ is unramified over $F$; its trace is the coefficientwise
finite-field trace. In $E$, $s=u(t)$ has positive order $e$, and
Lemma~\ref{B:free-basis} gives $[E:F']=e$. Separability implies
$u'(t)\ne0$: otherwise $u\in k_E((t^p))$, and the intermediate
extension $k_E((t))/k_E((t^p))$ would be purely inseparable of degree
$p$, contradicting separability of $E/F'$. Therefore $ds\ne0$ and
every formal differential on $E$ can be written as $x\,ds$.

Apply Lemma~\ref{B:trace-substitution} to $E/F'$. It gives
$\Res_{F'}(\Tr_{E/F'}(x)\,ds)=\Res_E(x\,ds)$ in $k_E$.
Apply $\Tr_{k_E/k_F}$ and commute it with extraction of the
$s^{-1}$ coefficient. Transitivity of matrix traces proves
\eqref{B:residue-trace}. Residues do not depend on the chosen
uniformizers by Lemma~\ref{B:substitution}.
\end{proof}

\begin{lemma}[Logarithmic derivative of a norm]\label{B:dlog-norm}
For $v\in E^\times$ in the same separable extension,
\begin{equation}\label{B:dlog}
 \frac{d\Nm_{E/F}(v)}{\Nm_{E/F}(v)}
                      =\Tr_{E/F}\left(\frac{dv}{v}\right).
\end{equation}
\end{lemma}
\begin{proof}
A derivation extends uniquely across a finite separable algebraic
extension: for a root $z$ of a separable polynomial $P$, differentiating
$P(z)=0$ gives $dz=- (dP)(z)/P'(z)$. This formula also shows that
differentiation commutes with every $F$-embedding into a normal closure.
Differentiate the finite product
$\Nm_{E/F}(v)=\prod_\sigma\sigma(v)$ and divide by that nonzero
product. The result is
$\sum_\sigma\sigma(dv/v)$, which is the trace of the differential.
Thus \eqref{B:dlog} is elementary algebra, independent of
\eqref{B:residue-trace}.
\end{proof}

Theorem~\ref{B:trace-residue} and Lemma~\ref{B:dlog-norm} supply all
local residue statements formerly cited in the Artin--Schreier
norm-character calculation. The Cartier identities used there are
proved there by their coefficient formulas; no statement from Tate's
paper is required to complete that argument.


\section{Elementary descent and ramification arguments}\label{app:descent}
We include the short algebraic arguments used in character descent, so
Hilbert~90 need not be imported from a reference. The additional
prime-square norm and ramification facts needed for the diamond are
proved in Appendix~\ref{app:diamond-foundations}.

\begin{lemma}[Cyclic Hilbert 90, in both forms]\label{C:hilbert90}
Let $E/F$ be cyclic of degree $n$, with generator $\sigma$.
Then
\[
 \ker\Nm_{E/F}=\{y/\sigma(y):y\in E^\times\},\qquad
 \ker\Tr_{E/F}=\{\sigma(y)-y:y\in E\}.
\]
\end{lemma}
\begin{proof}
Distinct field embeddings are linearly independent as maps. Here is the
short argument: take a nonzero relation of shortest length. Evaluating
at $tx$ and subtracting one embedding's value at $t$ times the relation
at $x$ removes that embedding. Choose $t$ on which another embedding
differs from it. The resulting shorter nonzero relation is a
contradiction.

If $a\in E^\times$ has norm one, put $b_0=1$ and
$b_i=\prod_{j=0}^{i-1}\sigma^j(a)$ for $1\le i\le n$.
By the independence just proved the map
$x\mapsto\sum_{i=0}^{n-1}b_i\sigma^i(x)$ is not identically zero.
Choose $x$ for which its value $y$ is nonzero. Since $b_n=1$ and
$a\sigma(b_i)=b_{i+1}$, one has $a\sigma(y)=y$.
This gives $a=y/\sigma(y)$. The reverse inclusion telescopes under
the norm.

For the additive statement, the trace map is nonzero by the same
independence, and hence is a surjective $F$-linear map to $F$.
Its kernel has dimension $n-1$. The map $\sigma-1$ has kernel $F$,
so its image also has dimension $n-1$. Its image lies in the trace
kernel by telescoping, proving equality. This proof applies when
the characteristic divides $n$.
\end{proof}

\begin{lemma}[Continuous extension of the characters used here]
\label{C:character-extension}
Let $H$ be a subgroup of $E^\times$ containing $U_E^m$ for some
$m\ge1$, and let $\chi:H\to\CC^\times$ be trivial on $U_E^m$.
Then $\chi$ extends to a continuous character of $E^\times$.
The same assertion holds if $\chi$ is first prescribed on several
subgroups, agrees on their intersections, and defines a character on
the subgroup they generate containing such a $U_E^m$.
\end{lemma}
\begin{proof}
The group $E^\times/U_E^m$ is a product of an infinite cyclic group
and the finite group $\OO_E^\times/U_E^m$. Extend a character of a
subgroup one generator at a time. If the first relation of a new
generator $g$ with the existing subgroup is $g^n=h$, choose a nonzero
complex $n$th root of the prescribed value of $h$. If there is no
such relation, choose any nonzero complex value. This defines an
extension at each step, and finitely many steps suffice. Pullback to
$E^\times$ is continuous because $U_E^m$ is open and remains in the
kernel. Agreement and multiplicativity on the generated subgroup are
explicit hypotheses in the last assertion, not conclusions of the
extension argument.
\end{proof}

\begin{lemma}[The first wild ramification term]\label{C:leading-ramification}
Let $\sigma$ be an inertia automorphism of a finite Galois extension
of local fields, let $\pi$ be an upstairs uniformizer, and suppose
\[
 \sigma(\pi)=\pi(1+c\pi^u),\qquad v(c)=0,\qquad u\ge1.
\]
For $x\ne0$ with normalized upstairs valuation $m$, one has
$v(\sigma x-x)\ge m+u$. If $m$ is prime to the residue characteristic,
then equality holds.
\end{lemma}
\begin{proof}
Expand $x=\sum_{j\ge m}a_j\pi^j$ using the multiplicative lifts of
residue representatives. Inertia fixes those lifts: they are the
unique lifts of their residues among roots of unity of order prime
to the residue characteristic (and zero). The series converges in
the local field, with only finitely many negative powers. For any
integer $j$, the binomial expansion at the principal unit gives
\[
 (1+c\pi^u)^j-1\equiv jc\pi^u\pmod{\pp^{u+1}}.
\]
This is also valid for negative $j$, by the convergent inverse series.
Consequently the $j$th summand changes by an element of valuation at
least $j+u$. When $m$ is prime to the residue characteristic, the
first summand changes with exact valuation $m+u$, and all following
summands have strictly larger valuation. This proves the assertions.
The hypothesis $u\ge1$ matters: in a tame row the first coefficient
is instead the residue of $(\sigma\pi/\pi)^m-1$ and may vanish.
\end{proof}

\begin{lemma}[Transitivity of the different]\label{C:different-tower}
For a tower of finite separable local extensions $L/E/F$,
\[
 \mathfrak D_{L/F}
      =\mathfrak D_{L/E}\mathfrak D_{E/F}\OO_L.
\]
\end{lemma}
\begin{proof}
Write $A=\OO_F$, $B=\OO_E$ and $C=\OO_L$. These are finite free
modules along the tower. Via the nondegenerate trace pairings their
dual modules identify with the inverse differents. The elementary
adjunction
\[
 \operatorname{Hom}_A(C,A)
   \simeq\operatorname{Hom}_B(C,\operatorname{Hom}_A(B,A))
\]
sends a functional $\ell$ to $c\mapsto(b\mapsto\ell(bc))$; its
inverse evaluates at $b=1$. Write the principal fractional $B$-ideal
$\mathfrak D_{E/F}^{-1}=dB$. Under the trace identifications, the
right side is $d\operatorname{Hom}_B(C,B)
=d\mathfrak D_{L/E}^{-1}$. Transitivity of field trace shows that
these identifications send $x\in L$ to the same functional
$c\mapsto\Tr_{L/F}(xc)$, with no scalar discrepancy. Thus
$\mathfrak D_{L/F}^{-1}=d\mathfrak D_{L/E}^{-1}$ as fractional
$C$-ideals. Inversion gives the assertion.
\end{proof}

\begin{lemma}[The discriminant of an integral order]\label{C:order-discriminant}
Let $E/F$ be finite separable and let $B\subset\OO_E$ be an
$\OO_F$-order of full rank. Then its discriminant ideal is the field
discriminant ideal multiplied by the square of its index ideal.
In particular the field discriminant divides the order discriminant.
For a totally ramified extension, the normalized different exponent
is at most the $F$-valuation of the order discriminant.
\end{lemma}
\begin{proof}
Choose $\OO_F$-bases of $B$ and $\OO_E$, and let $M$ be the
inclusion matrix. The two trace Gram matrices satisfy
$G_B=M^{\mathsf T}G_{\OO_E}M$, so
$\det G_B=(\det M)^2\det G_{\OO_E}$. This proves the first
assertion and divisibility. The valuation of the field discriminant
is the length over $\OO_F$ of the trace-dual lattice modulo
$\OO_E$: this follows by diagonalizing the trace matrix over the
discrete valuation ring. The trace dual is
$\mathfrak D_{E/F}^{-1}$ by the definition of the inverse different.
Its cyclic-prime specialization underlies the trace-ideal formula of
\cite[Theorem~3.7]{Ueda}. If the different exponent is $d$, this
quotient has $\OO_F$-length $[k_E:k_F]d$, since each successive
$\OO_E$-lattice quotient is $k_E$. In the totally ramified case
$[k_E:k_F]=1$, proving the stated bound.
\end{proof}


\addtocontents{toc}{\protect\enlargethispage{12pt}}
\section{The diamond-level local facts, without local reciprocity}
\label{app:diamond-foundations}
The general local reciprocity theorem and the general quotient theorem
for ramification groups are unnecessary here. We prove their particular
consequences used in this paper. Serre~\cite[Chapters~III--V, XIII--XIV]{Serre}
remains a reference for the classical background and its more general
formulations. The inputs below are the cyclic-prime results of
\cite[Theorem~3.7, Corollary~3.9, and Proposition~3.3]{Ueda}, elementary
finite-field and Galois theory, Hensel lifting, and the usual monogenic
different identity. The latter says
$\mathfrak D_{E/F}=(f'(\pi))$ when
$\OO_E=\OO_F[\pi]$ and $f$ is the minimal polynomial of $\pi$.
It is the specialization of the trace-dual description of the different
to a power basis. Its general algebraic form, simple-root Hensel lifting,
and finite-field cyclicity are already available in Mathlib~\cite{Mathlib}.
For completeness we give the short quantitative Newton argument below;
this avoids confusing a theorem for $\mathbf Z_p$ with one for every
local field.
The declaration and hypothesis matches are recorded separately, not
inserted into the mathematical argument.

We also use the elementary degree formula $[E:F]=e(E/F)f(E/F)$,
which is part of the local-field infrastructure of the FML development.
For a totally ramified extension, any uniformizer $\pi$ generates the
integer ring: $F(\pi)=E$ follows from the degree formula, and in the
basis $1,\pi,\ldots,\pi^{n-1}$, the valuations
$n v_F(a_i)+i$ of the terms $a_i\pi^i$ are distinct modulo $n$.
Thus an integral sum has every $a_i$ integral. This supplies the
monogenic hypothesis whenever the displayed different identity is used.

The cyclic coordinates used below need no further classification
input. In characteristic $p$, additive Hilbert~90 applied to $1$ in a
cyclic degree-$p$ extension gives $\sigma z-z=1$. Then $z^p-z\in F$,
and $z$ generates the extension. In degree two and characteristic
other than two, a nonzero $w=u-\sigma u$ has $\sigma w=-w$;
therefore $w^2\in F$ and $E=F(w)$. The reductions of these coordinates
at a specified valuation are proved in the corresponding case calculations.

\begin{lemma}[Quantitative lifting over the local fields used here]
\label{D:newton}
Let $R$ be a complete discrete valuation ring with valuation $v$, let
$f\in R[X]$, and let $a\in R$. Suppose $s=v(f'(a))<\infty$ and
$v(f(a))>2s$. Then there is a unique root $b$ in the ball
$v(b-a)>s$. Moreover $v(b-a)\ge v(f(a))-s$.
\end{lemma}
\begin{proof}
The case $f(a)=0$ is immediate, including uniqueness by the argument
below. Otherwise iterate $a_{n+1}=a_n-f(a_n)/f'(a_n)$.
Taylor's polynomial identity with integral coefficients shows
inductively that $v(f'(a_n))=s$, and, writing
$r_n=v(f(a_n))-2s>0$, that
\[
 v(a_{n+1}-a_n)=s+r_n,\qquad r_{n+1}\ge2r_n.
\]
If an iterate is a root, stop. The derivative differences have valuation
at least $s+r_n>s$, so they cannot change $s$; the linear Taylor term
cancels and the remaining terms have valuation at least $2(s+r_n)$.
Thus the iterates are integral and converge to a root with the asserted
error bound. For $b,c$ in the stated ball, divided differences give
$f(b)-f(c)=(b-c)(f'(a)+\varepsilon)$ with $v(\varepsilon)>s$.
The second factor is nonzero; hence two roots there are equal.
The proof uses only completeness, polynomial identities, and the
nonarchimedean valuation inequalities. Rescaling the variable and the
polynomial gives the fractional-coordinate applications wherever used.
\end{proof}

\begin{lemma}[Inertia, tame action, and subgroup numbering]\label{D:inertia}
For a finite Galois extension $K/F$ of local fields, let $G=\Gal(K/F)$,
let $I$ be the kernel of its residue action, and put $q=|k_F|$.
The residue action identifies $G/I$ with the cyclic group
$\Gal(k_K/k_F)$. Its fixed field $K^I/F$ is unramified.
The quotient $I/G_1$ is cyclic of order prime to the residue
characteristic $p$, and $G_1$ is a $p$-group.
A lift of residue Frobenius conjugates $I/G_1$ by the $q$th power map.
For $H\subset G$, the lower ramification groups on $K/K^H$ are
$H_j=H\cap G_j$ for every integer $j\ge0$.
\end{lemma}
\begin{proof}
Choose a generator $\bar a$ of $k_K/k_F$, lift its monic irreducible
polynomial to $\OO_F[X]$, and lift $\bar a$ to a root $a\in\OO_K$
by Hensel's lemma. The field $F(a)$ has degree at most
$f=[k_K:k_F]$ and residue degree at least $f$, hence degree exactly
$f$ and ramification index one. All residue conjugates of $\bar a$
lift uniquely inside $F(a)$, so $F(a)/F$ is Galois and reduction
identifies its group with $\Gal(k_K/k_F)$. Normality of $K/F$
extends its automorphisms to $K$. An element of $G$ acts trivially on
the residue field precisely when it fixes $a$, by uniqueness of the
simple-root lifts. This proves the assertions about $I$ and $G/I$.

Choose a uniformizer $\pi$ of $K$. Inertia fixes the multiplicative
lifts of every residue element. Expanding integers in powers of $\pi$
therefore tests lower ramification on $\pi$ alone. In particular,
\[
 I\longrightarrow k_K^\times,\qquad
 g\longmapsto\overline{g\pi/\pi}
\]
is a homomorphism with kernel $G_1$. For $j\ge1$ the map
\[
 G_j\longrightarrow k_K,\qquad
 g\longmapsto\overline{(g\pi/\pi-1)/\pi^j}
\]
is additive with kernel $G_{j+1}$. These statements follow by
multiplying the two uniformizer ratios and reducing the first
nonzero coefficient. The finite subgroup of $k_K^\times$ is cyclic
and has order prime to $p$; each additive image is a $p$-group.
The filtration eventually equals $1$, since every nonidentity inertia
automorphism moves $\pi$. This proves the assertion about $G_1$.

The tame residue is unchanged on replacing $\pi$ by $u\pi$ for a
unit $u$, since inertia fixes its residue. For a Frobenius lift $\phi$,
compute using $\pi'=\phi^{-1}\pi$:
\[
 \overline{(\phi g\phi^{-1})(\pi)/\pi}
   =\overline{\phi}\bigl(\overline{g(\pi')/\pi'}\bigr)
   =\overline{g\pi/\pi}^{\,q}.
\]
Finally, membership in a lower ramification group is the condition
$v_K(gx-x)\ge j+1$ for all $x\in\OO_K$. For an element of $H$
this condition is identical over $F$ and over $K^H$: both use the same
normalized valuation on $K$. This proves the subgroup assertion.
\end{proof}

\begin{lemma}[The breaks of a totally ramified prime-square diamond]
\label{D:break-ledger}
Let $K/F$ be totally ramified with $G=C_p^2$, where $p$ is the residue
characteristic. There are integers $1\le t\le b$ such that the
filtration is $G$ through $t$, possibly one line $H_0$ of order $p$
through $b$, and then $1$. Choose $L_0=K^{H_0}$ if $b>t$, and
choose any degree-$p$ field if $b=t$. The break of $L_0/F$ is $t$;
the break of $K/L_0$ is $b$. Every other lower field $L/F$ has break
\[
 s=t+(b-t)/p,
\]
and $K/L$ has break $t$. In particular $s$ is an integer, and
$b=t+p(s-t)$. Inertia of order $p$, rather than $p^2$, gives instead
one unramified lower field and ramified lower fields whose breaks are
unchanged by the unramified base extension.
\end{lemma}
\begin{proof}
Lemma~\ref{D:inertia} gives $G_1=G$. A decreasing filtration of
$C_p^2$ can contain at most one proper nontrivial subgroup.
The subgroup assertion gives the breaks on the two upper edges.
For a top uniformizer $\pi$, the monogenic different formula gives
\[
 d_{K/F}=v_K\prod_{g\ne1}(\pi-g\pi)
       =\sum_{j\ge0}(|G_j|-1)
       =(p^2-1)(t+1)+(p-1)(b-t).
\]
Here each nonidentity $g$ is counted exactly
$v_K(\pi-g\pi)$ times. Transitivity of the different and the
cyclic-prime formula give
\[
 d_{K/F}=(p-1)(b+1)+p\,d_{L_0/F}
         =(p-1)(t+1)+p\,d_{L/F}.
\]
Solving these equalities proves the stated lower-edge breaks,
including integrality of $s$, without a quotient-numbering theorem.

If inertia has order $p$, its fixed field $U$ is the unique unramified
lower field, and the other lower fields $E$ are totally ramified.
In $K=EU$, a uniformizer of $E$ is still a uniformizer of $K$.
The generator on $K/U$ restricts to the generator on $E/F$, and the
normalized valuation restricts without a multiplier on $E$.
Its uniformizer difference thus has the same valuation on both edges.
\end{proof}

\begin{proposition}[Separation of the norm kernels needed for SML]
\label{D:norm-separation}
Let $K/F$ have group $C_\ell^2$. The norm kernels
$N_{L_0/F}L_0^\times$ and $N_{L/F}L^\times$ are distinct if
$L_0$ is the unramified lower field, or, in total ramification,
$L_0$ is the distinguished field of Lemma~\ref{D:break-ledger} and
$L\ne L_0$. In the totally ramified one-break case all lower norm
kernels are pairwise distinct.
\end{proposition}
\begin{proof}
Every lower norm subgroup has index $\ell$ by the cyclic-prime results
of FML. With an unramified lower field its nontrivial norm characters
have conductor zero, whereas those of every ramified lower field
have positive conductor. Equal norm kernels would give equal character
groups, contrary to those conductors. In the totally ramified two-break
case, Lemma~\ref{D:break-ledger} gives the distinct lower conductors
$t+1<s+1$, proving the assertion in the same way.

It remains to distinguish the kernels when every edge has break $t$.
Here $\ell=p$. Fix $\pi\in K$ of valuation one and put
\[
 c_g=\overline{(g\pi/\pi-1)/\pi^t}\in k,
       \qquad k=k_F=k_K.
\]
Lemma~\ref{D:inertia} makes $g\mapsto c_g$ an injective additive
map, with two-dimensional $\F_p$-image $V$. For a line
$H=\langle\sigma\rangle$, set $L_H=K^H$, $c=c_\sigma$,
$\Pi_H=N_{K/L_H}\pi$, and $\varpi=N_{K/F}\pi$.
For $\tau\notin H$ write $e=c_\tau$. FML's critical norm polynomial
(\cite[Proposition~3.3]{Ueda}) gives
\[
 P_H(X)=X^p-c^{p-1}X,\qquad
 \overline{(\tau\Pi_H/\Pi_H-1)/\Pi_H^t}
       =e^p-c^{p-1}e=:d_H\ne0.
\]
The second equality is the first applied to
$N_{K/L_H}(\tau\pi/\pi)=\tau\Pi_H/\Pi_H$.
Apply the same critical norm theorem now to $L_H/F$ with the
compatible uniformizers $\Pi_H,\varpi$. Its critical norm image is
\[
 d_H^p\ker\Tr_{k/\F_p},
\]
whose annihilator is the line $d_H^{-p}\F_p$ in the residue trace
pairing. This line is precisely the restriction of the norm-character
group to the last unit layer, by the norm-filtration theorem and the
exact norm-character conductor.

Choose a basis $c_1,c_2$ of $V$ and put
$D=c_1c_2^p-c_1^pc_2\ne0$. For any ordered basis $c,e$ of $V$,
\[
 c e^p-c^pe=aD\quad\text{for some }a\in\F_p^\times,
 \qquad d_H=aD/c.
\]
Thus the annihilator line is $(c^p/D^p)\F_p$.
Distinct lines $H$ give distinct lines $c\F_p$, and Frobenius is
bijective, so these annihilators are distinct. The norm-character
groups, and hence their norm kernels in $F^\times$, are distinct.
This argument identifies whole last-layer characters, not just one
trivial character value.
\end{proof}

\begin{lemma}[Crossed norm characters from distinct norm kernels]
\label{D:crossed-norm}
Let $L_i,L_j$ be distinct lower fields in a $C_\ell^2$ diamond, and
suppose their lower norm kernels are distinct. Then
\[
 S(L_j/F)\longrightarrow S(K/L_i),\qquad
       \omega\longmapsto\omega\circ N_{L_i/F}
\]
is an isomorphism. In particular this applies to every distinguished
comparison in Proposition~\ref{D:norm-separation}.
\end{lemma}
\begin{proof}
Norm transitivity shows that every displayed pullback kills
$N_{K/L_i}K^\times$. A nontrivial $\omega$ has kernel
$N_{L_j/F}L_j^\times$, since that quotient has prime order $\ell$.
Its pullback cannot be trivial: otherwise one index-$\ell$ lower norm
subgroup would be contained in the other and hence equal to it.
Both character groups have order $\ell$ by FML, so the displayed
homomorphism is an isomorphism. This is the required base-change
identity proved directly, not an appeal to local reciprocity.
\end{proof}

\begin{lemma}[All three quadratic norm characters]\label{D:quadratic-norms}
In a biquadratic local extension, the three nontrivial lower norm
characters are pairwise distinct and, after labeling, satisfy
$\omega_3=\omega_1\omega_2$. Consequently
Lemma~\ref{D:crossed-norm} applies to every pair of its lower fields.
\end{lemma}
\begin{proof}
If $\operatorname{char}F\ne2$, associate to a square class $a$ the
norm character of $F(\sqrt a)/F$, and associate the trivial character
to the square class of $1$. Denote its value at $b$ by $(a,b)$.
The existence of a nonzero solution of
\[
 X^2=aY^2+bZ^2
\]
is equivalent to $(a,b)=1$. If $a$ is nonsquare, any such solution
has $Z\ne0$ and gives $b=N(X/Z+(Y/Z)\sqrt a)$; the split case is
immediate. The conic is symmetric in $a,b$, so $(a,b)=(b,a)$.
Multiplicativity in the second variable is the norm-character
property; symmetry gives it in the first variable. A nonsquare class
has nontrivial character by FML's index-two norm theorem. Thus the
map from square classes to quadratic characters is an injective
homomorphism. The three fields have classes $a,b,ab$, proving the claim.

If $\operatorname{char}F=2$, use the explicitly proved formula
of Lemma~\ref{D:EQ:ASsymbol}:
\[
 \omega_f(u)=(-1)^{\Tr_{k/\F_2}\operatorname{Res}_F(f\,du/u)}.
\]
Its proof uses only the single quadratic norm index, the Cartier
calculation, and Appendix~\ref{app:local-residues}; in particular it
uses no crossed norm-character assertion or SML comparison.
The formula is additive in the Artin--Schreier class $f$, and the same
proof establishes nontriviality for every nonzero class, ramified or
unramified. The classes of the three fields are $f,g,f+g$, giving
pairwise distinctness and the product identity.
\end{proof}

\begin{remark}[The scope of the replacements]\label{D:scope}
For odd $\ell$ the proof of SML compares a distinguished lower field
with every other one and then uses transitivity of equality. It does
not require norm separation for an arbitrary pair of larger-break
lower fields, or the general norm correspondence for all abelian
extensions. For $\ell=2$, Lemma~\ref{D:quadratic-norms} supplies every
pair as well as the product of the three lower characters. All the
ramification numbers used in the final case classification follow
from Lemmas~\ref{D:inertia}--\ref{D:break-ledger}. Thus neither full
local class field theory nor the general Herbrand quotient theorem is
an unstated prerequisite. The proofs above give the additional
local arguments; the elementary matrix, field, and finite-group
bookkeeping need not be reproduced in greater generality here.
\end{remark}


\begin{thebibliography}{99}
\bibitem{Dwork} B.~M.~Dwork,
\emph{On the Root Number in the Functional Equation of the Artin--Weil
$L$-Series}, Ph.D. thesis, Columbia University, 1954.
\bibitem{Lakkis} Konstantin Lakkis,
\emph{Die Galoisschen Gau{\ss}schen Summen von Hasse}, original source.
Historical and methodological reference only.
\bibitem{Langlands} R.~P.~Langlands,
\emph{On the Functional Equation of the Artin $L$-functions},
Yale University preprint (1970), IAS electronic edition.
\bibitem{Mathlib} The mathlib Community, \emph{Mathlib}, the Lean~4 mathematical
library. The source revision and exact algebraic declarations used as
background are recorded in the companion dependency inventory.
\bibitem{Serre} Jean-Pierre Serre, \emph{Local Fields}, translated by
Marvin Jay Greenberg, Graduate Texts in Mathematics, vol.~67,
Springer-Verlag, 1979.
\bibitem{Stickelberger} L. Stickelberger,
\emph{Ueber eine Verallgemeinerung der Kreistheilung},
Mathematische Annalen \textbf{37} (1890), 321--367.
Historical reference for the leading congruence proved in
Appendix~\ref{app:leading-gauss}; no theorem is quoted from this paper.
\bibitem{TateResidue} John Tate, \emph{Residues of differentials on curves},
Annales scientifiques de l'\'Ecole Normale Sup\'erieure (4)
\textbf{1} (1968), 149--159. Historical reference for local residues;
the formal substitution and trace-compatibility statements used here
are proved in Appendix~\ref{app:local-residues}.
\bibitem{Ueda} Fukuhiro Ueda,
\emph{A Local Proof of Langlands's First Main Lemma for Local Epsilon
Factors over Nonarchimedean Local Fields}, mathematical manuscript.
August 2026. The theorem and lemma numbers cited in the text refer to
this version.
\end{thebibliography}


\end{document}
